\section{The two methods: precise transfers and repairs}
\label{sec:methods}

\subsection*{MT-A01. What the adopted Dirichlet theorem says adelically}
The supplied declarations \texttt{riemannHypothesis},
\texttt{LFunction\_ne\_zero\_of\_half\_lt\_re}, and
\texttt{LFunction\_eq\_zero\_re} in \texttt{ALPOGE\_Main.lean} have the
definitions printed in \texttt{ALPOGE\_Statements.lean}.  We use these as the
cited input $A_D$, as requested.  The supplied archive contains these exported
declarations and the scan report, but not the imported tree
\texttt{Supplements.Kernels.MainTheoremsAll}; a kernel build of that tree is
not a check performed in this run.  This distinction concerns our evidence,
not the stipulated status of the input theorem.

Every unitary idele-class character of $\mathbb Q$ has the form
$\chi_0|\cdot|^{it}$, where $\chi_0$ is a finite-order Dirichlet character
and $t\in\mathbb R$.  Consequently $A_D$ gives the rank-one Hecke GRH over
$\mathbb Q$ for these characters, not just its finite-order specialization.
Indeed the norm map has a continuous section and compact kernel, which for
$\mathbb Q$ is $\widehat{\mathbb Z}^{\times}$ after the sign convention at
infinity.  A continuous character of this profinite group into the circle has
finite image: its image is a compact quotient of a profinite group, hence
totally disconnected, whereas an infinite closed subgroup of the circle is
the circle.  It therefore factors through some
$(\mathbb Z/N\mathbb Z)^\times$.  The remaining unitary character of
$\mathbb R_{>0}$ is $x^{it}$.  In the convention
$L(s,\chi_0|\cdot|^{it})=L(s+it,\chi_0)$, zeros are translated only in
their imaginary coordinate.  This proves the assertion.  Geometrically,
the finite-order part is a rank-one representation of the arithmetic
fundamental group; the continuous norm twist is an analytic normalization,
not an additional finite-monodromy motive.  The result does not cover
arbitrary characters over a nontrivial number field.

\subsection*{MT-M01. A nonzero counterexample to the draft's general zero-count lemma}
Masullo, Lemma 9.3.2.4, PDF p.~425, claims a local $O(\log |T|)$ zero count
for the Mellin transform of every $\varphi\in\mathcal S_E(G)$.
We checked the image of the original PDF page, including the displayed
Jensen inequality and its last inference.  The quantifier is present in the
PDF; this is not an inference from the model transcription.

Take $G=\mathbb R_{>0}$, trivial character, Haar measure $dt/t$, and set
\[
 g_a(u)=\frac{1}{\sqrt{4\pi a}}\exp\left(-\frac{u^2}{4a}\right),
 \qquad a=1+i,\qquad
 \varphi(e^u)=\frac{g_{1+i}(u)-g_{1-i}(u)}{2i}.
\]
The square roots use the branch with positive real part.  This is a real,
nonzero smooth function.  For every real $b$ and every nonnegative integer
$j$, the $j$th derivative of $e^{bu}g_{1\pm i}(u)$ is a polynomial in $u$
times $e^{bu-u^2/(4(1\pm i))}$.  Its absolute value is bounded by a
polynomial times $e^{bu-u^2/8}$, hence is rapidly decreasing.  Thus
$t^b\varphi(t)$ is Schwartz in the logarithmic coordinate for every $b$;
this is exactly the condition of Definition D.3.3.1, PDF p.~804 (also
\S1.1.2, PDF p.~11).  In particular $\varphi\in\mathcal S_E(G)$.

For $\operatorname{Re}a>0$, the Gaussian integral and analytic continuation
in $a$ give
\[
 \int_{\mathbb R}g_a(u)e^{su}\,du=e^{as^2},\qquad s\in\mathbb C.
\]
For completeness, the identity first follows by completing the square for
positive real $a$ and real $s$; both sides are entire in $s$ by Gaussian
domination on compact sets, and holomorphic in $a$ on
$\operatorname{Re}a>0$, where $\operatorname{Re}(1/a)>0$.  The identity
theorem proves the claimed continuation.  Therefore
\[
 \mathcal M\varphi(s)=e^{s^2}\sin(s^2).
\]
Its zeros on the positive imaginary axis are
$i\sqrt{k\pi}$, $k=1,2,\ldots$, and are simple.  For $I=[-1,1]$ and
$T\ge2$, their number in $I+i[T,T+1]$ is
\[
 \#\{k:T^2/\pi\le k\le(T+1)^2/\pi\}
   =\frac{2T+1}{\pi}+O(1).
\]
This is not $O(\log T)$.  The function is nonzero, so removing the zero
function from the lemma does not repair it.  The proof on p.~425 needs a
\emph{lower} bound for the modulus at a Jensen centre, whereas rapid decay
supplies upper bounds.  No statement about the zeros of a Hecke L-function
is contradicted by this example.

\subsection*{MT-M02. The displayed approximate identities are proper kernel projections}
Theorem 9.3.2.6, equation (9.3.2.10), PDF p.~426, defines subspaces by
intersecting kernels of all jet evaluations up to height $T$, then claims
their orthogonal projections tend strongly to the identity.  The printed
formula and the asserted inclusion in its proof have both been checked on
the page image.

Here is a counterexample even with a single zero, so the zero-count issue
in MT-M01 is irrelevant.  Put
\[
 g(u)=(4\pi)^{-1/2}e^{-u^2/4},\qquad
 \psi(e^u)=-g'(u).
\]
Then $\psi\in\mathcal S_E(\mathbb R_{>0})$, and integration by parts gives
$\mathcal M\psi(z)=z e^{z^2}$.  Its only zero is $0$, with multiplicity one,
strictly inside the strip for $I=[-1,1]$.  Choose Sobolev index $0$.
The draft's Hilbert norm is then
\[
 \|f\|_H^2=\int_{\mathbb R}|f(e^u)|^2e^{2|u|}\,du.
\]
The functional $\delta_0(f)=\int_{\mathbb R}f(e^u)\,du$ is continuous,
because Cauchy--Schwarz bounds its modulus by
$\|f\|_H(\int e^{-2|u|}\,du)^{1/2}=\|f\|_H$.
It is nonzero, since $\delta_0(g)=1$.  For every $T>0$ the subspace displayed
in (9.3.2.10) is the same proper closed subspace $\ker\delta_0$.
Its orthogonal projector cannot converge to the identity.  In particular,
the Riesz vector $r(e^u)=e^{-2|u|}$ has $\|r\|_H=1$ and is killed by every
one of these projectors.  It belongs to the indicated completion because
$C_c^\infty(\mathbb R)$ is contained in $\mathcal S_E$ in logarithmic
coordinates and is dense in this weighted $L^2$ space: truncate first and
then smooth on compact intervals, where the positive weight is bounded
above and below.  The ranges decrease when more evaluation
conditions are added, whereas the proof on p.~426 uses the opposite
inclusion.

\subsection*{MT-M03. The exact projection repair}
Let $H$ be a Hilbert space, let $(\ell_j)_{j\ge1}$ be continuous linear
functionals, and put
\[
 V_N=\bigcap_{j\le N}\ker\ell_j,\qquad V=\bigcap_{j\ge1}\ker\ell_j.
\]
Then $P_{V_N}\to P_V$ strongly.  The complementary projectors converge
strongly to $P_{V^\perp}$, the projector onto the closed span of the Riesz
vectors of the $\ell_j$.  Thus they give an approximate identity on the
spectral quotient $H/V\simeq V^\perp$.

Here is a full proof.  For fixed $x\in H$, put $y_N=P_{V_N}x$.  For
$M\ge N$, nestedness gives $P_{V_M}y_N=y_M$ and hence
\[
 \|y_N-y_M\|^2=\|y_N\|^2-\|y_M\|^2.
\]
The nonincreasing nonnegative sequence $\|y_N\|^2$ converges, so $(y_N)$ is
Cauchy.  Write its limit as $y$.  For each fixed $N$, all sufficiently late
$y_M$ lie in the closed space $V_N$, so $y\in V$.  Also
$x-y_N\perp V$ for every $N$; passing to the limit gives $x-y\perp V$.
Therefore $y=P_Vx$.  Taking complements proves the second assertion.
Finally the orthogonal complement of the closed span of the Riesz vectors
is precisely $V$, which proves the last identification.  Exhaustion by
height instead of an integer enumeration gives the same proof.

This repairs a functional-analytic direction error.  It does not identify
the common kernel with the closed adelic theta image, prove a Schauder
basis, establish nuclearity of the actual convolution operators, or justify
the geometric-kernel decomposition.  Those are separate statements and
must be checked in the quotient on which they are used.  In Grothendieck's
trace language this is the distinction between an annihilator and its
spectral quotient, before any trace or Fredholm determinant is taken.

\subsection*{MT-M04. An arithmetic replacement for the false general lemma}
For every fixed unitary Hecke character, the number of nontrivial zeros
in a unit-height window is $O_{K,\eta}(\log(|T|+3))$, unconditionally.
The full proof for arbitrary infinity type is NT-Z10 above: polynomial
vertical growth and a uniform lower bound at the Euler-product centre
give Jensen's estimate.  It includes pure norm twists, whose single
translated pole is excluded before applying Jensen at large height.
For a fixed primitive finite-order character a conductor-explicit
primary locator is Greni\'e--Molteni,
\emph{An explicit Chebotarev density theorem under GRH}, Appendix A,
Theorem A.1, author PDF p.~26: that appendix's zero count is unconditional.
Subtract the counts at consecutive heights.  The main term has derivative
$O(\log(|T|+3))$ for a fixed field and character and each remainder has the
same bound.  The arithmetic restriction therefore closes the local
zero-count gate, including infinite-order characters, without using GRH.
It supplies a route around MT-M01 for arithmetic divisors, but does not
establish Lemma 9.3.2.5 for arbitrary Schwartz convolution ideals.  The
constant in the infinite-order statement is fixed-character, not a
claimed uniform archimedean-conductor constant.

\subsection*{MT-ART01. A nonmonomial GL\texorpdfstring{$_2$}{2} consequence: odd Artin forms}
Assume $M_H$.  Let $\rho:G_{\mathbb Q}\to\operatorname{GL}_2(\mathbb C)$
be continuous, irreducible, odd and of finite image.  Then its Artin
L-function is entire and all of its nontrivial zeros lie on
$\operatorname{Re}s=1/2$.

The holomorphy input is already a theorem: Khare--Wintenberger,
\emph{Serre's modularity conjecture (I)}, Corollary 10.2(ii), author PDF
p.~21 (with Theorem 10.1(ii), pp.~20--21), associates $\rho$ to a weight-one
cuspidal newform.  Its standard L-function is entire.  Let $L$ be the finite
Galois extension cut out by $\rho$.  Brauer induction writes its character
as an integer linear combination of characters induced from
one-dimensional characters of subgroups of $\operatorname{Gal}(L/\mathbb Q)$.
Artin induction and class field theory therefore express $L(s,\rho)$ as
a finite product of integer powers of finite-order Hecke L-functions over
the corresponding fixed fields.  In the open critical strip, $M_H$ makes
each factor holomorphic and nonzero away from the critical line.  The
product, including its negative exponents, is consequently holomorphic and
nonzero there.  Combining this with the independent holomorphy theorem
proves GRH for $L(s,\rho)$.

This includes projective image $A_5$: no monomial expression with only
nonnegative exponents is required.  It is a new transfer in this catalogue,
not a new proof of the classical modularity theorem.  In motivic terms it
concerns odd rank-two Artin motives and their weight-one automorphic
realizations.  It gives no GRH for the weight-two L-function of an arbitrary
non-CM elliptic curve.  Nor does it prove holomorphy of arbitrary even
icosahedral or higher-dimensional Artin representations.  The inherited
Brauer divisor-support statements in RETURN\_ZETA\_TWO are the general
input behind this specialization.

\subsection*{MT-ART02. The missing positivity is a finite matrix at each zero}
Let $L/K$ be a finite Galois extension with group $G$, and let
$0<\operatorname{Re}s_0<1$.  By Brauer meromorphy define
\[
 n_\chi=\operatorname{ord}_{s=s_0}L(s,\chi,L/K)\in\mathbb Z,
 \qquad\Theta_{s_0}=\sum_{\chi\in\operatorname{Irr}(G)}n_\chi\chi.
\]
Form the $|G|$ by $|G|$ matrix
\[
 K_{g,h}=\Theta_{s_0}(g h^{-1}).
\]
Then all Artin L-functions for $L/K$ are holomorphic at $s_0$ if and only
if $K$ is positive semidefinite.  The matrix is Hermitian because
$\chi(g^{-1})=\overline{\chi(g)}$ and $n_\chi$ is real.

To prove the equivalence, consider the central convolution operator
$T=|G|^{-1}K$ on $\mathbb C[G]$.  On the matrix-coefficient space of an
irreducible character $\chi$, Schur orthogonality makes $T$ scalar with
eigenvalue $n_\chi/\chi(1)$ (interchanging $\chi$ and its dual under the
opposite convolution convention only permutes the list).  Each such space
has dimension $\chi(1)^2$ and together they exhaust the regular
representation.  Thus $K\succeq0$ is equivalent to $n_\chi\ge0$ for every
$\chi$, which is exactly absence of a pole of every factor.

The restrictions of $\Theta_{s_0}$ to cyclic subgroups are already actual
characters.  Indeed, for a cyclic subgroup $C$ and a one-dimensional
character $\eta$ of $C$, Frobenius reciprocity and Artin induction give
\[
 \langle\Theta_{s_0}|_C,\eta\rangle_C
   =\operatorname{ord}_{s=s_0}L(s,\eta,L/L^C)\ge0,
\]
because the right-hand side is an order of a holomorphic Hecke function in
the strip.  This proves positivity on every such restriction; it does not
prove the full matrix is positive.  Also
\[
 \Theta_{s_0}(1)=\sum_\chi n_\chi\chi(1)
   =\operatorname{ord}_{s=s_0}\zeta_L(s).
\]
The point of the criterion is the exact remaining object: a central
positive kernel in the Artin character ring.  $M_H$ removes all tests off
the critical line, while neither zero location nor cyclic positivity
supplies the full critical-line kernel.  An Alp\"oge-type positivity
argument would have to produce this kernel or an equivalent positive
decomposition, with the actual Artin orders as coefficients.

\subsection*{MT-ART04. All induced-linear positivity tests still leave a gap}
There is a virtual character which passes every positivity inequality
obtained by inducing a one-dimensional character of a subgroup, but is not
an actual character.  This refutes an abstract reconstruction rule; we do
not claim that this virtual character occurs as the orders of actual
Artin L-functions at a point.

Let $G=A_5$, and let $\chi$ be its four-dimensional augmentation character
for the natural action on five letters.  Define
\[
             \Theta=\operatorname{reg}_G-5\chi.
\]
The natural permutation action is doubly transitive, so its permutation
character has inner-product norm squared two.  Its invariant subspace is
one-dimensional; the orthogonal complement therefore has character
$\chi$ with norm squared one, proving that $\chi$ is irreducible.
The coefficient of $\chi$ in $\Theta$ is $4-5=-1$, so $\Theta$ is not an
actual character and its full matrix from MT-ART02 is not positive.

Nevertheless, for every $H\le G$ and every linear character $\eta$ of $H$,
\[
       \left\langle\Theta,\operatorname{Ind}_H^G\eta\right\rangle_G
                    \ge0.                                      \tag{MT.1}
\]
Here are all the details.  The conjugacy classes of $A_5$ have sizes
$1,15,20,12,12$, as follows by counting the double transpositions,
three-cycles and the two five-cycle classes.  A normal subgroup is a union
of these classes containing the identity, and its size divides $60$.
No proper nontrivial sum of these class sizes containing $1$ divides $60$;
thus $A_5$ is simple.  An action on cosets of a proper subgroup of index
$2,3,$ or $4$ would be an injective homomorphism into the corresponding
symmetric group, which is impossible by order.  Lagrange's theorem leaves
only the possible proper indices
\[
                    k\in\{5,6,10,12,15,20,30,60\}.
\]
Put $m=\langle\chi|_H,\eta\rangle_H$.  Since $\chi|_H$ has dimension four,
$0\le m\le4$.  Frobenius reciprocity says that
$\operatorname{Ind}_H^G\eta$, of dimension $k$, contains $m$ copies of
$\chi$, so $4m\le k$.  These two bounds imply $k-5m\ge0$ for every index
in the list except possibly $k=12$.  At $k=12$, $|H|=5$, and $H$ acts
regularly on the five letters.  The restriction of $\chi$ is the sum of
the four distinct nontrivial characters of $C_5$, so $m\le1$ and again
$k-5m\ge0$.  Finally
\[
 \left\langle\operatorname{reg}_G,
       \operatorname{Ind}_H^G\eta\right\rangle=k,
\]
proving (MT.1) for proper $H$.  For $H=G$, simplicity and nonabelianity
make the trivial character the only linear character, and its coefficient
in $\Theta$ is one.  This proves (MT.1) in all cases.

For comparison with the methods, Artin induction and Hecke holomorphy
give exactly inequalities of the form (MT.1) for an actual local Artin
order character.  The example proves that these inequalities alone do not
force a positive character, even before any question of zero location.
Thus a transfer of determinant positivity from the Alp\"oge method would
need additional arithmetic information, not just all the existing
induced-linear nonnegative orders.  The elementary exact computation
\texttt{checks/artin\_cone\_certificate.py} checks the negative matrix
quadratic form: with $v(g)=\chi(g)$ one has $v^*Kv=-900$.

\subsection*{MT-ART03. The classical simple-zero reduction}
Stark, \emph{Some effective cases of the Brauer--Siegel theorem},
Inventiones Mathematicae \textbf{23} (1974), Theorem 3, p.~144, proves
that if $\operatorname{ord}_{s_0}\zeta_L\le1$ and $s_0\ne1$, then all
Artin factors for $L/K$ are holomorphic at $s_0$.  We cite this classical
result rather than count another proof as a finding.  Applied with the
divisor-support theorem under $M_H$, it confines any still-possible pole
inside the strip to a multiple zero of the top-field zeta on the critical
line.  No simplicity assertion for these Dedekind zetas is being made.

\subsection*{MT-Q01. Exact path quadrics in every characteristic}
For $k\ge2$ let
\[
 V_k=\left\{(x_1,\ldots,x_k):\sum_{i=1}^{k-1}x_ix_{i+1}=0\right\}
       \subset\mathbb A^k_{\mathbb Z},\qquad
 r=\lfloor k/2\rfloor,\quad d=\lceil k/2\rceil.
\]
Then for every finite field $\mathbb F_q$, including characteristic two,
\[
 \#V_k(\mathbb F_q)=q^{k-1}+(q-1)q^{d-1}.                 \tag{MT.2}
\]
In the relative Grothendieck ring of varieties over $\mathbb Z$ one has
\[
                 [V_k]=\mathbb L^{k-1}+\mathbb L^d-\mathbb L^{d-1}.
                                                               \tag{MT.3}
\]
Here $\mathbb L=[\mathbb A^1_{\mathbb Z}]$; the identity is an identity of
classes, not a claim of a direct-sum splitting of all integral realizations.

To prove these formulas, group the odd coordinates into $x\in\mathbb A^d$
and the even coordinates into $y\in\mathbb A^r$.  The quadratic form
becomes $x^{\mathsf T}Ay$, where the integral $d$ by $r$ matrix $A$ has an
$r$ by $r$ triangular minor with diagonal entries one.  It has a left
inverse over $\mathbb Z$ and its columns extend to a unimodular basis of
$\mathbb Z^d$.  Thus after an invertible integral change of the $x$
coordinates the equation is $\sum_{j=1}^r x_jy_j=0$, with $d-r$ free
$x$ coordinates.  Over $y=0$ its fibre is $\mathbb A^d$.  Over
$\mathbb A^r\setminus\{0\}$, the map of vector bundles
$\mathcal O^d\to\mathcal O$ given by $(y_1,\ldots,y_r,0,\ldots,0)$ is
surjective, and its kernel is a vector bundle of rank $d-1$.
It is Zariski locally trivial on the opens where a coordinate $y_j$ is
invertible.  Additivity and the vector-bundle relation give
\[
 [V_k]=\mathbb L^d+(\mathbb L^r-1)\mathbb L^{d-1},
\]
which is (MT.3).  Counting the same fibres over $\mathbb F_q$ proves
(MT.2) directly, without a purity theorem.

Consequently the exact error term in the unweighted orthogonality identity
of the supplied Fable sketch is
\[
             q\,\#V_k(\mathbb F_q)=q^k+(q-1)q^d.
\]
For a fixed nonzero additive-character parameter its quadratic sum is
$q^d$, by summing first over $x$.  This gives square-root size when $k$ is
even; for odd $k$ the extra factor $q^{1/2}$ reflects the one-dimensional
radical.  It explains the sign and the exact constant for this family.
The general weighted operator estimate in
\texttt{OURS\_FABLE\_PROGRAM\_SUGGESTIONS\_20261120\_EXPLAINED.md},
\S3, remains a separate inherited result.  This calculation does not prove
the sign or norm assertions for arbitrary sheaves in Alp\"oge's steps 3
and 5, nor identify this elementary model with all the constructions in
the omitted Lean dependency tree.

\subsection*{MT-Q02. Its arithmetic zeta function and shifted zero lines}
The exact all-prime formula in MT-Q01 gives
\[
 Z(V_k/\mathbb F_q,T)
     =\frac{1-q^{d-1}T}{(1-q^{k-1}T)(1-q^dT)},
 \qquad
 \zeta_{V_k}(z)
     =\frac{\zeta(z-k+1)\zeta(z-d)}{\zeta(z-d+1)}.       \tag{MT.4}
\]
Indeed apply (MT.2) over every $\mathbb F_{q^m}$ in
$Z=\exp(\sum_{m\ge1}\#V_k(\mathbb F_{q^m})T^m/m)$ and use
$-\log(1-aT)=\sum_{m\ge1}a^mT^m/m$.  Multiplying the resulting factors
over rational primes proves the second identity first in a right
half-plane of absolute convergence, and then meromorphically.
Under $A_D$ the nontrivial divisor of this function is supported on
\[
    \operatorname{Re}z=k-\tfrac12,\qquad
    \operatorname{Re}z=d+\tfrac12,\qquad
    \operatorname{Re}z=d-\tfrac12.
\]
Divisor orders are signed and cancellations are kept exactly as in
(MT.4); the denominator's zeros are possible poles, not additional zeros.
The shifted pole of $\zeta$ at $1$ and shifted trivial zeros are recorded
by the same explicit formula, separately from this nontrivial-divisor
assertion.  This is an all-characteristic, concrete Artin--Tate example
linking the path model of the method with arithmetic zeta functions and
the Grothendieck ring.  It illustrates why a mixed-weight scheme has
several shifted lines rather than a single universal critical line.

\subsection*{MT-L01. Damped Euler products and Adams operations}
Let $\rho$ be a finite-image complex representation of $G_K$ with no
invariant vectors, of dimension $d$.  At each finite place $v$ put
$A_v=\rho(\operatorname{Frob}_v)|_{V^{I_v}}$.  Its eigenvalues have
absolute value one.  For $|c|\le1$ define, initially for $\operatorname{Re}s>1$,
\[
 D_c(s,\rho)=\prod_v\det(1-cA_v(Nv)^{-s})^{-1}.
\]
If $L(s,\rho)$ is holomorphic and nonzero on
$H=\{\operatorname{Re}s>1/2\}$, then $D_c$ has a holomorphic nonvanishing
continuation to $H$, given by
\begin{align*}
 D_c(s,\rho)&=\exp\{c\log L(s,\rho)+E_c(s)\},\\
 E_c(s)&=\sum_v\sum_{k\ge2}
   \frac{c^k-c}{k}\operatorname{tr}(A_v^k)(Nv)^{-ks}.
\end{align*}
The branch of $\log L$ is fixed by the Euler series in
$\operatorname{Re}s>1$.

For a full proof, expand $-\log\det(1-X)$ in powers of $X$ in the initial
half-plane.  The terms with $k=1$ cancel in
$\log D_c-c\log L$, giving the displayed $E_c$.  On a compact subset with
$\operatorname{Re}s\ge1/2+\varepsilon$, its terms in absolute value sum to
at most
\[
 2d\sum_v\sum_{k\ge2}(Nv)^{-k(1/2+\varepsilon)}
 \le \frac{2d}{1-2^{-1/2-\varepsilon}}
       \sum_v(Nv)^{-1-2\varepsilon}<\infty.
\]
Normal convergence, also uniform in $|c|\le1$, proves that $E_c$ is
holomorphic in $s$.  A nonvanishing holomorphic function on the simply
connected half-plane $H$ has a holomorphic logarithm, uniquely normalized
by its existing branch on the right.  The exponential above is therefore
holomorphic, nonzero, and agrees with the product where it converges.
This proves the continuation.

Under $M_H$, Brauer induction supplies the needed holomorphic nonvanishing
on $H$ for $\rho$ with no invariants: all Hecke factors have no zeros
there, and their orders at $1$ combine to
$-\dim V^{G_K}=0$.  This does not require Artin holomorphy on the boundary
line.  Under $A_D$ the same argument applies to sums of nonprincipal
Dirichlet characters over $\mathbb Q$.

The traces $\operatorname{tr}(A_v^k)$ are the local Adams-power characters;
the proof separates the first prime term from all higher Adams terms,
whose series converge in the larger half-plane.  This is a concrete
$\lambda$-ring link, rather than a claim that positivity of coefficients
follows from the formal operations.  For $c=1/m$ the phenomenon is closely
related to Gun--Hazra--Sahu, \emph{On holomorphy and non-vanishing of Artin
L-functions} (2025), Theorem 5.  The displayed formula covers all complex
$|c|\le1$ and makes its hypotheses explicit.

\subsection*{MT-L02. One irrational parameter detects the zero-free half-plane}
Fix an irrational real number $0<c<1$ and a nontrivial irreducible Artin
representation $\rho$.  Then the following are equivalent:
\begin{enumerate}
 \item $L(s,\rho)$ is holomorphic and nonzero on $H$;
 \item the Euler product $D_c(s,\rho)$ admits a single-valued meromorphic
       continuation to $H$.
\end{enumerate}
In the second assertion meromorphic is a global, single-valued condition;
a branch on a slit domain is not enough.

MT-L01 proves the forward implication.  Conversely $L$ is meromorphic by
Brauer induction and $E_c$ is holomorphic on $H$ by the same normally
convergent series without any hypothesis about the zeros of $L$.
Logarithmic differentiation in $\operatorname{Re}s>1$ gives
\[
 \frac{D_c'}{D_c}=c\frac{L'}L+E_c'.
\]
Both sides are meromorphic on connected $H$ if the asserted continuation
of $D_c$ exists, so the identity holds on $H$.  Taking residues at any
$s_0\in H$ gives
\[
 \operatorname{ord}_{s_0}D_c
      =c\operatorname{ord}_{s_0}L.
\]
Both orders are integers, and $c$ is irrational.  Both must be zero.
Thus $L$ is holomorphic and nonvanishing on $H$, proving the equivalence.
If instead $c=a/b$ in lowest terms, the same equality proves
$b\mid\operatorname{ord}_{s_0}L$ everywhere in $H$.

GRH in the meromorphic-divisor sense for the nontrivial irreducible Artin
factors is therefore equivalent, by the functional equation, to these
irrational-parameter continuation tests.  To include the trivial factor one
must separately add RH for $\zeta_K$; that input is already part of $M_H$.
This is an analytic multiplicity test.  It neither constructs the
continuation from arithmetic lattices nor removes poles on the critical
line.  With a trivial summand of multiplicity $r$, the known pole at $1$
contributes local monodromy $e^{-2\pi i c r}$; this is why the no-invariants
hypothesis was retained rather than suppressed.

\subsection*{MT-ART05. The exact induced-linear cone for \texorpdfstring{$A_5$}{A5}}
The coarse counterexample MT-ART04 can be sharpened to a degree-three
virtual character. More usefully, the same calculation gives explicit
lower bounds for any pole of an actual $A_5$ Artin function. These are
group-theoretic deductions, not realizations of the virtual examples as
arithmetic divisor orders.

Write the irreducible characters in the order
$1,\chi_3,\chi'_3,\chi_4,\chi_5$, and put
\[
 \Theta=a+b\chi_3+c\chi'_3+d\chi_4+e\chi_5,
 \qquad a,b,c,d,e\in\mathbb Z.
\]
Then the inequalities
$\langle\Theta,\operatorname{Ind}_H^{A_5}\eta\rangle\ge0$ for
every subgroup $H$ and every linear character $\eta$ are equivalent to
the following seven inequalities:
\[
 a\ge0,\quad e\ge0,\quad a+d\ge0,\quad b+c\ge0,\quad
 b+d+e\ge0,\quad c+d+e\ge0,\quad b+c+d\ge0.       \tag{MT-A5-cone}
\]
If $r=\Theta(1)=a+3b+3c+4d+5e$, then $d<0$ forces $r\ge3$,
and $b<0$ or $c<0$ forces $r\ge4$. These thresholds are sharp for
this abstract cone: respectively
\[
       1+\chi_3+\chi'_3-\chi_4,
       \qquad -\chi_3+\chi'_3+\chi_4               \tag{MT-A5-ex}
\]
have degrees three and four and satisfy every inequality.

\begin{proof}
We give the character and subgroup details so the finite assertion has
no unreported character-table oracle. The five conjugacy classes have
orders $1,2,3,5,5$ and sizes $1,15,20,12,12$. The natural four-dimensional
augmentation character has values $(4,0,1,-1,-1)$ and norm one, as in
MT-ART04. Its real representation is oriented because all permutations
are even. The Hodge-star decomposition of its exterior square has two
invariant real three-dimensional spaces. The exterior-square character
has values $(6,-2,0,1,1)$ by
\[
 \chi_{\wedge^2 V}(g)=\tfrac12(\chi_V(g)^2-\chi_V(g^2)).
\]
Its norm squared is two and it has no invariant vectors, so the two
three-dimensional summands are distinct irreducibles. Simplicity of
$A_5$ makes each faithful and its determinant trivial. An order-two or
order-three element in either corresponding $\mathrm{SO}_3$ representation
has trace $-1$ or $0$. At order five the two traces are
$\varphi=(1+\sqrt5)/2$ and $\bar\varphi=(1-\sqrt5)/2$; their sum is
one and each is the trace of a nontrivial order-five rotation.

There are six Sylow-five subgroups. Their normalizers have order ten
and are dihedral, since the centralizer of a five-cycle is its own
cyclic subgroup. The permutation action on the six subgroups has
values $(6,2,0,1,1)$: an order-five element fixes its unique Sylow
subgroup, an order-three element fixes none, and double counting gives
$6\cdot5/15=2$ fixed subgroups for an involution. Removing the constants
gives an irreducible character $\chi_5=(5,1,-1,0,0)$, of norm one.
The squared dimensions $1+9+9+16+25=60$ exhaust the character table.

Every proper subgroup is contained in an $A_4$, a $D_5$ of order ten,
or an $S_3$ of order six. Here is an elementary verification. Simplicity
excludes proper indices two, three and four. The remaining possible
proper subgroup orders divide sixty and belong to
$\{1,2,3,4,5,6,10,12\}$. There are no elements of order four, six or
ten. Thus subgroups of orders four, six and ten have types
$V_4,S_3,D_5$. The centralizer of an involution is its unique $V_4$,
whose normalizer is $A_4$. In a subgroup of order twelve, a normal
Sylow-three subgroup is impossible since the ambient normalizer of
$C_3$ has order six. There are therefore four Sylow-three subgroups,
contributing eight nonidentity elements. Three Sylow-two subgroups
would contribute nine further involutions, because distinct $V_4$'s
have disjoint nonidentity sets. Hence the Sylow-two subgroup is normal
and the order-twelve subgroup is its $A_4$ normalizer. The smaller
cyclic and Klein subgroups are contained in these listed groups.

Let $u,\bar u$ be the nontrivial linear characters of $A_4$,
$v$ its irreducible character of degree three, and let $\epsilon$
denote a sign character of either $D_5$ or $S_3$. The restriction
decompositions obtained from the displayed character values are
\[
\begin{array}{c|ccc}
 & A_4 & D_5 & S_3\\\hline
1&1&1&1\\
\chi_3&v&\epsilon+v_1&\epsilon+v_2\\
\chi'_3&v&\epsilon+v'_1&\epsilon+v_2\\
\chi_4&1+v&v_1+v'_1&1+\epsilon+v_2\\
\chi_5&u+\bar u+v&1+v_1+v'_1&1+2v_2
\end{array}                                                   \tag{MT-A5-res}
\]
Here $v_1,v'_1$ are the two dihedral characters of degree two and
$v_2$ is the standard $S_3$ character. The nonlinear characters in
these three groups are themselves induced from linear characters:
$v$ from a nontrivial character of $V_4$, $v_1,v'_1$ from $C_5$,
and $v_2$ from $C_3$. Transitivity of induction and Frobenius reciprocity
therefore show that all induced-linear tests imply nonnegative
multiplicities in every column of (MT-A5-res). Together with the
global trivial-character test $a\ge0$, these are exactly the seven
inequalities (MT-A5-cone), after deleting redundant inequalities.
Conversely those seven inequalities make every column an actual
character. Restricting further gives actual characters on all proper
subgroups, hence every required linear inner product is nonnegative.
The global linear test is again $a\ge0$. This proves the equivalence.

For the degree bounds, first put $d=-k<0$. Then $a\ge k$,
$b+c\ge k$, and $b,c\ge k-e$. If $e=0$, these imply
$r\ge k+6k-4k=3k\ge3$. If $e\ge1$, the first two inequalities give
$r\ge k+3k-4k+5e\ge5$. Next put $b=-k<0$. If $d\ge0$, we have
$c\ge k$, $d+e\ge k$, whence
$r\ge4d+5e\ge4k$. If $d=-j<0$, then
$a\ge j$, $c\ge k+j$, and $e\ge k+j$, giving
$r\ge j+3j-4j+5(k+j)\ge4k$. The argument for $c<0$ is identical.
Finally direct substitution verifies both examples (MT-A5-ex).
\end{proof}

\subsection*{MT-ART06. Pole multiplicity gates for actual \texorpdfstring{$A_5$}{A5} extensions}
Let $L/K$ be an $A_5$ extension and $s_0\ne1$. Its Heilbronn character
has coefficients equal to the local Artin orders. Hecke holomorphy and
induction give all the tests in MT-ART05, unconditionally. Hence every
Artin factor is holomorphic at a zero of $\zeta_L$ of order at most two.
A pole of the four-dimensional factor requires order at least three
and a zero of $\zeta_K$ at the same point. A pole of either
three-dimensional factor requires order at least four. Conditional on
$M_H$, all remaining nontrivial candidate poles lie on the critical line.

At the central point the lower bounds strengthen to six and eight,
respectively. To prove this final refinement, the permutation actions
of degrees five and six give
\[
 L(s,\chi_4)=\frac{\zeta_{L^{A_4}}(s)}{\zeta_K(s)},\qquad
 L(s,\chi_5)=\frac{\zeta_{L^{D_5}}(s)}{\zeta_K(s)}.
\]
The completed Dedekind zeta functional equation has sign $+1$, and
its completion factors are finite and nonzero at $1/2$. Consequently
the central orders $a,d,e$ are even integers, with $a,e\ge0$ by
(MT-A5-cone). If $d=-k<0$, then $k\ge2$; the preceding proof gives
$r\ge3k\ge6$ when $e=0$ and $r\ge5e\ge10$ otherwise.
If $b=-k<0$ and $d\ge0$, the nonnegative even integers $d,e$
satisfy $d+e\ge k\ge1$, so $r\ge4d+5e\ge8$.
If $d=-j<0$, then $j\ge2$ and the preceding bound
$r\ge5(k+j)$ is stronger. The same holds for $c<0$.
This proves all assertions; no assumption on simplicity of zeros has
been introduced. The degree-three and degree-four virtual examples
only show the limits of the indicated character tests away from the
additional central parity constraints.

These refinements put the residual Artin question into explicit local
inequalities on a finite Galois quotient, within the SGA 1 V \S7
framework. They sharpen ORDER item 42 and ZETA\_TWO ZT02 without
claiming general Artin holomorphy or constructing actual negative
Artin divisor orders.

\subsection*{MT-H01. A unitary Hecke family outside the algebraic-CM normalization}
Let $K=\mathbb Q(\sqrt5)$, $\epsilon=(1+\sqrt5)/2$, and let $\sigma$
be its nontrivial automorphism. Put $t_0=\pi/\log\epsilon$.
For each nonzero integer $k$ there is a conductor-one unitary Hecke
character defined on fractional ideals by
\[
       \eta_k((\alpha))=
       \left|\frac{\alpha}{\sigma(\alpha)}\right|^{ik t_0}.
                                                               \tag{MT-H1}
\]
It has infinite order, $\eta_k^\sigma=\eta_{-k}$, and is trivial
on the rational ideles. No positive power of $\eta_k$ becomes an
algebraic Hecke character after multiplying by a complex norm power
and a finite-order character. Thus the full unitary-Hecke input $M_H$
has an explicit scope beyond the algebraic Hecke characters used in
the potentially CM motive transfer.

\begin{proof}
The ring of integers is $\mathbb Z[\epsilon]$, with discriminant five.
Minkowski's ideal-class bound is $\sqrt5/2<2$, so every fractional
ideal is principal. An exact classical locator for the bound is
Milne, \emph{Algebraic Number Theory}, Theorem 4.3 (v3.08), PDF p.~72.
The units are $\{\pm\epsilon^j:j\in\mathbb Z\}$. For a direct
verification, a positive unit $u$ with $1<u<\epsilon$ would have
integer trace $u+u^{-1}$ in $(2,\sqrt5)$ if its norm were one, or
$u-u^{-1}$ in $(0,1)$ if its norm were minus one; neither interval
contains an integer. Multiplying an arbitrary positive unit by a
power of $\epsilon$ reduces it to $[1,\epsilon)$, so it is a power
of $\epsilon$, and signs give the assertion.

Multiplying $\alpha$ by a unit changes the logarithm of the absolute
ratio in (MT-H1) by $2j\log\epsilon$. Its phase therefore changes by
$2\pi i k j$, so (MT-H1) is well-defined and multiplicative. It is
unitary. The classical ideal/infinity-type description of Hecke
characters makes it an unramified finite-place character, with
archimedean exponents $\pm ikt_0$ (their simultaneous sign depends
on the ideal-to-idele convention). The unit compatibility just checked
is exactly what is required for this description. Rational ideles
have equal archimedean components and equal valuations at conjugate
places; their images are trivial. For a ramified rational place this
also follows here from the generator $\sqrt5$, whose absolute ratio
is one. Conjugating (MT-H1) inverts the ratio.

On the positive archimedean subgroup the exponents of any algebraic
Hecke character are integers, up to an overall convention sign.
A norm power adds the same complex number to the two exponents;
a finite-order character is trivial on this connected subgroup.
If a positive power $m$ of (MT-H1) had the asserted algebraic
normalization, subtraction of its two exponents would give
$2mikt_0$ equal to the difference of two integers. This is impossible.
The same nonconstant archimedean restriction proves infinite order.
\end{proof}

This does not assert a transcendence theorem for individual Hecke
values, or rule out every conceivable generalized motive. It precisely
separates the character class covered by CM01 from the broader unitary
character class quantified in Masullo's endpoint. A proposed transport
of the CM theta-value and algebraic-height steps must therefore address
this extra archimedean data; norm twisting alone cannot do so.

\subsection*{MT-H02. An explicit nonholomorphic \texorpdfstring{$\mathrm{GL}_2$}{GL2} family}
For $k\ge1$, the classical Hecke--Maa\ss\ correspondence associates to
$\eta_k$ a cuspidal dihedral Maa\ss\ newform $g_k$ of weight zero,
level five and nebentypus $\chi_5$. Its spectral parameter is $kt_0$,
its Laplace eigenvalue is $1/4+k^2t_0^2$, and
\[
       L(s,g_k)=L_K(s,\eta_k).
                                                               \tag{MT-H2}
\]
Consequently $M_H$ gives GRH for every member of this family. Its
entireness is already the classical Hecke theorem. The exact primary
locator for the correspondence, level, spectral parameter and coefficient
identity is Humphries--Khan, \emph{On the Random Wave Conjecture for
Dihedral Maa\ss\ Forms}, GAFA 30 (2020), Appendix A.1, printed pp.~106--107 (PDF pp.~73--74), the displayed
formula for $\psi((\alpha))$ and the subsequent formulas for
$t_g$ and $\lambda_{g_\psi}(n)$.

For clarity, the transfer proof is short and complete. MT-H01 gives
a non-norm character: its two archimedean exponents are unequal.
The cited correspondence is therefore cuspidal, and its formula
$\lambda_{g_k}(n)=\sum_{N\mathfrak a=n}\eta_k(\mathfrak a)$ identifies
the absolutely convergent Dirichlet series with (MT-H2), including all
finite primes. Analytic continuation gives the equality everywhere.
The character is unitary, so its nontrivial critical line is
$\operatorname{Re}s=1/2$, which is exactly the assertion supplied by
$M_H$. This proves the stated implication without a generic
non-dihedral Maa\ss\-form hypothesis.

These forms live on the analytic modular curve of level five and in
the automorphic-induction part of the adelic framework. They are an
explicit additional $\mathrm{GL}_2$ branch beside the holomorphic
CM and finite-image weight-one branches, not the missing general
non-CM $\mathrm{GL}_2$ theorem in ORDER item 42.

\subsection*{MT-H03. Exact tensor and symmetric-power products in that family}
Write $F_j(s)=L_K(s,\eta_j)$ for $j\ne0$ and
$F_0(s)=\zeta_K(s)=\zeta(s)L(s,\chi_5)$. Then $F_{-j}=F_j$.
For the actual local two-dimensional parameters $V_j$ of $g_j$, the
following finite Euler-product identities hold:
\[
 L(s,V_m\otimes V_n)=F_{m+n}(s)F_{m-n}(s),\qquad m,n\ge1,
                                                               \tag{MT-H3}
\]
and, for $n\ge0$ and $k\ge1$,
\[
 L(s,\operatorname{Sym}^n V_k)=
 \prod_{j=0}^{\lfloor(n-1)/2\rfloor}F_{(n-2j)k}(s)
 \begin{cases}\zeta(s),&n\text{ even},\\1,&n\text{ odd}.
 \end{cases}                                                   \tag{MT-H4}
\]
The product is empty when $n=0$. Conditional on $M_H$, all nontrivial
zeros of these products lie on $\operatorname{Re}s=1/2$. They are
entire except for a simple pole at one when $m=n$ in (MT-H3), or
when $n$ is even in (MT-H4). This is a covered arbitrarily large
degree family, not a use of an unknown general symmetric-power lift.

\begin{proof}
Conjugating ideals proves $F_{-j}=F_j$. At a split rational prime
the two local parameters of $V_k$ are $z^k,z^{-k}$, where
$z=\eta_1(\mathfrak p)$. The tensor weights are
$z^{m+n},z^{m-n},z^{-m+n},z^{-m-n}$, which gives (MT-H3).
The symmetric-power weights are $z^{(n-2j)k}$ for $0\le j\le n$;
pairing opposite weights leaves one middle weight $1$ precisely for
even $n$, proving (MT-H4) there.

At an inert prime the unique prime ideal is $(p)$ and has character
value one. The two induced Frobenius eigenvalues are $1,-1$ for
every nonzero $k$. Tensoring gives two of each sign. On the
symmetric power, the even-index monomials give eigenvalue $1$ and
the odd-index monomials give $-1$. If $n=2r$, these multiplicities
are $r+1,r$; if $n=2r+1$, they are $r+1,r+1$. These are exactly
the factors on the right sides. If $m=n$, the factor $F_0$ at an
inert prime is $(1-p^{-2s})^{-1}$, so the same check includes that
case.

The only ramified prime is five. The ideal $(\sqrt5)$ has
$\eta_k$-value one and the character is unramified at every finite
place. Thus its local character at that quadratic completion is
trivial, and its induced Weil parameter is $1\oplus\chi_{5,5}$.
The inertia-fixed dimension is one. A tensor product has two
trivial constituents and two $\chi_{5,5}$ constituents; a symmetric
power has respectively $r+1,r$ or $r+1,r+1$ of them. The latter
constituents have no inertia invariants. Hence the ramified Euler
factors also agree exactly, rather than only away from five.
Equivalently the proof is the weight-pair decomposition for diagonal
matrices $\operatorname{diag}(z,z^{-1})$ and an outside coset
acting by interchanging the two basis vectors.

Absolute convergence initially gives the identities for
$\operatorname{Re}s>1$. The right sides give their meromorphic
continuations. All $F_j$, $j\ne0$, are entire Hecke functions with
nonzero boundary value at one; the latter is Ramakrishnan--Valenza,
Theorem 7-28, printed p.~289. The same holds for $L(1,\chi_5)$.
Thus the displayed zeta factors give exactly the asserted simple
poles, with no cancellation at one. Finally $M_H$ locates all the
nontrivial zeros of each unitary Hecke factor. The identities contain
no negative exponents, so they introduce no hidden poles on the
critical line. This proves every assertion.
\end{proof}

The middle summand in (MT-H4) is the trivial character for this
specific family. In a differently normalized dihedral family it
depends on the central character; replacing it universally by either
$1$ or $\chi_K$ would be incorrect. This is the representation-ring
and lambda-operation part of the adelic framework, beside SGA 6 V
\S\S2,7, with the explicit analytic parameters of MT-H01 retained.

\subsection*{MT-ART07. The entire non-real $A_5$ Artin category}
Let $L/\mathbb Q$ be Galois with group $A_5$, and suppose $L$ is not
totally real. Every nontrivial irreducible Artin L-function for this
extension is entire, unconditionally. Under $M_H$ each satisfies GRH.
This closes the non-real $A_5$ specialization of the holomorphy gate;
MT-ART05/06 remain useful for totally real extensions and other base fields.

The holomorphy assertion is classical, not a new finding. An exact primary
locator is Calegari, \emph{The Artin conjecture for some $S_5$-extensions},
\S1, publisher-author PDF p.~2, the paragraph immediately before
Theorem 1.1: it states the stronger automorphy assertion for $A_5$ when
complex conjugation is nontrivial, citing Khare--Wintenberger and Kim's
icosahedral functorial constructions. The two conditions are equivalent
here because a Galois extension has a real embedding exactly when every
embedding is real, equivalently complex conjugation acts trivially.
We cite that classical result rather than repeat its proof.

Here is the complete GRH transfer and its tensor-category form.
For every finite-dimensional representation $V$ of $A_5$, put
$m=\dim V^{A_5}$. Complete reducibility gives
\[
 L(s,V)=\zeta(s)^m
          \prod_{\chi\ne1}L(s,\chi)^{\langle V,\chi\rangle},
 \qquad \langle V,\chi\rangle\in\mathbb Z_{\ge0}.
\]
The product after removing $\zeta^m$ is entire by the cited theorem.
For each irreducible factor, Brauer induction expresses its L-function
as a finite product of integer powers of finite-order Hecke L-functions.
At a point of the open critical strip off its middle line, $M_H$ makes
every such factor finite and nonzero. The product is consequently nonzero
there, including when some Brauer exponents are negative. This proves
the claimed GRH without making a positivity assumption on those exponents.

At $s=1$ every nontrivial irreducible factor is also nonzero. Indeed in a
Brauer expression $\chi=\sum_j a_j\operatorname{Ind}_{H_j}^{A_5}\eta_j$,
the order there is $-\sum_{\eta_j=1}a_j$, by the simple Dedekind pole and
Hecke boundary nonvanishing. Frobenius reciprocity identifies the sum with
$\langle\chi,1\rangle=0$. Thus $L(s,V)$ has a pole of order exactly $m$
at $1$, and has no other poles. Tensor products, duals, exterior and
symmetric powers, and every Schur functor produce representations of the
same finite group, so the same conclusions hold for them after removing
their explicitly determined invariant summands. This is a statement about
the whole Artin tensor category attached to this cover, not about a
general positive-weight motive. The companions in Grothendieck's language
are SGA~1, Expos\'e~V, \S7 (finite monodromy), and SGA~6, Expos\'e~V,
\S\S2,7 (representation-ring operations), together with AG-LAMBDA.

\subsection*{MT-ART08. A finite exceptional set is enough for a holomorphy transfer}
Here is a useful precise comparison lemma. Let $V$ be a self-dual,
finite-image complex Galois representation over a number field $K$.
Suppose an Euler product $F(s)$ converges in a right half-plane and
continues to an entire function, with the same finite local factors as $L(s,V)$ outside a
finite set $S$. Suppose also that, for each $v\in S$, the reciprocal of
the local factor of $F$ is a polynomial $Q_v((Nv)^{-s})$.
Then $L(s,V)$ is entire. No equality of ramified local parameters is
needed for this conclusion.

To prove this, write the Artin factor at $v$ as
$P_v((Nv)^{-s})^{-1}$. Every eigenvalue of Frobenius on the inertia
invariants is a root of unity, since $V$ has finite image. Hence the zeros
of $P_v((Nv)^{-s})$ all lie on $\operatorname{Re}s=0$.
Equality in a right half-plane, followed by meromorphic continuation,
gives
\[
 L(s,V)=F(s)\prod_{v\in S}
             \frac{Q_v((Nv)^{-s})}{P_v((Nv)^{-s})}.
\]
The right side is holomorphic on $\operatorname{Re}s>0$.
The classical completed Artin functional equation is
\[
 A^{s/2}\Gamma_V(s)L(s,V)
   =W A^{(1-s)/2}\Gamma_V(1-s)L(1-s,V),
\]
with $A>0$ and $W\ne0$; the gamma factor is a product of
$\Gamma_{\mathbb R}(s)$, $\Gamma_{\mathbb R}(s+1)$ and
$\Gamma_{\mathbb C}(s)$ with nonnegative integral multiplicities.
An exact locator for these formulas is Prasad--Yogananda,
\emph{A Report on Artin's holomorphy conjecture}, \S3.3,
PDF pp.~8--9, with Brauer induction in Theorem 3.3.
For $\operatorname{Re}s<1$, the numerator
$\Gamma_V(1-s)L(1-s,V)$ is holomorphic, since
$\operatorname{Re}(1-s)>0$. The reciprocal $1/\Gamma_V(s)$ is entire.
The functional equation therefore makes $L(s,V)$ holomorphic on this
second half-plane too. The two half-planes cover $\mathbb C$, proving
the assertion. For a non-self-dual $V$ the identical argument works if
the hypotheses are also supplied for $V^\vee$.

The finiteness of $S$ is an actual hypothesis. Agreement at a set of
places of density one does not give the finite product in the proof.
In particular this lemma cannot by itself upgrade the weak
five-dimensional correspondence in Calegari's Theorem 1.1.
In geometric terms this separates modification at finitely many closed
points of an arithmetic base from a density statement about infinitely
many Frobenius fibres.

\subsection*{MT-S501 and MT-S502. The $S_5$ transfer and its one-value gate}
Let $L/\mathbb Q$ be an $S_5$ extension, with complex conjugation of cycle
type $(12)(34)$. Assume $5$ is unramified and its Frobenius also has this
cycle type. Then the irreducible factors of dimensions $4$ and $6$ are
entire and, under $M_H$, satisfy GRH (MT-S501).

Let $E=L^{A_5}$, let $F\subset L$ have degree $6$ over $\mathbb Q$, and
put $H=EF$. Under $M_H$, the additional single condition
\[
                         \zeta_H(1/2)\ne0
\]
implies that every nontrivial irreducible factor is entire and satisfies
GRH (MT-S502). This is a proved reduction with the displayed extra input,
not a proof of that central nonvanishing.

The classical inputs are Calegari's Theorems 1.1 and 1.2,
publisher-author PDF pp.~2--3; the definition $H=EF$ is explicit in
Theorem 1.2 of arXiv:1112.1152v1, p.~2. The former supplies the
dimension-$4$ and dimension-$6$ holomorphy, and the latter supplies the
remaining case if $\zeta_H$ has no zero in the real interval $(0,1)$.
Under $M_H$ the only possible real zero in that interval is $1/2$.
Thus its interval condition is exactly equivalent to the displayed
single-value condition. Once holomorphy is available, the Brauer
zero-location argument of MT-ART07 proves GRH for every factor.
This is the full proof of both transfers. The one-dimensional factors
are Dirichlet functions and already fall under $A_D$.

The totally real quadratic field $E$ is the base for the odd projective
two-dimensional representations used in the cited theorem. Accordingly
this is a higher-rank consequence of a specific automorphic descent
theorem. It is not obtained by treating every $S_5$ factor as a single
Hecke character. In the Artin-motive programme, the remaining gate in
this family is one central Dedekind value. AG-SIGN controls a nonzero
leading coefficient after a central zero is removed; it does not prove
that the central order is zero.

\subsection*{MT-S503. One unconditional holomorphy example becomes a full GRH example}
Take the splitting field $L$ of
\[
                        x^5-x^3-x^2+x+1.
\]
Calegari, Theorem 5.1, publisher-author PDF p.~15, proves automorphy of
every representation of $\operatorname{Gal}(L/\mathbb Q)=S_5$;
the quintic field has discriminant $1609$. Consequently under $M_H$
every nontrivial irreducible Artin L-function of this explicit field
satisfies GRH, by the same signed Brauer argument. This example uses
that cited theorem, not the local-at-$5$ hypotheses of MT-S501, which
the source explicitly says this polynomial does not satisfy.
The historical bounded-height computation in Theorem 5.3 is part of
the classical proof's evidence; it is not a computation performed in
this run. Geometrically this supplies a concrete finite \(S_5\)-cover
of the arithmetic base whose complete Artin category is covered.

\subsection*{MT-M05. The stated Schauder quotient loses more than its zero evaluations see}
Lemma 9.3.2.5 on original PDF pp.~425--426 asserts a basis for the
quotient by the actual convolution ideal $\varphi*\mathcal S_E(G)$,
with coefficient functionals the Mellin jets at the zeros of
$\mathcal M\varphi$. We inspected the original p.~426 image: there
is no closure over this ideal. The assertion is false even when the
Mellin divisor consists of one simple zero. Thus the counterexample
does not rely on the zero-count defect in MT-M01.

Use $G=\mathbb R_{>0}$, trivial character and $u=\log t$; convolution
becomes ordinary additive convolution. Write
\[
 g_a(u)=(4\pi a)^{-1/2}e^{-u^2/(4a)},\qquad a>0.
\]
Choose $0<b<a$ and set $\varphi(e^u)=-g_a'(u)$,
$h(e^u)=-g_b'(u)$. Both functions belong to the exponential-Schwartz
space, and
\[
 \mathcal M\varphi(s)=s e^{as^2},\qquad
 \mathcal Mh(s)=s e^{bs^2}.
\]
The first function has just the simple zero $0$, lying in the interior
of the strip for $I=[-1,1]$. The only claimed coefficient functional is
$\delta_0$, and $\delta_0(h)=0$.
But $h\notin\varphi*\mathcal S_E$: if $h=\varphi*f$, multiplicativity
of the absolutely convergent Mellin transform gives
$\mathcal Mf(s)=e^{(b-a)s^2}$ for $s\ne0$, and then everywhere by
analyticity. On the imaginary axis this has modulus
$e^{(a-b)\tau^2}$. In contrast, $f(e^u)$ is integrable and therefore
$|\mathcal Mf(i\tau)|\le\int|f(e^u)|\,du$, a contradiction.
Thus the nonzero quotient class of $h$ has every asserted coefficient
zero, which is incompatible with uniqueness in a Schauder expansion.
The still simpler choice $\varphi=g_a$, $h=g_b$ has no Mellin zeros
and a nonzero quotient, so even the empty-divisor case fails.

\subsection*{MT-M06. The closed-ideal repair for every polynomial Gaussian divisor}
In logarithmic coordinates put
\[
 \mathcal E=\{f\in C^\infty(\mathbb R,\mathbb C):
   p_{N,j}(f):=\sup_x e^{N|x|}|f^{(j)}(x)|<\infty
   \text{ for every integer }N,j\ge0\}.
\]
This is $\mathcal S_E(\mathbb R_{>0})$ with its exponential-Schwartz
Fr\'echet topology. Polynomial factors in the usual Schwartz seminorms
are absorbed by increasing $N$. Let
\[
 P(s)=\prod_{\nu=1}^r(s-\rho_\nu)^{m_\nu},\quad
 d=\sum_\nu m_\nu,\qquad
 \varphi=P(-D)g_a,\quad D=d/dx.
\]
Distinct $\rho_\nu\in\mathbb C$, positive integral $m_\nu$, and
$a>0$ are arbitrary; $P=1,d=0$ is allowed. Then
\[
 \overline{\varphi*\mathcal E}
   =J_P:=\{h\in\mathcal E:
          (\mathcal Mh)^{(j)}(\rho_\nu)=0,
          \ 0\le j<m_\nu\},                         \tag{MT-M6}
\]
and the quotient by this \emph{closed} ideal is topologically
$\mathbb C^d$, with exactly the displayed Mellin jets as coordinates.
In particular its coordinate basis is absolute. This repairs the
finite-divisor case after making the closure operation explicit.

\begin{proof}
First, $g_a*\mathcal E$ is dense in $\mathcal E$. We give the full
argument, since density is distinct from the false surjectivity in
MT-M05. Convolution is continuous on $\mathcal E$: the inequality
$e^{N|x|}\le e^{N|x-y|}e^{N|y|}$ bounds each differentiated
convolution by an exponential supremum seminorm times an exponential
$L^1$ seminorm; the latter is bounded by a constant times
$p_{N+1,0}$.

Suppose a continuous linear functional $T$ vanishes on
$g_a*\mathcal E$. For every real $y$ and $\epsilon>0$ the function
$g_{a+\epsilon}(x-y)=g_a*g_\epsilon(\,\cdot-y)$ lies in this range.
It converges to $g_a(x-y)$ in every $p_{N,j}$ as
$\epsilon\downarrow0$, by uniform Gaussian estimates with the
parameter in a compact subinterval of $(0,\infty)$. Hence
$T(g_a(\,\cdot-y))=0$. The function
\[
 B(z)=T\bigl(e^{-x^2/(4a)}e^{zx/(2a)}\bigr)
\]
is entire: on every compact set of $z$, all its proposed complex
derivatives converge in every $p_{N,j}$ by Gaussian domination.
Completing the square gives $B(y)=0$ for real $y$, so the identity
theorem gives $B=0$.

For any Schwartz function $\psi$, Fourier inversion now gives
\[
 T(e^{-x^2/(4a)}\psi(x))
  =\frac1{2\pi}\int_{\mathbb R}\widehat\psi(\xi)
        B(2ai\xi)\,d\xi=0.
\]
This passage under $T$ is justified in the Fr\'echet space itself:
the $p_{N,j}$ seminorm of
$e^{-x^2/(4a)}e^{i\xi x}$ is at most
$C_{N,j,a}(1+|\xi|)^j$, whereas $\widehat\psi$ is rapidly
decreasing. For $h\in C_c^\infty$, take
$\psi=e^{x^2/(4a)}h$, which is again compactly supported and smooth.
Then $T(h)=0$. Finally $C_c^\infty$ is dense in $\mathcal E$:
multiply by a smooth cutoff equal to one on $[-R,R]$ and supported
in $[-2R,2R]$; the product rule and the bound from $p_{N+1,j}$ make
every $p_{N,j}$ of the discarded tail tend to zero. Thus $T=0$.
Hahn--Banach separation proves the asserted density.

Second, the range of $P(-D)$ on $\mathcal E$ is exactly $J_P$.
Integration by parts gives $\mathcal M(-Df)(s)=s\mathcal Mf(s)$,
so inclusion in $J_P$ follows. For the converse, if
$\mathcal Mh(\rho)=0$, define
\[
 f(x)=e^{-\rho x}\int_{-\infty}^x e^{\rho t}h(t)\,dt
     =-e^{-\rho x}\int_x^\infty e^{\rho t}h(t)\,dt.
\]
The equality is precisely the vanishing of the Mellin value.
The first expression controls the negative tail and the second the
positive tail. For any prescribed $N$, using an exponential bound
on $h$ with exponent larger than $N+|\operatorname{Re}\rho|+1$
proves exponential decay of $f$ on both tails; the equation
$f'=h-\rho f$ then gives the same bounds for every derivative.
Therefore $f\in\mathcal E$, and $(D+\rho)f=h$.
Equivalently $-f$ solves $(-D-\rho)(-f)=h$, and its Mellin transform
is $\mathcal Mh(s)/(s-\rho)$. Repeating this division through all
roots, including multiplicities, proves $P(-D)\mathcal E=J_P$.

The Mellin jet functionals are continuous, so $J_P$ is closed.
Since $g_a*\mathcal E$ is dense and $P(-D)$ is continuous,
$P(-D)(g_a*\mathcal E)$ is dense in $P(-D)\mathcal E=J_P$.
Differentiation commutes with convolution, identifying that range
with $\varphi*\mathcal E$. This proves (MT-M6).

For $d=0$ the quotient is zero and no interpolation is required.
For $d>0$, it remains to verify that no coordinate is missing. Given arbitrary
$d$ target jets, ordinary polynomial Hermite interpolation constructs
a polynomial $Q$ of degree below $d$ for which
$Q(s)e^{bs^2}$ has exactly those jets, with any fixed $b>0$.
Indeed multiplication by the nowhere-zero entire function $e^{bs^2}$
is invertible on each finite jet space, so interpolate the resulting
jets of $Q$. The function $Q(-D)g_b$ belongs to $\mathcal E$ and
realizes them. Choosing the targets to be the coordinate vectors
gives a continuous linear section from the finite-dimensional space
$\mathbb C^d$ to $\mathcal E$. The jet map consequently induces
a continuous isomorphism of the quotient with $\mathbb C^d$,
with continuous inverse. Finite-dimensional coordinate expansions are
absolute, completing the proof.
\end{proof}

For $P=s$, MT-M05 therefore exhibits a proper, nonclosed convolution
ideal whose closure is $\ker\delta_0$. The closed quotient really
does have the one coordinate the draft intended. For $P=1$ the
principal ideal is dense but proper, and the closed quotient is zero.
This is a concrete repair rather than a certification of the general
infinite-divisor statement. For arithmetic Hecke divisors the remaining
gate is controlled infinite interpolation and the claimed absolute
basis/nuclear estimates on the correct separated quotient. A local
$O(\log T)$ zero count alone does not prove those estimates.
In cote~1 trace language, Hausdorff completion and the identification
of the spectral quotient precede taking its nuclear trace.

\subsection*{MT-M07. An infinite closed quotient with an absolute Mellin basis}
The closed-ideal repair also admits a complete infinite-divisor example.
Fix $a,\ell>0$ in the logarithmic space $\mathcal E$ of MT-M06, and put
\[
 \varphi(x)=g_a(x-\ell)-g_a(x),\qquad
 \rho_n=\frac{2\pi i n}{\ell}\quad(n\in\mathbb Z).
\]
Then $\mathcal M\varphi(s)=e^{as^2}(e^{\ell s}-1)$, with precisely
the simple zeros $\rho_n$. Let
\[
 (\mathcal P_\ell f)(t)=\sum_{k\in\mathbb Z}f(t+k\ell).
\]
The quotient $\mathcal E/\overline{\varphi*\mathcal E}$ is
topologically $C^\infty(\mathbb R/\ell\mathbb Z)$ under
$\mathcal P_\ell$, and has an absolute Schauder basis whose coefficient
functionals are exactly $f\mapsto\mathcal Mf(\rho_n)$.

\begin{proof}
The periodization series and every derivative converge uniformly on a
fundamental interval, by exponential decay. This gives a continuous map
to the indicated smooth periodic space. Choose
$\chi\in C_c^\infty(\mathbb R)$ nonnegative and positive on
$[-\ell,\ell]$, and set
$\theta=\chi/\mathcal P_\ell\chi$. The denominator is positive,
smooth and periodic, so $\theta$ is smooth, compactly supported and
$\mathcal P_\ell\theta=1$. The map
$R(F)(x)=\theta(x)F(x)$ is a continuous section of periodization,
proving surjectivity.

Let $\Delta_\ell h(x)=h(x-\ell)-h(x)$. Telescoping shows that its
range lies in $\ker\mathcal P_\ell$. Conversely, if
$\mathcal P_\ell f=0$, define
\[
 h(x)=\sum_{k\ge1}f(x+k\ell)
     =-\sum_{k\le0}f(x+k\ell).
\]
These locally normally convergent series define a smooth function.
For $x\ge0$ use the first expression and for $x\le0$ the second.
For every positive integer $N$ and every $j$, summing the geometric
series of bounds $p_{N,j}(f)e^{-N|x+k\ell|}$ proves
$p_{N,j}(h)<\infty$; the $N=0$ case follows from $N=1$.
Thus $h\in\mathcal E$, and direct index shifting gives
$\Delta_\ell h=f$. The range of $\Delta_\ell$ is therefore the
closed space $\ker\mathcal P_\ell$.

By MT-M06, $g_a*\mathcal E$ is dense in $\mathcal E$.
Continuity of $\Delta_\ell$ and commutation with convolution imply
\[
 \overline{\varphi*\mathcal E}
 =\overline{\Delta_\ell(g_a*\mathcal E)}
 =\ker\mathcal P_\ell.
\]
The continuous section $R$ identifies the quotient topology with
the usual smooth periodic topology, rather than merely giving a
vector-space identification.

The Fourier expansion in the convention used here is
\[
 \mathcal P_\ell f(t)=\frac1\ell
       \sum_{n\in\mathbb Z}\mathcal Mf(\rho_n)e^{-\rho_n t}.
\]
Its coefficients are obtained by integrating over a fundamental
interval and summing translates, justified by absolute convergence.
They decrease faster than every power of $|n|$: integrate by parts
arbitrarily many times and use the integrability of every derivative
of $f$. The Fourier series consequently converges absolutely in
every smooth seminorm and equals $\mathcal P_\ell f$ by Fourier
uniqueness. Representatives for the basis are
\[
                     b_n(x)=\ell^{-1}\theta(x)e^{-\rho_n x}.
\]
Indeed $\mathcal Mb_n(\rho_m)=\delta_{m,n}$, by
$\mathcal P_\ell\theta=1$ and integration on one period. Moreover,
$p_{N,j}(b_n)\le C_{N,j}(1+|n|)^j$ since $\theta$ has compact
support. Hence the expansion
$\sum_n\mathcal Mf(\rho_n)b_n$ is absolutely convergent even in
$\mathcal E$, and represents the class of $f$ in the quotient.
The jets determine that class uniquely by the Fourier identification.
Ordering the indices as $0,1,-1,2,-2,\ldots$ gives an absolute
Schauder basis, as asserted.
\end{proof}

The actual ideal still cannot replace its closure: with $0<b<a$,
$h=\Delta_\ell g_b$ has all these Mellin values zero but is not
$\varphi*f$ for any $f\in\mathcal E$, by exactly the Fourier-growth
division argument of MT-M05. Thus this example tests an infinite
divisor with controlled spacing, not just the empty or finite case.

\subsection*{MT-M08. The corresponding trace and positive convolution formula}
For $f\in\mathcal E$, convolution by $f$ descends to the quotient of
MT-M07. On $L^2(\mathbb R/\ell\mathbb Z,dt)$ it is trace class and
\[
 \operatorname{Tr}C_f
   =\sum_{n\in\mathbb Z}\mathcal Mf(\rho_n)
   =\ell\sum_{k\in\mathbb Z}f(k\ell).                 \tag{MT-M8}
\]
All sums converge absolutely. If
$f=h*\widetilde h$ with $h\in\mathcal E$ and
$\widetilde h(x)=\overline{h(-x)}$, this operator is positive and
\[
                  \operatorname{Tr}C_f
                   =\sum_n|\mathcal Mh(\rho_n)|^2\ge0.
\]

To prove the assertions, first use absolute convergence and changes of
variables to periodize a convolution:
\[
 \mathcal P_\ell(f*g)(t)
  =\int_0^\ell(\mathcal P_\ell f)(t-u)
                     (\mathcal P_\ell g)(u)\,du.
\]
This proves descent and identifies the smooth kernel
$K_f(t,u)=(\mathcal P_\ell f)(t-u)$.
The orthonormal Fourier vectors
$\ell^{-1/2}e^{-\rho_n t}$ are eigenvectors with eigenvalues
$\mathcal Mf(\rho_n)$. These eigenvalues decrease faster than every
power, as proved in MT-M07, so their absolute sum is finite.
The diagonal expansion is thus a trace-class expansion and gives
the first trace sum. Evaluating the absolutely convergent Fourier
series of $\mathcal P_\ell f$ at zero gives the second equality
in (MT-M8), including its factor $\ell$. Finally the Mellin transform
of $\widetilde h$ at a purely imaginary argument is the conjugate
of that of $h$; the eigenvalues are therefore the displayed squares.

There is also a direct nuclear realization on the smooth quotient.
Explicitly, put $c_n(F)=\ell^{-1}\int_0^\ell F(t)e^{\rho_nt}\,dt$
and $e_n(t)=e^{-\rho_nt}$. The $c_n$ are an equicontinuous family,
bounded by the uniform norm. In
\[
 C_f=\sum_n(1+|n|)^{-2}\,
       c_n\otimes\bigl((1+|n|)^2\mathcal Mf(\rho_n)e_n\bigr),
\]
the scalars are summable and the vector family is bounded in every
$C^\infty$ seminorm. This is a nuclear decomposition. Fourier partial
sums converge uniformly on bounded sets of $C^\infty$ (apply a
uniform higher-derivative bound to the coefficient tail), so the
finite-rank trace limit is the same absolute sum above.

This supplies a fully checked spectral-quotient, basis and trace
model beside cote~1's nuclear trace language. It does not identify
an arithmetic prime distribution with this lattice or prove the
corresponding basis theorem for the actual Hecke zero set. In that
case the exact infinite interpolation and operator estimates remain
to be supplied on the correct closed quotient.

\subsection*{MT-M09. Even a bounded local zero count does not give the asserted point basis}
Passing to the closed ideal is necessary, but does not prove the
general basis claim. There exists $\varphi\in\mathcal E$ whose Mellin
transform has only simple zeros on the imaginary axis and only
$O(1)$ zeros in every unit-height interval, for which
\[
        Q=\mathcal E/\overline{\varphi*\mathcal E}
\]
has no Schauder basis with the individual zero evaluations as its
coefficient functionals. This holds for every enumeration of the
zeros, and therefore also rules out an absolute Schauder basis of
that specified kind.

\begin{proof}
Put $\epsilon_n=e^{-n^4}$ for $n\ge1$ and define
\[
 H(s)=\prod_{n=1}^{\infty}
       \left(1+\frac{s^2}{n^2}\right)
       \left(1+\frac{s^2}{(n+\epsilon_n)^2}\right),
 \qquad F(s)=e^{as^2}H(s),\quad a>0.
\]
The product converges uniformly on compact sets because the sum of
the absolute deviations of its factors from one does. Away from the
listed zeros the tail logarithms converge locally uniformly, so
there are no additional zeros. Its zeros are exactly
$\pm in,\ \pm i(n+\epsilon_n)$, all simple. Since
$0<\epsilon_n<1/2$, each unit-height interval contains a bounded
number of them.

For $r=|s|$ the product satisfies
\[
 \log|H(s)|\le2\sum_{n\ge1}\log(1+r^2/n^2)\le2\pi r.
\]
For the last bound compare the decreasing function of $x$ with its
integral on $(0,\infty)$; integration by parts gives
$\int_0^\infty\log(1+r^2/x^2)\,dx=\pi r$.
Consequently, on every strip $|\sigma|\le A$,
\[
 |F(\sigma+it)|
 \le \exp\{aA^2-at^2+2\pi\sqrt{A^2+t^2}\}.
\]
Define
\[
 \varphi(x)=\frac1{2\pi}\int_{\mathbb R}F(it)e^{-itx}\,dt.
\]
The Gaussian bound permits every derivative in $x$ under the integral.
Shifting the contour to $\operatorname{Re}s=\sigma$ gives, for
every real $\sigma$ and every $j\ge0$,
\[
 \varphi^{(j)}(x)
   =\frac{e^{-\sigma x}}{2\pi}
      \int_{\mathbb R}(-\sigma-it)^jF(\sigma+it)e^{-itx}\,dt.
\]
The horizontal sides of the contour tend to zero by the displayed
Gaussian estimate, uniformly for $\sigma$ in the finite interval
used in that shift. Choosing $\sigma=N+1$ for $x\ge0$ and
$\sigma=-N-1$ for $x\le0$ proves every seminorm $p_{N,j}$ is finite.
Thus $\varphi\in\mathcal E$. Cauchy estimates applied on slightly
larger strips show that $F(it)$ and all its derivatives are Schwartz.
Fourier inversion gives $\mathcal M\varphi(it)=F(it)$;
both sides are entire, so $\mathcal M\varphi=F$ everywhere.

Every evaluation at a listed zero annihilates the convolution ideal
and its closure, and therefore descends continuously to $Q$.
Let $p=p_{1,0}$ and let
$\overline p([f])=\inf_{j\in\overline{\varphi*\mathcal E}}p(f+j)$
be the corresponding quotient seminorm. For any $f\in\mathcal E$,
\[
 |(\delta_{in}-\delta_{i(n+\epsilon_n)})(f)|
  \le\epsilon_n\int_{\mathbb R}|x||f(x)|\,dx
  \le2\epsilon_n p(f).
\]
The first inequality follows from
$|e^{i\epsilon_n x}-1|\le\epsilon_n|x|$.
Because the difference functional annihilates the closed ideal,
the same bound holds with $p$ replaced by $\overline p$ on $Q$.

Suppose the claimed Schauder basis existed, and denote the basis
vector associated to $\delta_{in}$ by $b_n$. Biorthogonality of
basis vectors and coefficient functionals gives
$\delta_{in}(b_n)=1$ and
$\delta_{i(n+\epsilon_n)}(b_n)=0$. Hence
\[
                 \overline p(b_n)\ge(2\epsilon_n)^{-1}.
\]
Take the single test function $g_b\in\mathcal E$ for any fixed
$b>0$. Its coefficient at $\delta_{in}$ is $e^{-bn^2}$.
The corresponding terms of its proposed basis expansion satisfy
\[
 \overline p\bigl(e^{-bn^2}b_n\bigr)
        \ge\frac12 e^{n^4-bn^2}\longrightarrow\infty.
\]
Terms of a convergent series in a locally convex space must tend
to zero. In every enumeration these infinitely many distinct terms
occur arbitrarily far along the series, giving a contradiction.
This proves the assertion.
\end{proof}

This is a counterexample for general exponential-Schwartz Mellin
divisors, not a claim that such clustered zeros occur for an actual
Hecke L-function. It pinpoints why NT-Z10/MT-M04 cannot by themselves
establish the arithmetic Schauder step: even an $O(1)$ local count
does not control the norm of individual interpolation vectors.
The next possible repair must incorporate quantitative separation,
or group nearby evaluations into divided-difference blocks and prove
convergence of the resulting block expansion. Neither property is
supplied by a zero-location hypothesis alone in the present argument.

\subsection*{MT-M10. A point-jet basis forces a polynomial cluster-product bound}
The counterexample MT-M09 has a quantitative converse obstruction.
Let $\mathcal E$ be the exponential-Schwartz space of MT-M06,
let $J\subset\mathcal E$ be a closed linear subspace, and put
$Q=\mathcal E/J$. Fix a real $\sigma$ and a discrete set
$Z\subset\sigma+i\mathbb R$, with positive integer multiplicities
$m_\rho$. Suppose every functional
\[
 \delta_{\rho,k}(f)=\frac{(\mathcal Mf)^{(k)}(\rho)}{k!},
 \qquad 0\le k<m_\rho,
\]
annihilates $J$ and occurs as a coefficient functional of a
Schauder basis of $Q$ with continuous coordinate functionals.
Other basis coordinates are permitted.
Then there are constants $c>0$ and $A\ge0$, independent of
$\rho=\sigma+it$ and of the finite subset
$W\subset Z\setminus\{\rho\}$, such that
\[
       \prod_{\nu\in W}|t-\operatorname{Im}\nu|^{m_\nu}
                    \ge c(1+|t|)^{-A}.                \tag{MT-M10}
\]
In particular the distinct ordinates have a uniform polynomial
separation bound. The conclusion already follows from an ordinary
Schauder basis; absolute convergence is not needed.

\begin{proof}
The seminorms $p_{N,j}(f)=\sup_x e^{N|x|}|f^{(j)}(x)|$ make
$\mathcal E$ a Fr\'echet space: a Cauchy sequence converges uniformly,
with all derivatives, on every compact interval, and the weighted
uniform limits give the same smooth limit and all the seminorm
limits. The closed quotient $Q$ is Fr\'echet as well.
For a continuous seminorm $p$ on $\mathcal E$, write
\[
 \overline p([f])=\inf_{u\in J}p(f+u).
\]
The quotient seminorms of finite maxima of the $p_{N,j}$ generate
the quotient topology.

First fix any such $\overline p$. Let $S_n$ be the partial-sum
operators for the proposed basis, and let $R_n=S_n-S_{n-1}$ be
its rank-one coordinate projections, with $S_0=0$.
Convergence of every basis expansion makes
$\sup_n\overline p(R_nq)$ finite for each $q\in Q$.
For completeness, the uniform-boundedness step follows directly
from Baire's theorem. The closed sets
\[
 A_k=\{q:\sup_n\overline p(R_nq)\le k\},\qquad k=1,2,\ldots,
\]
cover $Q$. One contains $q_0+U$, where $U$ is a neighborhood of
zero, and $q_0$ itself may be chosen in that set. Subtraction
gives $\overline p(R_nu)\le2k$ for $u\in U$, uniformly in $n$.
Choose $U$ to contain a unit ball for a finite maximum of
quotient seminorms and rescale. This proves that for a continuous
quotient seminorm $\overline q$ and a constant $C$,
\[
             \overline p(R_nq)\le C\overline q(q)
                         \quad\hbox{for every }n,q.    \tag{1}
\]
The same conclusion holds when the seminorm on the right is
zero, by applying the bound to arbitrary scalar multiples.

Let $b_{\rho,0}$ be the basis vector with coefficient
$\delta_{\rho,0}$. For a fixed Gaussian $g_a$ with
$\mathcal Mg_a(s)=e^{as^2}$, use the test function
\[
                f_t(x)=e^{-\sigma x-itx}g_a(x).
\]
It satisfies $\delta_{\rho,0}(f_t)=1$.
Every fixed seminorm of $f_t$ is bounded by a fixed power of
$1+|t|$, by the Leibniz formula and Gaussian decay.
Apply (1) to $[f_t]$ and its coordinate projection. For each
fixed $\overline p$ there are constants $C_p,A_p$ such that
\[
       \overline p(b_{\rho,0})\le C_p(1+|t|)^{A_p}
                                      \quad(\rho\in Z). \tag{2}
\]
The constants do not depend on the enumeration position of
$b_{\rho,0}$.

We next give a uniform estimate for confluent divided differences.
For real nodes $t_0,\ldots,t_m$, repeated nodes allowed, the
divided difference of $u\mapsto e^{ixu}$ is
\[
 [t_0,\ldots,t_m]e^{ixu}
   =\int_{\Delta_m}(ix)^m
           e^{ix\sum_{j=0}^m\lambda_jt_j}\,d\lambda,
 \qquad
 \Delta_m=\{\lambda_j\ge0:\textstyle\sum\lambda_j=1\},
                                                        \tag{3}
\]
where $d\lambda$ has total mass $1/m!$.
One elementary verification suffices also for repeated nodes.
For distinct nodes, the divided difference of $u^k$ is zero
when $k<m$ and is the complete homogeneous polynomial
$h_{k-m}(t_0,\ldots,t_m)$ when $k\ge m$. This follows by
expanding the partial-fraction identity
\[
 \sum_{j=0}^m
 \frac{1}{(1-zt_j)\prod_{\ell\ne j}(t_j-t_\ell)}
          =\frac{z^m}{\prod_{j=0}^m(1-zt_j)}.
\]
On the simplex,
$\int_{\Delta_m}\lambda_0^{a_0}\cdots\lambda_m^{a_m}\,d\lambda
=\prod_j a_j!/(m+\sum a_j)!$; successive integration of
monomials proves this formula. Expanding the exponential and
using this integral gives exactly the same power series,
proving (3). Continuity as nodes coalesce gives the confluent
version, whose functionals are derivatives divided by factorials.
In particular the absolute value in (3) is at most $|x|^m/m!$.

Take $t_0=t$ and repeat each $\operatorname{Im}\nu$, $\nu\in W$,
exactly $m_\nu$ times among the other nodes. Apply the divided
difference to $u\mapsto\mathcal Mf(\sigma+iu)$.
All resulting point-jet functionals annihilate $J$.
Choose an integer $N\ge|\sigma|+1$ and put $p=p_{N,0}$.
Formula (3), Fubini's theorem and exponential decay give
\[
 \left|[t_0,\ldots,t_m]\mathcal Mf(\sigma+iu)\right|
 \le\int_{\mathbb R}\frac{|x|^m}{m!}e^{\sigma x}|f(x)|\,dx
 \le2p(f).                                             \tag{4}
\]
The last constant is independent of both the nodes and $m$,
since $\int_{\mathbb R}|x|^m e^{-|x|}\,dx=2m!$.
The same inequality holds on $Q$ with $\overline p$.

Biorthogonality says that on $b_{\rho,0}$ all the other
point jets in this divided difference vanish. The coefficient
of the simple node $t$ is
$\prod_{\nu\in W}(t-\operatorname{Im}\nu)^{-m_\nu}$.
For repeated other nodes this follows either from partial
fractions or by coalescing distinct nodes away from $t$.
Consequently (4) and (2) imply
\[
 \prod_{\nu\in W}|t-\operatorname{Im}\nu|^{-m_\nu}
       \le2\overline p(b_{\rho,0})
       \le2C_p(1+|t|)^{A_p},
\]
which is (MT-M10).
\end{proof}

For any fixed $0<r<1$, choosing all other zeros within distance
$r$ also bounds their total multiplicity by
$(A\log(1+|t|)-\log c)/|\log r|$.
The product condition is stronger than such a local count:
it controls the joint cost of arbitrarily close clusters.
Taking one other node recovers polynomial separation.
Taking the clusters in MT-M09 contradicts (MT-M10) directly.
This identifies an explicit necessary analytic input for
Masullo 9.3.2.5, PDF~425--426, even after the closed-ideal repair.
It neither asserts that arithmetic zeros violate that condition
nor proves that the condition suffices for a basis.

\subsection*{MT-M11. The same basis would bound arithmetic leading coefficients}
Fix a primitive unitary Hecke $L$-function $L(s,\eta)$ and assume
its nontrivial zeros lie on $\operatorname{Re}s=1/2$.
Suppose a closed quotient as in MT-M10 has all its nontrivial
zero jets as coefficient functionals, with their actual
multiplicities. Then there are constants $c,C,T_0>0$, depending
on the fixed function and the proposed basis, such that at every
nontrivial zero $\rho=1/2+iT$ with $|T|\ge T_0$ and order $m$,
\[
               \left|\frac{L^{(m)}(\rho,\eta)}{m!}\right|
                        \ge c(1+|T|)^{-C}.             \tag{MT-M11}
\]
In particular, at simple zeros the reciprocal of $|L'(\rho,\eta)|$
would have at most polynomial growth. The zero-location input is
$A_D$ for Dirichlet functions and $M_H$ for general Hecke functions;
the asserted quotient basis is an additional hypothesis.

\begin{proof}
We use the fixed-function analytic facts proved in NT-Z10:
polynomial growth on a fixed vertical strip away from its
finitely located poles, and $O(\log(2+|T|))$ zeros, with
multiplicity, in each bounded-height window.
The absolutely convergent Euler product gives
$|L(2+iT,\eta)|\ge c_\eta>0$: each local inverse factor
has modulus at least $(1+(\operatorname N\mathfrak p)^{-2})^{-1}$,
and the product of these lower bounds is positive.

Put $c_T=2+iT$. Choose $4\le R\le5$ so that the boundary of
$D=\{s:|s-c_T|<R\}$ contains no zero. At sufficiently large
$|T|$, it contains no pole or trivial zero either, because
the imaginary parts of their finitely many arithmetic
progressions parallel to the real axis are fixed by the infinity type.
Its zeros are therefore precisely the relevant critical-line
zeros in that disk, say $\rho_j$ with multiplicities $m_j$.
There are
$M=\sum_jm_j=O_\eta(\log(2+|T|))$ of them.
In the normalized coordinate $z=(s-c_T)/R$, put
$a_j=(\rho_j-c_T)/R$ and form the finite Blaschke product
\[
 B(s)=\prod_j
       \left(\frac{z-a_j}{1-\overline{a_j}z}\right)^{m_j}.
\]
Constant unimodular normalizations of the factors are
irrelevant. The function $G=L/B$ is holomorphic and zero-free
on the disk and its boundary. On the boundary $|B|=1$.
The polynomial growth bound and the maximum principle give
$|G(s)|\le M_T=C_1(1+|T|)^{C_2}$ throughout $D$, enlarging
$M_T$ if necessary to make it at least one.
At the center, $|B(c_T)|\le1$, so
$|G(c_T)|\ge c_\eta$.

The nonnegative harmonic function
$u(s)=\log M_T-\log|G(s)|$ satisfies the disk Harnack bound
\[
 u(\rho)\le
 \frac{R+|\rho-c_T|}{R-|\rho-c_T|}\,u(c_T)
 \le\frac{11}{5}\,u(c_T),
\]
because $|\rho-c_T|=3/2$ and $R\ge4$.
This elementary bound follows by comparing the maximum and
minimum of the Poisson kernel on each concentric boundary
circle and then letting that circle increase to radius $R$.
It implies $|G(\rho)|\ge c_2(1+|T|)^{-C_3}$.

Let $a=(\rho-c_T)/R$, so $|a|\le3/8$.
The leading coefficient of its own factor in $B$ contributes
$[R(1-|a|^2)]^{-m}\ge5^{-m}$.
Every other factor at $\rho$ has modulus
\[
 \frac{|a-a_j|}{|1-\overline{a_j}a|}
 \ge\frac{|\rho-\rho_j|}{10}.
\]
Thus, with $W$ the distinct other zeros in $D$,
\[
 \left|\lim_{s\to\rho}\frac{B(s)}{(s-\rho)^m}\right|
 \ge 10^{-M}\prod_{\rho_j\in W}|\rho-\rho_j|^{m_j}.
\]
All these points are on the same vertical line. Apply
(MT-M10) to this finite set $W$. Since
$10^{-M}\ge c_3(1+|T|)^{-C_4}$, the displayed expression
has a polynomial lower bound. Multiplying by the lower bound
for $|G(\rho)|$ proves (MT-M11).
\end{proof}

This is a necessary consequence of the proposed arithmetic
point-jet basis, not a derivative estimate obtained from GRH
alone. It converts the remaining interpolation problem into
a concrete additional test on arithmetic leading coefficients.
The work placement is the Masullo 9.3.2.5 repair, NT-Z10,
and the cote~1 trace/Fredholm programme reported in SA23
geo-H-lidskii/geo-H-genus. Positivity of the trace cannot
replace this topological prerequisite.

\subsection*{MT-M12. An exact simple-zero gate for the canonical arithmetic quotient}
Fix a primitive unitary Hecke character $\eta$ over a number field
$K$, and assume GRH for its $L$-function. Assume additionally that
all its nontrivial zeros are simple. Let
\[
 \Xi_\eta(s)=P_\eta(s)\Lambda(s,\eta)
\]
be the entire completed function with its possible poles removed.
For a non-pure-norm character take $P_\eta=1$. If
$L(s,\eta)=\zeta_K(s+i\tau)$, take
$P_\eta(s)=(s+i\tau)(s-1+i\tau)$.
The zeros of $\Xi_\eta$ are precisely the nontrivial zeros, denoted
$Z$. Put
\[
 \psi_\eta(x)=\frac1{2\pi}\int_{\mathbb R}\Xi_\eta(it)e^{-itx}\,dt,
 \qquad Q_\eta=\mathcal E/\overline{\psi_\eta*\mathcal E}.
\]
Then $\psi_\eta\in\mathcal E$, with
$\mathcal M\psi_\eta=\Xi_\eta$, and the following are equivalent:
\begin{enumerate}
\item There are $c,C>0$ such that
$|L'(\rho,\eta)|\ge c(1+|\operatorname{Im}\rho|)^{-C}$ for
every $\rho\in Z$.
\item $Q_\eta$ has an absolute Schauder basis with coefficient
functionals $[f]\mapsto\mathcal Mf(\rho)$, $\rho\in Z$.
\item $Q_\eta$ has a Schauder basis with these continuous
coefficient functionals.
\end{enumerate}
Under these conditions the closed convolution ideal is exactly
\[
 K_Z=\{f\in\mathcal E:\mathcal Mf(\rho)=0
                              \text{ for every }\rho\in Z\}. \tag{MT-M12a}
\]
Moreover $Q_\eta$ is topologically the rapidly decreasing sequence
space
\[
 s(Z)=\{(a_\rho):\sup_{\rho\in Z}
       (1+|\operatorname{Im}\rho|)^k|a_\rho|<\infty
                                  \text{ for every }k\ge0\}. \tag{MT-M12b}
\]
This equivalence uses $A_D$ for Dirichlet functions, or $M_H$
conditionally for general unitary Hecke functions. Simplicity is
an extra assumption throughout; condition (1) is the remaining
gate, not a consequence of those GRH inputs proved here.

\begin{proof}
\emph{Analytic preliminaries.}
Write $n=[K:\mathbb Q]$ and $\kappa=\pi n/4$.
The fixed-function polynomial strip bounds of NT-Z10 and
Stirling's formula for the $n$ factors $\Gamma_{\mathbb R}$
give, for every fixed $A>0$, constants $C_A,B_A$ such that
\[
 |\Xi_\eta(\sigma+it)|
       \le C_A(1+|t|)^{B_A}e^{-\kappa|t|}
                            \qquad(|\sigma|\le A).     \tag{5}
\]
On bounded $t$ the entire function is bounded on each compact
strip rectangle, so enlarging $C_A$ covers possible gamma-pole
cancellations there. Derivative bounds follow by Cauchy's formula
on a slightly larger strip. Thus $\Xi_\eta(it)$ is Schwartz.
The contour-shifting and Fourier-inversion argument of MT-M09,
using (5) in place of the Gaussian bound, shows that the displayed
inverse transform belongs to $\mathcal E$ and has Mellin transform
$\Xi_\eta$. More explicitly, shift to $\sigma=N+1$ on $x\ge0$
and to $\sigma=-N-1$ on $x\le0$; every integral with the extra
factor $(-\sigma-it)^j$ is absolutely convergent by (5).
This proves all $p_{N,j}$ bounds. The horizontal contour edges
tend to zero by the same estimate, and Fourier uniqueness then
proves the transform identity.

Assume condition (1). At a zero $\rho=1/2+iT$, differentiation
of the completed function gives
\[
 \Xi_\eta'(\rho)
  =P_\eta(\rho)q_\eta^{\rho/2}
       \prod_j\Gamma_{\mathbb R}(\rho+\mu_j)L'(\rho,\eta),
\]
where $q_\eta=|D_K|\operatorname N\mathfrak f_\eta$
is the finite conductor.
Stirling's formula and (1) imply, after enlarging constants to
cover the finitely many low zeros,
\[
        |\Xi_\eta'(\rho)|
           \ge c_1(1+|T|)^{-C_1}e^{-\kappa|T|}.         \tag{6}
\]

\emph{Polynomially bounded cardinal functions.}
Fix $a>0$. For each zero $\rho$ define the entire function
\[
 H_\rho(s)=
     \frac{\Xi_\eta(s)}
          {\Xi_\eta'(\rho)(s-\rho)}\,e^{a(s-\rho)^2}.    \tag{7}
\]
It takes the value one at $\rho$ and zero at every other zero.
For each fixed strip $|\operatorname{Re}s|\le A$, (5)--(6)
imply a bound
\[
 |H_\rho(\sigma+it)|
       \le C_A'(1+|T|)^{B_A'}
                              e^{-a(t-T)^2/2}.         \tag{8}
\]
Here are the details near the apparent denominator. If
$|s-\rho|\ge1$, use (5) and (6), together with
$\big||T|-|t|\big|\le|T-t|$ and
$1+|t|\le(1+|T|)(1+|t-T|)$.
The remaining fixed polynomial and
$e^{\kappa|t-T|}$ are absorbed by half the Gaussian.
If $|s-\rho|<1$, use
\[
 \frac{\Xi_\eta(s)}{s-\rho}
       =\int_0^1\Xi_\eta'(\rho+u(s-\rho))\,du
\]
and the Cauchy derivative bound from (5) on a slightly larger
strip. Since $|t-T|<1$, the same estimate (8) follows by
adjusting the constants.

Invert $H_\rho$ by Fourier transform and call the inverse
$h_\rho$. The same contour argument gives $h_\rho\in\mathcal E$,
$\mathcal Mh_\rho=H_\rho$, and, for each $N,j$,
\[
             p_{N,j}(h_\rho)
                    \le C_{N,j}(1+|T|)^{B_{N,j}}.      \tag{9}
\]
Indeed, on the shifted contour used for $x\ge0$ or $x\le0$,
integrating $(1+|t|)^j$ against the Gaussian in (8) costs
only another fixed power of $1+|T|$.

For any $f\in\mathcal E$, its values $\mathcal Mf(1/2+iT)$
decrease faster than every power of $|T|$. This follows by
integrating by parts arbitrarily many times in
$\int e^{x/2}f(x)e^{iTx}\,dx$; each resulting $L^1$ norm
is bounded by finitely many defining seminorms of $f$.
The local count NT-Z10 implies
$\sum_{\rho\in Z}(1+|\operatorname{Im}\rho|)^{-2}<\infty$.
Consequently
\[
              I(f)=\sum_{\rho\in Z}\mathcal Mf(\rho)h_\rho
                                                               \tag{10}
\]
converges absolutely in every seminorm of $\mathcal E$,
defines a continuous linear map, and has exactly the same
values as $f$ at all the zeros. Likewise, for every
$(a_\rho)\in s(Z)$ the series $\sum_\rho a_\rho h_\rho$
converges in $\mathcal E$ and depends continuously on that
sequence. These assertions use (9) and choose the coefficient
decay power separately for each target seminorm.

\emph{Division on vertical strips.}
We now show the ideal equality, which does not follow merely
from constructing the cardinal functions.
Condition (1) implies a uniform, fixed-strip minimum-modulus
estimate
\[
 |L(\sigma+it,\eta)|
   \ge c_A(1+|t|)^{-C_A}
            \min\{1,\operatorname{dist}(\sigma+it,Z)\},
 \qquad |\sigma|\le A,\quad |t|\ge T_A.                 \tag{11}
\]
Here is a complete proof. Put $r=A+2$ and center a disk at
$c_t=2+it$. Choose $R\in[4r+4,4r+5]$ with zero-free boundary.
At large $|t|$ there are no poles or trivial zeros in that disk.
Its $M=O_A(\log(2+|t|))$ zeros are all simple and on the
critical line. Remove them by the finite Blaschke product
$B$ for this disk, as in MT-M11, and put $G=L/B$.
Polynomial strip growth gives $|G|\le C(1+|t|)^C$ on the
whole disk, while the Euler lower bound at $c_t$ and Harnack's
inequality give
$|G|\ge c(1+|t|)^{-C}$ on the disk of radius $R/2$.
The Harnack factor on that smaller disk is at most three.
The point $s=\sigma+it$ is within radius $r<R/4$.

If the disk has no zeros, this already proves (11).
Otherwise choose a zero $\rho_0$ in it nearest to $s$.
If $|\rho_0-c_t|>R/2$, every zero in the disk is at
distance at least $R/4$ from $s$. Each Blaschke factor
therefore has modulus at least $1/8$, and
$|B(s)|\ge8^{-M}$, a polynomial lower bound.
If $|\rho_0-c_t|\le R/2$, simplicity gives
$L'(\rho_0)=G(\rho_0)B'(\rho_0)$.
In normalized coordinates $|a_0|\le1/2$, the derivative
of its own Blaschke factor is at most $4/(3R)$, and
every other factor at $\rho_0$ is at most
$2|\rho_0-\rho_j|/R$. Thus
\[
 |B'(\rho_0)|\le C_A^M
                     \prod_{\rho_j\ne\rho_0}|\rho_0-\rho_j|.
\]
The assumed lower bound for $L'$, the upper bound for $G$,
and $M=O_A(\log(2+|t|))$ imply
\[
          \prod_{\rho_j\ne\rho_0}|\rho_0-\rho_j|
                          \ge c(1+|t|)^{-C}.           \tag{12}
\]
Since $\rho_0$ is nearest to $s$,
$|\rho_0-\rho_j|\le2|s-\rho_j|$. Each denominator in a
Blaschke factor at $s$ is at most two. Hence
\[
 |B(s)|\ge(4R)^{-M}|s-\rho_0|
                       \prod_{\rho_j\ne\rho_0}|\rho_0-\rho_j|,
\]
and (12) proves (11). Zeros outside the large disk have
distance greater than one from $s$, so using the distance
to the full set $Z$ makes no difference in the minimum.
The estimate at a zero is interpreted as zero on both
sides and follows by continuity.

\emph{The closed ideal is the full zero kernel.}
The inclusion
$\overline{\psi_\eta*\mathcal E}\subset K_Z$ is immediate
from multiplication of Mellin transforms.
Take $f\in K_Z$ and write $F=\mathcal Mf$.
The quotient $F/\Xi_\eta$ is entire because all zeros
are simple and have been cancelled. On every fixed strip,
(11) and the rapid vertical decay of $F$ and $F'$ imply
rapid decay of $F/L$ for large $|t|$. If the distance to
$Z$ is at least one, divide the rapid bound for $F$
by (11). If it is less than one, choose a nearest zero
$\rho$ and use
$F(s)=\int_\rho^s F'(w)\,dw$ on the straight segment;
that segment lies in a slightly larger fixed strip and
a bounded-height window. The distance factor cancels
the corresponding factor in (11).

Inverting the gamma factors now costs at most
$C_A(1+|t|)^{C_A}e^{\kappa|t|}$ on the strip.
For any $\epsilon>0$ the entire function
\[
             J_\epsilon(s)=e^{\epsilon s^2}
                                  \frac{F(s)}{\Xi_\eta(s)}
\]
therefore decays faster than every power on every vertical
strip, in fact with a Gaussian bound. On compact rectangles
it is bounded by entireness. Contour shifting and Fourier
inversion give $j_\epsilon\in\mathcal E$ with this Mellin
transform. Uniqueness of the transform gives
\[
                  \psi_\eta*j_\epsilon=g_\epsilon*f.
\]
Finally $g_\epsilon*f\to f$ in $\mathcal E$ as
$\epsilon\downarrow0$. To verify this in each $p_{N,j}$,
write the difference as the integral of
$g_\epsilon(y)(f^{(j)}(x-y)-f^{(j)}(x))$.
For $|y|\le1$, the weighted mean-value bound is
$|y|e^{N|y|}p_{N,j+1}(f)$; outside any fixed
neighborhood of zero, use
$(e^{N|y|}+1)p_{N,j}(f)$ and the Gaussian tail.
Both integrated bounds tend to zero.
Thus $f$ belongs to the closed convolution ideal,
proving (MT-M12a).

The coordinate map $Q_\eta\to s(Z)$ is now injective.
It is continuous by the integration-by-parts bounds above,
which descend to quotient seminorms because all coordinates
annihilate the ideal. The continuous series map
$(a_\rho)\mapsto[\sum_\rho a_\rho h_\rho]$ is its inverse.
This proves (MT-M12b), and the usual coordinate vectors
give (2), with representatives $h_\rho$.
Statement (2) implies (3). Statement (3) implies (1)
by MT-M11 with $m=1$, extending its high-zero estimate
over the finitely many remaining zeros. This completes
the equivalence.
\end{proof}

Replacing $\psi_\eta$ by
$g_b*\psi_\eta$ for any $b>0$ leaves the closed ideal unchanged:
Gaussian convolution has dense range by MT-M06 and convolution
by $\psi_\eta$ is continuous.

The theorem settles a precise conditional interpolation and
spectral-synthesis interface, rather than identifying every
quotient occurring in Masullo's draft. For this canonical
completed-Hecke quotient with simple zeros, the Schauder and
absolute-Schauder existence gates are exactly the polynomial
reciprocal-derivative gate. The extension to arbitrary
multiplicities requires bounds for the whole local Hermite
interpolation data and is not asserted here.

\subsection*{MT-M13. Nuclear arithmetic convolution under the exact basis gate}
Under the equivalent conditions of MT-M12, convolution by
$f\in\mathcal E$ acts on $Q_\eta\simeq s(Z)$ as
\[
            (a_\rho)\longmapsto
                         (\mathcal Mf(\rho)a_\rho).
\]
It has an explicit nuclear representation on that Fr\'echet
space and a trace-class extension to $\ell^2(Z)$, with
\[
                   \operatorname{Tr}C_f
                         =\sum_{\rho\in Z}\mathcal Mf(\rho).
                                                               \tag{MT-M13}
\]
Here $\operatorname{Tr}$ denotes the trace of that Hilbert-space extension.
The sum converges absolutely. For
$h^\sharp(x)=e^{-x}\overline{h(-x)}$ and
$f=h*h^\sharp$, this Hilbert-space operator is positive and
\[
                   \operatorname{Tr}C_f
                         =\sum_{\rho\in Z}|\mathcal Mh(\rho)|^2
                         \ge0.
\]

\begin{proof}
Convolution descends because it preserves the closed
convolution ideal. Multiplication of Mellin transforms
gives the diagonal action. The eigenvalues decrease
faster than every power of the ordinate, so their
absolute sum is finite by NT-Z10. A diagonal operator
with these eigenvalues on the orthonormal coordinate
basis of $\ell^2(Z)$ is trace class and has the stated
trace.

To see nuclearity without an unspecified basis estimate,
write $b_\rho=[h_\rho]$ and
$\delta_\rho([g])=\mathcal Mg(\rho)$. The functionals
$\delta_\rho$ are equicontinuous: for a fixed
$N>1/2$, their absolute values are uniformly bounded
by a constant times the quotient seminorm of $p_{N,0}$.
The representation
\[
 C_f=\sum_{\rho\in Z}(1+|\operatorname{Im}\rho|)^{-2}
       \delta_\rho\otimes
       \bigl((1+|\operatorname{Im}\rho|)^2
                                  \mathcal Mf(\rho)b_\rho\bigr)
\]
has summable scalar coefficients, an equicontinuous
family of functionals, and a bounded family of vectors
in every quotient seminorm by (9) and rapid decay.
It is therefore an explicit nuclear representation.
Its diagonal sum is the same absolutely convergent
trace displayed above; the unambiguous Hilbert
trace needs no additional abstract trace convention.

The weighted involution preserves $\mathcal E$ and satisfies
\[
      \mathcal M h^\sharp(s)
                =\overline{\mathcal Mh(1-\overline s)}.
\]
At every $\rho$ on the critical line,
$1-\overline\rho=\rho$. Thus the diagonal eigenvalues
for $h*h^\sharp$ are the nonnegative squares stated,
which proves positivity and its trace formula.
\end{proof}

This is a rigorous operator model conditional on the exact
simple-zero basis gate. Its Hilbert completion is constructed
from the zero coordinates; it is not an independently
constructed geometric or prime-side operator proving GRH.
Its location in the programme is Masullo 9.3.2.5--9.3.3,
the Tate-adjoint involution, and cote~1's nuclear trace
framework as carried in SA23 geo-H-lidskii/geo-H-genus.

\subsection*{MT-M14. Bounded multiplicities and the full Hermite basis}
Retain the canonical completion and quotient of MT-M12, assume
GRH, and replace simplicity by the hypothesis
\[
                  1\le m_\rho\le m_0<\infty
                              \qquad(\rho\in Z),       \tag{13}
\]
where $m_\rho$ is the actual zero multiplicity.
Then the following are equivalent:
\begin{enumerate}
\item For some $c,C>0$,
$|L^{(m_\rho)}(\rho,\eta)/m_\rho!|
       \ge c(1+|\operatorname{Im}\rho|)^{-C}$ for all $\rho$.
\item The canonical quotient has an absolute Schauder basis
whose coefficient functionals are
$[f]\mapsto(\mathcal Mf)^{(k)}(\rho)/k!$ for all
$0\le k<m_\rho$.
\item It has an ordinary Schauder basis with precisely those
continuous coefficient functionals.
\end{enumerate}
Under these conditions the closed ideal is the full zero-jet
kernel, and the quotient is the rapid sequence space on
$\{(\rho,k):0\le k<m_\rho\}$, weighted by
$1+|\operatorname{Im}\rho|$. Assumption (13) and condition (1)
are additional arithmetic hypotheses, not consequences of GRH
established in this run.

\begin{proof}
The necessity of (1) from (3) is exactly MT-M11.
We give all changes needed for sufficiency; they also show why
the fixed bound $m_0$ is used.
Write $z=s-\rho$, $m=m_\rho$, and
\[
                  G_\rho(s)=\Xi_\eta(s)/(s-\rho)^m .
\]
The leading-coefficient hypothesis, multiplied by the gamma
and conductor factors, gives
\[
 |G_\rho(\rho)|\ge c(1+|T|)^{-C}e^{-\kappa|T|},
           \qquad T=\operatorname{Im}\rho.
\]
Cauchy's formula on a fixed larger disk around $\rho$ and
the strip bound (5) give
$|G_\rho^{(j)}(\rho)/j!|
 \le C(1+|T|)^C e^{-\kappa|T|}$ for $0\le j\le m_0$.
To see the division estimate directly, on a circle $|z|=2$
divide $\Xi_\eta(\rho+z)$ by $z^m$ and apply Cauchy's
formula inside; $m\le m_0$ keeps all constants fixed.

Fix $a>0$. Let $P_{\rho,k}(z)$ be the Taylor polynomial
through degree $m-1-k$ of the analytic germ
$e^{-az^2}/G_\rho(\rho+z)$.
The reciprocal-coefficient recursion starts with
$1/G_\rho(\rho)$ and at each subsequent degree divides
a finite sum by $G_\rho(\rho)$. The displayed upper and
lower bounds therefore bound every coefficient of
$P_{\rho,k}$ by
$C(1+|T|)^C e^{\kappa|T|}$, uniformly for
$0\le k<m\le m_0$. Only finitely many recursion steps
are involved. Define
\[
 H_{\rho,k}(s)=
      z^kP_{\rho,k}(z)e^{az^2}G_\rho(s).                \tag{14}
\]
Its Taylor expansion at $\rho$ is
$z^k+O(z^m)$, and it vanishes to the required multiplicity
at every other zero. On every fixed vertical strip its
bound has the same Gaussian form (8), with constants
depending also on $m_0$: away from $|z|<1$, use (5)
and the fixed polynomial degree in (14); inside,
use the Cauchy bound for $G_\rho$.
Inverse Mellin transforms $h_{\rho,k}\in\mathcal E$
therefore satisfy polynomial seminorm bounds (9)
uniformly over $k<m_\rho$.

The coefficients $(\mathcal Mf)^{(k)}(\rho)/k!$ are
uniformly rapidly decreasing in $|T|$ for
$k<m_0$. Indeed, they are Mellin transforms of the
fixed finite list $x^k f(x)/k!$, and the
integration-by-parts argument of MT-M12 applies to
each. The series with coefficients these jets and
vectors $h_{\rho,k}$ consequently converges absolutely
in $\mathcal E$, is continuous in $f$, and interpolates
all its jets. Conversely every rapid sequence of jets
is obtained from the corresponding convergent series.
Multiplicity is included in the local zero count,
so the same summability estimate applies.

It remains to verify the full ideal equality.
Repeat the disk argument proving (11), now with
Blaschke factors raised to the actual multiplicities.
The number of factors counted with multiplicity is still
$M=O_A(\log(2+|t|))$. If a nearest zero $\rho_0$
is in the inner half disk, the leading coefficient
identity is
\[
 \frac{L^{(m_{\rho_0})}(\rho_0)}{m_{\rho_0}!}
   =G(\rho_0)
       \lim_{s\to\rho_0}
                    \frac{B(s)}{(s-\rho_0)^{m_{\rho_0}}}.
\]
The upper bound for that leading Blaschke coefficient
is $C_A^M$ times
$\prod_{\rho_j\ne\rho_0}
          |\rho_0-\rho_j|^{m_{\rho_j}}$.
Condition (1) therefore supplies the polynomial lower
bound (12) for this multiplicity-weighted product.
Nearest-zero comparison gives
\[
 |L(s,\eta)|\ge c_A(1+|t|)^{-C_A}
                         |s-\rho_0|^{m_{\rho_0}}
\]
when the nearest distance is less than one.
When that distance is at least one, or the nearest
disk zero is outside the inner half disk, the same
argument gives a polynomial lower bound without
the distance factor. All constants are uniform
because $m_{\rho_0}\le m_0$ and $M=O_A(\log|t|)$.

If $F=\mathcal Mf$ vanishes to the required orders,
Taylor's formula at a nearest zero gives
\[
 F(s)=\frac{(s-\rho)^m}{(m-1)!}
       \int_0^1(1-u)^{m-1}
                         F^{(m)}(\rho+u(s-\rho))\,du .
\]
For $m\le m_0$ its derivative bound is uniformly
rapid in $|t|$. The distance factors cancel.
Thus $F/L$ is rapidly decreasing on each fixed
strip at large height. Exactly the Gaussian-regularized
division and approximate-identity proof in MT-M12
shows that $f$ lies in
$\overline{\psi_\eta*\mathcal E}$.
The reverse inclusion is immediate.
The interpolation series now identifies the quotient
topologically with the rapid jet-sequence space,
as in MT-M12. This proves (2), hence (3), and
completes the equivalence.
\end{proof}

No unspecified infinite Hermite estimate has been hidden in
the argument: the reciprocal-coefficient recursion has at most
$m_0$ steps. With unbounded multiplicities, the exponents of
the polynomial bounds produced by that recursion need not be
uniform. That case remains a separate interpolation gate.

\subsection*{MT-M15. Multiple zeros: positive traces and nilpotent operators}
Under the conditions of MT-M14, convolution has an explicit
nuclear representation on the rapid jet-sequence quotient
and a trace-class realization on the standard $\ell^2$
space of those jets, with
\[
          \operatorname{Tr}C_f
                    =\sum_{\rho\in Z}m_\rho\mathcal Mf(\rho).
                                                               \tag{MT-M15}
\]
For $f=h*h^\sharp$ this trace is
$\sum_\rho m_\rho|\mathcal Mh(\rho)|^2\ge0$.
However, if any $m_\rho>1$, there is no positive-definite
inner product on the full faithful jet quotient for which
every such adjoint-square convolution is a positive
operator. In particular, the trace positivity does not
upgrade to positivity of the entire nonsemisimple action.

\begin{proof}
At a zero $\rho$ let $a_j(f)=(\mathcal Mf)^{(j)}(\rho)/j!$.
The product rule in normalized jets is
\[
       a_k(f*g)=\sum_{j=0}^k a_j(f)a_{k-j}(g).
\]
Thus the block of $C_f$ at $\rho$ is a lower triangular
Toeplitz matrix with diagonal $a_0(f)$ and size $m_\rho$.
All its entries decrease rapidly with the ordinate,
uniformly in the at most $m_0$ indices.
The trace norm of a finite matrix is at most the sum
of the trace norms of its individual matrix units,
so the block trace norms are summable. This proves
the Hilbert trace-class assertion and its trace.

For nuclearity, use the rank-one expansion in these
matrix entries. The jet functionals are equicontinuous:
for a fixed $N>3/2$,
\[
 |(\mathcal Mg)^{(k)}(\rho)/k!|
 \le \frac1{k!}\int_{\mathbb R}|x|^k e^{x/2}|g(x)|\,dx
 \le2p_{N,0}(g),
\]
uniformly even before using $k<m_0$.
The interpolation vectors have polynomial seminorm
bounds, and the matrix entries have rapid decay.
Insert the summable scalar weight $(1+|T|)^{-2}$
as in MT-M13. With at most $m_0^2$ matrix units per
zero, this is an explicit nuclear representation.
The weighted involution has jet coefficients
\[
                   a_j(h^\sharp)=(-1)^j\overline{a_j(h)}.
\]
Its zeroth coefficient gives the asserted nonnegative
trace for an adjoint square.

Finally suppose $m=m_\rho>1$. The interpolating functions
in MT-M14 realize every local jet and zero jets elsewhere.
The local algebra is $\mathbb C[z]/(z^m)$, with
$z^\sharp=-z$ and conjugation on scalars.
Let $p$ denote its identity, extended by zero in every
other zero block, and let $n=z$ in this block and zero
elsewhere. Both are represented by actual functions
of $\mathcal E$, by (14). On the $m$-dimensional local
block, their convolution operators are $I$ and a
nonzero nilpotent $N$.
For every real $u$, the class $h=p+iu\,n$ satisfies
$h^\sharp=h$, and hence its adjoint square acts there as
\[
                         I+2iuN-u^2N^2.
\]
If all adjoint-square operators were positive in a
positive-definite inner product, these operators and
the corresponding operators with $-u$ would be
self-adjoint on the invariant local block.
Subtracting them for $u\ne0$ shows that $iN$ is
self-adjoint, equivalently $N^*=-N$.
A skew-adjoint operator on a finite-dimensional positive
inner-product space is normal and diagonalizable;
nilpotence forces all eigenvalues to vanish, so it is
zero. This contradicts the regular action of $z$
on $\mathbb C[z]/(z^m)$.
The argument is internal to a finite block and
does not depend on any choice of Hilbert completion.
\end{proof}

Passing to the semisimple quotient that retains only the
zeroth coordinate of each block removes the nilpotents.
Using $m_\rho$ copies of each scalar coordinate then
recovers the positive spectral trace with multiplicity.
It does not faithfully recover the discarded higher
jets. This is the precise distinction needed when
transferring Alpöge-style positivity to an orbital
jet quotient in Masullo 9.3.2.5--9.3.3, or to the
cote~1/SA23 trace programme. No assertion that the
draft explicitly claims a positive faithful Hilbert
realization of this entire quotient is made here.

\subsection*{MT-M16. The full arithmetic jet basis has one exact quantitative gate}
The bounded-multiplicity hypothesis in MT-M14 can be removed.
Fix a primitive unitary Hecke function, assume GRH, and use
the canonical $\Xi_\eta,\psi_\eta,Q_\eta$ of MT-M12.
Allow all actual multiplicities $m_\rho$, without a uniform
bound. Write $\rho=1/2+iT$ and set
\[
 P_\rho=
  \prod_{\substack{\nu\in Z,\ \nu\ne\rho\\
              |\operatorname{Im}\nu-T|\le1}}
              |T-\operatorname{Im}\nu|^{m_\nu},
\]
with empty product one. The following conditions are equivalent:
\begin{enumerate}
\item The first nonzero finite-$L$ Taylor coefficients satisfy
$|L^{(m_\rho)}(\rho,\eta)/m_\rho!|
                       \ge c(1+|T|)^{-C}$ uniformly in $\rho$.
\item The local products satisfy
$P_\rho\ge c(1+|T|)^{-C}$ uniformly in $\rho$.
\item The canonical closed convolution quotient has a
Schauder basis with the prescribed normalized Mellin
zero-jet coefficient functionals.
\item It has such an absolute Schauder basis, its closed
ideal is the full zero-jet kernel, and it is topologically
the rapid sequence space on
$\mathcal J=\{(\rho,k):0\le k<m_\rho\}$, weighted by $1+|T|$.
\end{enumerate}
Constants may differ between these equivalent conditions.
Only GRH for this fixed function is used analytically:
$A_D$ supplies it for Dirichlet functions and $M_H$
supplies it conditionally in general. The equivalent
quantitative conditions remain an extra gate; none is
being deduced here from zero location alone.

\begin{proof}
All constants below may depend on the fixed function and
on the indicated fixed strip. NT-Z10 gives both polynomial
strip growth of the finite $L$-function and a local count
including multiplicities. In particular
\[
       m_\rho=O(\log(2+|T|)),\qquad
       \sum_{\rho\in Z}m_\rho(1+|T|)^{-2}<\infty.       \tag{15}
\]
The finitely many low zeros can be included by adjusting
positive constants in every estimate.

\emph{Equivalence of the two quantitative conditions.}
Take the disk centered at $2+iT$ with a zero-free boundary
of radius $R\in[4,5]$. At large height it contains no
poles or trivial zeros. Its total zero multiplicity is
$M=O(\log(2+|T|))$. Remove its zeros by the finite
Blaschke product, with multiplicities, and write $L=GB$.
As proved in MT-M11, $G$ is bounded above by a polynomial
in $|T|$ and below by an inverse polynomial on the
inner half disk; this uses the Euler lower bound at
the center and Harnack's inequality.

The point $\rho$ lies in that inner disk. Its own leading
Blaschke coefficient is
$[R(1-|a|^2)]^{-m_\rho}$ with $a=-3/(2R)$.
Every other factor at $\rho$ is its distance from $\rho$
times a factor bounded above and below by fixed positive
constants. These statements follow directly from
$(z-a_j)/(1-\overline{a_j}z)$ and $|a|\le3/8$.
All other zeros of distance at most one are in the disk.
The distances of the remaining disk zeros are between
one and $2R$. Since there are $M$ factors with multiplicity,
all constant factors together cost at most a fixed power
of $2+|T|$. Therefore
\[
 c(1+|T|)^{-C}P_\rho
   \le\left|\frac{L^{(m_\rho)}(\rho,\eta)}{m_\rho!}\right|
   \le C(1+|T|)^C P_\rho.                              \tag{16}
\]
This proves (1) if and only if (2).

\emph{A radius adapted to the full multiplicity.}
Assume (2), and let $\delta_\rho$ be the lesser of one
and the distance to the nearest distinct zero.
If a distinct zero $\nu$ has this distance less than one,
the product $P_\nu$ contains the factor
$\delta_\rho^{m_\rho}$. Every other factor in $P_\nu$
is at most one. Hence
\[
 \delta_\rho^{m_\rho}\ge P_\nu
                         \ge c(1+|T|)^{-C}.
\]
When $\delta_\rho=1$ the same conclusion is immediate.
Put $r_\rho=\delta_\rho/4$. Equation (15) now gives
\[
                        r_\rho^{-m_\rho}
                                      \le C(1+|T|)^C. \tag{17}
\]
This is the quantitative replacement for the fixed number
of reciprocal-Taylor steps used in MT-M14.

\emph{Minimum modulus without bounded multiplicities.}
The disk proof of (11) works with all multiplicities once
(2) is available. We give the detail that makes the
constants independent of the largest multiplicity.
For a point $s=\sigma+it$ in a fixed strip, choose the
larger disk from MT-M12 and write again $L=GB$.
If that disk has no zeros, the lower bound for $G=L$
already gives the required polynomial lower bound.
If its nearest zero is outside the inner half disk,
every Blaschke factor at $s$ is at least $1/8$.
Otherwise let $\rho_0$ be a nearest disk zero in that
inner half disk. For every other zero,
$|\rho_0-\nu|\le2|s-\nu|$.
The product of the other distances, with multiplicity,
is at least $P_{\rho_0}$: omitted factors of distance
at most one only increase that product, while all
included distances greater than one are at least one.
The normalized Blaschke denominators and radii cost
only a constant to the total multiplicity
$M=O(\log(2+|t|))$. The lower bound for $G$ therefore
gives
\[
 |L(s,\eta)|\ge c_A(1+|t|)^{-C_A}
             \operatorname{dist}(s,Z)^{m_{\rho_0}}
                       \quad\hbox{if }\operatorname{dist}(s,Z)<1,
                                                               \tag{18}
\]
and a polynomial lower bound when that distance is
at least one. The choice of a nearest zero in (18)
is harmless in a tie. Zeros outside the larger disk
have distance greater than one from $s$.

\emph{Uniform Hermite interpolation.}
Put $m=m_\rho$, $z=s-\rho$, and
$G_\rho(s)=\Xi_\eta(s)/z^m$.
On the circle $|z|=r_\rho$, $\rho$ is the nearest zero.
Divide (18) by $|z|^m$ and use the fixed-strip lower
Stirling bound for the gamma factors. This gives
\[
             |G_\rho(s)|\ge
                 c(1+|T|)^{-C}e^{-\kappa|T|}
                            \quad(|z|=r_\rho).
\]
The function $G_\rho$ has no zero in that disk.
Consequently Cauchy's formula bounds the $j$th Taylor
coefficient of $e^{-az^2}/G_\rho(\rho+z)$ by
\[
       C(1+|T|)^C e^{\kappa|T|}r_\rho^{-j}.
                                                               \tag{19}
\]
For every $j<m$, (17) makes this a polynomial bound
times $e^{\kappa|T|}$, with constants independent of $j,m$.
Use exactly the Taylor truncation $P_{\rho,k}$ and
cardinal function (14) of MT-M14, now for all $k<m$.

Their fixed-strip bounds are uniform. When $|z|\ge1$,
the degree of $z^kP_{\rho,k}$ is at most $m-1$, so
\[
 |z^kP_{\rho,k}(z)/z^m|
       \le m\,C(1+|T|)^C e^{\kappa|T|}.
 \]
There is no growing positive power of $|z|$.
The factor $m$ is absorbed using (15), and the Gaussian
bound (8) follows from (5), just as before.
When $|z|<1$, the maximum principle applied to
$G_\rho$ on $|z|=2$ gives
$|G_\rho(s)|\le C(1+|T|)^C e^{-\kappa|T|}2^{-m}$.
The polynomial $z^kP_{\rho,k}$ is bounded by
$m$ times the coefficient bound (19), with (17).
This proves (8) on the remaining region as well,
after increasing a fixed constant.
Inverse transforms $h_{\rho,k}$ therefore satisfy
the polynomial seminorm bounds (9) uniformly over
every jet index $k<m_\rho$.

The normalized jets of any $f\in\mathcal E$ also have
uniform rapid height decay for \emph{all} $k\ge0$:
for each fixed $j$ and each fixed real-part strip,
\[
 \sup_{k\ge0}
 (1+|t|)^j\left|\frac{(\mathcal Mf)^{(k)}(\sigma+it)}{k!}\right|
                    \le C_{A,j}\,p_{N,j}'(f),          \tag{20}
\]
where $p_{N,j}'$ is a finite maximum of defining
seminorms, with sufficiently large fixed $N$.
Indeed, integrate by parts $j$ times in the transform
of $e^{\sigma x}x^k f(x)/k!$.
Every derivative of $x^k/k!$ is either zero or
$x^{k-\ell}/(k-\ell)!$; the remaining coefficients
depend only on $A,j$. For $N-A>1$, its weighted
$L^1$ norm is at most
$2/(N-A)^{k-\ell+1}\le2$, times the corresponding
seminorm of $f$. The region $|t|\le1$ follows from
the same integral bound without integration by parts.
This proves (20).

Equations (15), (20), and the uniform cardinal bounds
give absolute interpolation and its continuous inverse
for every rapid sequence on $\mathcal J$.
Finally, division proves the full ideal equality.
If $F=\mathcal Mf$ has all prescribed zero jets, its
Taylor remainder at a nearest zero is
\[
 F(s)=z^m\int_0^1m(1-u)^{m-1}
                  \frac{F^{(m)}(\rho+uz)}{m!}\,du .
 \]
The measure in this integral has total mass one.
Estimate (20) thus bounds the normalized remainder
uniformly in $m$, and its $|z|^m$ cancels the distance
factor in (18). The quotient $F/L$ is rapidly
decreasing on every fixed strip at large height.
The entire quotient $F/\Xi_\eta$, multiplied by
$e^{\epsilon s^2}$, is the Mellin transform of an
element of $\mathcal E$ for every $\epsilon>0$.
The Gaussian approximate-identity argument of
MT-M12 proves that $f$ lies in the closed ideal.
The reverse inclusion is immediate.
This yields the asserted topological rapid-sequence
identification and (4). Statement (4) implies (3),
and (3) implies (1) by MT-M11. All four conditions
are therefore equivalent.
\end{proof}

This removes both the simplicity assumption of MT-M12 and the
uniform multiplicity assumption of MT-M14. It does not prove the
remaining leading-coefficient or cluster-product bound.
The local $O(\log|T|)$ count is used together with that
quantitative condition; MT-M09 shows why the count alone
cannot replace it.

\subsection*{MT-M17. Nuclear trace with arbitrary zero multiplicities}
Under the equivalent conditions of MT-M16, every
$f\in\mathcal E$ acts by nuclear convolution on $Q_\eta$
and by a trace-class operator on $\ell^2(\mathcal J)$.
Its Hilbert trace is
\[
       \operatorname{Tr}C_f
                  =\sum_{\rho\in Z}m_\rho\mathcal Mf(\rho).
\]
For $f=h*h^\sharp$ this is the nonnegative sum
$\sum_\rho m_\rho|\mathcal Mh(\rho)|^2$.
The faithful-operator positivity obstruction of MT-M15
persists whenever any zero is multiple.

To prove these assertions without a uniform block-size
hypothesis, note that (15) bounds the total multiplicity
in each unit-height interval by $O(\log(2+|T|))$.
The number of Toeplitz matrix units in that interval is
therefore at most $O(\log^2(2+|T|))$.
Estimate (20) makes every matrix entry of $C_f$ uniformly
rapidly decreasing in height, independently of its jet
index. Thus the sum of the trace norms of all matrix
units is finite, giving the trace-class assertion and
the displayed diagonal trace. For nuclearity use the
equicontinuous normalized-jet functionals, the uniform
polynomial seminorm bounds for $h_{\rho,k}$, and the
scalar weight $(1+|T|)^{-2}$ for each matrix unit.
These weights are summable even with the
$O(\log^2(2+|T|))$ count, and the resulting vector
family is bounded in every seminorm by (20).
This is the same explicit nuclear decomposition as
MT-M15, now with all indices included.
The involution formula and the finite local nilpotent
argument in MT-M15 make no bounded-multiplicity
assumption, so they prove the remaining statements.

This is an exact reduction for the prescribed normalized-jet
coordinates. MT-M18 and MT-M19 below supply a different,
block-adapted absolute basis, spectral synthesis and nuclearity
without that quantitative condition. Identification with a
particular independently defined adelic quotient is a separate
map; MT-M21 begins the comparison with Meyer's construction.

\subsection*{MT-M18. Unconditional arithmetic block synthesis and an absolute adapted basis}

Fix a primitive unitary Hecke character \(\eta\) over a number field.
Use the entire completion \(\Xi=\Xi_\eta=P_\eta\Lambda(s,\eta)\)
and its inverse Mellin transform \(\psi=\psi_\eta\) from MT-M12,
but make no GRH or simplicity assumption.  The polynomial \(P_\eta\)
removes the possible pure-norm poles exactly as there.  Let \(Z\)
be the set of distinct zeros of \(\Xi\), with their actual
multiplicities \(m_\rho\).  These are the nontrivial zeros and
lie in \(0\le\operatorname{Re}s\le1\).
Put
\[
 \begin{gathered}
 \mathcal E=\{f\in C^\infty(\mathbb R,\mathbb C):
   p_{N,j}(f)=\sup_x e^{N|x|}|f^{(j)}(x)|<\infty
                    \text{ for all integers }N,j\ge0\},\\
 Q=\mathcal E/\overline{\psi*\mathcal E}.
 \end{gathered}
\]
All constants below concern this one fixed Hecke function; no
uniformity in its field, conductor, or infinity type is asserted.

There are real numbers
\[
 n+\tfrac14<h_n<n+\tfrac34\quad(n\in\mathbb Z)
\]
not equal to any zero ordinate, with the following properties.
Let \(\Omega_n\) be the interior of the rectangle
\([-1,2]+i[h_n,h_{n+1}]\), and orient its boundary
\(\Gamma_n\) positively.  Let \(A(\Omega_n)\) be the Banach
algebra of functions holomorphic in \(\Omega_n\) and continuous
on its closure, with the supremum norm.  Write
\[
 \begin{split}
 I_n&=\{F\in A(\Omega_n):
        F^{(j)}(\rho)=0\text{ for }\rho\in Z\cap\Omega_n,\
                                  0\le j<m_\rho\},\\
 J_n&=A(\Omega_n)/I_n,\qquad
 \|a\|_n=\inf_{F\bmod I_n=a}\sup_{\Gamma_n}|F|,\qquad
 d_n=\dim_{\mathbb C}J_n=\sum_{\rho\in Z\cap\Omega_n}m_\rho .
 \end{split}                                                    \tag{M18.1}
\]
An empty block has \(J_n=0\).  Then
\[
 d_n=O(\log(2+|n|)),\qquad
 s(J)=\{(a_n):q_k(a):=\sup_n(2+|n|)^k\|a_n\|_n<\infty
                                  \text{ for every }k\ge0\}      \tag{M18.2}
\]
is a Fr\'echet space, and the Mellin block restriction gives a
topological isomorphism
\[
 Q\simeq s(J).                                                    \tag{M18.3}
\]
In particular
\[
 \overline{\psi*\mathcal E}
 =\{f\in\mathcal E:(\mathcal Mf)^{(j)}(\rho)=0
                         \text{ for every }\rho,\ 0\le j<m_\rho\}.
                                                                  \tag{M18.4}
\]
The block expansion in \(Q\) converges absolutely in every
continuous seminorm.  Choosing an Auerbach basis in each
\((J_n,\|\cdot\|_n)\) gives an absolute Schauder basis of \(Q\)
with block-adapted continuous coefficient functionals.

Neither (M18.3) nor this last basis assertion prescribes the
individual zero jets as its coordinate functionals.  Thus the
result does not assert the quantitative condition of MT-M16.
The block partition and its adapted bases are choices; the
closed quotient in (M18.3) is the canonical one.

\begin{proof}
\emph{Analytic inputs and the transform topology.}
NT-Z10 proves unconditionally, including for a fixed
infinite-order unitary character, the polynomial upper bound
for the finite \(L(s,\eta)\) on every fixed vertical strip and
the bound \(O(\log(2+|T|))\) for the number of nontrivial
zeros, with multiplicities, in any fixed-height window.
It also gives
\[
                   |L(2+iT,\eta)|\ge\zeta_K(2)^{-1}.              \tag{M18.5}
\]
The standard completed Hecke functional equation locates all
the zeros of the entire completion in the fixed critical strip.
Exact classical source locators are Ramakrishnan--Valenza,
\emph{Fourier Analysis on Number Fields}, \S7.4,
Theorem 7-19, printed p.~279, and \S7.7, Theorem 7-28,
printed p.~289; the fixed-strip and local-count deductions
are proved in NT-Z10.  We use these unconditional facts,
not a zero-location consequence of either GRH input.

Put \(\kappa=\pi[K:\mathbb Q]/4\).  Stirling's formula and
the preceding strip bound give, for every fixed \(A>0\),
\[
 |\Xi(\sigma+it)|\le C_A(2+|t|)^{B_A}e^{-\kappa|t|}
                         \quad(|\sigma|\le A).                  \tag{M18.6}
\]
Cauchy's formula gives the same type of bound for \(\Xi'\)
on a smaller strip.  For large \(|t|\), the factor
\(\Xi(s)/L(s,\eta)\) is nonzero and has upper and lower
bounds by \(e^{-\kappa|t|}\) times fixed positive or negative
powers of \(2+|t|\), uniformly on each fixed strip.  Its
finitely located gamma exceptions are immaterial at large
height.

For clarity, the transform fact used below is elementary.
If an entire function \(H\) decreases faster than every
power of \(|t|\), uniformly on each strip \(|\sigma|\le A\),
its inverse Fourier integral on \(\sigma=0\) is the Mellin
transform inverse of an element of \(\mathcal E\).
For \(x\ge0\), shift the inverse contour to
\(\sigma=N+1\), and for \(x\le0\) shift it to
\(\sigma=-N-1\).  The horizontal edges tend to zero by the
uniform strip decay.  The \(j\)th derivative is bounded
by the integral of \(|s|^j|H(s)|\) on the shifted line,
times \(e^{-(N+1)|x|}\); these integrals are finite.
This proves all \(p_{N,j}\) bounds and continuity from the
corresponding strip seminorms.  Fourier inversion proves
the transform identity.  Conversely, integration by parts
in \(\mathcal Mf(\sigma+it)=\int e^{(\sigma+it)x}f(x)\,dx\)
shows rapid vertical decay, with each strip seminorm
bounded by finitely many \(p_{N,j}\).  Entireness follows
by exponential domination.  In particular (M18.6) proves
\(\psi\in\mathcal E\) and \(\mathcal M\psi=\Xi\), without GRH.

\emph{Selecting all horizontal boundaries at once.}
For \(n\in\mathbb Z\), let \(M_n\) count, with multiplicities,
the zeros whose ordinates lie in \([n-2,n+3]\).  Thus
\(M_n\le C\log(2+|n|)\).  On
\(I_n^*=(n+1/4,n+3/4)\), set
\[
 U_n(h)=\sum_{\substack{\rho\in Z\\
                    \operatorname{Im}\rho\in[n-2,n+3]}}
       m_\rho\log^+\frac1{|h-\operatorname{Im}\rho|}.
\]
Every summand has integral at most
\(\int_{-1}^1-\log|u|\,du=2\) over \(I_n^*\).
Averaging over its length \(1/2\) therefore permits a
choice \(h_n\in I_n^*\), avoiding every zero ordinate,
with \(U_n(h_n)\le 5M_n\); if \(M_n=0\), use its midpoint.
(The harmless slack from \(4M_n\) to \(5M_n\) gives a
positive-measure set when \(M_n>0\).)
Every zero with \(|h_n-\operatorname{Im}\rho|<1\)
is included in this sum.  Consequently
\[
 \prod_{\rho\in Z}
   \min\{1,|h_n-\operatorname{Im}\rho|\}^{m_\rho}
      \ge e^{-5M_n}\ge c(2+|n|)^{-C}.                            \tag{M18.7}
\]
Only finitely many factors in this product differ from one.
Also \(1/2<h_{n+1}-h_n<3/2\), so the rectangles have
uniformly bounded perimeter and cover all zero ordinates.

We claim, for every fixed \(A>0\), that
\[
 \min_{|\sigma|\le A}|\Xi(\sigma+i h_n)|
       \ge c_A(2+|n|)^{-C_A}e^{-\kappa|n|}.                      \tag{M18.8}
\]
Here is a complete minimum-modulus argument.  For large
\(|h_n|\), center a disk at \(c_n=2+i h_n\) and take a
zero-free boundary of radius
\(R\in[4(A+3)+4,4(A+3)+5]\).  Such a radius exists because
there are only finitely many zeros in that compact annulus.
At sufficiently large height the disk contains no finite-\(L\)
pole or trivial zero, since their possible ordinates are fixed
by the infinity type.  Its total zero multiplicity is
\(M=O_A(\log(2+|n|))\), by NT-Z10.  Remove its zeros by the
finite disk Blaschke product \(B\), with multiplicities,
and write \(L=GB\).  The factor \(G\) is holomorphic and
zero-free on the disk and its boundary.  Since \(|B|=1\)
on the boundary, (M18.6)'s finite-\(L\) antecedent and the
maximum principle give \(|G|\le C(2+|n|)^C\) there and
inside.  At the center, \(|B(c_n)|\le1\), so (M18.5)
gives \(|G(c_n)|\ge c>0\).
Apply Harnack's inequality to the nonnegative harmonic
function \(\log(C(2+|n|)^C)-\log|G|\).  On the inner
half disk its value is at most three times its value
at the center.  Thus
\[
                        |G(s)|\ge c(2+|n|)^{-C}
\]
on that inner disk.  This form of Harnack follows by
bounding the ratios of the Poisson kernels on concentric
circles; the ratio is \((R+r)/(R-r)\le3\) for \(r\le R/2\).

For \(s=\sigma+i h_n\), \(|\sigma|\le A\), the normalized
coordinate \(z=(s-c_n)/R\) has \(|z|<1/4\).
For a disk zero \(\rho\), with \(a_\rho=(\rho-c_n)/R\),
\[
 \left|\frac{z-a_\rho}{1-\overline{a_\rho}z}\right|
    \ge\frac{|s-\rho|}{2R}
    \ge\frac{\min\{1,|h_n-\operatorname{Im}\rho|\}}{2R}.
\]
Products of omitted factors in (M18.7) are at most one.
It follows that
\[
 |B(s)|\ge(2R)^{-M}
   \prod_{\rho\in Z}
      \min\{1,|h_n-\operatorname{Im}\rho|\}^{m_\rho}
             \ge c_A(2+|n|)^{-C_A}.
\]
Multiplication by \(G\) and by the lower Stirling bound
for \(\Xi/L\) proves (M18.8) at large height.
There are only finitely many remaining \(n\); each
horizontal segment is compact and zero-free for \(\Xi\),
so reducing \(c_A>0\) covers them.

The same disk proof, now centered at \(2+it\), gives a
polynomial lower bound for \(L(\sigma+it,\eta)\) on either
fixed vertical line \(\sigma=-A\) or \(\sigma=A\), for
every \(A\ge2\) and large \(|t|\).  Indeed their distances
from all disk zeros in the critical strip are at least one,
so every normalized Blaschke factor is at least \(1/(2R)\).
There is no logarithmic-distance product to control.
The lower Stirling bound and the compact low-height range
give the corresponding lower bound for \(\Xi\).
The same argument applies to the line \(\sigma=-1\)
and to \(\sigma=2\).  Combining it with (M18.8) yields
\[
 \inf_{w\in\Gamma_n}|\Xi(w)|
       \ge c(2+|n|)^{-C}e^{-\kappa|n|}.                          \tag{M18.9}
\]
All constants here are independent of separations between
zeros inside a rectangle.

\emph{The finite block algebras and an entire interpolation kernel.}
All zeros of a block are interior points.  Each of its
finitely many jet functionals is continuous on
\(A(\Omega_n)\), by Cauchy's formula on a disk contained
in \(\Omega_n\).  The map to these jets is onto, by
ordinary polynomial Hermite interpolation.  Its kernel
is \(I_n\), a closed ideal, proving that (M18.1) is a
genuine finite-dimensional quotient norm, with dimension
\(d_n\).  The maximum principle identifies the supremum
over the closure with that over \(\Gamma_n\).
The quotient norm is submultiplicative because the
supremum norm is.  The local zero count gives (M18.2)'s
dimension bound.

Fix \(a>0\).  For \(F\in A(\Omega_n)\) define
\[
 H_nF(z)=\frac1{2\pi i}\int_{\Gamma_n}
    F(w)\frac{\Xi(z)-\Xi(w)}{\Xi(w)(z-w)}
                      e^{a(z-w)^2}\,dw .                       \tag{M18.10}
\]
The apparent singularity at \(z=w\) is removable.
Consequently this is an entire function of \(z\).
If \(z\) is inside the rectangle, Cauchy's formula
for the second term gives
\[
 H_nF(z)=F(z)+\Xi(z)\frac1{2\pi i}
             \int_{\Gamma_n}
               \frac{F(w)e^{a(z-w)^2}}{\Xi(w)(z-w)}\,dw .
                                                                  \tag{M18.11}
\]
The integral on the right is holomorphic in the interior.
If \(z\) is outside the closed rectangle, that second
Cauchy integral is zero and \(H_nF(z)\) is divisible
locally by \(\Xi(z)\).
These statements follow also for continuous boundary
values by applying the formulas on interior rectangles
and taking limits.  Thus \(H_nF\) matches every
prescribed block jet of \(F\), and vanishes to the
full multiplicity at every zero outside the block.

Moreover, if all the block jets of \(F\) vanish, then
\(F/\Xi\) is holomorphic in \(\Omega_n\) and continuous
on its closure.  For exterior \(z\), both terms in
(M18.10) then integrate to zero by Cauchy's theorem.
Entireness implies \(H_nF=0\) everywhere.  It follows
that \(H_n\) depends only on \(a_n=F\bmod I_n\).
Denote this entire function by \(H_n a_n\).

For every fixed \(A>0\) there are constants \(C_A,B_A'>0\)
and \(c_a>0\) such that
\[
 |H_n a_n(\sigma+it)|
       \le C_A(2+|n|)^{B_A'}\|a_n\|_n
                         e^{-c_a(t-n)^2},
                 \qquad |\sigma|\le A.                         \tag{M18.12}
\]
To check the estimate before taking the quotient infimum,
bound \(|F|\) by its boundary supremum in (M18.10).
If \(|z-w|\ge1\), use (M18.6) and (M18.9) to bound
the difference quotient by
\[
 C_A(2+|n|)^C(2+|t|)^C
                    e^{\kappa\,\lvert t-\operatorname{Im}w\rvert}
       +1 .
\]
Here \(|\operatorname{Im}w-n|<2\), while
\[
 |e^{a(z-w)^2}|
       \le C_A e^{-a(t-\operatorname{Im}w)^2}.
\]
A fixed polynomial and the fixed linear exponential
are absorbed by a smaller Gaussian, and
\(2+|t|\le(2+|n|)(2+|t-n|)\).
If \(|z-w|<1\), the identity
\[
 \frac{\Xi(z)-\Xi(w)}{z-w}
      =\int_0^1\Xi'(w+u(z-w))\,du
\]
and the strip derivative bound give a polynomial in
\(2+|n|\) after division by (M18.9).
In this region \(|t-n|<3\), so it too satisfies (M18.12)
after a constant adjustment.  The perimeter is bounded
independently of \(n\).  Taking the infimum over all
representatives of \(a_n\) proves (M18.12).

The transform fact above now gives a unique
\(\mathcal T_n(a_n)\in\mathcal E\) whose Mellin transform is
\(H_n a_n\).  It is linear in \(a_n\), and for every
defining seminorm
\[
 p_{N,j}(\mathcal T_n(a_n))
       \le C_{N,j}(2+|n|)^{B_{N,j}}\|a_n\|_n .                  \tag{M18.13}
\]
Indeed on either shifted vertical line, the integral
of \((1+|t|)^j\) against (M18.12)'s Gaussian costs
only another fixed power of \(2+|n|\).
In particular no individual reciprocal derivative
or Hermite interpolation matrix norm has entered
these estimates.

\emph{Spectral synthesis for the closed ideal.}
Write \(J\) for the right side of (M18.4).
It is a closed subspace, since each Mellin jet is
continuous.  Multiplication of Mellin transforms
immediately gives
\(\overline{\psi*\mathcal E}\subseteq J\).
For the reverse inclusion take \(f\in J\) and
\(F=\mathcal Mf\).  The quotient \(F/\Xi\) is entire:
every zero, including its full multiplicity, has
been cancelled.

Fix \(A\ge2\).  Use the wider rectangles
\([-A,A]+i[h_n,h_{n+1}]\).  On their horizontal
edges (M18.8) applies; on their vertical edges
the vertical-line argument following (M18.8)
applies.  The rapid fixed-strip decay of \(F\)
and these lower bounds show, on the whole boundary,
\[
 \left|\frac{F(s)}{\Xi(s)}\right|
       \le C_{A,f}(2+|n|)^{C_A}e^{\kappa|n|}.
\]
The maximum principle for the entire quotient
extends this bound through the rectangle.
Since their heights have gaps between \(1/2\)
and \(3/2\), it follows that on every fixed strip
\[
 \left|\frac{F(\sigma+it)}{\Xi(\sigma+it)}\right|
       \le C_{A,f}(2+|t|)^{C_A}e^{\kappa|t|}.                  \tag{M18.14}
\]
For every \(\epsilon>0\), the entire function
\[
                   e^{\epsilon s^2}F(s)/\Xi(s)
\]
therefore has rapid decay on every fixed strip.
It is the Mellin transform of some
\(j_\epsilon\in\mathcal E\).
With \(g_\epsilon(x)=(4\pi\epsilon)^{-1/2}
e^{-x^2/(4\epsilon)}\), Fourier uniqueness gives
\[
                         \psi*j_\epsilon=g_\epsilon*f.
\]
The right side tends to \(f\) in \(\mathcal E\).
For completeness, on \(|y|\le1\) the weighted
mean-value bound for
\(f^{(j)}(x-y)-f^{(j)}(x)\) is
\(|y|e^{N|y|}p_{N,j+1}(f)\).
On the complement of any fixed neighborhood of
zero use
\((e^{N|y|}+1)p_{N,j}(f)\).
Integration against \(g_\epsilon(y)\) makes both
bounds tend to zero, by its first moment and
Gaussian tails.  This proves the reverse inclusion
and hence (M18.4), independently of the interpolation
construction.

\emph{The topological decomposition and the adapted basis.}
The restriction map
\[
 R:\mathcal E\longrightarrow s(J),\qquad
        (Rf)_n=(\mathcal Mf)|_{\Omega_n}\bmod I_n
\]
is continuous.  Indeed the quotient norm is at most
\(\sup_{\Gamma_n}|\mathcal Mf|\), and these boundary
suprema decrease faster than every power of \(|n|\),
with each such bound controlled by finitely many
seminorms of \(f\).  Define
\[
                    S(a)=\sum_{n\in\mathbb Z}\mathcal T_n(a_n).
\]
For each seminorm \(p_{N,j}\), (M18.13) and
\(q_{B_{N,j}+2}(a)\), with the exponent rounded up
if necessary, bound the sum of its term seminorms
by a constant times \(\sum_n(2+|n|)^{-2}\).
Thus this series converges absolutely in every
seminorm of \(\mathcal E\), and defines a continuous
linear map \(S:s(J)\to\mathcal E\).
Completeness of \(s(J)\) follows directly from the
complete finite-dimensional norms in its weighted
supremum seminorms; completeness of \(\mathcal E\)
follows from convergence of all weighted derivative
seminorms on compact intervals.

Every Mellin jet of \(S(a)\) may be computed
termwise by continuity.  The interpolation property
of (M18.11) gives \(RS=\mathrm{id}\).
Also \(\ker R=J=\overline{\psi*\mathcal E}\).
Therefore \(R\) induces the continuous bijection
(M18.3), with continuous inverse
\(a\mapsto[S(a)]\).
The block operators \(\mathcal T_n R_n\) on \(\mathcal E\)
are finite-rank continuous projections, mutually
annihilating for distinct \(n\).
Their absolute sum is the continuous projection
\(SR\); on \(Q\) their sum is the identity.

Finally every finite-dimensional complex normed
space has an Auerbach basis \(v_{n,1},\ldots,v_{n,d_n}\):
its vectors and dual coordinate functionals
\(\phi_{n,j}\) all have norm one.
Here is the elementary existence argument.  In
any fixed coordinates maximize the absolute
determinant of \(d_n\) vectors in the compact unit
ball.  A maximum is positive.  Replacing its
\(j\)th column by any unit-ball vector multiplies
the determinant by the corresponding dual
coordinate.  Maximality makes that coordinate
have modulus at most one; its value on its own
column is one.  This proves the claimed norms.

For \(a=(a_n)\in s(J)\), its resulting scalar
coefficients satisfy
\(|\phi_{n,j}(a_n)|\le\|a_n\|_n\).
Writing \(e_{n,j}\) for \(v_{n,j}\) in its block
and zero elsewhere, one has
\[
 \sum_{n,j}q_k\bigl(\phi_{n,j}(a_n)e_{n,j}\bigr)
 \le\sum_n d_n(2+|n|)^k\|a_n\|_n
 \le q_{k+3}(a)\sum_n d_n(2+|n|)^{-3}<\infty .
\]
This proves an absolute expansion in every
seminorm, in any enumeration of these vectors.
The coordinate functionals are continuous and
unique by the block isomorphism and the finite
basis in each block.  Transport through (M18.3)
proves the asserted absolute Schauder basis.
\end{proof}

The spectral block \(J_n\), as an algebra, is the
finite-length complex algebra
\[
 \mathbb C[s]\Big/
       \prod_{\rho\in Z\cap\Omega_n}(s-\rho)^{m_\rho}
 \simeq
 \prod_{\rho\in Z\cap\Omega_n}
                  \mathbb C[z]/(z^{m_\rho}).
\]
Its stable norm is the quotient holomorphic supremum
norm in (M18.1), not a norm obtained by declaring its
individual jet coordinates orthonormal.
Thus the construction is a finite-length spectral
divisor decomposition compatible with the Mellin
convolution algebra.  It directly refines the repair
of Masullo's Lemma 9.3.2.5, original PDF pp.~425--426:
closure of the convolution ideal is essential, and
the adapted basis has different coordinates from
that lemma's prescribed jets.  The inherited programme
locators are cote~1's nuclear-trace framework and
SA23 geo-H-lidskii/geo-H-genus, as carried in the
attached returns; no new inspection of original
cote scans is claimed here.  Nuclearity and the
trace action of these finite spectral algebras
are treated separately in MT-M19.

No first-priority claim is made for spectral synthesis in related
function spaces; the explicit block norms and bounds above supply
the precise topology and interpolation statement used here.



\subsection*{MT-M19. Unconditional nuclear convolution and its exact positivity criterion}
Use the canonical quotient, stable block algebras \(J_n\), and
topological isomorphism \(Q\simeq s(J)\) of MT-M18. No GRH,
simplicity, or reciprocal-derivative bound is assumed initially.
For every \(f\in\mathcal E\), convolution \(C_f\) is nuclear on
\(Q\). There is a Hilbert completion \(\mathcal H_B\) obtained
by making an Auerbach basis in each block orthonormal, on which
every \(C_f\) is trace class, with
\[
 \operatorname{Tr}_{\mathcal H_B} C_f
       =\sum_{\rho\in Z}m_\rho\mathcal Mf(\rho).       \tag{M19.1}
\]
The series is absolutely convergent and the trace is independent
of all block and Auerbach choices. Nuclearity is uniform when
\(f\) ranges over a bounded subset of \(\mathcal E\).

Define \(h^\sharp(x)=e^{-x}\overline{h(-x)}\). Then
\[
 \operatorname{Tr}_{\mathcal H_B} C_{h*h^\sharp}\ge0
          \quad\hbox{for every }h\in\mathcal E
 \quad\Longleftrightarrow\quad
 \text{GRH for }L(s,\eta).                              \tag{M19.2}
\]
Under GRH the trace equals
\(\sum_\rho m_\rho|\mathcal Mh(\rho)|^2\).
This concerns the trace; the chosen block Hilbert norm need
not make the operator \(C_{h*h^\sharp}\) positive.

\begin{proof}
Put \(F=\mathcal Mf\) and \(b_n(f)=\sup_{w\in\Gamma_n}|F(w)|\).
These numbers decrease faster than every power of \(2+|n|\),
with each decay bound controlled by finitely many defining
seminorms of \(f\). This follows from the integration-by-parts
estimate in MT-M18 on the one fixed strip containing all
\(\Gamma_n\). Convolution preserves the closed ideal: the
identity \(\mathcal M(f*g)=\mathcal Mf\,\mathcal Mg\) holds
by absolute integrability, and convolution is continuous on
\(\mathcal E\). Explicitly its seminorms are bounded by
\[
 p_{N,j}(f*g)
 \le p_{N,j}(f)\int_{\mathbb R}e^{N|y|}|g(y)|\,dy
 \le 2p_{N,j}(f)p_{N+1,0}(g).
\]
On \(J_n\), the induced map \(B_n(f)\) is multiplication by
the restriction of \(F\). The quotient Banach-algebra norm
therefore gives
\[
                      \|B_n(f)\|\le b_n(f).            \tag{M19.3}
\]

Let \(v_{n,j},\phi_{n,j}\), \(1\le j\le d_n\), be the
Auerbach vectors and dual functionals of MT-M18, all of norm
one. Regard each vector as supported in its own block and
write \(\ell_{n,j}(a)=\phi_{n,j}(a_n)\). The family
\(\ell_{n,j}\) is equicontinuous, since
\(|\ell_{n,j}(a)|\le q_0(a)\). Set
\[
 \lambda_{n,j}=(2+|n|)^{-2},\qquad
 u_{n,j}=(2+|n|)^2 B_n(f)v_{n,j}.
\]
The sum of \(|\lambda_{n,j}|\) is finite because
\(d_n=O(\log(2+|n|))\). For each \(k\),
\[
 q_k(u_{n,j})\le(2+|n|)^{k+2}b_n(f),
\]
so \(u_{n,j}\) is a bounded family in \(s(J)\). The series
\[
 C_f a=\sum_{n,j}\lambda_{n,j}\ell_{n,j}(a)u_{n,j}       \tag{M19.4}
\]
is thus an explicit nuclear decomposition, absolutely convergent
in every seminorm. Its value is the block multiplication map
because the finite Auerbach expansion in each block is exact.
The same estimates are uniform on a bounded set of \(f\)'s,
which proves the asserted uniformity.

Give the scalar Auerbach coordinates their ordinary square-sum
norm, and let \(\mathcal H_B\) be its completion. For each
nonempty block this Euclidean norm lies between
\(d_n^{-1/2}\|a_n\|_n\) and
\(d_n^{1/2}\|a_n\|_n\), since the vectors and coordinate
functionals have norm one. Consequently the inclusion
\(s(J)\to\mathcal H_B\) is continuous and injective, and
has dense image because it contains all finitely supported
block vectors. Every matrix entry of \(B_n(f)\) in these
orthonormal coordinates has modulus at most \(b_n(f)\),
by (M19.3). Expressing it as a sum of rank-one matrix units
therefore gives
\[
       \|B_n(f)\|_1\le d_n^2b_n(f),\qquad
       \sum_n\|B_n(f)\|_1<\infty.                     \tag{M19.5}
\]
The block direct sum is trace class and extends \(C_f\).

Algebraically, polynomial Hermite interpolation identifies
\(J_n\) with
\(\prod_{\rho\in Z\cap\Omega_n}\mathbb C[z]/(z^{m_\rho})\).
In the monomial basis of a local factor, multiplication by
\(F(\rho+z)\) is triangular with its constant term
\(F(\rho)\) on every diagonal entry. Its finite-dimensional
trace is \(m_\rho F(\rho)\). Trace is invariant under a
change of basis, so (M19.5) proves (M19.1). It also proves
the choice-independence stated there. Absolute convergence
follows directly from
\(\sum_n d_n b_n(f)<\infty\). All traces in the statement
are the ordinary Hilbert traces just constructed; no
unjustified trace on an arbitrary locally convex space is used.

Direct substitution gives
\[
 \mathcal Mh^\sharp(s)
       =\overline{\mathcal Mh(1-\bar s)}.              \tag{M19.6}
\]
The unitary Hecke functional equation implies that the zero
involution \(\rho\mapsto\rho^\dagger=1-\bar\rho\)
preserves \(Z\) and its multiplicities. Indeed it relates
\(\Lambda(s,\eta)\) to \(\Lambda(1-s,\bar\eta)\),
while complex conjugation relates the latter character to
\(\eta\). The pole-removal polynomial introduces no zeros
into \(Z\). Formula (M19.6) therefore preserves the full
zero-jet ideal, so \(\sharp\) descends to an antilinear
algebra involution of \(Q\). If GRH holds then \(\rho^\dagger=\rho\),
and (M19.1), (M19.6) give the stated nonnegative sum.

Conversely, suppose \(\rho\ne\rho^\dagger\) is a pair
off the critical line, of common multiplicity \(m\).
They have the same ordinate and hence lie in the same
block. In that block prescribe the constant jet 1 at
\(\rho\), the constant jet \(-1\) at \(\rho^\dagger\),
and the zero jet at every other zero; prescribe zero in
all other blocks. This is a finitely supported element
of \(s(J)\). The right inverse in MT-M18 realizes it as
\(\mathcal Mh\) for some \(h\in\mathcal E\). Formula
(M19.1) then has just two nonzero contributions for
\(h*h^\sharp\), both equal to \(-m\). Its trace is
\(-2m<0\), contradicting the left side of (M19.2).
This proves the equivalence.
\end{proof}

This is a precise zero-side Weil positivity criterion in
the canonical quotient, with an explicit unconditional
nuclear and trace-class realization. It supplies no proof
of positivity independent of the zeros. The trace is the
sum of traces on finite-length spectral algebras; multiple
zeros contribute their lengths. In the programme of
cote~1 and SA23 geo-H-lidskii/geo-H-genus, MT-M21--MT-M22 identify the specified Meyer adelic summation
quotient with this model, and MT-M24 identifies the actual
Masullo smooth quotient. MT-M25 treats the actual line Hilbert
completion and its lost jets. A cohomological realization and
a positivity argument defined independently of the zeros remain separate. Adopted \(A_D\)
supplies (M19.2)'s positive conclusion for Dirichlet
functions; \(M_H\) supplies it conditionally in general.

\subsection*{MT-M20. A weaker positive Hilbert completion and its exact obstruction}
Keep the unconditional canonical quotient of MT-M18.
If all nontrivial zeros are simple, individual Mellin
evaluation induces a continuous injective map with dense
image
\[
       Q\longrightarrow\ell^2(Z),\qquad
       [g]\longmapsto(\mathcal Mg(\rho))_{\rho\in Z}.   \tag{M20.1}
\]
It therefore defines a positive-definite continuous Hilbert
norm on \(Q\). Every convolution extends diagonally to a
trace-class operator on this completion. Under GRH,
\(C_{h^\sharp}=C_h^*\), and \(C_{h*h^\sharp}\) is a
positive operator. No reciprocal-derivative bound is
needed for this weaker Hilbert norm.

More precisely, the following are equivalent:
\begin{enumerate}
\item GRH holds for \(L(s,\eta)\), and every nontrivial
zero is simple.
\item There is a positive-definite continuous inner
product on \(Q\) such that every convolution extends to
a bounded operator on its Hilbert completion and every
\(C_{h*h^\sharp}\) is positive there.
\end{enumerate}
The inner product is required to be faithful on the
\emph{full jet quotient}. Semisimplifying its multiple-zero
local algebras changes this requirement.

\begin{proof}
Point evaluation at any zero in \(\Omega_n\) has norm at
most one on the quotient \(J_n\), by the maximum principle
and the definition of its quotient norm. Thus
\[
 \sum_{\rho\in Z}|\mathcal Mg(\rho)|^2
 \le\sum_n d_n\|(Rg)_n\|_n^2
 \le q_2(Rg)^2\sum_n d_n(2+|n|)^{-4}<\infty.           \tag{M20.2}
\]
This proves continuity. When all zeros are simple, zero
values mean all zero jets, and MT-M18 identifies the kernel
exactly with the closed ideal; hence the map is injective.
Every finitely supported zero-value vector lies in its
image, since it specifies elements of finitely many
\(J_n\)'s and the interpolation right inverse realizes
them. This proves density and (M20.1).

Convolution by \(f\) has diagonal entries
\(\mathcal Mf(\rho)\). Their absolute sum is finite by
the local zero count and rapid strip decay, so its extension
is trace class. Under GRH (M19.6) makes the entries for
\(h^\sharp\) the complex conjugates of those for \(h\).
The adjoint and positivity statements follow. This proves
the implication (1) to (2).

For the converse, finite-support elements of the block
algebras are actual elements of \(Q\), by MT-M18, and
convolution realizes their algebra multiplication. If
\(\rho\ne\rho^\dagger\), let \(e,e^\dagger\) be the
local idempotents equal to 1 in the full local algebras
at those two zeros and zero in every other local algebra.
The involution interchanges them. For \(h=e-e^\dagger\)
one has
\[
                  h*h^\sharp=-(e+e^\dagger).
\]
On the nonzero finite-dimensional invariant subspace
\((e+e^\dagger)Q\), convolution by this element is
\(-\mathrm{id}\). Its quadratic form is strictly negative
on every nonzero vector in any positive-definite inner
product. This contradicts (2), and proves GRH.

Now suppose a critical-line zero has multiplicity
\(m\ge2\). In its local algebra
\(\mathbb C[z]/(z^m)\), the involution is
\(z^\sharp=-z\), by (M19.6). Let \(e\) be that local
idempotent and \(n\) the element represented by \(z\),
both zero in all other local factors. They are realized
by elements of \(Q\). On the finite-dimensional invariant
space \(eQ\), write \(N=C_n\). It is a nonzero nilpotent
operator. For real \(u\), the element \(h_u=e+iun\)
satisfies \(h_u^\sharp=h_u\), and
\[
 C_{h_u*h_u^\sharp}|_{eQ}
                     =I+2iuN-u^2N^2.                \tag{M20.3}
\]
By (2) these operators are positive, in particular
self-adjoint, for both \(u=1\) and \(u=-1\).
Subtracting them shows that \(iN\) is self-adjoint
in the restricted positive-definite inner product.
A self-adjoint nilpotent on a finite-dimensional
inner-product space is zero: diagonalize the Hermitian
matrix, whose real eigenvalues must all vanish.
This contradicts \(N\ne0\). Thus all zeros are simple,
completing the converse.
\end{proof}

The Hilbert norm in (M20.1) can be strictly weaker than
the original Fr\'echet topology. Its existence therefore
does not make point evaluations a Schauder coordinate
system for that topology: MT-M16 still gives the exact
extra leading-coefficient condition for that assertion.
With multiple zeros, evaluation kills the higher-jet
nilpotent directions. A semisimple Hilbert model with
\(m_\rho\) copies of each scalar eigenvalue retains
the multiplicity-weighted trace, but is not a faithful
Hilbert completion of this full jet quotient. No
simplicity consequence of \(A_D\) or \(M_H\) is claimed.


\subsection*{MT-M21. Identification with Meyer's zeta-summation quotient over \(\mathbb Q\)}
For the Riemann zeta function, the canonical quotient in MT-M18
is naturally isomorphic, as a topological convolution module,
to the zero-space in Meyer's classical spectral realization.
This identification is unconditional and includes all zero
multiplicities.

Here are the precise source and conventions. In Ralf Meyer,
\emph{A spectral interpretation for the zeros of the Riemann
zeta function}, G\"ottingen Seminars 2004--05, pp.~117--137,
let \(H_- =\mathcal O(\mathbb R_+^\times)\) be the strong
Schwartz space of Section~2, let \(H_+\) be the even Schwartz
space on \(\mathbb R\), and let
\[
 H_\cap=\{u\in H_+:u(0)=\widehat u(0)=0\},\qquad
 (Zu)(x)=\sum_{k\ge1}u(kx).
\]
His Theorem~3.3, printed p.~124 (PDF p.~8), proves the
closed-image assertion for the zeta summation map; Section~4,
printed pp.~125--126, defines
\(H_-^0=H_-/ZH_\cap\).
Theorem~4.1, printed p.~126 (PDF p.~10), characterizes its
Mellin-transform kernel by two half-plane division conditions.
Under the logarithmic coordinate map
\[
       \mathcal L:\mathcal E\longrightarrow H_-,
       (\mathcal Lf)(x)=f(\log x),
\]
one has the exact equality of closed subspaces
\[
          \mathcal L\bigl(\overline{\psi_\zeta*\mathcal E}\bigr)
                         = ZH_\cap.                  \tag{M21.1}
\]
Consequently \(Q_\zeta\simeq H_-^0\), and the adapted
absolute basis, nuclear convolution, and trace constructions
of MT-M18--MT-M20 apply to this independently defined
zeta-summation quotient with their stated positivity hypotheses.

\begin{proof}
The map \(\mathcal L\) is a topological isomorphism:
the invariant derivatives \((x\,d/dx)^j\) become
\(d^j/dt^j\) for \(t=\log x\), and requiring rapid
decay after every real power weight \(x^\sigma\)
is equivalent to the defining seminorms
\(\sup_t e^{N|t|}|f^{(j)}(t)|\). Polynomial factors in
\(t\), if included in the Schwartz definition, are
absorbed by increasing \(N\). Its Mellin transform is
exactly \(\mathcal Mf(s)\), since \(d^\times x=dt\).
Multiplicative convolution becomes additive convolution.

Let \(F=\mathcal Mf\). Meyer's Theorem~4.1 says that
\(\mathcal Lf\in ZH_\cap\) precisely when
\[
 t\longmapsto F(\sigma+it)/\zeta(\sigma+it)
          \text{ is Schwartz for every }\sigma\ge\tfrac12,
                                                               \tag{M21.2}
\]
and
\[
 t\longmapsto F(\sigma+it)/\zeta(1-\sigma-it)
          \text{ is Schwartz for every }\sigma\le\tfrac12.
                                                               \tag{M21.3}
\]
The conditions include holomorphic removal of the corresponding
zeros on those lines. If they hold, every nontrivial zero
\(\rho\) of real part at least \(1/2\) is a zero of \(F\)
to at least its full multiplicity, by (M21.2). For zeros in
the other half-plane, the zeta functional equation says that
\(\zeta(1-s)\) has the same nontrivial zero divisor as
\(\zeta(s)\); condition (M21.3) gives the same conclusion.
Thus \(f\) belongs to the full zero-jet kernel in MT-M18.

Conversely, suppose \(F\) vanishes on this entire divisor.
The function \(F/\zeta\) is holomorphic on
\(\operatorname{Re}s>-2\). All nontrivial zero singularities
are removed, the pole of \(\zeta\) at 1 gives a removable
zero of its reciprocal, and its first trivial zero is at
\(-2\), outside this open half-plane.

Fix \(A\ge2\). Use rectangles
\([-1,A+1]+i[h_n,h_{n+1}]\), with the horizontal levels
constructed in MT-M18 for \(\zeta\). The finite-\(L\)
Blaschke estimates in that proof give a fixed inverse
polynomial lower bound for \(|\zeta|\) on the horizontal
edges at large height. The two vertical edges are a fixed
positive distance from every nontrivial zero, so the same
disk argument gives their polynomial lower bound. They
also avoid all trivial zeros. The finitely many low-height
boundaries cause no problem: the levels \(h_n\) are never
zero, and both vertical lines are zero-free.

For every \(N\), rapid fixed-strip decay of \(F\)
therefore bounds \(|F/\zeta|\) on these boundaries by
\(C_{A,N,f}(2+|n|)^{-N}\), after increasing the decay
power used for \(F\). The maximum principle for the
holomorphic quotient gives the same bound inside each
rectangle. Since adjacent heights have uniformly bounded
gaps, \(F/\zeta\) decreases rapidly on the whole strip
\([-1,A+1]\). Cauchy's formula on small fixed disks
inside that strip gives rapid decay of all its vertical
derivatives on \([1/2,A]\). This proves (M21.2).

Apply this same argument to \(F^\vee(s)=F(1-s)\).
It is an entire strip-Schwartz function and vanishes on
the full nontrivial divisor, because that divisor is
invariant under \(s\mapsto1-s\). The result for
\(F^\vee/\zeta\), followed by \(s\mapsto1-s\),
is exactly (M21.3). Meyer's image criterion now gives
\(\mathcal Lf\in ZH_\cap\).

We have proved that \(\mathcal L^{-1}(ZH_\cap)\)
is precisely the full zero-jet kernel. MT-M18 identifies
that kernel with \(\overline{\psi_\zeta*\mathcal E}\),
proving (M21.1). Both are closed subspaces of the
topologically identified test-function spaces, so their
quotients, with quotient topology, are naturally
topologically isomorphic. The logarithmic identification
commutes with scaling and integrated convolution, giving
the stated module assertion.
\end{proof}

This settles a specified arithmetic comparison map, beyond
a zero-built abstract model. It is a comparison with
Meyer's Poisson/zeta-summation construction, not an
identification with an unspecified adelic quotient in
Masullo's draft or with an \(\ell\)-adic cohomology group.
The classical spectral realization itself is Meyer's
result; the proof here supplies the explicit closed-ideal
and stable-block crosswalk used in this investigation.
The full number-field idele-class construction is located
in Meyer, \emph{On a representation of the idele class
group related to primes and zeros of L-functions},
Duke Math. J.~127 (2005), 519--595, author version
\texttt{math/0311468v3}, Theorem~5.8 (pp.~43--44) and
Theorem~5.11 (p.~47). A sectorwise topological comparison
with that construction requires its precise image
criterion and is not inferred merely from equal spectra.


\subsection*{MT-M22. The fixed Hecke sector of Meyer's adelic zero-space}

For every fixed primitive unitary Hecke \(L(s,\eta)\), the
canonical quotient \(Q_\eta\) of MT-M18 is isomorphic to the
corresponding fixed compact-character sector of Meyer's
idele-class zero-space.  The isomorphism is unconditional
and compatible with convolution, including all zero
multiplicities.  The sector is first given its quotient
Fr\'echet topology from \(H_-^\chi\); the proof also
identifies the quotient bornology by an explicit splitting.

We fix the character convention explicitly.  Write
\[
 L(s,\eta)=\prod_{\mathfrak p}
       (1-c_{\mathfrak p}(N\mathfrak p)^{-s})^{-1},
\]
where \(c_{\mathfrak p}=0\) at a ramified place.  Let
\(\beta\) be the associated unitary idele-class character
in Meyer's convention: for his \(S\)-prime element
\(p_v\in C_S\), at every unramified \(v\notin S\),
\[
                      \beta(p_v)^{-1}=c_{\mathfrak p_v}.
                                                               \tag{M22.1}
\]
Thus the component is parameterized by
\(\omega_s=\beta|\cdot|^s\).
This convention is dictated by the inverse in Meyer's
Euler factor, rather than by a choice of the name
\(\eta\) or \(\eta^{-1}\) for its idele realization.
For example, \(L(s,\eta)=\zeta_K(s+i\tau)\) corresponds
to \(\beta=|\cdot|^{i\tau}\).
If an idele character has instead been labeled by
its direct values \(c_{\mathfrak p_v}\) on Meyer's
\(p_v\), it must be inverted to obtain \(\beta\).

Let \(C=C_K\), \(C^1=\ker|\cdot|\), and choose a
continuous smooth splitting \(a:\mathbb R\to C\)
with \(|a(t)|=e^t\).
Normalize Haar measure on \(C^1\) to one, so
\(d^\times x=dc\,dt\) in \(x=ca(t)\).
Put \(\chi=\beta|_{C^1}\).
Use Meyer's spaces and embeddings
\[
 H_-=\mathcal S(C)_{\mathbb R},\qquad
 H_+=\mathcal S(\mathbb A_K)/K^\times,\qquad
 H_-^0=H_-/(H_+\cap H_-),
\]
where the intersection uses his embeddings into
\(\mathcal S(C)_{]1,\infty[}\oplus
  \mathcal S(C)_{]-\infty,0[}\).
Superscript \(\chi\) denotes the subspace on which
left translation by \(c\in C^1\) acts by \(\chi(c)\).
Then
\[
 (\mathcal L_\beta f)(x)=\beta(x)^{-1}f(\log|x|)
\]
induces an isomorphism
\[
              Q_\eta \ \simeq\ (H_-^0)^\chi.                  \tag{M22.2}
\]
It is topological for the specified quotient topology.
Norm translation by \(t\) corresponds to idele
translation by \(a(t)\), with the explicit unitary
factor \(\beta(a(t))\).

The primary source is Ralf Meyer, \emph{On a representation
of the idele class group related to primes and zeros of
\(L\)-functions}, Duke Math. J.~127 (2005), 519--595,
author version \texttt{math/0311468v3}.
The exact input is the weighted spaces in Definition 4.1
and (22), author pp.~24--25; the embeddings in Lemma 5.4,
pp.~40--41; and the image criterion in \S5.7, pp.~46--47,
following Lemma 5.10.  Theorem 5.1, pp.~38--39, supplies
the finite-level description used there.
Theorems 5.8, pp.~43--44, and 5.11, p.~47, are the
classical summability and spectral-multiplicity statements.
The comparison below uses the image criterion,
rather than inferring an isomorphism from equal spectra.

\begin{proof}
\emph{The exact test-function sector.}
A norm splitting exists by varying a positive component
at one archimedean place with the appropriate local
norm exponent.  In the resulting product
\(C=C^1\times a(\mathbb R)\), every \(\chi\)-homogeneous
function has the form \(\chi(c)^{-1}u(t)\).
Since \(\beta\) is unitary,
\(\beta(a(t))=e^{i\tau t}\) for a fixed real \(\tau\).
Multiplication by this fixed modulation preserves
every exponential-Schwartz seminorm, in both
directions, by the Leibniz formula.

Consequently \(\mathcal L_\beta\) is a topological
isomorphism from \(\mathcal E\) onto \(H_-^\chi\).
Indeed the norm-direction invariant derivative becomes
\(d/dt\); derivatives in the compact direction act
on its one fixed smooth character by constants.
All norm weights are \(e^{\alpha t}\), and the
requirement for every real \(\alpha\) is equivalent
to the seminorms \(e^{N|t|}|f^{(j)}(t)|\).
Polynomial factors in \(t\) are absorbed by increasing
\(N\).  The character is trivial on a suitable
compact subgroup with elementary compact quotient,
so these functions occupy one fixed Bruhat level.
There is no limit over varying compact characters.

Haar normalization as in Meyer's \S3.1, p.~16,
gives exactly
\[
 \widehat{\mathcal L_\beta f}(\omega_s)
       =\int_C\beta(x)^{-1}f(\log|x|)
                      \beta(x)|x|^s\,d^\times x
       =\mathcal Mf(s).                                        \tag{M22.3}
\]
The transform is zero on all other compact-character
components by orthogonality on \(C^1\).
A direct convolution calculation, using
\(\beta(y)^{-1}\beta(y^{-1}x)^{-1}=\beta(x)^{-1}\),
shows
\[
  (\mathcal L_\beta f)*(\mathcal L_\beta g)
                         =\mathcal L_\beta(f*g).
\]
Also
\[
 \lambda_b\mathcal L_\beta f
   =\beta(b)\mathcal L_\beta
         \bigl(f(\,\cdot-\log|b|)\bigr),\qquad b\in C.            \tag{M22.4}
\]
These identities specify the representation convention.

\emph{Meyer's image condition on this component.}
Choose once and for all a sufficiently large finite
set \(S\), containing the archimedean places, all
ramification of \(\beta\), and the places needed for
his normalized additive character and differential
idele.  His multiplier \(\mathcal L_S\) has transform
\[
 \widehat{\mathcal L_S}(\omega)
       =\prod_{v\notin S}(1-\omega(p_v)^{-1})^{-1}.
\]
By (M22.1), on \(\omega_s\) this is the partial finite
Hecke function
\[
 L_\eta^S(s)
       =L(s,\eta)D_{\eta,S}(s),\qquad
 D_{\eta,S}(s)=
        \prod_{v\in S_{\mathrm{fin}}}
             (1-c_{\mathfrak p_v}(N\mathfrak p_v)^{-s}).         \tag{M22.5}
\]
A ramified factor in \(D_{\eta,S}\) is one.
The reflected quasi-character is
\[
 \omega_s^*=|\cdot|\omega_s^{-1}
                    =\beta^{-1}|\cdot|^{1-s},
\]
and therefore
\[
              \widehat{\mathcal L_S}(\omega_s^*)
                       =L_{\bar\eta}^S(1-s).                   \tag{M22.6}
\]
Here \(\bar\eta\) denotes the dual unitary Hecke
character, with conjugate Euler coefficients.

For an entire strip-Schwartz function \(F(s)\)
on this component, the equivalent reduced
image criterion printed on Meyer p.~47 is
\[
 \begin{array}{ll}
 F(\alpha+it)/L_\eta^S(\alpha+it)
        \text{ is Schwartz in }t,&0<\alpha\le1,\\[2mm]
 F(it)/L_{\bar\eta}^S(1-it)
        \text{ is Schwartz in }t.&
 \end{array}                                                    \tag{M22.7}
\]
All apparent zero singularities must be removable.
The source's remaining conditions for
\(\alpha>1\) and \(\alpha<0\) follow from \(F\)
being the transform of \(H_-\).
Zero extension to other compact-character
components causes no additional Schwartz condition.
Thus (M22.7) characterizes membership of
\(\mathcal L_\beta f\) in \(\Sigma(H_+)\),
where \(F=\mathcal Mf\).

This is also the relevant intersection in
the quotient, rather than only its first-coordinate
image.  Meyer's Poisson formula gives
\[
 \Sigma u-J\Sigma\mathfrak F^*u
                         =c_0-c_1|\cdot|^{-1}.
\]
The second term lies in
\(\mathcal S(C)_{]-\infty,0[}\).
If \(\Sigma u\in H_-\), so does the right side.
Its decay as \(|x|\to0\) forces \(c_0=c_1=0\):
these distinct powers cannot decay after every
negative weight.  Hence
\(i_+u=i_-(\Sigma u)\).
Conversely, an element of the intersection has
its first coordinate in \(\Sigma(H_+)\cap H_-\).
Compact-character projection commutes with the
maps.  Therefore (M22.7) characterizes
\((H_+\cap H_-)^\chi\).

\emph{Removing the image criterion's growth gate.}
If (M22.7) holds, every zero of \(\Xi_\eta\) in
\(0<\operatorname{Re}s\le1\) is a zero of \(F\)
with at least its full multiplicity, by its first
condition.  The omitted finite Euler factors have
no zeros in that right half-plane, and the gamma
factors have neither zeros nor poles there.
The possible pure-norm pole is cancelled in
\(\Xi_\eta\) and is not a zero of it.
If a nontrivial zero is on \(\operatorname{Re}s=0\),
the functional equation identifies its multiplicity
with that of \(L_{\bar\eta}(1-s)\); the second
condition gives the same conclusion.
Thus membership implies all the zero-jet
vanishings of MT-M18.

Conversely suppose \(F=\mathcal Mf\) vanishes
on the full divisor of \(\Xi_\eta\).
Set \(G=F/\Xi_\eta\), which is entire.
The rectangle proof of (M18.14) gives,
for every \(A,N>0\), the sharpened estimate
\[
 \sup_{|\sigma|\le A}|G(\sigma+it)|
       \le C_{A,N,f}e^{\kappa|t|}(2+|t|)^{-N},
 \qquad \kappa=\pi[K:\mathbb Q]/4.                              \tag{M22.8}
\]
Indeed the boundary lower bound for \(\Xi_\eta\)
costs \(e^{\kappa|n|}(2+|n|)^{B_A}\).
Use the decay power \(N+B_A\), rounded up and
increased if needed, for \(F\) on that boundary.
The maximum principle and bounded block heights
give (M22.8), with no zero-spacing hypothesis.

Write the standard completion as
\[
 \Xi_\eta(s)=P_\eta(s)q_\eta^{s/2}
                  \Gamma_\eta(s)L(s,\eta).
\]
On each closed strip in \(\operatorname{Re}s>0\),
the unitary archimedean gamma product is
holomorphic and bounded above by
\(e^{-\kappa|t|}\) times a fixed power
of \(2+|t|\), by Stirling's formula.
Consequently
\[
        F(s)/L(s,\eta)
          =G(s)P_\eta(s)q_\eta^{s/2}\Gamma_\eta(s)
\]
is holomorphic and rapidly decreasing there,
by (M22.8).  The reciprocal of a finite-\(L\)
pole is a zero, so pure-norm poles cause
no exception.

The completed functional equation gives
\[
 \Xi_\eta(s)
   =\epsilon_\eta P_\eta(s)q_\eta^{(1-s)/2}
            \Gamma_{\bar\eta}(1-s)L(1-s,\bar\eta),
                    \qquad|\epsilon_\eta|=1.
\]
It follows that \(F(s)/L(1-s,\bar\eta)\)
is holomorphic and rapidly decreasing
on each closed strip in
\(\operatorname{Re}s<1\).
The reflected gamma arguments have positive
real part there.  Thus neither half-plane
division asks \(F\) to cancel trivial
finite-\(L\) zeros outside the appropriate half.

For \(\delta>0\), every nontrivial omitted
Euler factor satisfies
\[
 |1-c_{\mathfrak p}(N\mathfrak p)^{-s}|
                         \ge1-(N\mathfrak p)^{-\delta}
       \quad(\operatorname{Re}s\ge\delta).
\]
The ramified factors are one.
Thus \(D_{\eta,S}^{-1}\) is bounded on each
such right half-plane, as is
\(D_{\bar\eta,S}(1-s)^{-1}\) when
\(\operatorname{Re}s\le1-\delta\).
Their derivatives are bounded on smaller
fixed strips.  Therefore
\[
 \begin{split}
 F/L_\eta^S&\text{ is strip-Schwartz on }\operatorname{Re}s>0,\\
 F/L_{\bar\eta}^S(1-s)&\text{ is strip-Schwartz on }
                                      \operatorname{Re}s<1.
 \end{split}                                                    \tag{M22.9}
\]
Cauchy's formula on slightly larger strips
gives every vertical derivative estimate.
These statements imply (M22.7).
We have proved
\[
 \mathcal L_\beta^{-1}\bigl((H_+\cap H_-)^\chi\bigr)
   =\{f:(\mathcal Mf)^{(j)}(\rho)=0,\
          \Xi_\eta(\rho)=0,\ j<m_\rho\}
   =\overline{\psi_\eta*\mathcal E},                             \tag{M22.10}
\]
where the last equality is MT-M18.

\emph{The precise quotient topology and bounded lifting.}
Give \(H_-^\chi\simeq\mathcal E\) its
just-identified Fr\'echet topology and its
bounded-set bornology.  At this fixed character
these are precisely the source's weighted
Schwartz structure: it lies at one Bruhat
level, and the countable family of real
weight endpoints in (22) generates its
seminorms and bounded sets.
Here bounded and precompact sets coincide,
so the convention of using the precompact
bornology gives the same result.  Indeed,
for a bounded sequence in \(\mathcal E\),
the bounds on one additional derivative
give equicontinuity of every derivative
on each compact interval.  Diagonal
Arzel\`a--Ascoli gives a subsequence
converging with every derivative on every
compact interval.  For a prescribed
\(p_{N,j}\), the uniform \(p_{N+1,j}\)
bound makes the weighted tails outside
\([-R,R]\) at most \(Ce^{-R}\).
The subsequence is therefore Cauchy in
every defining seminorm and converges
in the complete space \(\mathcal E\).
This proves relative compactness of every
bounded set in its metrizable topology;
the converse is immediate.
The subspace in (M22.10) is closed.
Hence (M22.10) immediately gives a
topological isomorphism from \(Q_\eta\)
onto the sector quotient furnished with
the quotient topology from \(H_-^\chi\).

Here is why the original quotient bornology
agrees as well.  The maps \(i_+\) and \(i_-\)
are bornological embeddings by the cited
source.  Therefore the bornology of their
intersection equals the subspace bornology
inherited from \(H_-^\chi\): a subset of the
intersection is bounded in either embedded
space exactly when it is bounded in the
common ambient space.  No topology on
\(H_+^\chi\), or on its quotient, is being
assumed in this argument.

Compact-character averaging is a bounded
projection preserving \(H_+\) and \(H_-\).
At each finite level it is ordinary integration
over a compact group against its smooth
character, and these projections agree under
the level inclusions.
It preserves their intersection and splits
both spaces into its image and kernel.
The algebraic quotient and its quotient
bornology therefore split accordingly.
In particular \((H_-^0)^\chi\) is the
bornological quotient
\(H_-^\chi/(H_+\cap H_-)^\chi\);
representatives can be projected before
passing to the quotient.

Finally the continuous linear section
\(S:s(J)\to\mathcal E\) in MT-M18 supplies
an explicit lifting map for this quotient,
after transport by \(\mathcal L_\beta\).
Every bounded set in the quotient
Fr\'echet space lifts to a bounded set
through this section.  Conversely the
quotient map is continuous, so it takes
bounded sets to bounded sets.
The same two assertions hold for
precompact sets, using continuity of
the section and quotient map.
The quotient itself has bounded sets
precompact, since its continuous section
maps them to bounded, hence precompact,
subsets of \(\mathcal E\), and its
quotient map sends these back to the
original sets.
Thus the source's quotient bornology
is exactly the bornology of the specified
Fr\'echet quotient.  This is an explicit
splitting argument, not an inference
from equality of spectra.
Together with (M22.3)--(M22.4) and
multiplicativity, it proves all parts
of (M22.2).
\end{proof}

The stable blocks and nuclear traces in
MT-M18--MT-M19 therefore apply to this
independently defined adelic summation
quotient in each fixed Hecke sector.
All constants may depend on the character
and its infinity type; no uniform
decomposition over all characters is
asserted.  The trace positivity and
faithful-operator distinctions of
MT-M19--MT-M20 retain their exact GRH
and multiplicity hypotheses.
This identifies Meyer's specified
quotient, rather than an unspecified
quotient in Masullo's draft.
Its geometric location is the norm
direction of the idele-class action,
Tate's Poisson summation, and cote~1's
nuclear-trace framework as carried in
SA23 geo-H-lidskii/geo-H-genus.
The primary checks cover the named
normalization, image, and embedding
passages; no complete-paper audit
is claimed.


\subsection*{MT-M23. Every nonabelian Galois top field obstructs weak Mertens and faithful positivity}

Let $L/K$ be a nonabelian Galois extension of number fields. The
classical theorem of Hu--Kaneko--Martin--Schildkraut (HKMS),
\emph{Order of zeros of Dedekind zeta functions}, Theorem1.1,
author version2107.03269v2 p.2, proves unconditionally that
$\zeta_L$ has infinitely many nontrivial zeros of multiplicity
at least two. The complete first proof is \S3 p.5; a second
proof occupies \S4 pp.5--6. Their Theorem1.2, p.2 with proof
\S5 pp.6--8, gives infinitely many zeros of multiplicity at
least three whenever $\operatorname{Gal}(L/K)$ has an irreducible
representation of degree at least three. These are cited classical
results; no new multiplicity theorem is claimed here.

Set
\[
 M_L(x)=\sum_{N\mathfrak a\le x}\mu_L(\mathfrak a),\qquad
 I_L(X)=\int_1^X|M_L(x)/x|^2\,dx.
\]
The following consequences hold without either adopted GRH input:
\[
 \limsup_{X\to\infty}\frac{I_L(X)}{(\log X)^3}>0,
 \qquad
 \limsup_{X\to\infty}\frac{I_L(X)}{\log X}=+\infty.
 \tag{MT-M23a}
\]
If the group has an irreducible representation of degree at least
three, the exponent $3$ in the first assertion can be replaced by
$5$. The canonical full-jet quotient of $\zeta_L$ has no faithful
continuous Hilbert completion on which every convolution extends
boundedly and every $C_{h*h^\sharp}$ is a positive operator.
Nevertheless its nuclear trace construction and adapted absolute
basis exist unconditionally by MT-M18--MT-M19, and $M_H$ makes
the traces of these adjoint-square convolutions nonnegative.

\paragraph{Complete proof of the consequences.}
Choose a multiple zero supplied by HKMS, of actual multiplicity
$m\ge2$, or $m\ge3$ in the stronger case. If $\zeta_L$ has any
off-critical zero, the functional equation and conjugation give
a zero with real part greater than $1/2$. The last assertion of
NT-L16 then gives infinite limsup after division by every fixed
nonnegative power of $\log X$; all displayed conclusions follow.
Otherwise all nontrivial zeros are critical, and our chosen zero
$\rho$ satisfies the exact NT-L16 bound
\[
 \limsup_{X\to\infty}\frac{I_L(X)}{(\log X)^{2m-1}}
 \ge
 \frac{1}{(2m-1)((m-1)!)^2
       |\rho\,\zeta_L^{(m)}(\rho)/m!|^2}>0.
\]
When $2m-1$ equals the claimed exponent this is the first assertion.
When it is greater, take an unbounded sequence on which the
displayed ratio is bounded below by a fixed positive constant;
multiplying by the remaining positive power of $\log X$ gives
infinite limsup. The same argument with exponent one proves
the second assertion. This dichotomy uses no unproved zero
location. Finally MT-M20 proves that a faithful positive Hilbert
realization with the stated properties forces every nontrivial
zero to be simple. The actual multiple zeros contradict that
necessary condition. MT-M18--MT-M19 give the separate basis and
trace statements, and their GRH equivalence gives the conditional
trace positivity under $M_H$.

For the programme this is a finite-monodromy Artin-motive
obstruction to discarding nilpotent spectral blocks. It sharpens
the Kummer examples NT-L17 to all nonabelian Galois top fields.
The cote1/SA23 \texttt{geo-H-lidskii} trace route must retain the
difference between positivity of a scalar trace and positivity
of every operator in a faithful full-multiplicity Hilbert model.
HKMS's general multiplicity prediction for representations of
degree at least four remains their Conjecture6.2, not a theorem
used here. A non-Galois field is covered only when it is itself
the top field of a nonabelian Galois extension of some subfield.


\subsection*{MT-ART09. A smaller relative-zero gate for solvable Artin holomorphy}

Let $L/K$ be a finite solvable Galois extension, with group $G$,
and let $s_0\ne1$. Write
\[
 q(s_0)=\operatorname{ord}_{s=s_0}
             \frac{\zeta_L(s)}{\zeta_K(s)}\ge0.
\]
The nonnegativity follows from the classical Aramata--Brauer
theorem. When $|G|$ has at least two distinct prime divisors,
let $p_2$ be the second smallest and put
\[
                 B(G)=\max\{2,p_2-2\}.
\]
Then every Artin $L(s,\chi,L/K)$ is holomorphic at $s_0$ if
$q(s_0)\le B(G)$. This is the combination of Peng-Jie Wong,
\emph{Zeros of Dedekind zeta functions and holomorphy of Artin
$L$-functions}, Theorems1.3 and1.4, printed p.254, with their
complete proofs on pp.259 and260. The solvability hypothesis
is essential to the statements as cited here. In a $p$-group
all irreducible characters are monomial, so Artin holomorphy
already follows globally from Hecke holomorphy.

For $4\nmid |G|$ the sufficient bound $q(s_0)\le3$ holds,
without adding a solvability assumption: Wong Corollary5.3,
printed pp.259--260, includes the group-theoretic implication
that these groups are solvable. His Corollary5.1, p.259, is the
odd-degree special case. One may therefore replace $B(G)$ by
$\max\{B(G),3\}$ when applicable. The stronger Sylow-structure
variants in his Theorems5.2 and5.4 are separately located there;
no unstated prime-factor hypothesis from them is being used.

\paragraph{Complete transfer to the adopted Hecke hypothesis.}
Under $M_H$, a possible pole of an Artin function in
$0<\Re s<1$ must lie on $\Re s=1/2$: Brauer induction writes
that function as a finite quotient of Hecke functions, whose
possible nontrivial denominator zeros are all on that line.
This is the MT-ART01 confinement, and does not assert Artin
holomorphy. The contrapositive of the just-cited local theorem
now proves that a possible pole must satisfy
\[
 \Re s_0=\tfrac12,\qquad q(s_0)\ge B(G)+1.
 \tag{MT-ART09a}
\]
If $4\nmid |G|$, use the strengthened bound at least four.
Thus, if the relative zeta function has order at most $B(G)$
at every point of the critical line, all Artin functions are
holomorphic throughout the open critical strip. This last
zero-order condition is an extra input, not supplied by $M_H$.
For an unconditional local conclusion at a particular $s_0$,
only the corresponding relative-order inequality is needed.

This is a strictly finer criterion than measuring the order
of $\zeta_L$ alone: by the definition of a meromorphic order,
\[
 q(s_0)=\operatorname{ord}_{s_0}\zeta_L
                 -\operatorname{ord}_{s_0}\zeta_K.
\]
The base contribution is removed exactly, including when the
base has a multiple zero. It is not necessary to posit any
noncoincidence of zeros. The argument is a complete input
transfer to the cited classical local theorem, not a new proof
of that theorem or of general Artin holomorphy.

Geometrically, $\zeta_L/\zeta_K$ is the $L$-function of the
augmentation Artin motive attached to the finite etale cover
after deleting ramified primes, with the standard inertia
factors restored. The relative virtual permutation class is
$[\mathbb Q[G]]-[\mathbb Q]$ in its representation $K_0$.
This belongs to the cote1 trace and SGA5 finite-correspondence
framework already fixed in AG-C1, and sharpens the
solvable part of Q-ARTIN-POS beyond MT-ART02. It does not turn a
virtual-character inequality into coefficientwise positivity.


\subsection*{MT-M24. The specified Masullo smooth quotient is the Meyer quotient}
This is a comparison of the actual test spaces and summation maps. It does
not assume Masullo's global converse theorem, his spectral basis claims,
or $M$. Fix the primitive unitary Hecke function and the convention
$\beta(\mathfrak p_v)^{-1}=c_{\mathfrak p_v}$ from MT-M22; put
$C=C_K$, $C^1=\ker|\cdot|$, and $\chi=\beta|_{C^1}$.
Use the norm splitting and volume-one measure on $C^1$ chosen there.
Write $\mathcal E$ for the exponential Schwartz space of MT-M18,
$\mathcal I_\eta=\overline{\psi_\eta*\mathcal E}$ for its closed
completed-Hecke ideal, and $Q_\eta=\mathcal E/\mathcal I_\eta$.
For a real number $a$, let $M_a u(g)=|g|^a u(g)$, and define
\[
 \Theta_\chi=\Phi_{K,\chi}\bigl(\mathcal S_0(\mathbf A_K)\bigr),
 \qquad
 Q^{\mathrm{Mas}}_{\chi,a}
 =\mathcal S_{\mathrm E}(C)(\chi)\big/
       \overline{M_a\Theta_\chi}^{\,\mathcal S_{\mathrm E}(C)(\chi)}.
 \tag{M24.1}
\]
Here the symbols on the right have their literal meanings in Masullo:
\[
 \mathcal S_0(\mathbf A_K)
  =\{f\in\mathcal S(\mathbf A_K):f(0)=\widehat f(0)=0\},
 \qquad
 \Phi_K(f)(g)=\sum_{\xi\in K^\times}f(\xi g),
 \tag{M24.2}
\]
and $\Phi_{K,\chi}$ is compact-character averaging of $\Phi_K$.
The closure in (M24.1) does not enlarge $M_a\Theta_\chi$.
There is an explicit topological isomorphism
\[
 Q_\eta\xrightarrow{\ \sim\ }Q^{\mathrm{Mas}}_{\chi,a},
 \qquad
 [f]\longmapsto
 \bigl[g\longmapsto |g|^a\beta(g)^{-1}f(\log|g|)\bigr].
 \tag{M24.3}
\]
For $a=0$ this identifies the quotient in Masullo (7.1.4.7),
original PDF p.~327. For $a=1/2$ it identifies the centered quotient
in (1.8.3.1), original PDF p.~45. These two source displays use
different twists; (M24.3) keeps the twist explicit.

\begin{proof}
First verify the test-space topology rather than just its elements.
Masullo Definition D.3.3.1, original PDF p.~804, requires
$|\cdot|^b u\in\mathcal S(C)$ for every real $b$.
Definition D.3.3.3, p.~805, gives each fixed compact-invariance
level the projective topology of all these weighted Schwartz spaces,
and gives their union the locally convex inductive topology.
On the $\chi$-sector the map
\[
 \mathcal L_\beta f(g)=\beta(g)^{-1}f(\log|g|)
 \tag{M24.4}
\]
is a linear bijection from $\mathcal E$ onto that sector.
Indeed, the $\chi$-eigencondition for left translation says
$u(cg)=\chi(c)^{-1}u(g)$ for $c\in C^1$; its restriction to
the norm splitting therefore determines $u$. The unitary value
of $\beta$ on that splitting is $e^{i\tau x}$ for a fixed real
$\tau$, so removing it only changes derivatives by fixed finite
linear combinations.

Every function in this sector is invariant under $\ker\chi$.
The compact quotient $C^1/\ker\chi$ is a closed subgroup of the
circle, hence is either finite or a circle, and this is one permitted
elementary compact-invariance level. At this level the weighted
Schwartz seminorms, after $x=\log|g|$, are equivalent to
\[
 p_{N,j}(f)=\sup_x e^{N|x|}|f^{(j)}(x)|,
 \qquad N,j\ge0.
 \tag{M24.5}
\]
To see the equivalence directly, the weights $b=N$ and $b=-N$
bound (M24.5); conversely every polynomial factor in $x$ occurring
in a weighted Schwartz seminorm is absorbed by one larger
exponential weight. Derivatives on the compact factor multiply
this fixed character by constants. Inclusion of this level in the
inductive limit is continuous. The inverse of (M24.4) on the
sector is also continuous: compact-character averaging followed
by restriction to the norm splitting is continuous into
$\mathcal E$ on each elementary level, by these same seminorm
estimates. It is therefore continuous on the inductive limit.
This establishes the stated sector topology directly, without
using a topological claim about a different quotient.

Next identify the denominator as a set of actual functions.
Use Meyer's common ambient space and embeddings from MT-M22:
\[
 i_+(v)=(\Sigma v,J\Sigma\mathcal F_*v),
 \qquad i_-(u)=(u,u),
 \qquad H_+=\mathcal S(\mathbf A_K)/K^\times,
 \quad H_-=\mathcal S_{\mathrm E}(C).
\]
The summation maps agree term by term:
$\Sigma f(g)=\Phi_K(f)(g)$ by (M24.2).
Their Poisson identity is
\[
 J\Sigma\widehat f-\Sigma f
      =f(0)-|\cdot|^{-1}\widehat f(0).
 \tag{M24.6}
\]
For $f\in\mathcal S_0(\mathbf A_K)$ the two entries agree.
The first entry has all the right-half-plane Schwartz weights,
and the second all the left-half-plane weights, by the defining
Meyer embeddings. Their common value has all real weights:
between any two admitted weights the weighted derivative is
bounded by the two endpoint weighted derivatives, and polynomial
Schwartz factors are treated in the same way. Thus this common
value belongs to $H_-$, and
$\Phi_K(\mathcal S_0)\subset H_+\cap H_-$.

Conversely, if $(u,u)=i_+(v)$ lies in this intersection, choose
an adelic Schwartz representative $f$ of the coinvariant class
$v$. Equation (M24.6) then vanishes identically. The functions
$1$ and $|\cdot|^{-1}$ are linearly independent because the norm
map is onto $\mathbf R_{>0}$, so $f(0)=\widehat f(0)=0$.
Consequently
\[
       \Phi_K(\mathcal S_0(\mathbf A_K))=H_+\cap H_-.
 \tag{M24.7}
\]
This argument uses an actual representative of the coinvariant
class, not a claim that every completion class has an adelic
representative. $H_+$ here is Meyer's quotient of the adelic
Schwartz space.

Compact-character averaging preserves both embedded spaces
and their intersection. On functions, the volume-one projector is
\[
 P_\chi u(g)=\int_{C^1}u(cg)\chi(c)\,dc.
\]
Masullo (2.2.5.2), p.~95, is precisely this formula applied to
$\Phi_K f$, except that his chosen compact Haar measure can have
volume $c_K\ne1$. This multiplies the operator by the fixed
nonzero scalar $c_K$ and leaves its image unchanged. Hence
\[
       \Theta_\chi=(H_+\cap H_-)^\chi.
 \tag{M24.8}
\]
One can also verify surjectivity in (M24.8) directly: an element
of the intersection in the $\chi$-sector already has an
$\mathcal S_0$ representative by (M24.7), and applying
$\Phi_{K,\chi}$ to that representative returns $c_K$ times the
element. No choice of a compact subgroup lift in the idele group
is required.

MT-M22, whose image criterion and topology were proved using
Meyer's specified maps, now gives
\[
 \mathcal L_\beta^{-1}\Theta_\chi
 =\{f\in\mathcal E:(\mathcal Mf)^{(j)}(\rho)=0
         \text{ for all }0\le j<m_\rho\}
 =\mathcal I_\eta.
 \tag{M24.9}
\]
The middle space is closed, since each indicated evaluation is
continuous. Thus $\Theta_\chi$ is closed in the actual sector
topology just checked. Multiplication by $|\cdot|^a$ is an
automorphism of that topology, with inverse multiplication by
$|\cdot|^{-a}$; this follows either from (M24.5), using an
integer larger than $|a|$, or directly by translating the family
of all real weights. This proves the closure assertion and the
continuous two-sided quotient isomorphism (M24.3).
\end{proof}

The action is also explicit. For $b\in C$,
\[
 \lambda_b M_a\mathcal L_\beta f
  =\beta(b)|b|^{-a}M_a\mathcal L_\beta
       \bigl(f(\,\cdot-\log|b|)\bigr).
 \tag{M24.10}
\]
For an idele-class convolution test $k$, set
\[
 k_\beta(y)=\int_{C^1}k(c\tau(e^y))\beta(c\tau(e^y))\,dc.
\]
Changing variables in the Haar integral gives
\[
 C_k M_a\mathcal L_\beta f
   =M_a\mathcal L_\beta\bigl((e^{-a\cdot}k_\beta)*f\bigr).
 \tag{M24.11}
\]
In particular the spectral multiplier is evaluated at $\rho-a$,
with the compact-character factor $\beta$ retained. At $a=0$ this
is precisely the MT-M22 module comparison; at $a=1/2$ it is the
centered version, not an assertion that the same convolution
operator is unchanged by the twist.
An alternative Haar normalization rescales the scalar convolution
integral by the same fixed positive constant; it changes neither
the denominator nor the quotient topology.

\paragraph{Exact source scope and the remaining completion question.}
The original Masullo PDF pages used here are: the theta-sum
definition (2.2.3.1), pp.~81--82, and its underlying sum (2.2.2.1),
p.~77; the polar hyperplanes (2.2.4.23), p.~91;
character averaging (2.2.5.2), p.~95; the actual smooth quotient
(7.1.4.7), p.~327; its centered introductory version (1.8.3.1),
p.~45; and Definitions D.3.3.1/D.3.3.3, pp.~804--805.
The original p.~327 quotient formula was rendered and read;
the other defining formulas were checked against extraction
from the original PDF. The Meyer inputs are his summation map
and Lemma~5.3, pp.~39--40; the embeddings in Lemma~5.4,
pp.~40--41; and Poisson formula (31), p.~41, with the complete
provided arguments already read for MT-M22. No new proof of
the full Masullo draft is claimed.

Masullo Definition~10.1.1.1, p.~434, instead closes
$M_a\Theta_\chi$ in a weighted Hilbert norm, and
(10.2.1.1)--(10.2.1.2), p.~439, and (10.4.1.3), p.~456,
form the resulting Hilbert quotients. Those are different
topological completions of the identified smooth quotient.
Their kernel, retained jets and trace-class properties must be
proved in that norm; MT-M25 and MT-M26 address those maps.
The quotient identification itself is unconditional. Any
subsequent use of universal Hecke GRH is still the conditional
input $M$.

The geometric crosswalk is now an actual Tate--Poisson
summation quotient, rather than an unspecified space sharing
its zero multiplicities. It advances the inherited cote~1 /
SA23 geo-H-lidskii and geo-H-genus interfaces at the precise
test-space and closed-denominator level. It makes no claim
to construct an $\ell$-adic or motivic realization.

\subsection*{MT-M25. Exact line-completion kernel and the strict jet cutoff}
Fix the Hecke function, $\mathcal E$, $\Xi=\Xi_\eta$,
$\mathcal I=\mathcal I_\eta$ and $Q=\mathcal E/\mathcal I$
of MT-M18. For $s\ge0$ and $a\in\mathbf R$, put
\[
 H_s=L^2\bigl(\mathbf R,(1+x^2)^{s/2}\,dx\bigr),
 \qquad T_af(x)=e^{ax}f(x),
 \qquad J_{a,s}=\overline{T_a\mathcal I}^{\,H_s},
 \quad Z_{a,s}=H_s/J_{a,s}.
 \tag{M25.1}
\]
These are exactly the fixed-character versions of Masullo's
line norm (9.1.2.3), original PDF p.~399, his closure
(10.1.1.3), p.~434 with $I=\{0\}$, and his quotient
(10.2.1.2), p.~439, by MT-M24. The fixed unitary norm
modulation has absolute value one and is harmless for this norm.
Let $m_\rho$ be the multiplicity of a zero of $\Xi$, and set
\[
       q_s=\max\{0,\lceil(s-1)/2\rceil\}.
 \tag{M25.2}
\]
Then
\[
 J_{a,s}=
 \bigcap_{\substack{\rho=a+i\gamma:\Xi(\rho)=0\\
                    0\le k<\min(m_\rho,q_s)}}
 \ker\left[u\longmapsto\int_{\mathbf R}
                  u(x)x^k e^{i\gamma x}\,dx\right].
 \tag{M25.3}
\]
Every functional displayed here is continuous on $H_s$.
The natural map
\[
 q_{a,s}:Q\longrightarrow Z_{a,s},\qquad[f]\longmapsto[T_af]
 \tag{M25.4}
\]
is continuous with dense image, and its kernel consists exactly
of the classes whose jets $F^{(k)}(\rho)$, $F=\mathcal Mf$,
vanish at all the indices in (M25.3). Thus the line completion
retains only the zeros on $\Re\rho=a$ and only their derivatives
with the \emph{strict} inequality
\[
                         2k<s-1.
 \tag{M25.5}
\]
It is faithful on the full smooth quotient if and only if all
zeros lie on that line and $m_\rho\le q_s$ for every zero.
No GRH or simplicity is assumed in these assertions.

\begin{proof}
The continuous bilinear dual of $H_s$ can be written as
$H_{-s}$ under $\ell_g(u)=\int u(x)g(x)\,dx$.
This is just the Riesz representation theorem after multiplying
by $(1+x^2)^{s/4}$; complex conjugation can be absorbed into
the representing $g$.
For a polynomial $P$ of degree $k$ with nonzero leading
coefficient,
\[
 P(x)e^{i\gamma x}\in H_{-s}
 \quad\Longleftrightarrow\quad
 \int_{|x|\ge1}|x|^{2k-s}\,dx<\infty
 \quad\Longleftrightarrow\quad 2k<s-1.
 \tag{M25.6}
\]
This proves continuity of the displayed functionals and the
necessity of their strict numerical threshold. It also shows
that the number of admissible nonnegative integers $k$ is
exactly $q_s$, including $0\le s\le1$ when there are none.
Each such functional vanishes on $T_a\mathcal I$ by MT-M18's
full-jet ideal equality, hence on its Hilbert closure.

We prove there are no other annihilators. The elementary
inequality
\[
 1+(x+y)^2\le2(1+x^2)(1+y^2)
\]
implies, for either sign of $s$,
\[
 \|u(\,\cdot-y)\|_{H_{\pm s}}
       \le 2^{s/4}(1+y^2)^{s/4}\|u\|_{H_{\pm s}}.
 \tag{M25.7}
\]
For the negative sign use the same inequality with $x$ and
$x-y$ interchanged before taking reciprocals. Minkowski's
integral inequality therefore gives
\[
 \|k*u\|_{H_{\pm s}}
 \le2^{s/4}\|u\|_{H_{\pm s}}
       \int|k(y)|(1+y^2)^{s/4}\,dy.
 \tag{M25.8}
\]
In particular convolution by a Schwartz function is bounded
on both spaces. Smooth compactly supported functions are
dense in both weighted $L^2$ spaces: truncate first, and then
approximate on the remaining compact interval, where the
weight and its inverse are bounded. These facts also prove
strong continuity of translation, first on that dense subspace
and then by (M25.7).

Let $g\in H_{-s}$ annihilate $T_a\mathcal I$.
It defines a tempered distribution, since Cauchy--Schwarz
bounds $|\int ug|$ by the continuous Schwartz seminorm
$\|u\|_{H_s}\|g\|_{H_{-s}}$.
Use the Fourier convention
$\mathcal F_+u(t)=\int u(x)e^{itx}\,dx$ and its inverse
$(2\pi)^{-1}\int\phi(t)e^{-itx}\,dt$.
Define the tempered distribution $D$ by
$D(\phi)=\ell_g(\mathcal F_+^{-1}\phi)$.
Since $\psi_a=T_a\psi_\eta\in\mathcal E$ and
$\mathcal F_+\psi_a(t)=\Xi(a+it)$, we obtain
\[
                         \Xi(a+it)D=0.
 \tag{M25.9}
\]
Here is why the identity holds on every Schwartz test,
not just an entire Fourier transform. For $u\in\mathcal E$,
$\psi_a*u\in T_a\mathcal I$ because
$T_a(f_1*f_2)=T_af_1*T_af_2$ and $T_a\mathcal E=\mathcal E$.
Thus $\ell_g(\psi_a*u)=0$. The space $\mathcal E$ contains
$C_c^\infty$ and is dense in the usual Schwartz space;
convolution by $\psi_a$ is continuous there, and
$\ell_g$ is continuous there. This extends the equality to
all Schwartz $u$, which is precisely (M25.9).

The zeros of $t\mapsto\Xi(a+it)$ on the real line form a
locally finite set. Away from them, a compactly supported
smooth test can be divided by this nonvanishing multiplier,
so (M25.9) shows $\operatorname{supp}D$ is contained in that
set. Near a zero $\gamma$ of order $m_\rho$, multiplication
by a smooth invertible function reduces (M25.9) to
$(t-\gamma)^{m_\rho}D=0$. A distribution supported there
with this property depends only on the first $m_\rho$ jets
of a test: subtract its Taylor polynomial of degree
$m_\rho-1$ times a cutoff equal to one near $\gamma$;
the remainder is divisible by $(t-\gamma)^{m_\rho}$ there.
It is therefore a linear combination of
$\delta_\gamma^{(k)}$, $0\le k<m_\rho$.

Choose a smooth frequency cutoff supported near this single
zero and equal to one near it. Multiplying $D$ by that
cutoff corresponds to convolution of $g$ with a Schwartz
kernel, with an inessential reflection dictated by the
chosen Fourier convention. By (M25.8) the resulting
function still belongs to $H_{-s}$. Its inverse transform
is a polynomial times $e^{i\gamma x}$. Formula (M25.6)
therefore excludes every nonzero coefficient of degree
$k$ with $2k\ge s-1$. Locally only the functionals listed
in (M25.3) occur.

Finally these local pieces exhaust the annihilator in its
Hilbert norm. Choose $\vartheta\in C_c^\infty(\mathbf R)$
equal to one near zero and put
$\vartheta_R(t)=\vartheta(t/R)$, $R\ge1$.
The associated convolution kernels have the form
$Rk(Rx)$, possibly reflected, where $k$ is Schwartz and
$\int k=1$. Their weighted $L^1$ norms in (M25.8) are
uniformly bounded, because
\[
 \int R|k(Ry)|(1+y^2)^{s/4}\,dy
       =\int|k(v)|(1+v^2/R^2)^{s/4}\,dv
       \le\int|k(v)|(1+v^2)^{s/4}\,dv.
\]
Translation continuity and dominated convergence in the
same integral show $Rk(R\cdot)*g\to g$ in $H_{-s}$.
The distribution $\vartheta_RD$ is supported on finitely
many zeros, so each such convolution is a finite linear
combination of the permitted polynomial waves. Hence
the full annihilator of $T_a\mathcal I$ is the Hilbert
closure of the span of precisely the functionals in
(M25.3). Taking its annihilator proves (M25.3).

The inclusion $T_a:\mathcal E\to H_s$ is continuous,
since one exponential seminorm with weight greater than
$|a|$ absorbs the polynomial factor in (M25.1). Its range
is dense because $T_a\mathcal E=\mathcal E$ contains
$C_c^\infty$. It vanishes on $\mathcal I$ after passing
to the quotient, and so induces (M25.4), continuously
with dense image. The functional indexed by $(\rho,k)$
evaluates $T_af$ to
\[
       \int e^{ax}f(x)x^ke^{i\gamma x}\,dx
                    =(\mathcal Mf)^{(k)}(\rho).
\]
This gives its exact kernel. If all zero jets are retained,
MT-M18 makes that kernel zero. If any zero jet is omitted,
the finite-block interpolation of MT-M18 realizes a class
with just that jet nonzero and all other zero jets zero.
It is a nonzero class in the kernel of (M25.4). This proves
the stated necessary and sufficient condition for faithfulness.
\end{proof}

\paragraph{A literal boundary correction with an arithmetic witness.}
The original PDF (10.2.1.6), p.~440, displays
\[
 N_s(\rho)=\min\{N(\rho),[(s-1)/2+1]\}.
\]
At $s=3$ this allows the first derivative at any double zero.
But its proposed representing function is
$x e^{i\gamma x}$, whose squared $H_{-3}$ norm contains
\[
              \int_{|x|\ge1}\frac{x^2}{(1+x^2)^{3/2}}\,dx
                         =+\infty.
 \tag{M25.10}
\]
Thus the endpoint inclusion is false as written; the
correct retained multiplicity is
$\min\{m_\rho,q_s\}$. Both the source norm (9.1.2.3),
p.~399, and the displayed cutoff (10.2.1.6), p.~440, were
rendered and visually read from the original PDF.
This is not a hypothetical objection requiring an unknown
double zero of the Riemann zeta function. By NT-L17,
the actual degree-six field
$K=\mathbf Q(\sqrt[3]{2},\sqrt{-3})$ has nontrivial zeros
of $\zeta_K$ of order at least two. Choose such a zero
$\rho$, put $a=\Re\rho$, and use its trivial Hecke
character. Then (M25.10) is a concrete counterexample
to the stated membership at its permitted $a$ and $s=3$.
Under $M$, one may take $a=1/2$.

More generally, for every fixed finite $s$, choose an odd
prime $r$ with $r-1>q_s$. The Kummer field of NT-L17
has zeros of multiplicity at least $r-1$. Its line quotient
at $a=1/2$ cannot faithfully retain the full smooth
quotient: under $M$ it loses the higher jets, and without
$M$ either that loss occurs or an off-line zero is lost.
The strict cutoff is a property of this particular
Hilbert completion, not a failure of the unconditional
adapted-block synthesis in MT-M18.

This gives the precise analytic map at the inherited
cote~1 / SA23 geo-H-lidskii and geo-H-genus interface:
the Hilbert norm can discard part of the finite-length
spectral divisor. It does not prove or refute the
conditional endpoint $M$, and it invokes none of the
draft's unverified general basis or nuclearity arguments.

\subsection*{MT-ART10. An integral Galois-module class forces a central zero}

Let $N/\mathbb Q$ be a tamely ramified $Q_8$-extension and let
$\chi$ be the unique irreducible complex character of degree two.
Fr\"ohlich's classical theorem identifies the two signs
\[
 W(\chi)=+1\quad\Longleftrightarrow\quad
       \mathcal O_N\text{ is free over }\mathbb Z[Q_8].
 \tag{MT-ART10a}
\]
The exact primary statement is Fr\"ohlich, \emph{The Galois
module structure of algebraic integer rings in fields with
generalised quaternion group} (1974), Theorem1, printed p.82,
with the locally free class invariant defined on pp.81--82.
The original proof is his \emph{Artin-root numbers and normal
integral bases for quaternion fields}, Inventiones17(1972),
143--166. Its full proof was not retrieved here; (MT-ART10a)
is used as the classical theorem, not claimed as a fresh proof.
Martinet's \emph{Fonctions z\^eta au point $1/2$}, \S VI,
printed114--115, gives the equivalent normal-basis formulation.

Consequently, if $\mathcal O_N$ is not free then
\[
 L(\tfrac12,\chi)=0,
 \qquad \operatorname{ord}_{s=1/2}\zeta_N(s)\ge2.
 \tag{MT-ART10b}
\]
These assertions are unconditional and are compatible with
both adopted GRH inputs.

\paragraph{Complete short deduction.}
Choose a cyclic subgroup $C_4\subset Q_8$ and a faithful
degree-one character $\lambda$ of $C_4$. Its induction is
$\chi$: its values are $2$ at the identity, $-2$ at the central
involution, and zero elsewhere, which is the irreducible
degree-two character. Artin induction-invariance therefore
identifies $L(s,\chi)$ with a nontrivial finite-order Hecke
function over $N^{C_4}$, so it is entire. Since $\chi$ is
selfdual, its completed function satisfies
\[
                 \Lambda(s,\chi)=W(\chi)\Lambda(1-s,\chi).
\]
Nonfreeness and (MT-ART10a) give $W(\chi)=-1$. Evaluation at
$s=1/2$ forces $\Lambda(1/2,\chi)=0$. The archimedean factors
and conductor power are finite and nonzero there, so the
finite-part $L$-function also vanishes. If $K=N^{\{1,-1\}}$
is the biquadratic subfield, the regular-character identity
gives the all-prime factorization
\[
                         \zeta_N(s)=\zeta_K(s)L(s,\chi)^2.
\]
The first factor is holomorphic at $1/2$, so the order is at
least two. More precisely the order of $L(s,\chi)$ there is
positive and odd, by the same functional equation. This proves
(MT-ART10b) without a general Artin-holomorphy hypothesis.

\paragraph{A completely specified classical example.}
Martinet, \emph{Modules sur l'alg\`ebre du groupe quaternionien}
(1971), TheoremIV.1 printed406 and Example3 printed407, gives
\[
 K=\mathbb Q(\sqrt5,\sqrt{41}),\qquad
 M=\frac{5+\sqrt5}{2}\frac{41+\sqrt{205}}{2},\qquad
 N=K(\sqrt M).
 \tag{MT-ART10c}
\]
This is a totally real tame $Q_8$-extension, ramified only at
$5$ and $41$, with nonfree integer ring. In his trace criterion
the two values are $205$ and $(1+5+41+205)/4=63$; they differ
modulo four, so the nonfree case applies. The radicand was
checked directly in the original page407 image, where the
second square root is $\sqrt{5\cdot41}$.
Thus this particular field has $\zeta_N(1/2)=0$ and satisfies
(MT-ART10b). For a convenient polynomial description, multiplying
$X^2-M$ over the four conjugates of $M$ gives
\[
 X^8-205X^6+13940X^4-378225X^2+3404025.
 \tag{MT-ART10d}
\]
An exact rational coefficient calculation of this norm is
recorded separately; the Galois group, ramification and integral
module classification are the cited classical example.

The programme bridge here is precise: tameness makes
$\mathcal O_N$ a locally free integral Galois module, and its
class in the locally free class group of $\mathbb Z[Q_8]$
controls one functional-equation sign. Rationalizing erases
this integral class while retaining the usual regular
representation. This is an example of the integral $K_0$
information distinguished in AG-LEFSCH-LATTICE and
AG-LEFSCH-DEFECT, not a general formula for all arithmetic
$K$-groups. It also shows that a residual condition such as
the nonvanishing of a Dedekind zeta function at $1/2$ cannot
be discarded merely because $M_H$ holds. It is not an example
inside the particular $S_5$ family of MT-S502; that family's
central nonvanishing gate remains its own arithmetic question.


\subsection*{MT-M26. Trace class on the specified strip Hilbert quotient}
Fix one primitive unitary Hecke function and the notation of
MT-M18 and MT-M24. Let $s\ge0$, $a\in\mathbf R$, and let
$I=[u,v]$ be a compact interval with $u<v$. Define
\[
 \begin{split}
 w_{I,s}(x)&=\max\{e^{2ux},e^{2vx}\}(1+x^2)^{s/2},\\
 H_{I,s}&=L^2(\mathbf R,w_{I,s}(x)\,dx),\\
 J_{a,I,s}&=\overline{T_a\mathcal I_\eta}^{\,H_{I,s}},
 \qquad Z_{a,I,s}=H_{I,s}/J_{a,I,s},
 \qquad T_af(x)=e^{ax}f(x).
 \end{split}                                                    \tag{M26.1}
\]
Assume there is a closed strip $[c_0,c_1]$ containing the real
parts of all zeros of $\Xi_\eta$, with
\[
                 [c_0-a,c_1-a]\subset(u,v).
                                                               \tag{M26.2}
\]
Then the natural map $Q_\eta\to Z_{a,I,s}$ is continuous,
injective and has dense image. More precisely,
\[
 J_{a,I,s}
 =\{h\in H_{I,s}:(\mathcal Mh)^{(k)}(\rho-a)=0
       \text{ for all zeros }\rho\text{ and }0\le k<m_\rho\}.
                                                               \tag{M26.3}
\]
The Mellin transforms and derivatives in (M26.3) are defined
by absolutely convergent integrals inside the strip $I$.
For every $g\in\mathcal E$, convolution induces a trace-class
operator on this \emph{actual} Hilbert quotient, and
\[
       \operatorname{Tr}_{Z_{a,I,s}}C_g
          =\sum_{\Xi_\eta(\rho)=0}
                        m_\rho(\mathcal Mg)(\rho-a),
                                                               \tag{M26.4}
\]
with absolute convergence. The trace-class estimates are
uniform when $g$ ranges over a bounded subset of $\mathcal E$.

Unconditionally one may take $[c_0,c_1]=[0,1]$, so the result
holds whenever $[0,1]\subset a+\operatorname{int}I$.
Under GRH for this one function, one may take $c_0=c_1=1/2$.
In particular, under the conditional input $M$, the result
holds at $a=1/2$ for every $I=[-\alpha,\alpha]$ with
$\alpha>0$, for each fixed Hecke character separately.

\begin{proof}
We first specify the adapted blocks needed inside the available
strip. Choose fixed real numbers
\[
        a+u<b_-<c_0\le c_1<b_+<a+v.
\]
Use the selected horizontal heights $h_n$ from MT-M18 and
rectangles $\Omega_n=(b_-,b_+)\times(h_n,h_{n+1})$.
The proof of MT-M18 gives the same block theorem for these
vertical edges, as the following estimates make explicit.
The horizontal minimum-modulus estimate there was proved
on every fixed horizontal segment, so it applies on
$[b_-,b_+]$. On either vertical edge the distance from every
zero is at least
$\delta=\min(c_0-b_-,b_+-c_1)>0$.
In the Blaschke-product proof of MT-M18, each factor on these
edges is therefore at least
$\min(1,\delta)/(2R)$, where $R$ is the fixed disk radius
used there. There are $O(\log(2+|n|))$ such factors.
Their product has a polynomial lower bound. The same
zero-free factor estimate and the two-sided Stirling bound
then give
\[
 \inf_{w\in\partial\Omega_n}|\Xi_\eta(w)|
    \ge c(2+|n|)^{-B}e^{-\kappa|n|}.
                                                               \tag{M26.5}
\]
At bounded height the edges contain no zeros and compactness
supplies the bound. The fixed width and bounded height of the
rectangles give uniformly bounded perimeter. The entire
difference-kernel construction in MT-M18, its Gaussian estimate,
and its quotient-infimum argument consequently give finite
Banach algebras $J_n$ and lifts $\mathcal T_n:J_n\to\mathcal E$
with
\[
 \begin{gathered}
 d_n=\dim J_n=O(\log(2+|n|)),\qquad
 R_m\mathcal T_n=\begin{cases}\operatorname{id}_{J_n},&m=n,\\0,&m\ne n,
                         \end{cases}\\
 p_{N,j}(\mathcal T_n z)
       \le C_{N,j}(2+|n|)^{B_{N,j}}\|z\|_{J_n}.
 \end{gathered}                                                 \tag{M26.6}
\]
Here $J_n$ is the quotient of the holomorphic supremum-norm
algebra on the closed rectangle by its full block jets.
The absolute synthesis, closed-ideal equality and two-sided
block isomorphism of MT-M18 also hold: their proof uses
(M26.5), bounded perimeter and inclusion of every zero in
exactly one rectangle, all just verified. In particular the
lifts are not unstable individual Hermite interpolation
vectors. This establishes the precise version of MT-M18
being used below.

For $h\in H_{I,s}$ its Mellin integral is holomorphic on
$u<\Re z<v$. If $\sigma$ ranges over a compact subinterval
of $(u,v)$, then for every integer $k\ge0$,
\[
 \begin{split}
 |(\mathcal Mh)^{(k)}(\sigma+it)|
 &\le\|h\|_{H_{I,s}}
   \left(\int_{\mathbf R}
     \frac{|x|^{2k}e^{2\sigma x}}{w_{I,s}(x)}\,dx\right)^{1/2}\\
 &\le C_k\|h\|_{H_{I,s}},
 \end{split}                                                    \tag{M26.7}
\]
uniformly in $t$. The integral converges by the positive
distance to both endpoints, uniformly in that compact
subinterval. The same estimate justifies differentiation
under the integral on compact subsets of the strip.
Consequently
\[
 R_n^H(h)=\bigl[w\longmapsto(\mathcal Mh)(w-a)\bigr]\in J_n
 \quad\hbox{satisfies}\quad
       \|R_n^H(h)\|_{J_n}\le C\|h\|_{H_{I,s}},
                                                               \tag{M26.8}
\]
where $C$ is independent of $n$. This follows by bounding
the quotient norm by the supremum of this representative
on the contour. For $h=T_af$ the representative is
$\mathcal Mf(w)$, so $R_n^H$ vanishes on $T_a\mathcal I_\eta$,
and by continuity on its Hilbert closure. It induces a
bounded map $r_n:Z_{a,I,s}\to J_n$ with norm at most $C$.

Let $\pi:H_{I,s}\to Z_{a,I,s}$ be the quotient map, and put
$s_n=\pi T_a\mathcal T_n$. The map $T_a:\mathcal E\to H_{I,s}$
is continuous: a seminorm $p_{N,0}$ with
$N>|a|+\max(|u|,|v|)+1$ bounds its norm by an integrable
exponential times the fixed polynomial weight. Thus (M26.6)
implies, for some fixed constants $C,B$,
\[
 \|s_n\|_{J_n\to Z_{a,I,s}}\le C(2+|n|)^B,
 \qquad r_m s_n=\delta_{mn}\operatorname{id}_{J_n}.
                                                               \tag{M26.9}
\]
The last identity follows directly from the jets of the
lifts. In particular each block is actually realized in
the Hilbert quotient.

Convolution by $g\in\mathcal E$ is bounded on $H_{I,s}$.
Indeed, with $B_I=\max(|u|,|v|)$, the elementary polynomial
weight inequality used in MT-M25 gives
\[
 \|h(\,\cdot-y)\|_{H_{I,s}}
 \le2^{s/4}e^{B_I|y|}(1+y^2)^{s/4}\|h\|_{H_{I,s}}.
                                                               \tag{M26.10}
\]
Minkowski's inequality integrates this bound against
$|g(y)|$, and the resulting integral is finite.
The subspace $T_a\mathcal I_\eta$ is invariant under this
convolution because
\[
           g*T_af=T_a\bigl((e^{-a\cdot}g)*f\bigr)
                                                               \tag{M26.11}
\]
and $\mathcal I_\eta$ is a convolution ideal in $\mathcal E$.
Its Hilbert closure is therefore invariant as well. Hence
$C_g$ induces a bounded operator on $Z_{a,I,s}$.

Let $M_{g,n}$ be multiplication in $J_n$ by the block class
of $w\mapsto\mathcal Mg(w-a)$. Submultiplicativity of the
quotient algebra norm shows
\[
 \|M_{g,n}\|\le b_n(g):=
       \sup_{w\in\partial\Omega_n}|\mathcal Mg(w-a)|.
                                                               \tag{M26.12}
\]
For each $A>0$, $b_n(g)\le C_A(g)(2+|n|)^{-A}$,
with $C_A(g)$ bounded on bounded subsets of $\mathcal E$,
by its strip-Schwartz Mellin transform. These constants
are bounded by finitely many $\mathcal E$ seminorms:
integrating by parts sufficiently many times on the
fixed strip bounds the high ordinates, and direct
integration bounds the remaining compact set. The finite-rank
operators $s_nM_{g,n}r_n$ have rank at most $d_n$ and
satisfy the Hilbert trace-norm bound
\[
 \|s_nM_{g,n}r_n\|_1
       \le C d_n(2+|n|)^B b_n(g).
                                                               \tag{M26.13}
\]
Here we used the elementary inequality
$\|T\|_1\le\operatorname{rank}(T)\|T\|$ for a finite-rank
operator on a Hilbert space. Thus the series
\[
                     \sum_n s_nM_{g,n}r_n
                                                               \tag{M26.14}
\]
converges in trace norm, uniformly for $g$ in bounded test
sets. These same estimates show that
$g\mapsto C_g$ is continuous into the trace-class norm.

It equals the actual induced convolution. To verify this
without assuming spectral completeness in the Hilbert
space, first take a vector $\pi T_af$ with $f\in\mathcal E$.
Equation (M26.11) and the absolute block synthesis in the
\emph{already proved smooth quotient} give the identity
with (M26.14) on this vector. Applying $\pi T_a$ is
legitimate because it is continuous from that smooth
quotient into the Hilbert quotient; no injectivity has
been assumed. These vectors are dense: $T_a\mathcal E$
contains $C_c^\infty$, which is dense in the weighted
$L^2$ space by truncation and compact approximation.
Both operators are bounded, so the identity holds on
all of $Z_{a,I,s}$.

We can now prove the exact closure statement (M26.3).
One inclusion follows already from (M26.7). For the other,
if every zero jet of $h\in H_{I,s}$ vanishes, then
$r_n[ h ]=0$ for every $n$. The operator identity just
proved gives $C_g[h]=0$ for every $g\in\mathcal E$.
Take a smooth compactly supported approximate identity
$g_\epsilon$ of integral one. Translation is strongly
continuous on $H_{I,s}$: this holds first for compactly
supported smooth functions, then for all functions by
their density and (M26.10). That estimate is uniformly
bounded on small $y$, so
$g_\epsilon*h\to h$ in $H_{I,s}$.
Passing to the quotient gives $[h]=0$, proving (M26.3).
In particular $Q_\eta\to Z_{a,I,s}$ is injective, since
all zero jets of a smooth vector still determine its
class by MT-M18. Continuity and density were proved above.

Finally the finite-rank trace identity and (M26.9) give
\[
 \operatorname{Tr}(s_nM_{g,n}r_n)
       =\operatorname{Tr}_{J_n}(M_{g,n}r_ns_n)
       =\operatorname{Tr}_{J_n}M_{g,n}.
\]
In the local algebra at a zero $\rho$ of multiplicity
$m_\rho$, multiplication by a holomorphic function is
triangular in the powers of $w-\rho$, with that function's
value at $\rho$ repeated $m_\rho$ times on the diagonal.
This computes the trace as
$\sum_{\rho\in\Omega_n}m_\rho\mathcal Mg(\rho-a)$.
Summing is legitimate by (M26.13), and the scalar sum
is also absolutely convergent since $d_n=O(\log(2+|n|))$
and the Mellin transform is rapidly decreasing on the
fixed strip. This proves (M26.4).
\end{proof}

The injection from the smooth quotient has proper dense
image. Indeed $\Xi_\eta$ has infinitely many nontrivial
zeros, hence infinitely many nonzero $J_n$. The rapid
block space of MT-M18 is not normable: if one norm gave
its topology, it would be bounded above by a constant
times some $q_K$, while $q_{K+1}$ would be bounded by
a constant times that norm. A unit vector in a single
block at height $n$, scaled by $(2+|n|)^{-K}$, contradicts
these two inequalities as $|n|\to\infty$. Were the
continuous injection onto the Hilbert quotient surjective,
the Fr\'echet open mapping theorem would make the smooth
quotient normable. Thus this is a faithful Hilbert
completion, not an identification of its norm topology
with the stronger smooth topology.

At $a=1/2$ put $h^*(x)=\overline{h(-x)}$. Under GRH,
(M26.4) yields
\[
 \operatorname{Tr}C_{h*h^*}
      =\sum_\rho m_\rho|\mathcal Mh(\rho-1/2)|^2\ge0.
                                                               \tag{M26.15}
\]
On any strip satisfying the unconditional sufficient
condition $[0,1]\subset1/2+\operatorname{int}I$, positivity
of these traces for every $h\in\mathcal E$ is equivalent
to GRH for the fixed function. For if
$\rho^\dagger=1-\overline\rho\ne\rho$ is an off-line pair,
the unitary functional equation gives equal multiplicities.
MT-M18 finite-block interpolation, followed by $T_{1/2}$,
produces $h$ with Mellin values $1,-1$ at
$\rho-1/2,\rho^\dagger-1/2$, respectively, and zero at
all other zeros. Formula (M26.4) then gives
$\operatorname{Tr}C_{h*h^*}=-2m_\rho$.
No positivity of the operator $C_{h*h^*}$ itself is claimed;
the faithful full-jet obstruction of MT-M20 still applies.

\paragraph{Literal source interface and remaining scope.}
Masullo's norm (9.2.2.31), original PDF p.~416, becomes
(M26.1) after $x=\log|g|$ in the fixed compact-character
sector; the supremum over subintervals is attained at
the endpoints of the compact interval $I$. His closure
Definition~10.1.1.1, p.~434, becomes $J_{a,I,s}$ by
MT-M24. For $I=[-\alpha,\alpha]$ and $a>0$, his stated
parameter range, the resulting space
is his specified quotient (10.4.1.3), original PDF p.~456,
in that sector. The convolution convention is exactly
(M24.11); an idele-class test uses the scalar function
$k_\beta$, so its trace is evaluated at the correspondingly
shifted quasi-character. Compactly supported smooth tests
are included in $\mathcal E$.

This independently supplies fixed-character spectral
trace class on those literal quotients; it does not invoke
Theorem~10.3.2.1 or the draft's general Schauder claims.
All constants may depend on the fixed field and character.
A simultaneous summation over varying compact characters
would require additional uniform conductor and archimedean
parameter estimates; that is not inferred from a
fixed-character result. Neither the trace formula nor
the space identification proves $M$. The remaining
positivity input is still conditional.

In the inherited cote~1 / SA23 geo-H-lidskii and geo-H-genus
crosswalk, the map from the completed-Hecke spectral
divisor to a specified adelic Hilbert quotient is now
explicit and faithful whenever the containing strip has
interior room. The line quotient MT-M25 explains why
that interior condition cannot simply be discarded.

\subsection*{MT-M27. Full adelic trace class on a wide strip}
Let $K$ be a fixed number field, $C=C_K$, and
$C^1=\ker|\cdot|$. Choose a norm splitting and normalize
the compact Haar measure to have volume one. Put
\[
 \begin{split}
 E_K&=\mathcal S_{\mathrm E}(C),\qquad
 \Theta_K=\Phi_K(\mathcal S_0(\mathbf A_K)),\\
 H_{J,s}&=L^2\!\left(C,
       \max\{|g|^{2c},|g|^{2d}\}
       (1+|\log|g||^2)^{s/2}\,dg\right),\\
 \mathcal J_{J,s}&=\overline{\Theta_K}^{\,H_{J,s}},
 \qquad Z_{J,s}=H_{J,s}/\mathcal J_{J,s},
 \end{split}                                                    \tag{M27.1}
\]
where $s\ge0$ and $J=[c,d]$ satisfies $c<0<1<d$.
The topology on $E_K$ is the one specified in Masullo
Definitions D.3.3.1/D.3.3.3, original PDF pp.~804--805.
The Hilbert space and denominator in (M27.1) are his
specified norm and Hilbert closure, before the harmless
centering twist of MT-M24.

For each $\chi\in\widehat{C^1}$ choose a unitary extension
$\beta_\chi$ to $C$ and use the Hecke convention of MT-M22.
Let $\Xi_\chi$ have zeros $\rho$ with multiplicities
$m_{\chi,\rho}$. Define
\[
 \widehat f(\beta_\chi|\cdot|^z)
       =\int_C f(g)\beta_\chi(g)|g|^z\,dg.
                                                               \tag{M27.2}
\]
Then every $f\in\mathcal D(C)$ induces a trace-class
convolution operator on $Z_{J,s}$, and
\[
 \operatorname{Tr}_{Z_{J,s}}C_f
     =\sum_{\chi\in\widehat{C^1}}
        \sum_{\Xi_\chi(\rho)=0}
                   m_{\chi,\rho}
                     \widehat f(\beta_\chi|\cdot|^\rho).
                                                               \tag{M27.3}
\]
The series is absolutely convergent, uniformly on bounded
subsets of $\mathcal D(C)$. The trace norms are uniformly
bounded on those same sets. This includes the infinitely
many archimedean characters at any fixed finite level.
No GRH is required.

There are also exact kernel identities:
\[
 \begin{split}
 \overline{\Theta_K}^{\,E_K}
   &=\{u\in E_K:
      \partial_z^k\widehat u(\beta_\chi|\cdot|^z)|_{z=\rho}=0
            \text{ for all }\chi,\rho,0\le k<m_{\chi,\rho}\},\\
 \mathcal J_{J,s}
   &=\{u\in H_{J,s}:
      \partial_z^k\widehat u(\beta_\chi|\cdot|^z)|_{z=\rho}=0
            \text{ for all }\chi,\rho,0\le k<m_{\chi,\rho}\}.
 \end{split}                                                    \tag{M27.4}
\]
Thus the corrected full smooth quotient
$E_K/\overline{\Theta_K}^{E_K}$ maps continuously and
injectively, with proper dense image, into the specified
Hilbert quotient. We do not identify its stronger smooth
topology with the Hilbert norm topology.

\paragraph{Classical input and exact scope.}
The additional uniform input here is the published Meyer
Lemma~5.6, author PDF pp.~41--42, in
\emph{On a representation of the idele class group related
to primes and zeros of $L$-functions},
\url{https://arxiv.org/pdf/math/0311468} (v3).
For sufficiently large finite $S$, restricting its
stated $C_K$-valued map to $U_S$-invariants, as in
its proof, gives
\[
 (P_--P_+^S)C_f:
      L^2(C_S)_\alpha\oplus L^2(C_S)_{1-\alpha}
           \longrightarrow
          \mathcal S(C_S)_{>}\oplus\mathcal S(C_S)_{<}
                                                               \tag{M27.5}
\]
extends to a nuclear operator for $\alpha>1$,
uniformly for $f$ in bounded test sets.
Here $>$ means all weights $>1$, $<$ all weights $<0$,
$P_-$ is the diagonal projection, and $P_+^S$ has range
in the embedded adelic summation space $H_+^S$.
The supplied proof of this lemma was read in full,
including its Poisson and cutoff reduction; its cited
upstream local kernel lemmas are classical inputs, not
claimed newly fully audited. The rank-one meaning of
``nuclear'' is explicitly given in Meyer \S2.1, formula
(5), author pp.~6--7. We prove the transfer from (M27.5)
to the particular Hilbert quotient; it does not follow
merely from equality of spectral multiplicities.

\begin{proof}
\emph{Character projections and the two closed kernels.}
The compact abelian Fourier decomposition gives an
orthogonal Hilbert sum
\[
                     H_{J,s}=\mathop{\widehat\bigoplus}_\chi
                                      H_{J,s}(\chi).
                                                               \tag{M27.6}
\]
The weight depends only on the norm, so the compact
character projectors are orthogonal contractions.
The subspace $\Theta_K$ is invariant under these
projectors: by the global set identity (M24.7), it is
the intersection of the two embedded Meyer spaces,
both preserved by compact averaging. Its $\chi$ image
is exactly $\Theta_\chi$ of MT-M24.
Consequently $\mathcal J_{J,s}\cap H_{J,s}(\chi)$ is
the Hilbert closure of $\Theta_\chi$: apply the
contractive projector to a convergent approximating
sequence from $\Theta_K$. MT-M26 identifies this sector
closure with its full zero-jet kernel, because
$[0,1]\subset\operatorname{int}J$.
All these jet functionals are continuous on the full
Hilbert space: the compact projection has norm one,
and the Cauchy--Schwarz estimate (M26.7) applies in
that sector. Finite sums of character projections
converge strongly in (M27.6). It follows that vanishing
of all character/zero jets is equivalent to membership
in $\mathcal J_{J,s}$. This proves the second identity
in (M27.4).

For completeness the corresponding smooth closure
requires no uniform inverse estimates. An element
$u\in E_K$ lies at one elementary compact-invariance
level, with compact factor a finite abelian group
times a torus. At this level its seminorms are the
weighted suprema of all derivatives in the norm and
compact directions. Integration by parts on that
torus bounds every Fourier coefficient, in each
fixed norm-direction seminorm, by arbitrarily high
negative powers of its lattice frequency. For
example, apply $(1-\Delta)^N$ in the compact variable;
the character eigenvalue grows as $(1+|\nu|^2)^N$.
Taking $2N$ larger than the torus dimension makes
the Fourier tails summable, and increasing $N$
treats any specified compact derivative as well.
Thus finite character truncations converge to $u$
in every seminorm of that elementary level, hence
in $E_K$.
If all its zero jets vanish, each single-character
truncation belongs to the actual image $\Theta_\chi$
by MT-M24. Every finite truncation is therefore in
$\Theta_K$, and $u\in\overline{\Theta_K}^{E_K}$.
The reverse inclusion follows from continuity of
each zero-jet functional and its vanishing on
$\Theta_K$. This proves the first identity in
(M27.4). The natural map from the corrected smooth
quotient into the Hilbert quotient is consequently
injective and continuous. Its image is dense:
finite character sums times smooth compactly
supported functions of $\log|g|$ belong to $E_K$
and are dense in (M27.6). It is proper, since
surjectivity would, after a single compact-character
projection, contradict the proper-image assertion
of MT-M26 in that fixed sector.

\emph{A finite level for each test set, not a finite set of characters.}
Choose $\alpha>1$ so that
$[1-\alpha,\alpha]\subset(c,d)$.
A test $f\in\mathcal D(C)$ is invariant under
$U_S=\prod_{v\notin S}\mathcal O_v^\times$
for sufficiently large finite $S$ containing all
infinite places; $S$ can be enlarged to make the
$S$-ideal class group trivial. One common $S$ works
for a bounded set of tests. This follows directly
from the compact-invariance-level definition of
the test-function bornology, and is stated
explicitly in Meyer's proof of Theorem~5.8,
author p.~43. Enlarging $S$ preserves invariance.
For such $S$,
$C/U_S\simeq C_S=\mathbf A_S^\times/\mathcal O_{K,S}^\times$:
the trivial $S$-class group lets one adjust all
valuations outside $S$ by a principal idele; the
kernel is precisely the $S$-unit group.

Averaging over $U_S$ is an orthogonal contraction
on $H_{J,s}$, preserves $\mathcal J_{J,s}$, and
therefore gives an orthogonal projection on
$Z_{J,s}$. Since $f$ is $U_S$-invariant,
$C_f=C_fP_{U_S}=P_{U_S}C_f$.
It is enough to prove trace class on this invariant
Hilbert subspace, identified with functions on
$C_S$, and then extend it by zero on the orthogonal
complement. The invariant denominator is the
closure of $P_{U_S}\Theta_K$, by the same averaging
argument as above. No truncation of the infinitely
many compact Lie characters of $C_S^1$ is made.

\emph{An explicit bounded map from Meyer's ambient pair.}
Write
\[
 A_S=\mathcal S(C_S)_{>}\oplus\mathcal S(C_S)_{<},
 \qquad E_S=\mathcal S(C_S)_{\mathbf R},
 \qquad i_-(h)=(h,h).
\]
Choose a smooth norm cutoff $\phi$ equal to one
near zero and zero near infinity. Meyer's map
\[
        p_-(u_0,u_1)=(1-\phi)u_0+\phi u_1
                                                               \tag{M27.7}
\]
is bounded from $A_S$ to $E_S$.
This also follows immediately from weighted
seminorms: on the upper tail use a right weight
larger than the desired weight, on the lower
tail a smaller left weight, and treat derivatives
of the cutoff on a compact norm interval.
It satisfies $p_-i_-h=h$ and $P_-=i_-p_-$.

Put $\mathcal P=\operatorname{span}\{1,|\cdot|^{-1}\}$.
Poisson gives
\[
 X_S=H_+^S+i_-(E_S)
      =\{(u_0,u_1)\in A_S:u_0-u_1\in\mathcal P\}.
                                                               \tag{M27.8}
\]
Here is the explicit splitting needed for this identity
and for the subsequent bounded map. Choose two adelic
Schwartz functions, invariant outside $S$, whose pairs
$(f(0),\widehat f(0))$ are linearly independent.
They exist by taking fixed standard finite factors
and archimedean smooth functions: one vanishes at
zero but has nonzero integral, and another has
nonzero value at zero. Normalize their two pairs.
Their embedded images in $H_+^S$ give a linear map
$V:\mathcal P\to H_+^S$ such that
$dV=\operatorname{id}_{\mathcal P}$, where
$d(u_0,u_1)=u_0-u_1$. This uses exactly Poisson
(M24.6), including its two signs.
For any pair in $A_S$, evaluate its difference at
norms $1$ and $e$, and let $\ell(u_0,u_1)$ be the
unique element of $\mathcal P$ with those two
values. The matrix
$\left(\begin{smallmatrix}1&1\\1&e^{-1}\end{smallmatrix}\right)$
is invertible, so $\ell:A_S\to\mathcal P$ is a
bounded linear map and equals $d$ on $X_S$.
Subtracting $Vd$ makes the two entries equal;
their common function has all real Schwartz
weights. This proves (M27.8) directly.

Let $q:H_{J,s}^{U_S}\to Z_{J,s}^{U_S}$ be the
Hilbert quotient map. Inclusion $E_S\to H_{J,s}^{U_S}$
is bounded, by the exponential weights and compactness
of the norm-one factor. Therefore
\[
 B:A_S\longrightarrow Z_{J,s}^{U_S},\qquad
 B(y)=q\bigl(p_-(y-V\ell(y))\bigr)
                                                               \tag{M27.9}
\]
is bounded. It satisfies $B(i_-h)=q(h)$.
It vanishes on $H_+^S$: if $y\in H_+^S$, then
$y-V\ell(y)$ is still in $H_+^S$ and is diagonal;
its common function belongs to the theta image
by the Poisson intersection identity (M24.7),
so its Hilbert quotient class is zero.
Thus on $X_S$, (M27.9) is the actual canonical
quotient map, with an explicitly defined bounded
extension to the whole ambient pair. No assertion
about an unspecified quotient topology is used.

\emph{Transfer of the nuclear extension.}
There is a bounded diagonal inclusion
\[
 \delta:H_{J,s}^{U_S}\longrightarrow
       L^2(C_S)_\alpha\oplus L^2(C_S)_{1-\alpha},
       \qquad h\longmapsto(h,h).
                                                               \tag{M27.10}
\]
Indeed both endpoint exponential weights on the
right are bounded by the maximum of the weights
at $c,d$; the polynomial factor in (M27.1) is
at least one. The inclusion has norm at most
$\sqrt2$ for the chosen normalizations.
Compose (M27.5) with this inclusion and (M27.9):
\[
              A_f=B(P_--P_+^S)C_f\delta:
                     H_{J,s}^{U_S}\to Z_{J,s}^{U_S}.
                                                               \tag{M27.11}
\]
It is a nuclear map between Hilbert spaces,
uniformly for the bounded test sets under consideration.
To make the bornological-to-Hilbert step explicit,
Meyer's definition writes the middle map as
$\sum_j\lambda_j w_j\ell_j$, where
$\sum|\lambda_j|<\infty$, $(w_j)$ is bounded in
$A_S$ and $(\ell_j)$ is equibounded in the source
dual. The images $Bw_j$ are norm bounded.
The functionals $\ell_j\delta$ are uniformly
norm bounded. If a precompact rather than the
norm-bounded convention is used on the Hilbert
source, the same assertion follows from uniform
boundedness: equiboundedness on its compact
subsets implies pointwise boundedness of these
continuous functionals, hence a common operator
norm bound. A linear functional bounded on all
compact subsets of a Banach space is continuous,
since otherwise a sequence tending to zero
could have unbounded images. This also verifies
that no bornology convention inserts a hidden
topology assumption. The resulting rank-one
series has summable products of Hilbert norms.

On a smooth $h\in E_S$, the definitions give
$P_-C_fi_-h=i_-C_fh$, while
$P_+^SC_fi_-h\in H_+^S$ is killed by $B$.
Consequently $A_fh=q(C_fh)$ there.
Convolution by $f$ is bounded on $H_{J,s}$:
the radial weight has the translation bound
\[
 \|\lambda_b h\|_{H_{J,s}}
  \le2^{s/4}\max(|b|^{B_0},|b|^{-B_0})
       (1+|\log|b||^2)^{s/4}\|h\|_{H_{J,s}},
 \qquad B_0=\max(|c|,|d|),
\]
and $f$ is compactly supported and integrable.
It preserves the theta image and its closure,
so it induces a bounded operator $\overline C_f$
on the Hilbert quotient. Since $E_S$ is dense
in $H_{J,s}^{U_S}$, the preceding equality extends
to $A_f=\overline C_fq$ on the full Hilbert space.
The quotient has its isometric orthogonal
section $\sigma:Z_{J,s}^{U_S}\to H_{J,s}^{U_S}$.
Thus $\overline C_f=A_f\sigma$ is trace class.
The rank-one bounds just proved are uniform for
bounded test sets, and extension by zero to
$(Z_{J,s}^{U_S})^\perp$ preserves the trace norm.

\emph{The full trace and absolute character summation.}
All compact character projectors commute with
$C_f$, so (M27.6) descends to an orthogonal
sum of quotient sectors. Their operators are
exactly those of MT-M26, and their traces are
the inner sums in (M27.3). A trace-class operator
on an orthogonal sum has trace equal to the
sum of its sector traces, by computing in
orthonormal bases of the summands.
For clarity, absolute convergence of the full
\emph{zero} sum also follows directly from its
Hilbert trace norm. Take any finite collection
of the finite block subspaces of MT-M26 in
finitely many character sectors. Their direct
sum is finite dimensional and invariant under
$C_f$; its eigenvalues are exactly the indicated
Mellin values with their block multiplicities.
An orthonormal Schur basis for this invariant
subspace gives
\[
       \sum_{\text{these blocks}}m_{\chi,\rho}
            |\widehat f(\beta_\chi|\cdot|^\rho)|
                     \le\|C_f\|_1.
                                                               \tag{M27.12}
\]
Indeed choose diagonal phases so that the sum
of the absolute diagonal entries is a trace
against a finite-rank contraction; its absolute
value is at most the trace norm. Increasing the
finite set of blocks proves absolute convergence
of the whole spectral sum and the same bound.
Here are also the tails needed for uniform absolute
convergence, rather than only a uniform bound on the
total sum. A bounded test set lies in one space of
smooth functions with fixed compact support on an
elementary quotient. It is relatively compact there
by the derivative bounds and Arzel\`a--Ascoli.
The linear map from this Fr\'echet test space to
trace-class operators is continuous: its boundedness
on bounded sets, just proved, implies continuity
into a Banach space. Explicitly, failure of continuity
would give a sequence tending to zero in successively
more of the defining seminorms while its image norms
tend to infinity, contradicting bounded-set preservation.
The image of the test set is consequently relatively
compact in trace norm. If $P_F$ is a finite sum of
compact-character projections, then $P_F\to1$
strongly; hence
$\|T-P_FTP_F\|_1\to0$ uniformly for $T$ in this
relatively compact trace-class set. Pointwise this
follows by approximation by finite-rank operators;
uniformity follows from a finite trace-norm net and
the contraction property of the projections.
Applying (M27.12) on the complementary sectors bounds
all their absolute spectral terms by this tail norm.
For the remaining finitely many characters, the
fixed-sector estimates of MT-M26 give uniform
decay of the absolute sums outside finitely many
height blocks. The character projection of a
bounded test set is bounded in that sector's
$\mathcal E$, so those estimates apply. Combining
the two tails proves the asserted uniform absolute
convergence on bounded test sets.
This proves (M27.3).
\end{proof}

For the centered or otherwise twisted literal Masullo
quotient, put $J=a+I$. Multiplication by $|\cdot|^a$
is an isometry from (M27.1) to the Hilbert space with
base $I$ and denominator
$\overline{|\cdot|^a\Theta_K}$, and
$C_kM_a=M_aC_{|\cdot|^{-a}k}$.
Consequently, whenever $[0,1]\subset a+\operatorname{int}I$,
its full trace is
\[
       \sum_{\chi,\rho}m_{\chi,\rho}
                    \widehat k(\beta_\chi|\cdot|^{\rho-a}).
                                                               \tag{M27.13}
\]
For $a=1/2$ and $I=[-\alpha,\alpha]$, the unconditional
sufficient condition is $\alpha>1/2$.

At this centered wide-strip quotient, let
$h^*(g)=\overline{h(g^{-1})}$. Under $M$ every spectral
quasi-character in (M27.13) is unitary, and therefore
\[
       \operatorname{Tr}C_{h*h^*}
          =\sum_{\chi,\rho}m_{\chi,\rho}
               |\widehat h(\beta_\chi|\cdot|^{\rho-1/2})|^2\ge0.
                                                               \tag{M27.14}
\]
This uses $M$ as a condition; trace class alone does
not supply this positivity. Conversely, positivity
for every $h\in\mathcal D(C)$ would imply GRH for
every Hecke character of this fixed field: a single
off-line pair gives a negative fixed-character
Schwartz test by MT-M26; multiplying it by smooth
expanding norm cutoffs approximates it in that
sector's $\mathcal E$ topology by elements of
$\mathcal D(C)$, preserving a negative trace for
large cutoff. Continuity of the fixed-sector
trace-class map justifies this last step. No
positive-operator realization is inferred.

\paragraph{What remains after this transfer.}
This is the full adelic wide-strip result, not merely
a list of fixed-character trace formulas. The
varying-character summability is supplied by
Meyer's classical uniform nuclear extension, with
an explicit map into Masullo's specified Hilbert
quotient. The all-character test topology remains
the ordinary inductive union of finite invariance
levels; constants need not be uniform when that
level is allowed to vary without a boundedness
condition. The narrow centered intervals
$0<\alpha\le1/2$, even under $M$, do not embed
into the two Hilbert spaces in (M27.5).
MT-M26 treats each of their fixed-character
sectors; simultaneous trace-class control on
such a narrow interval remains a separate gate.

The original Masullo norm and closure locators
remain (9.2.2.31), p.~416; Definition~10.1.1.1,
p.~434; and (10.4.1.3), p.~456. The geometric
crosswalk to inherited cote~1 / SA23
geo-H-lidskii and geo-H-genus now includes the
actual full adelic Hilbert quotient on a wide
strip. Meyer's classical spectral realization
and uniform nuclearity retain their original
credit; this proof supplies the precise
comparison and completion transfer used here.

\subsection*{MT-ART11. The tame integral class is a central-parity invariant}
Let $N/K$ be a finite Galois extension of number fields, tamely
ramified at every finite place, with group $G$. All module classes
in this statement are over $\mathbb Z[G]$, after restriction of
scalars from $\mathcal O_K$. Put
\[
 U_{N/K}=[\mathcal O_N]-[\mathcal O_K[G]]
       \in\operatorname{Cl}(\mathbb Z[G]),\qquad
 m_\chi=\operatorname{ord}_{s=1/2}L(s,\chi,N/K)\in\mathbb Z.
                                                        \tag{ART11.1}
\]
Orders here are meromorphic orders, so a negative value denotes
a pole. Let $R_G^s$ be the virtual symplectic-character group and
$T(\chi)=\chi+\overline\chi$. The classical root-number map
\[
 t:\operatorname{Hom}_{\Omega_{\mathbb Q}}
       (R_G^s/T(R_G),\{\pm1\})
          \longrightarrow\operatorname{Cl}(\mathbb Z[G])
                                                        \tag{ART11.2}
\]
is additive after writing its sign-group law multiplicatively.
Then
\[
 U_{N/K}=t\bigl(\chi\longmapsto(-1)^{m_\chi}\bigr),
 \qquad 2U_{N/K}=0.                                    \tag{ART11.3}
\]
In particular,
\[
 \zeta_N(1/2)\ne0\ \Longrightarrow\ U_{N/K}=0,
 \qquad
 U_{N/K}\ne0\ \Longrightarrow\
        \operatorname{ord}_{s=1/2}\zeta_N(s)\ge2.          \tag{ART11.4}
\]
Neither implication uses Artin holomorphy or GRH. The first
means stable freeness as a $\mathbb Z[G]$-module of rank
$[K:\mathbb Q]$; it does not by itself assert a relative normal
integral basis over $\mathcal O_K[G]$.

\paragraph{Classical input.}
Taylor, \emph{On Fr\"ohlich's conjecture for rings of integers
of tame extensions}, Invent.\ Math.\ 63 (1981), 41--79,
Theorem~1, p.~41, gives $U=t(W)$ and $2U=0$.
Cassou-Nogu\`es--Taylor, Ann.\ Inst.\ Fourier 33(2) (1983),
1--17, (2.2)--(2.5), pp.~6--7, and Corollary~A, p.~10,
fix (ART11.2) and the relative normalization in (ART11.1).
The original Taylor theorem page was visually inspected;
the 1983 paper's full main proof, p.~11, was read with its
explicit upstream Gauss-sum theorems retained as citations.

\begin{proof}
The completed Artin function has a meromorphic functional
equation. If $\chi$ is symplectic it is self-dual, so
\[
 \Lambda(1/2+z,\chi)=W(\chi)\Lambda(1/2-z,\chi).
\]
Its archimedean gamma and conductor factors are holomorphic
and nonzero at $1/2$. Write its nonzero Laurent expansion as
$a z^{m_\chi}+O(z^{m_\chi+1})$, where $a\ne0$ and
$m_\chi$ may be negative. Comparing the two leading
coefficients gives $1=W(\chi)(-1)^{m_\chi}$.
The same calculation extends multiplicatively to virtual
symplectic characters. Thus the parity character in
(ART11.3) is exactly the classical root-number character;
in particular it has the Galois invariance and the
$T(R_G)$-quotient property needed in (ART11.2).
Taylor's theorem gives (ART11.3).

Here are full details of the nonvanishing step, so that no
unproved holomorphy assertion enters (ART11.4). Set
\[
             \Theta=\sum_{\chi\in\operatorname{Irr}(G)}
                         m_\chi\chi.
\]
Suppose $\zeta_N(1/2)\ne0$. For a cyclic subgroup
$C\le G$, put $F=N^C$. The abelian factorization is
\[
       \zeta_N(s)=\prod_{\eta\in\widehat C}
                       L(s,\eta,N/F).
\]
Every factor is holomorphic at $1/2$, including the trivial
factor $\zeta_F$. Their nonnegative orders sum to zero,
so each order is zero. Artin induction and Frobenius
reciprocity identify these orders with
$\langle\operatorname{Res}^G_C\Theta,\eta\rangle_C$.
Hence $\operatorname{Res}^G_C\Theta=0$ for every cyclic
$C$. Every $g\in G$ lies in $\langle g\rangle$, so
$\Theta(g)=0$ for every $g$, and irreducible-character
independence gives $m_\chi=0$ for every $\chi$.
The root-number character is therefore trivial and
$U_{N/K}=0$. This also proves that nonvanishing of
$\zeta_N$ excludes central Artin poles in this argument.

For the contrapositive, $\zeta_N$ is holomorphic at $1/2$.
Its completed functional equation has sign $+1$ and
nonzero archimedean factors there. A central zero
therefore has positive even order, proving the lower
bound two in (ART11.4).
\end{proof}

If $G$ has no irreducible symplectic characters, its
symplectic quotient $R_G^s/T(R_G)$ is zero and $U_{N/K}=0$.
For odd $|G|$, Taylor's Corollary, p.~41, with complete
proof p.~42, strengthens this to actual freeness over
$\mathbb Z[G]$, using the Eichler condition and Jacobinski
cancellation. This cancellation step is part of that
classical corollary. There is no general freeness claim
from a vanishing stable class in the preceding proof.

For use beside the positivity gate MT-ART02, there is the
following precise refinement. If all the central Artin
orders are nonnegative, then
\[
 \operatorname{ord}_{1/2}\zeta_N
   =\sum_{\chi\in\operatorname{Irr}(G)}\chi(1)m_\chi
   \ \ge\!
   \sum_{\substack{\chi\ {
m irreducible\ symplectic}\\W(\chi)=-1}}
                         \chi(1).                       \tag{ART11.5}
\]
Indeed each selected order is a nonnegative odd integer
and therefore at least one; every other term is
nonnegative. Thus an integral sign obstruction gives a
degree-weighted lower bound once the central Heilbronn
matrix is positive. Without that extra positivity, the
unconditional conclusion remains (ART11.4).

The placement is the locally free integral $K_0$ of the
finite Galois cover, in contrast with its regular rational
representation after the normal-basis theorem. It extends
the concrete integral distinction of MT-ART10 and
AG-LEFSCH-LATTICE/AG-LEFSCH-DEFECT. In the inherited
cote~16, scans~14--18, no.~7 integral-ring discussion,
it supplies an exact sign-valued invariant that survives
in the lattice but is absent from the rational module.
The associated Artin representation lies in the
SGA~6, V, \S2 and \S7, especially (7.1.1),
representation-$\lambda$-ring crosswalk fixed in AG-LAMBDA.
These are programme placements, not an assertion that
SGA~6 or the archive contains Taylor's arithmetic theorem.

\subsection*{MT-ART12. The hermitian trace class recovers every central parity}
Keep the tame extension and the right-module convention.
Equip $\mathcal O_N$ with its $G$-valued trace form
and the comparison module with its standard trace form:
\[
 \begin{split}
 h_N(m,n)&=\sum_{\gamma\in G}
                \operatorname{Tr}_{N/\mathbb Q}(m n^\gamma)
                                \gamma^{-1},\\
 h_0(m,n)&=\operatorname{Tr}_{K/\mathbb Q}(\overline m n),
       \qquad m,n\in K[G],\\
 X_{N/K}&=[\mathcal O_N,h_N]-[\mathcal O_K[G],h_0]
                   \in K_0H(\mathbb Z[G]).
 \end{split}                                             \tag{ART12.1}
\]
The bar in $h_0$ is the group-ring involution
$\gamma\mapsto\gamma^{-1}$, and the field trace is
applied coefficientwise. Here $K_0H$ is the Grothendieck
group for these locally free hermitian modules and
orthogonal sum. Its discriminant has the canonical
sign extraction $\theta$ of Cassou-Nogu\`es--Taylor,
(2.9), p.~8. Their main theorem, p.~10, and
Corollary~B, p.~11, give
\[
 \theta_\chi(X_{N/K})=W(\chi),\qquad
 \eta_\chi(X_{N/K})=
       \bigl(\operatorname{N}\mathfrak f(\chi)\bigr)^{1/2}
                                      W_\infty(\chi)
                                                        \tag{ART12.2}
\]
for every symplectic $\chi$. These exact normalizations
are (3.6)--(3.9), p.~9; the original page was inspected,
including the square root of the conductor and the bar
in $h_0$.

The consequence for the present spectral programme is
\[
       (-1)^{\operatorname{ord}_{1/2}L(s,\chi)}
                         =\theta_\chi(X_{N/K}).           \tag{ART12.3}
\]
Thus the hermitian trace class determines each
symplectic central parity individually, including
meromorphic orders, without GRH or Artin holomorphy.
The underlying locally free class is obtained by the
forgetful map and gives $t(W)$ as in MT-ART11; inversion
of $t$ is not needed in (ART12.3).

\begin{proof}
Apply the complete Laurent-series parity calculation
in MT-ART11 to each symplectic character and substitute
the classical first equality in (ART12.2). This proves
(ART12.3) with no positivity premise. If two tame
extensions over the same base have identified groups
and the same class $X$ under that identification, the
homomorphisms $\theta_\chi$ take the same values, so
all their symplectic central orders have the same
parity. Applying the forgetful map gives equal
underlying stable classes, by (ART11.3). Conversely,
this proof makes no injectivity assumption on $t$;
stable-class data alone have not been substituted
for the hermitian class.
\end{proof}

The definition of the reduced hermitian group in
the cited source uses extension to
$\overline{\mathbb Q}[G]$; the source's p.~5 and p.~10
were visually checked to retain that algebraic closure.
No isometry over $\mathbb Q[G]$ is silently inferred.
The construction supplies a precise hermitian refinement
of the integral $K_0$ comparison in MT-ART11 and of
the trace programme at cote~1. It gives root signs
and central parities, while nonnegative multiplicities
remain the distinct Heilbronn-positivity gate. The
extra conditional lower bound (ART11.5) states exactly
what that positivity would add.

\subsection*{MT-M28. Full adelic trace class on every positive-width critical strip}
Let \(K\) be a fixed number field and retain the actual summation
image and test-function conventions of MT-M24 and MT-M27. Choose
an archimedean norm splitting \(C_K=C_K^1\times\mathbf R_+^\times\)
and give \(C_K^1\) Haar volume one. For \(a\in\mathbf R\), \(s\ge0\)
and \(I=[u,v]\), \(u<v\), put
\[
 \begin{split}
 H_{I,s}&=L^2\!\left(C_K,
    \max\{|g|^{2u},|g|^{2v}\}
         (1+\log^2|g|)^{s/2}\,dg\right),\\
 Z_{a,I,s}&=H_{I,s}/
       \overline{|\cdot|^a\Theta_K}^{\,H_{I,s}},
 \qquad \Theta_K=\Phi_K(\mathcal S_0(\mathbf A_K)).
 \end{split}                                                    \tag{M28.1}
\]
Assume GRH for all unitary Hecke characters of this \(K\), as
supplied by the conditional input \(M\), and assume
\[
                         \tfrac12-a\in(u,v).                  \tag{M28.2}
\]
Every \(f\in\mathcal D(C_K)\) then induces a trace-class
convolution operator on the Hilbert quotient (M28.1). With the
character convention of MT-M22 and with each unitary extension
\(\beta_\chi\) chosen trivial on the norm splitting, its trace is
\[
 \operatorname{Tr}_{Z_{a,I,s}}C_f
   =\sum_{\chi\in\widehat{C_K^1}}\ \sum_{\Xi_\chi(\rho)=0}
          m_{\chi,\rho}\,
          \widehat f(\beta_\chi|\cdot|^{\,\rho-a}).
                                                               \tag{M28.3}
\]
The sum is absolutely convergent, uniformly on bounded test
sets, and the trace norms are bounded on those sets. At every
fixed compact-invariance and compact-support test level the
map \(f\mapsto C_f\) is continuous into trace norm.

The denominator in (M28.1) is exactly the kernel of all
character/zero jets
\[
 \partial_z^k\widehat h(\beta_\chi|\cdot|^z)
                  \big|_{z=\rho-a},\qquad 0\le k<m_{\chi,\rho}.
                                                               \tag{M28.4}
\]
Consequently the corrected full smooth quotient
\[
 |\cdot|^a\mathcal S_{\mathrm E}(C_K)\big/
           |\cdot|^a\overline{\Theta_K}^{\,\mathcal S_{\mathrm E}(C_K)}
                                                               \tag{M28.5}
\]
maps continuously and injectively into \(Z_{a,I,s}\), with
proper dense image. In particular, at \(a=1/2\), every
\(I=[-\alpha,\alpha]\) with \(\alpha>0\) is allowed.
The zero-width completion remains governed by MT-M25.

\paragraph{Classical input and scope of the uniform proof.}
The analytic input besides \(M\) is the usual Hecke
continuation and functional equation, with its finite-order
property, already used in NT-L03, NT-Z10 and NT-Z17.
The continuation/functional-equation locator is
Ramakrishnan--Valenza, \emph{Fourier Analysis on Number Fields},
Theorem~7-19, printed pp.~279--280; the real and complex
local factors are in \S7.1, printed pp.~244--245, with their
explicit integral calculations on pp.~250--252.
These local formulas and the entire supplied proof of
Theorem~7-19 were read, and the relevant local-factor pages
were inspected as original page images. We do not claim that
Theorem~7-19 separately states a uniform finite-order estimate:
only the classical finite-order property of each function is
used below, and all uniform estimates are proved here.
The source is
\url{https://webusers.imj-prg.fr/~julien.marche/Ref_GL_GT/Fourier_Analysis_on_Number_Fields.pdf}.
The exact summation-image identification is MT-M24, and
the already proved one-character Hilbert theorem is MT-M26.
Meyer's uniform nuclear lemma from MT-M27 is not needed in
the uniform estimates below.

\begin{proof}
\emph{1. Character parameters at one elementary test level.}
Fix a compact subgroup \(k\subset C_K^1\) such that
\(T=C_K^1/k\) is a compact abelian Lie group. Its dual is
\[
                  \Lambda=\widehat T\simeq\mathbf Z^d\times F,
                                                               \tag{M28.6}
\]
where \(F\) is finite. Every bounded subset of
\(\mathcal D(C_K)\) is contained in one such invariance
level, with a common compact support in \(x=\log|g|\)
and bounded derivatives of every order on \(T\times\mathbf R\).
This is the elementary-level definition of the Bruhat test
bornology/topology used in MT-M27. We prove uniform estimates
at this level, without replacing its infinitely many
archimedean characters by a finite set.

Choose generators of the free part in (M28.6), and extend
each character to \(C_K\) by making it trivial on the chosen
norm splitting. Then the extension at lattice index
\(\nu\in\mathbf Z^d\times F\) is the product of the fixed
generator extensions and one of finitely many torsion
extensions. Each fixed extension has finite conductor.
The conductor of a product divides the least common
multiple of the conductors of its factors. Hence all finite
conductors at this level divide one fixed ideal.

For the primitive Hecke function \(L_\nu\) in the convention
of MT-M22, write its completion as
\[
 \Xi_\nu(z)=Q_\nu^{z/2}
      \prod_{j=1}^{n_K}\Gamma_{\mathbf R}(z+\mu_{\nu,j})L_\nu(z),
 \qquad \Gamma_{\mathbf R}(z)=\pi^{-z/2}\Gamma(z/2),
                                                               \tag{M28.7}
\]
up to a nonzero scalar. Here \(n_K=[K:\mathbf Q]\).
For every \(\nu\ne0\) this completion is entire. The
single trivial compact character, with the chosen extension,
gives \(\zeta_K\); for it use the pole-removed completion
and the fixed-character result MT-M26 throughout. A pure
norm character trivial on the norm splitting is trivial,
so there are no other polar exceptions in this family.

The local formulas in the cited pages show that real
places contribute shifts \(\epsilon+i\tau\), with
\(\epsilon=0\) or \(1\), and complex places contribute
\[
      |m|/2+i\tau,\qquad |m|/2+1+i\tau.
 \]
Indeed
\(\Gamma_{\mathbf C}(z)=
\Gamma_{\mathbf R}(z)\Gamma_{\mathbf R}(z+1)\);
an alternative constant normalization cancels in all
ratios below. Angular integers and real norm parameters
of products grow at most linearly in the exponent vector.
Thus, writing \(H_\nu=2+|\nu|\) for any fixed lattice norm,
\[
 \begin{gathered}
  1\le Q_\nu\le Q_0,\qquad
  \operatorname{Re}\mu_{\nu,j}\in\tfrac12\mathbf Z_{\ge0},
  \qquad |\mu_{\nu,j}|\le C H_\nu .
 \end{gathered}                                                 \tag{M28.8}
\]
The constants depend only on the field and this test level.
They are independent of \(\nu\).

\emph{2. Uniform gamma ratios and finite-\(L\) bounds.}
We record the elementary gamma estimates needed to avoid
an uncanceled, possibly very large, completed-function
normalization. For \(\operatorname{Re}z\ge c>0\), the
digamma series gives
\[
 \psi(z)=-\gamma+
      \sum_{m\ge0}\left(\frac1{m+1}-\frac1{m+z}\right),
 \qquad
 |\psi(z)|\le C_c+2\log(2+|z|),
 \qquad |\operatorname{Im}\psi(z)|\le C_c.                    \tag{M28.9}
\]
Here is a direct bound. Split the series at
\(N=\lceil2|z|+2\rceil\). The finite part is bounded by
two harmonic sums, using \(|m+z|\ge m+c\); the tail is
bounded by
\(\sum_{m\ge N}|z-1|/((m+1)|m+z|)=O(1)\).
For \(z=x+iy\), the imaginary part is bounded by
\[
 |y|\sum_{m\ge0}\frac1{(m+x)^2+y^2}
 \le \frac1{2c}+\int_0^\infty
           \frac{|y|}{(t+x)^2+y^2}\,dt
 \le \frac1{2c}+\frac\pi2,
 \]
with the \(y=0\) case immediate.
Consequently a bounded change in real part, within a
right half-plane, changes \(|\Gamma_{\mathbf R}|\) by
at most a fixed power of \(2+|z|\), in either direction.
A change of imaginary part by \(h\) changes it by at
most \(e^{C_c|h|}\). These statements follow by
integrating the real or imaginary derivative of
\(\log|\Gamma_{\mathbf R}|\); the half-plane is free
of gamma poles and zeros.

Fix \(A>0\). The Euler product gives
\[
 \zeta_K(2)^{-1}\le |L_\nu(2+it)|\le\zeta_K(2),
 \qquad |L_\nu(A+2+it)|\le\zeta_K(2).
                                                               \tag{M28.10}
\]
The lower bound follows factor by factor from
\((1+x)^{-1}\ge1-x\), \(0\le x<1\), including the
omitted ramified factors. On \(\sigma=-A\), the
functional equation expresses \(L_\nu(\sigma+it)\)
as \(L_{\bar\nu}(1-\sigma-it)\), a bounded conductor
power and a product of ratios
\[
  \frac{\Gamma_{\mathbf R}(1-\sigma+
                   \operatorname{Re}\mu_{\nu,j}
                   +i(t+\operatorname{Im}\mu_{\nu,j}))}
       {\Gamma_{\mathbf R}(\sigma+\mu_{\nu,j}+it)}
                                                               \tag{M28.11}
\]
in absolute value. Conjugation was used in the
numerator; the root number has modulus one.
Use
\(\Gamma_{\mathbf R}(z+2)=
 (z/(2\pi))\Gamma_{\mathbf R}(z)\)
to move the denominator to a fixed right half-plane.
The reciprocal of the original denominator becomes
the reciprocal of that shifted gamma function times
a polynomial of fixed degree. This remains valid
at poles of the original denominator, by analytic
continuation. Now (M28.9) and (M28.8) bound every
ratio by a fixed power of \(H_\nu+|t|\).
It follows that both boundary lines of the strip
\([-A,A+2]\) satisfy such a polynomial bound.

For explicitness, the Phragm\'en--Lindel\"of step is
uniform despite using only finite order separately
for each \(\nu\). Divide \(L_\nu(z)\) by
\((H_\nu+C_A+z)^B\), with integer \(B\) sufficiently
large and \(C_A>A+2\). This denominator has no zero
in the strip, and on its boundaries its modulus
is comparable with \((H_\nu+|t|)^B\). The quotient
is uniformly bounded on both boundaries and has
finite order for each \(\nu\). Multiply it by
\[
          \exp\{-\epsilon\cos(\lambda(z-z_0))\},
 \]
where \(z_0\) is the real midpoint of the strip and
\(\lambda>0\) is small enough that
\(\cos(\lambda(\sigma-z_0))>c>0\) throughout it.
The multiplier has modulus at most one on the
vertical boundaries and tends to zero faster than
every finite-order growth on horizontal boundaries
as their heights tend to infinity. The maximum
principle on rectangles, followed first by their
height tending to infinity and then by
\(\epsilon\downarrow0\), gives the same uniform
bound throughout the strip. Thus, for every
fixed \(A\),
\[
       |L_\nu(\sigma+it)|
           \le C_A(H_\nu+|t|)^{B_A}
                 \quad (|\sigma|\le A,\ \nu\ne0).
                                                               \tag{M28.12}
\]
No GRH has been used in (M28.8)--(M28.12).

Jensen's formula, centered at \(2+iT\), now gives
for each fixed radius \(R\)
\[
  \#\{L_\nu\text{ zeros in }|z-(2+iT)|\le R\}
             \le C_R\log(H_\nu+|T|),                           \tag{M28.13}
\]
counting multiplicity. Indeed use a larger concentric
circle of radius at least \(2R\), (M28.12) for its
upper bound and the lower bound (M28.10) at the
center. A boundary avoiding zeros can be chosen in
a fixed bounded radius interval. The logarithmic
ratio in Jensen's inequality stays bounded below.
This counts trivial zeros as well as nontrivial
ones and is uniform at every height, not only
above a character-dependent exceptional range.

\emph{3. Uniform narrow block contours.}
Continue with \(\nu\ne0\); the single polar family
member is included using MT-M26 at the end of step~4.
Choose \(0<\delta<1/4\) so small that
\[
 [b_-,b_+]=[1/2-\delta,1/2+\delta]
                         \subset a+(u,v).
 \]
By the assumed GRH, all zeros of the entire
completion \(\Xi_\nu\) lie on \(\operatorname{Re}z=1/2\).
All other zeros of the finite \(L_\nu\) are the
trivial zeros canceling gamma poles, and have
nonpositive real part. Thus all nontrivial zeros
are separated by \(\delta\) from the two vertical
edges, and all trivial zeros are separated by at
least \(b_->1/4\) from every point of the narrow
closed strip.

Put \(X_{\nu,n}=H_\nu+|n|\). Count the nontrivial
zero ordinates in \([n-2,n+3]\), with multiplicity,
by \(M_{\nu,n}\). Equation (M28.13), with a fixed
sufficient radius, gives \(M_{\nu,n}\le C\log X_{\nu,n}\).
The elementary averaging argument of MT-M18 now
chooses
\[
        h_{\nu,n}\in(n+1/4,n+3/4)
 \]
avoiding every zero ordinate, such that
\[
 \prod_{\Xi_\nu(\rho)=0}
      \min(1,|h_{\nu,n}-\operatorname{Im}\rho|)^{m_\rho}
           \ge e^{-5M_{\nu,n}}\ge X_{\nu,n}^{-C}.             \tag{M28.14}
\]
For completeness, each function
\(\log^+(1/|h-\gamma|)\) has integral at most \(2\)
on an interval of length \(1/2\). Their sum
therefore has a value at most \(5M_{\nu,n}\);
the slack gives a positive-measure set when
\(M_{\nu,n}>0\). Every factor differing from one
comes from the specified ordinate window.

Let \(\Omega_{\nu,n}\) be the rectangle with
vertical edges \(b_\pm\) and horizontal edges
\(h_{\nu,n},h_{\nu,n+1}\); write
\(\Gamma_{\nu,n}=\partial\Omega_{\nu,n}\).
Its perimeter is uniformly bounded. The disk
Blaschke/Harnack proof of MT-M18, now using the
uniform bounds (M28.10)--(M28.13), gives
\[
        \inf_{w\in\Gamma_{\nu,n}}|L_\nu(w)|
                         \ge c X_{\nu,n}^{-B}.               \tag{M28.15}
\]
Here are the details needed for uniformity. Center
a disk at \(2+i\operatorname{Im}w\), with a
zero-free boundary of radius in a fixed interval
large enough that the narrow strip point lies in
the inner half disk. Remove all finite-\(L\)
zeros in this disk by its finite Blaschke product.
The remaining zero-free factor has modulus at
most \(C X_{\nu,n}^B\) on the boundary and at least
\(\zeta_K(2)^{-1}\) at the center. Harnack applied
to the logarithm of this upper bound divided by
its modulus gives a lower bound \(cX_{\nu,n}^{-B}\)
on the inner half disk. There are at most
\(C\log X_{\nu,n}\) Blaschke factors.

On a vertical edge every nontrivial factor has
distance at least \(\delta\) from its zero, and
every trivial factor distance at least \(b_-\).
Each normalized factor is bounded below by this
distance divided by twice the disk radius.
On a horizontal edge, use (M28.14) for the
nontrivial factors, and \(b_-\) for the trivial
factors. Omitting factors from the product in
(M28.14) only increases that product.
These lower bounds, each with \(O(\log X_{\nu,n})\)
factors, prove (M28.15).
The argument includes small heights and trivial
zeros whose ordinates grow with \(\nu\). It
does not absorb any such range into a
character-dependent compactness constant.

\emph{4. Uniform lifts with the completed gamma factors canceled.}
Let
\[
 J_{\nu,n}=A(\Omega_{\nu,n})/
       \{F:F^{(j)}(\rho)=0\
             \text{for every block zero and }j<m_\rho\},
                                                               \tag{M28.16}
\]
with the quotient of the supremum norm on the
closed rectangle. This is the same finite
Banach algebra as in MT-M18. Its dimension is
\[
                  d_{\nu,n}\le C\log X_{\nu,n}.                \tag{M28.17}
\]
If the block is empty its quotient is the zero
space. All assertions below then hold with
the corresponding maps zero.

Use the entire difference-kernel lift of MT-M18,
with its Gaussian parameter fixed to \(1\):
\[
 H_{\nu,n}F(z)=\frac1{2\pi i}
       \int_{\Gamma_{\nu,n}}
       F(w)\frac{\Xi_\nu(z)-\Xi_\nu(w)}
                    {\Xi_\nu(w)(z-w)}
                         e^{(z-w)^2}\,dw.                    \tag{M28.18}
\]
As proved there, it is entire, has the prescribed
jets in this block and zero full jets in all
other blocks, and depends only on the class of
\(F\) in \(J_{\nu,n}\). For each fixed \(\nu\),
the MT-M18 construction and the contour
adaptation in MT-M26 show that it is the
Mellin transform of an element of \(\mathcal E\).
We now prove a uniform estimate for that element.

For an integer \(N\ge0\), take the two inversion
lines
\[
                 \operatorname{Re}z=\pm(N+9/4).              \tag{M28.19}
\]
They lie outside all the rectangles. On them the
Cauchy integral of the \(-\Xi_\nu(w)\) term in
(M28.18) is zero. Only the ratio
\(\Xi_\nu(z)/\Xi_\nu(w)\) remains.
Move each numerator gamma factor
\(\Gamma_{\mathbf R}(z+\mu_{\nu,j})\) to a right
half-plane by a fixed even shift, if necessary.
Every recurrence denominator has real part in
\[
        \pm(N+9/4)+\tfrac12\mathbf Z+2\mathbf Z,
 \]
and therefore absolute value at least \(1/4\).
There is no gamma pole anywhere on either
entire inversion line, for any \(\nu\) or height.
The shifted numerator and the denominator
\(\Gamma_{\mathbf R}(w+\mu_{\nu,j})\) both lie
in a fixed right half-plane.

Use (M28.9) first to change the imaginary part,
then to change the real part by their bounded
difference. Uniformly for \(z=\pm(N+9/4)+it\)
and \(w\in\Gamma_{\nu,n}\), this gives
\[
 \left|
  \frac{\Gamma_{\mathbf R}(z+\mu_{\nu,j})}
       {\Gamma_{\mathbf R}(w+\mu_{\nu,j})}
 \right|
 \le C_N(H_\nu+|t|)^{B_N}
                      e^{C_N|t-\operatorname{Im}w|}.         \tag{M28.20}
\]
The fixed recurrence factors have been included
in \(C_N\). The finite conductor ratio is bounded
by (M28.8). The finite-\(L\) numerator is bounded
by (M28.12) and its denominator by (M28.15).
Thus (M28.7) gives the uniform completed ratio
\[
 \left|\frac{\Xi_\nu(z)}{\Xi_\nu(w)}\right|
 \le C_N X_{\nu,n}^{B_N}
            (H_\nu+|t|)^{B_N}
                      e^{C_N|t-\operatorname{Im}w|}.         \tag{M28.21}
\]
The normalization scalar in (M28.7) cancels
exactly. No bound for an uncanceled completed
function, and in particular no polynomial
bound for a potentially factorial-size gamma
normalization, is asserted or used.

The Gaussian in (M28.18) has modulus
\(\exp((\operatorname{Re}z-\operatorname{Re}w)^2
       -(t-\operatorname{Im}w)^2)\).
The real-part difference is bounded for fixed
\(N\), and
\(|\operatorname{Im}w-n|\le7/4\). The Gaussian
absorbs the exponential linear factor in
(M28.21) and every polynomial in \(|t-n|\).
Also \(H_\nu+|t|\le X_{\nu,n}(1+|t-n|)\).
The denominator \(|z-w|\) is bounded below
by a positive constant on the inversion lines.
Since the contour perimeter is bounded, we obtain
\[
 |H_{\nu,n}F(\pm(N+9/4)+it)|
 \le C_N X_{\nu,n}^{B_N}\|F\|_\infty
                              e^{-c(t-n)^2},                 \tag{M28.22}
\]
with \(c>0\) independent of \(\nu,n,N\), after
reducing it once.

Invert the Mellin transform on these lines.
For \(x\ge0\) use the positive line and for
\(x\le0\) the negative line. The \(j\)-th
derivative costs \(|z|^j\), whose Gaussian
integral is bounded by a constant times
\((2+|n|)^j\). Since the resulting function
is already in \(\mathcal E\) for each
individual \(\nu\), the contour shifts are
legitimate; (M28.22) now supplies their
uniform constants. Taking the infimum over
all representatives \(F\) of a fixed block
class yields linear lifts
\(\mathcal T_{\nu,n}:J_{\nu,n}\to\mathcal E\)
with
\[
 p_{N,j}(\mathcal T_{\nu,n}b)
       \le C_{N,j}X_{\nu,n}^{B_{N,j}}\|b\|_{J_{\nu,n}},
 \qquad R_{\nu,m}\mathcal T_{\nu,n}=\delta_{mn}\operatorname{id}.
                                                               \tag{M28.23}
\]
Here \(p_{N,j}(h)=\sup_x e^{N|x|}|h^{(j)}(x)|\).
The last identity means the identity in the
matching block and zero in the others.
The exceptional \(\nu=0\) has the corresponding
fixed-character polynomial bounds by MT-M26;
enlarging the constants once includes it.
These are stable quotient-block estimates,
with no separation condition on zeros inside
one rectangle.

\emph{5. The actual quotient and its full zero-jet kernel.}
The compact character projectors are orthogonal
contractions on \(H_{I,s}\). The denominator
of (M28.1) is invariant under them: MT-M24
identifies \(\Theta_K\) with the intersection
of the two summation spaces, and compact
averaging preserves that intersection. Applying
a projector to an approximating sequence
shows that its intersection with one character
sector is exactly the Hilbert closure of the
corresponding sector image.

Apply MT-M26 in each sector, using (M28.2).
It identifies that closure with all the
sector zero jets. Every such jet is a
continuous functional on the full Hilbert
space, by the compact projection and the
interior-strip Cauchy--Schwarz estimate.
Finite character truncations converge
strongly in the Hilbert direct sum.
Therefore the closed denominator is exactly
the simultaneous kernel (M28.4).

The full smooth closed-kernel identity of
MT-M27 was proved by compact Fourier
truncation and MT-M24, independently of
the Hilbert strip width. It continues to
apply here. It proves faithfulness and
continuity of (M28.5); compactly supported
smooth radial functions and finite
character sums prove density. If this
map were onto, applying any one compact
character projector would contradict
the proper-image assertion of MT-M26.
This proves the stated proper dense
smooth inclusion without equating
the smooth and Hilbert topologies.

\emph{6. Trace-norm summation over characters and blocks.}
Let \(\pi:H_{I,s}\to Z_{a,I,s}\) be the quotient map.
The block restriction map, after compact
character projection, is
\[
  r_{\nu,n}[h]
       =\left[w\longmapsto
          \widehat h(\beta_\nu|\cdot|^{\,w-a})\right]
           \in J_{\nu,n}.
 \]
The real parts \(w-a\) lie in a fixed compact
subinterval of \((u,v)\). Cauchy--Schwarz
against the radial weight, and the fact
that the compact character has modulus
one, therefore give
\[
                \|r_{\nu,n}\|\le C                         \tag{M28.24}
\]
uniformly in both indices. As in MT-M26,
this follows by bounding the quotient
norm by the supremum of the displayed
representative on the contour.

For \(b\in J_{\nu,n}\), define
\[
  s_{\nu,n}b
       =\pi\left[
          \beta_\nu(g)^{-1}|g|^a
               (\mathcal T_{\nu,n}b)(\log|g|)
              \right].
 \]
A single fixed seminorm \(p_{N,0}\), with
\(N>|a|+\max(|u|,|v|)+1\), bounds its
Hilbert norm. Equation (M28.23) gives
\[
 \|s_{\nu,n}\|\le C X_{\nu,n}^B,\qquad
 r_{\mu,m}s_{\nu,n}
      =\begin{cases}\operatorname{id},&(\mu,m)=(\nu,n),\\
                      0,&\text{otherwise}.
        \end{cases}                                          \tag{M28.25}
\]
No compact-direction derivatives occur
in this Hilbert norm, so no extra
character factor is concealed here.

For a test \(f\) invariant under \(k\), let
\[
 f_\nu(x)=\int_{C_K^1}
          f(c\,\tau(e^x))\,\beta_\nu(c)\,dc,\qquad
 b_{\nu,n}(f)=
       \sup_{w\in\Gamma_{\nu,n}}
                      |\mathcal M f_\nu(w-a)|.
 \]
For every \(A>0\),
\[
                   b_{\nu,n}(f)\le C_A(f)X_{\nu,n}^{-A},       \tag{M28.26}
\]
where \(C_A(f)\) is bounded by finitely
many seminorms at this test level and
is uniformly bounded on bounded test
sets. To see this directly, integration
by parts on the torus, using
\((1-\Delta_T)^r\), gives arbitrary
negative powers of \(1+|\nu|\).
Integration by parts in \(x\), on the
common compact support, gives arbitrary
negative powers of \(1+|\operatorname{Im}w|\),
uniformly for the real part in the
fixed contour interval. The derivatives
of \(e^{(\operatorname{Re}w-a)x}f\)
are bounded by finitely many test
seminorms. Since
\(|\operatorname{Im}w-n|\le7/4\),
combining these two arbitrary powers
gives (M28.26). The finite group in
(M28.6) only changes a constant.

Let \(M_{f,\nu,n}\) be multiplication
in \(J_{\nu,n}\) by
\([\mathcal Mf_\nu(w-a)]\). Its norm
is at most \(b_{\nu,n}(f)\), by the
submultiplicative quotient algebra
norm. The finite-rank operator
\[
                 A_{\nu,n}(f)
                    =s_{\nu,n}M_{f,\nu,n}r_{\nu,n}
 \]
has rank at most \(d_{\nu,n}\). Hence,
as an operator on the Hilbert quotient,
\[
 \|A_{\nu,n}(f)\|_1
       \le C\,d_{\nu,n}X_{\nu,n}^B b_{\nu,n}(f).
                                                               \tag{M28.27}
\]
This is the elementary Hilbert estimate
\(\|A\|_1\le\operatorname{rank}(A)\|A\|\);
the intervening finite-dimensional
spaces need not carry Hilbert norms.
Equations (M28.17) and (M28.26) show
that
\[
                  \sum_{\nu\in\Lambda}\sum_{n\in\mathbf Z}
                                      A_{\nu,n}(f)            \tag{M28.28}
\]
converges in trace norm. For example,
take \(A>B+d+3\) in (M28.26), since
the lattice plus the height index
has dimension \(d+1\). Its tails
converge uniformly on bounded test
sets, and the resulting trace norm
is bounded by finitely many test
seminorms.

This is the actual convolution, not
only a formal operator with the
same eigenvalues. Convolution by
\(f\) is bounded on \(H_{I,s}\):
the translation estimate (M26.10)
integrated over its compact support
proves this. It preserves the closed
denominator, since
\[
 C_f(|\cdot|^a h)
       =|\cdot|^a C_{|\cdot|^{-a}f}h
 \]
and the summation image is a
convolution submodule.
On a single-character smooth vector
the operator (M28.28) equals this
convolution by the already proved
one-character synthesis and dense
extension in MT-M26. On characters
not in \(\Lambda\), both operators
are zero: \(k\)-invariance of \(f\)
makes its Fourier coefficient zero.
Finite-character smooth vectors
are dense in the full Hilbert
quotient. Both operators are
bounded, so their equality extends
to that quotient. This completes
the full adelic trace-class proof.

\emph{7. Trace, uniform convergence, and the analytic scope.}
By (M28.25), the finite-rank trace
of \(A_{\nu,n}(f)\) is the trace
of multiplication by
\([\mathcal Mf_\nu(w-a)]\) in
\(J_{\nu,n}\). The local algebra
decomposes as the direct sum of
\(\mathbf C[z]/(z-\rho)^{m_\rho}\)
over its block zeros. Multiplication
is triangular in the power bases,
with diagonal value
\(\mathcal Mf_\nu(\rho-a)\) repeated
\(m_\rho\) times. Thus this trace
is the corresponding finite part
of (M28.3). Trace-norm convergence
allows summation of these traces.
The individual spectral terms
are absolutely summable as well:
\[
 \sum_{\rho\text{ in block}}
    m_\rho|\mathcal Mf_\nu(\rho-a)|
            \le d_{\nu,n}b_{\nu,n}(f),
 \]
by the maximum principle on the
block. Equations (M28.17) and
(M28.26) give uniform tails of
this double series directly.
This proves all the asserted
trace and continuity properties.

For a specified bounded test set
at the level \(k\), its trace-class
and trace-formula conclusions use
GRH only for the characters in
\(\widehat{C_K^1/k}\). Other
character sectors are annihilated
by these tests. The global
faithfulness statement (M28.4)
uses all character sectors.
Thus the full conditional input
\(M\) is sufficient, while the
operator conclusion at one level
has the stated smaller Hecke
family as its precise analytic
input.

At \(a=1/2\), define
\(h^*(g)=\overline{h(g^{-1})}\).
The group is abelian and
unimodular. Under the assumed
GRH, every
\(\beta_\chi|\cdot|^{\rho-1/2}\)
is unitary, and (M28.3) gives
\[
       \operatorname{Tr} C_{h*h^*}
          =\sum_{\chi,\rho}
                 m_{\chi,\rho}
                 |\widehat h(
                    \beta_\chi|\cdot|^{\,\rho-1/2})|^2
          \ge0.                                               \tag{M28.29}
\]
This is a trace identity. It
does not identify the group
involution with the operator
adjoint for the specified
Hilbert inner product.
All higher zero jets remain
present in (M28.4), consistently
with the obstruction in MT-M20
to a faithful positive
full-jet realization when
multiple zeros occur.
\end{proof}

\paragraph{Run and geometric crosswalk.}
This is a new uniform transfer proof in
\texttt{GRH\_NIGHT/complexity}, following the fixed-character
MT-M18/MT-M24/MT-M26 and the full wide-strip MT-M27.
Its literal source targets are Masullo
Definition~9.2.2.31, original PDF p.~416, the Hilbert
closure in Definition~10.1.1.1, p.~434, and the
strip quotient in Definition~10.4.1.3, p.~456,
with the centering twist from MT-M24. It closes
the all-character positive-width strip gate
conditionally on \(M\); it is not a deduction
of \(M\). In the attached geometric programme,
it supplies an actual trace-class realization
for the idele/test-function quotient interface
of the \texttt{cote1}/\texttt{SA23geo-H-lidskii}
crosswalk. The positivity statement remains
the conditional trace statement (M28.29);
the stronger faithful positive-inner-product
gate is still exactly MT-M20.

\subsection*{MT-NB01. Quantitative approximation after the adopted Riemann hypothesis}

The qualitative Nyman--Beurling and B\'aez-Duarte criteria are
already background in \texttt{RETURN\_ALPOGE\_BRIDGES\_cp02},
CC-P01--CC-P09. Here the additional issue is the rate, including
the contribution of multiple zeros. Put
\[
 d_N^2=\inf_{a_1,\ldots,a_N\in\mathbf C}
 \frac1{2\pi}\int_{\mathbf R}
 \left|1-\zeta(\tfrac12+it)
       \sum_{n\le N}a_n n^{-1/2-it}\right|^2
                 \frac{dt}{1/4+t^2}.
 \tag{NB01.1}
\]
The adopted input $A_D$ and the classical theorems below give,
for every $\epsilon>0$,
\[
 d_N^2\ll_\epsilon
 \frac{(\log\log N)^{5/2+\epsilon}}{\sqrt{\log N}},
 \qquad
 \liminf_{N\to\infty}d_N^2\log N
       \ge C_\zeta^{(2)}
       :=\sum_{\rho\ {
m distinct}}
                    \frac{m_\rho^2}{|\rho|^2}.
 \tag{NB01.2}
\]
The upper bound is Balazard--de Roton,
\emph{Sur un crit\`ere de B\'aez-Duarte pour l'hypoth\`ese
de Riemann}, Int. J. Number Theory 6(4) (2010), 883--903,
unnumbered main Theorem in \S1; author version
\texttt{arXiv:0812.1689v1}, p.~3. Its constructive choice
appears in \S2, Propositions~1--3, author pp.~4--5:
one can take
\[
 a_n=\mu(n)n^{-\varepsilon_N},\qquad
 \varepsilon_N=25(\log\log N)^{5/2+\epsilon}
                         (\log N)^{-1/2}
 \tag{NB01.3}
\]
for all sufficiently large $N$. The implied constants are
effective in the source's convention; no numerical constant
or numerical threshold is extracted here.
The lower bound is Burnol, \emph{A lower bound in an
approximation problem involving the zeros of the Riemann
zeta function}, Adv. Math. 170 (2002), 56--70,
Theorem~1.3, author version \texttt{math/0103058v2}, p.~4.
That theorem uses continuous dilations $\lambda\le\alpha\le1$;
the discrete space in (NB01.1) is a subspace at $\lambda=1/N$,
so its distance is at least the continuous distance.
The Mellin--Plancherel identification is stated and proved in
Balazard--de Roton, Proposition~1, author pp.~4--5.

The familiar sum over zeros with their ordinary multiplicities
is, under $A_D$,
\[
 C_\zeta^{(1)}:=\sum_\rho\frac{m_\rho}{|\rho|^2}
                    =2+\gamma-\log(4\pi).
 \tag{NB01.4}
\]
An exact locator for the identity in this normalization is
Balazard--de Roton, author p.~3, immediately before (3),
and Bettin--Conrey--Farmer, author p.~2.
Thus $C_\zeta^{(2)}\ge C_\zeta^{(1)}$, with equality precisely
when every nontrivial zero is simple. The adopted RH input
does not remove this distinction.

For comparison, Bettin--Conrey--Farmer,
\emph{An optimal choice of Dirichlet polynomials for the
Nyman--Beurling criterion}, Proc. Steklov Inst. Math.
280, Suppl.~2 (2013), S30--S36, Theorem~1,
\texttt{arXiv:1211.5191v1}, p.~2, proves the sharper conclusion
\[
 d_N^2\sim\frac{2+\gamma-\log(4\pi)}{\log N}
 \tag{NB01.5}
\]
under RH and the additional hypothesis
\[
 \sum_{|\operatorname{Im}\rho|\le T}
                    |\zeta'(\rho)|^{-2}
                  \ll T^{3/2-\delta}\quad\text{for some }\delta>0.
 \tag{NB01.6}
\]
The hypothesis includes simplicity. Their explicit polynomial is
$V_N(s)=\sum_{n\le N}(1-\log n/\log N)\mu(n)n^{-s}$.
The theorem gives this polynomial's error asymptotic; Burnol's
lower bound then gives (NB01.5) for the infimum. These are
classical results with their exact remaining inputs, not new
proofs of the cited analytic estimates.

\subsection*{MT-NB02. The multiplicity-sensitive Hecke approximation constant}

Let $K$ be a number field, and let $F(s)=L_K(s,\eta)$ be a
primitive unitary Hecke $L$-function. Allow $\eta=1$, but exclude
the nontrivial pure norm characters $|\cdot|^{it}$, whose pole
would lie away from $1$. For such a pure norm character the
following statement instead applies to its vertical translate
$\zeta_K(s)$. Write $F(s)=\sum_{n\ge1}a_n n^{-s}$ and define
\[
 \Psi_F(x)=\operatorname*{Res}_{s=1}
                   \frac{x^s F(s)}s-\sum_{n\le x}a_n,
 \qquad
 B_F^\lambda=\operatorname{span}\{
                  t\longmapsto\Psi_F(\alpha/t):
                               \lambda\le\alpha\le1\}.
 \tag{NB02.1}
\]
Assume GRH for this particular $F$, supplied by $A_D$ for a
primitive Dirichlet function and by the conditional input
$M_H$ in the general Hecke case. Then $\Psi_F(1/t)$ belongs
to $L^2((0,\infty),dt)$, the distance
\[
 D_F(\lambda)=\operatorname{dist}_{L^2(0,\infty)}
                      (\mathbf1_{(0,1)},B_F^\lambda)
 \tag{NB02.2}
\]
tends to zero, and
\[
 \liminf_{\lambda\downarrow0}D_F(\lambda)^2\log(1/\lambda)
    \ge C_F^{(2)}:=\sum_{\rho\ {
m distinct}}
                          \frac{m_\rho^2}{|\rho|^2}<\infty.
 \tag{NB02.3}
\]
The zeros in this formula are the nontrivial zeros of $F$,
all on the critical line under the stated input.
The precise classical theorem being specialized is
Anne de Roton, \emph{Une approche hilbertienne de
l'hypoth\`ese de Riemann g\'en\'eralis\'ee}, Bull. Soc.
Math. France 134(3) (2006), 417--445: Definition~1 p.~418,
Definition~2 and Theorem~1 p.~419, Theorem~2 p.~420.
The lower-limit argument is \S7, pp.~441--444, particularly
Propositions~7.1--7.2 and \S7.2. The original p.~420 image
and that complete final proof were read. Earlier technical
operator estimates remain cited prerequisites.

\begin{proof}[Complete proof of the stated Hecke specialization]
We check the source's Selberg-class hypotheses, rather than
assuming that an arbitrary meromorphic Artin quotient satisfies
them. Grouping the ideal Dirichlet series by absolute norm gives
the displayed coefficients. If $n_K=[K:\mathbf Q]$, their
absolute values are bounded by $d_{n_K}(n)$: at each rational
prime the local ideal-counting coefficients are bounded by those
of $(1-p^{-s})^{-n_K}$. Consequently $a_n\ll_\epsilon n^\epsilon$.
The unitary Euler factors have roots of modulus at most one.
For a rational prime $p$, the logarithmic coefficient at $p^r$
is a sum over prime ideals of residue degree $f\mid r$,
with each summand bounded by $f/r$. Since the sum of the
residue degrees is at most $n_K$, its absolute value is at
most $n_K/r$. All other logarithmic coefficients vanish.
This verifies the source's Euler-product bound with exponent
zero, strictly smaller than $1/2$.

The classical Hecke continuation, finite-order property and
functional equation supply the remaining axioms; see the exact
Ramakrishnan--Valenza Theorem~7-19 and local-factor locators
in MT-M28. The real gamma shifts have nonnegative real part,
and the complex shifts have real part $|m|/2\ge0$.
There is no pole for a non-pure-norm character. The allowed
trivial character gives $\zeta_K$, with its single pole at $1$.
Thus the permitted functions belong to the precise Selberg
class of de Roton's Definition~1. Her Theorem~1, with the
specified GRH, gives the $L^2$ assertion and convergence of
the distance defined in (NB02.1)--(NB02.2). Her Theorem~2
then gives the lower bound; squaring a nonnegative lower-limit
inequality yields (NB02.3).

For clarity, convergence of the multiplicity-square sum also
follows directly from the fixed-function local zero bound in
NT-Z10. In every unit height interval the sum of multiplicities
is $O_F(\log(2+|T|))$. The sum of their squares is at most
the square of that sum. The contribution of intervals of
height at least two is therefore bounded by a constant times
$\sum_{j\ge2}j^{-2}\log^2j$, which converges. The finitely many
bounded-height zeros cause no problem, since a critical-line
zero has $|\rho|\ge1/2$.
\end{proof}

\subsection*{MT-NB03. Integral Galois classes force an approximation penalty}

For a function in MT-NB02 set
\[
 C_F^{(1)}=\sum_\rho\frac{m_\rho}{|\rho|^2},\qquad
 \Delta_F=\sum_\rho\frac{m_\rho(m_\rho-1)}{|\rho|^2}.
 \tag{NB03.1}
\]
Then $C_F^{(2)}=C_F^{(1)}+\Delta_F$, and
$\Delta_F=0$ precisely when all its nontrivial zeros are simple.
For a tame finite Galois extension $N/K$ put
$U_{N/K}=[\mathcal O_N]-[\mathcal O_K[G]]$
in the locally free class group of $\mathbf Z[G]$, with the
restriction-of-scalars convention of MT-ART11. If
$U_{N/K}\ne0$, then, conditionally on GRH for $\zeta_N$,
\[
 \begin{split}
 \liminf_{\lambda\downarrow0}D_{\zeta_N}(\lambda)^2
                      \log(1/\lambda)&\ge C_{\zeta_N}^{(2)},\\
 C_{\zeta_N}^{(2)}&\ge16,\qquad
 C_{\zeta_N}^{(2)}\ge C_{\zeta_N}^{(1)}+8.
 \end{split}                                                   \tag{NB03.2}
\]
The same constants apply to the explicit tame quaternionic
field of MT-ART10, directly from its proved central order
$r\ge2$; this use does not identify nonfreeness with a nonzero
stable class. Independently of $U_{N/K}$, if $N/K$ is nonabelian
Galois, then $\Delta_{\zeta_N}>0$ under that same GRH input.
Thus a multiplicity-one constant is strictly too small in
all these nonabelian examples.

\begin{proof}
All the sums converge by the argument in MT-NB02. The identity
in (NB03.1) follows term by term from $m^2=m+m(m-1)$.
Every summand of $\Delta_F$ is nonnegative, and it vanishes
exactly when its positive integer multiplicity equals one.

MT-ART11 proves unconditionally, without Artin holomorphy,
that $U_{N/K}\ne0$ forces an even positive central order
$r=\operatorname{ord}_{s=1/2}\zeta_N(s)\ge2$. Its proof uses
Taylor's root-number class theorem and the cyclic-restriction
argument recorded there. The single central zero contributes
$4r^2\ge16$ to $C_{\zeta_N}^{(2)}$ and
$4r(r-1)\ge8$ to $\Delta_{\zeta_N}$. This proves (NB03.2)
by MT-NB02. The arithmetic central-zero assertion is
unconditional; GRH enters to obtain the approximation-space
conclusion and place every nontrivial zero in that formula.

For any nonabelian $N/K$, the classical
Hu--Kaneko--Martin--Schildkraut Theorem~1.1, already located
and used in MT-M23, gives infinitely many multiple nontrivial
zeros of $\zeta_N$. Under the present GRH each contributes a
strictly positive term to $\Delta_{\zeta_N}$; even one would
suffice. Hence the claimed strict inequality follows. No
uniform numerical lower bound is asserted in this last case.
\end{proof}

\subsection*{MT-NB04. A precise derivative exponent sufficient for the optimal rate}

Suppose in addition to $A_D$ that, for some $c>0$ and
$0\le\beta<1/4$, every nontrivial zero satisfies
\[
             |\zeta'(\rho)|\ge c(1+|\operatorname{Im}\rho|)^{-\beta}.
 \tag{NB04.1}
\]
Then all the zeros are simple, (NB01.6) holds for some
$\delta>0$, and the optimal asymptotic (NB01.5) follows.

\begin{proof}
The displayed positive lower bound implies simplicity. The
classical zero-count bound $N(T)=O(T\log(2+T))$, with the
same multiplicity convention as NT-Z10, gives for $T\ge2$
\[
 \sum_{|\operatorname{Im}\rho|\le T}|\zeta'(\rho)|^{-2}
                 \ll_{c,\beta}T^{1+2\beta}\log T.
\]
Take $\delta=1/4-\beta>0$. Since $\log T\ll_\delta T^\delta$,
the right side is $O(T^{5/4+\beta})=O(T^{3/2-\delta})$.
Bettin--Conrey--Farmer Theorem~1 gives the error asymptotic
for $V_N$, and Burnol gives the matching lower bound for
the infimum, as in MT-NB01. This proves the conclusion.
\end{proof}

This sufficient exponent condition is stronger than the
unspecified polynomial reciprocal-derivative bound forced by
the point-jet basis in MT-M11. It is not a consequence here
of $A_D$, of the repaired adapted-block basis, or of simple
zero locations alone. It makes the additional quantitative
condition needed by that particular route explicit.

\paragraph{Geometric and run placement of MT-NB01--04.}
The Hecke functions are the rank-one adelic objects of Tate's
thesis; for $\zeta_N$ the object is the pushforward of the unit
Artin motive along $\operatorname{Spec}N\to\operatorname{Spec}K$.
The functions $\Psi_F$ turn its coefficient-counting distribution
into a dilation subspace. The zero divisor, including its
orders, obstructs approximation by that subspace. In de Roton's
proof each order-$m$ zero supplies a full $m$-jet Gram block;
the inverse block's first diagonal entry is $m^2$ (\S7.2,
p.~444). This is the analytic reason that discarding higher
jets loses the constants in MT-NB03, consistent with MT-M20
and MT-M25.

For the archive programme, this is an approximation-theoretic
addition to the trace/Fredholm interface of cote~1, scans
79--80 and 111--112, inherited through
\texttt{OURS\_INTAKE\_RESEARCH\_SA23\_20261120.md} \S4.1,
\texttt{geo-H-lidskii}/\texttt{geo-H-genus}; the sheets themselves
are not cited as Nyman--Beurling theorems. MT-NB03 also uses
the integral $K_0$ comparison of MT-ART11 and
AG-LEFSCH-LATTICE/DEFECT, placed at cote~16 scans~14--18,
no.~7, and SGA6 Expos\'e~V \S2 and \S7 (7.1.1), with the
same inherited-locator qualification as MT-ART11. The precise
run extension is the quantitative completion of the qualitative
CC-P01--CC-P09 bridge, with the negative derivative moment
(NB01.6) retained as the remaining optimal-rate condition.

\subsection*{MT-M29. The determinant of the graded arithmetic spectral module}
Fix one unitary Hecke character and its pole-removed entire
completion \(\Xi\). Work first in the centered, actual wide-strip
Hilbert quotient of MT-M24/MT-M26, with \(a=1/2\) and
\([0,1]\subset1/2+\operatorname{int}I\). Its definition by adelic
summation and Hilbert closure is unconditional. Write \(U_y\)
for norm translation by \(\tau(e^y)\), and \(D\) for its
infinitesimal generator. Its dense smooth quotient is the
centered version of \(Q_\eta\).

There is a canonical algebraic submodule \(V_{\mathrm{fin}}\):
the vectors contained in finite-dimensional \(U_y\)-invariant
subspaces. It is dense in the smooth quotient and in the
specified Hilbert quotient. Its primary decomposition and
canonical filtration are
\[
 V_{\mathrm{fin}}=\bigoplus_\rho V_\rho,\qquad
 V_\rho\simeq\mathbf C[u]/(u^{m_\rho}),\qquad
 D|_{V_\rho}=(\rho-\tfrac12)+u,\qquad
 F^jV_\rho=(D-\rho+\tfrac12)^jV_\rho .
                                                               \tag{M29.1}
\]
The isomorphism uses the local Mellin coordinate
\(u=z-(\rho-1/2)\). Form the associated graded
\[
       V_{\mathrm{gr}}
          =\bigoplus_\rho\bigoplus_{j=0}^{m_\rho-1}
                         F^jV_\rho/F^{j+1}V_\rho.             \tag{M29.2}
\]
This construction is canonical and functorial at each finite
spectral stage for maps intertwining norm translation. It
retains \(m_\rho\) dimensions at \(\rho\); it replaces the
nilpotent extensions by scalar actions.

Assume now GRH for this function. It is supplied by \(A_D\)
for a Dirichlet function over \(\mathbf Q\), or by the
conditional input \(M\) for the general Hecke character.
Give the one-dimensional graded lines in (M29.2) the metric
in which the class of \(u^j\) has norm one, and make them
orthogonal. The resulting Hilbert completion
\(\mathcal H_{\mathrm{gr}}\) carries the unitary norm action
\[
        U_y^{\mathrm{gr}}=e^{iyT},\qquad
        T|_{(\rho,j)}=\gamma,\qquad \rho=\tfrac12+i\gamma.
                                                               \tag{M29.3}
\]
Every integrated \(\mathcal E\)-test operator is trace class,
with the same trace as on the actual wide-strip quotient:
\[
 \operatorname{Tr}_{\mathcal H_{\mathrm{gr}}}C_f^{\mathrm{gr}}
       =\sum_\rho m_\rho\,\mathcal Mf(\rho-\tfrac12)
       =\operatorname{Tr}_{Z_{1/2,I,s}} C_f .                  \tag{M29.4}
\]
The grading and this new Hilbert completion change the
representation when a multiple zero occurs; they do not
equip the original Jordan blocks with a positive invariant
metric.

Suppose in addition that
\[
 \Xi(\tfrac12+z)=\epsilon\Xi(\tfrac12-z),\qquad
       \epsilon\in\{1,-1\}.
                                                               \tag{M29.5}
\]
Delete the central spectral block, then take the
positive-frequency spectral subspace
\(\mathcal H_+=1_{(0,\infty)}(T)\mathcal H_{\mathrm{gr}}\).
On it,
\[
                          B=T_+^{-2}                         \tag{M29.6}
\]
is a bounded positive trace-class operator. If
\(r=\operatorname{ord}_{s=1/2}\Xi(s)\), set
\[
 H(t)=t^{-r}\Xi(\tfrac12+it),\qquad
              \mathfrak D(w)=H(\sqrt w)/H(0).
 \]
Then the operator from the actual arithmetic module has
\[
                       \det(1-wB)=\mathfrak D(w).             \tag{M29.7}
\]
This is the AG-DET operator, now identified as a functional
calculus of the norm generator on the graded finite
spectral module. Its numerical multiplicities are precisely
the original zero orders.

\begin{proof}
\emph{The finite spectral core comes from the actual quotient.}
MT-M18 and MT-M24 identify the smooth quotient with its full
zero jets. For a finite set \(S\) of zeros, there is a unique
class with any prescribed jets at \(S\) and zero jets at
every other zero: its existence follows from the finitely
many block lifts of MT-M18, and uniqueness from the exact
closed-ideal equality. These classes form a finite-dimensional
subalgebra \(V_S\), canonically isomorphic to
\[
                     \prod_{\rho\in S}
                              \mathbf C[u]/(u^{m_\rho}).
                                                               \tag{M29.8}
\]
It has a canonical unit \(e_S\), whose local class is one
at \(S\) and zero elsewhere. Multiplication by \(e_S\) is
the finite-jet quotient projector. Thus \(V_S\) is both
a subspace and a finite quotient of the actual smooth
arithmetic quotient; no choice of interpolation
representatives survives in (M29.8).

These finite projectors extend boundedly to the actual
wide-strip Hilbert quotient: its finite jet restrictions
are continuous by MT-M26, and the finitely many lifted
classes have finite Hilbert norm. This proves that \(V_S\)
is a genuine finite spectral submodule there as well.
The finite unions of the MT-M18 blocks exhaust all zeros,
and their partial sums converge in the smooth quotient.
Hence \(\bigcup_S V_S\) is dense in that quotient and,
by its dense inclusion, in the Hilbert quotient.

Norm translation acts in the Mellin transform by
multiplication by \(e^{zy}\) in the centered variable.
Consequently its action on the \(\rho\)-factor in (M29.8)
is multiplication by
\[
                  e^{(\rho-1/2+u)y}
                   =e^{(\rho-1/2)y}e^{uy}.
 \]
Differentiating this finite matrix exponential gives
the generator in (M29.1).

It remains to verify that this core is intrinsic, rather
than merely a specified collection of finite subspaces.
Let \(W\) be any finite-dimensional invariant subspace of
the Hilbert quotient. The continuous one-parameter group
on \(W\) is a matrix exponential; write \(D_W\) for its
matrix generator, and let \(P\) be its minimal polynomial.
Every finite block restriction intertwines the norm
group, and therefore intertwines its generator on \(W\).
On a local factor with \(P(\rho-1/2)\ne0\),
\[
                 P((\rho-1/2)+u)
 \]
is invertible, because its constant term is nonzero.
Thus the restriction of every vector of \(W\) to this
factor is zero. Only the finitely many zeros whose
centered parameters are roots of \(P\) can remain.
If \(S\) is that finite set, subtracting the finite
projector \(e_S\) kills all jets of a vector in \(W\).
The exact Hilbert kernel statement of MT-M26 forces
that difference to be zero. Therefore \(W\subset V_S\).
Conversely every \(V_S\) is finite dimensional and
invariant, so \(V_{\mathrm{fin}}=\bigcup_S V_S\).
This establishes the intrinsic description.

\emph{The filtration and its functoriality.}
For the primary summand \(V_\rho\), put
\(N_\rho=D-(\rho-1/2)\). Its image powers
\(N_\rho^j V_\rho\) are the filtration in (M29.1);
in the coordinate algebra they are \((u^j)\).
Any finite-stage intertwining map commutes with \(D\).
It preserves the primary decomposition and sends
\(N_\rho^j V_\rho\) into the corresponding image power
in its target. It therefore induces a map on every
graded quotient, compatibly with compositions.
The inclusions as \(S\) grows have the same property.
This proves the claimed algebraic functoriality.
No boundedness assertion for arbitrary infinite-stage
intertwiners is required.

On \(F^jV_\rho/F^{j+1}V_\rho\), the nilpotent part
\(N_\rho\) acts as zero, so \(U_y\) acts by
\(e^{(\rho-1/2)y}\). Each of these quotients has
dimension one, and there are exactly \(m_\rho\)
of them. The distinguished vector \([u^j]\)
uses the local unit in (M29.8) and the specified
global Mellin coordinate; it is independent of
any chosen theta-function representative.
Declaring these vectors orthonormal defines
the metric stated above.

\emph{The Hilbert action and its trace.}
Under GRH the centered parameter is \(i\gamma\).
Thus every norm-translation scalar on the associated
graded has modulus one. The Hilbert completion is
explicitly
\[
             \ell^2\{(\rho,j):0\le j<m_\rho\}.
 \]
The norm group acts by the scalars \(e^{i\gamma y}\).
It is strongly continuous, first on finite-support
vectors and then on all vectors by their density
and the unit norm of the action.
Its self-adjoint frequency generator is multiplication
by \(\gamma\), with domain
\[
       \left\{(a_{\rho,j}):
                    \sum_{\rho,j}\gamma^2|a_{\rho,j}|^2
                                      <\infty\right\}.
 \]
Self-adjointness follows directly: the adjoint of
this multiplication operator has exactly that
domain, since testing against the finite-support
coordinate vectors determines its coordinates.
This proves (M29.3).

For \(f\in\mathcal E\), the integrated action is
diagonal with entry \(\mathcal Mf(i\gamma)\)
on each of the \(m_\rho\) graded lines. The transform
decreases faster than every power of \(|\gamma|\).
The fixed-function local zero count gives
\(\sum_\rho m_\rho(2+|\gamma|)^{-2}<\infty\).
Hence this diagonal operator is trace class, with
trace and trace norm obtained by summing its
entries and their absolute values, respectively.
MT-M26 gives the identical trace on the actual
wide-strip Hilbert quotient. This proves (M29.4).
The proof also gives continuity into trace norm
for the \(\mathcal E\) test topology.

The central subspace has dimension \(r<\infty\).
After it is removed, positive zero ordinates
are bounded away from zero, since zeros of a
nonzero entire function have no finite accumulation
point. Thus \(T_+^{-2}\) is bounded and positive,
with eigenvalues \(\gamma^{-2}\), one for each of
the \(m_\rho\) graded lines at each positive
ordinate. Their sum converges by the same
zero count. This proves (M29.6), including trace
class and compactness. If there are no positive
zeros the space is zero and its determinant is one.

By (M29.5), the noncentral ordinates occur in
pairs \(\gamma,-\gamma\) with equal multiplicities.
AG-DET, whose complete factorization proof
applies to this entire order-at-most-one
completion, states precisely
\[
        \mathfrak D(w)=
          \prod_{\gamma>0}(1-w/\gamma^2)^{m_{1/2+i\gamma}}.
 \]
The spectrum of (M29.6), with its just-proved
graded multiplicities, gives this same locally
uniformly convergent Fredholm product.
This proves (M29.7) and the comparison with
AG-DET. It uses no assertion about individual
Hermite vectors forming a Schauder basis.
\end{proof}

\paragraph{The ordinary trace-GNS completion is different.}
The distinction is essential when zeros have multiplicity.
On the centered convolution algebra \(Q\), let
\[
       \tau(f)=\operatorname{Tr}_{Z_{1/2,I,s}}C_f,\qquad
       f^*(x)=\overline{f(-x)}
 \]
in the fixed-character radial notation. The Hecke
functional equation makes the centered ideal invariant
under this involution. Under the same GRH input,
\[
 \langle[f],[g]\rangle_\tau
          =\tau(f*g^*)
          =\sum_{\rho}m_\rho\,
                    \mathcal Mf(i\gamma)
                        \overline{\mathcal Mg(i\gamma)} .
                                                               \tag{M29.9}
\]
Its radical is exactly the ideal of classes whose
values, without their higher derivatives, vanish at
every zero. Indeed all weights in the sum for
\(f=g\) are positive, so zero norm is equivalent
to all those values vanishing. Finite interpolation
from (M29.8) gives every finitely supported list
of values. The completion of the quotient by
this radical is therefore
\[
       \mathcal H_{\mathrm{GNS}}
             \simeq\ell^2(\{\text{distinct }\rho\},m_\rho).
                                                               \tag{M29.10}
\]
The evaluation map is continuous from \(Q\):
rapid Mellin decay and the local zero count
bound the squared norm by finitely many
\(\mathcal E\) seminorms, and it vanishes on
the closed ideal. Thus the density assertion
in (M29.10) is a statement about the actual
arithmetic quotient, not only an abstract
diagonal model.

Multiplication by a test transform on (M29.10)
has one Hilbert dimension per distinct zero.
The normalized coordinate vector at \(\rho\)
is \(m_\rho^{-1/2}e_\rho\); its diagonal matrix
entry is still \(\mathcal Mf(i\gamma)\).
Consequently its ordinary Hilbert trace counts
each distinct zero once. The integer \(m_\rho\)
in (M29.9) is a weight in the spectral measure,
not a Hilbert spectral multiplicity.
It follows that the positive-frequency inverse
square on this GNS space has ordinary
Fredholm determinant
\[
                \prod_{\substack{\gamma>0\\\text{distinct}}}
                                      (1-w/\gamma^2).         \tag{M29.11}
\]
It agrees with (M29.7) if and only if every
positive zero is simple. This follows also
by comparing the orders of the isolated
zeros of the two entire determinants.
The associated graded construction (M29.2)
is what preserves the multiplicities required
by the arithmetic trace and the completed
function.

There are actual number-field examples in which
this distinction matters. For
\(E=\mathbf Q(\zeta_3,\sqrt[3]{2})\), NT-L17
proves unconditionally that \(\zeta_E\) has
infinitely many zeros of multiplicity at
least two. Under \(M\) these lie on the
critical line; infinitely many are
noncentral, and the selfdual symmetry
gives infinitely many positive ordinates.
Thus the ordinary GNS determinant (M29.11)
does not equal the pole-removed completed
\(\zeta_E\) determinant, whereas the
graded determinant (M29.7) does.
This example uses no simplicity conjecture.

\paragraph{The established family and the remaining geometric gate.}
For a selfdual primitive Dirichlet function,
including the separately pole-removed zeta
function, all of the positive conclusions
above use the designated theorem input
\(A_D\), without \(M\). A complex Dirichlet
character and its conjugate are handled
by the finite direct sum of their actual
arithmetic modules and the product of
their completions. Take the associated
graded in each summand, then the positive
frequency subspace of their direct sum.
The product completion is selfdual, its
zero multiplicities are the sums of the
summand multiplicities, and the resulting
Fredholm determinant is its AG-DET
determinant. Each of these assertions
follows from finite direct sums and the
already proved trace/product formulas.
The same construction applies under \(A_D\)
to a specified Hecke factor whose completed
function is a finite product of Dirichlet
completions with the same center, once
that classical factorization is supplied.
For general Hecke factors, and for their
selfdual finite products, the location
input remains conditional on \(M\).

This advances the comparison requested in
AG-Q08 and in the commentary after AG-DET:
the determinant operator is now obtained
from the canonical filtration of the
locally finite part of the actual
adelic summation quotient. Its positivity
still uses the stated zero-location
theorem. It is not an independently
prime-defined positive cohomology or
an identification with the
Gan--Gross--Prasad period determinants
in Alp\"oge's \emph{Remark}, steps~4
and~7 (PDF pp.~3--4). Those period and
integral-structure comparison maps
remain separate.

The precise Masullo source route is
different from an assertion of positivity
of the full Jordan representation.
Section~1.4.6, original PDF p.~30,
explicitly does not invoke Weil
positivity; \S1.4.8--1.4.10,
pp.~31--32, places the main burden
on computing the same geometric
trace intrinsically in arbitrarily
narrow strip completions.
Theorem~12.5.1.1, pp.~519--520,
equations (12.5.1.1)--(12.5.1.4),
asserts the comparison with the
zeros inside the chosen strip for
tests vanishing on \(C_K^1\).
Original pages~30,31,519,520
were rendered and read; the
complete displayed proof of
Theorem~12.5.1.1 was read as
a dependency map, not as an
audit of its cited upstream
kernel theorems. Equation
(10.3.1.17), p.~450, was
also inspected in the original
page image: its triangular
derivative block is consistent
with the local algebra above,
without supplying the missing
uniform orbital comparison.

MT-M28 proves the narrow-strip
trace class under \(M\). It
does not prove that geometric
comparison independently of
zero location. The remaining
gate is the zero-independent
geometric trace identity, or
an independent positive
cohomological/period comparison.
The new graded determinant
comparison closes the
particular module-identification
step and retains the full
numeric multiplicity, while
making the loss of nilpotent
extension data explicit.

\paragraph{Run locators.}
This new comparison is proved in
\texttt{GRH\_NIGHT/complexity},
\texttt{work/root/graded\_adelic\_determinant.tex}.
It connects MT-M18/MT-M24/MT-M26
to AG-DET and to the inherited
\texttt{ALPOGE\_BRIDGES} rows
BR-TR-01/02, the SA23
\texttt{geo-H-lidskii}/\texttt{geo-H-genus}
rows, and cote~1 scans~79--80
and~111--112 as recorded in
the supplied SA23 intake.
Those archive locators retain
their inherited status.

\subsection*{MT-M30. An ideal-valued splitting of restriction to norm one}
Let \(K\) be a fixed number field, choose the norm splitting
and Haar normalization of MT-M24, and put
\[
 E_K=\mathcal S_{\mathrm E}(C_K),\qquad
 \mathcal I_a=|\cdot|^a\overline{\Theta_K}^{\,E_K},\qquad
 E_K^0=\{f\in E_K:f|_{C_K^1}=0\},
                                                               \tag{M30.1}
\]
where \(a\in\mathbf R\) is fixed. Let
\(\mathcal S(C_K^1)\) denote the Bruhat smooth functions on the
compact group \(C_K^1\). Restriction to norm one is the
continuous map
\[
                 r:E_K\longrightarrow\mathcal S(C_K^1).
 \]
There is a continuous linear, compact-translation-equivariant
map
\[
         B_a:\mathcal S(C_K^1)\longrightarrow\mathcal I_a,
                   \qquad rB_a=\operatorname{id}.             \tag{M30.2}
\]
No GRH input is needed. In particular
\[
 P_a=\operatorname{id}-B_ar
 \quad\hbox{is a continuous projection onto }E_K^0,\qquad
                   P_af-f\in\mathcal I_a.                    \tag{M30.3}
\]
Inclusion consequently induces a topological isomorphism
\[
        E_K^0/(\mathcal I_a\cap E_K^0)
                          \ \simeq\ E_K/\mathcal I_a,         \tag{M30.4}
\]
whose inverse is induced by \(P_a\).
The target in (M30.2) is the \emph{closed} centered
theta ideal. Actual summation-image membership of
an infinite character sum is not asserted.

Write
\[
         \mathcal D^0(C_K)
              =\{f\in\mathcal D(C_K):f|_{C_K^1}=0\}.
 \]
This space is dense in \(E_K^0\).
Therefore two continuous linear functionals on \(E_K\)
that annihilate \(\mathcal I_a\) are equal everywhere
if and only if they agree on \(\mathcal D^0(C_K)\).
The map \(P_a\) need not preserve compact radial
support or positivity.

\begin{proof}
\emph{Uniform normalized elements of the ideal.}
At an elementary compact level, use the family of
characters \(\beta_\nu\), \(\nu\in\mathbf Z^d\times F\),
and the parameter \(H_\nu=2+|\nu|\) from step~1 of
MT-M28. Each extension is trivial on the fixed norm
splitting. For \(\nu\ne0\) use the entire completion
\[
 \Xi_\nu(z)=Q_\nu^{z/2}
        \prod_{j=1}^{n_K}\Gamma_{\mathbf R}(z+\mu_{\nu,j})
                         L_\nu(z).
 \]
For \(\nu=0\) use the fixed pole-removed completion
of \(\zeta_K\), for example its standard completion
multiplied by \(z(z-1)\). These choices have precisely
the nontrivial completed zero divisors used by MT-M24.
All satisfy \(\Xi_\nu(2)\ne0\), by the Euler product
and the finite nonzero gamma factors there.
Define
\[
          F_\nu(z)=
              \frac{\Xi_\nu(z)}{\Xi_\nu(2)}\,e^{(z-2)^2}.
                                                               \tag{M30.5}
\]
It is entire, vanishes to the full order at every
completed zero, and satisfies \(F_\nu(2)=1\).
Let \(\psi_\nu\) be its inverse Mellin transform.
For each fixed \(\nu\), the fixed-strip bounds and
Gaussian decay in MT-M18 show
\(\psi_\nu\in\mathcal E\).

We need, and now prove, bounds uniform at this
compact level:
\[
           p_{N,j}(\psi_\nu)
                        \le C_{N,j}H_\nu^{B_{N,j}}.            \tag{M30.6}
\]
The finite conductor is uniformly bounded and
\(\operatorname{Re}\mu_{\nu,j}\) is a nonnegative
half-integer with \(|\mu_{\nu,j}|\ll H_\nu\).
The Euler lower bound at \(2\) and the uniform
finite-\(L\) polynomial upper bound are precisely
(M28.10) and (M28.12). Their proofs were
unconditional. On either line
\(\operatorname{Re}z=\pm(N+9/4)\), the same
gamma-ratio proof as in (M28.20), now with the
single fixed denominator point \(w=2\), gives
\[
 \left|\frac{\Xi_\nu(z)}{\Xi_\nu(2)}\right|
       \le C_N(H_\nu+|\operatorname{Im}z|)^{B_N}
                              e^{C_N|\operatorname{Im}z|}.
                                                               \tag{M30.7}
\]
In detail, shift numerator gamma factors right
by a fixed even integer. Every recurrence
denominator has real part a quarter-integer
of absolute value at least \(1/4\). The
remaining gamma ratio is in a fixed right
half-plane. The bounds
\(|\operatorname{Im}\psi(z)|\le C_c\) and
\(|\psi(z)|\le C_c+2\log(2+|z|)\) there
give (M30.7) by integrating the logarithmic
derivative in imaginary and real directions.
This uses a canceled ratio throughout, so
no factorial-size completion normalization
remains. The single principal character
has its fixed polynomial pole-removal
factor and is included by enlarging
the constants once.

Multiplication by the Gaussian in (M30.5)
turns (M30.7) into
\[
 |F_\nu(\pm(N+9/4)+it)|
                 \le C_NH_\nu^{B_N}e^{-c t^2}.
 \]
Inverting on these lines, differentiating
under the absolutely convergent integrals,
and using the positive line for \(x\ge0\)
and negative line for \(x\le0\), proves
(M30.6), just as in (M28.22)--(M28.23).
This part of MT-M28 used only Euler
convergence, the classical functional
equation and per-function finite order.
Neither its GRH contour selection nor
its narrow-strip minimum modulus was
used here.

Set
\[
                 \varphi_\nu(x)=e^{ax}\psi_\nu(x).
 \]
Multiplication by \(e^{ax}\) is a continuous
automorphism of \(\mathcal E\), so it
preserves the polynomial-in-\(H_\nu\)
form of (M30.6), with changed indices
and constants. Its normalization is
explicitly
\[
                   \mathcal M\varphi_\nu(2-a)=1.              \tag{M30.8}
\]
The full-zero-jet identification in
MT-M24 shows that this is an element
of the centered \(\nu\)-sector
theta ideal \(\mathcal I_{a,\nu}\).
One can also see the ideal membership
before taking its closure: the
Gaussian factor in (M30.5) is itself
the Mellin transform of an element
of \(\mathcal E\), so \(\psi_\nu\)
lies in the principal convolution
ideal generated by the inverse
transform of \(\Xi_\nu\).

\emph{A polynomial lower bound for the normalizing norm.}
It is important not to divide by an
uncontrolled small number. Put \(t_0=2-a\)
and choose an integer \(L>|t_0|+1\).
By (M30.6),
\[
                  p_{L,0}(\varphi_\nu)\le C H_\nu^B.
 \]
For \(R\ge1\), the tail of the integral
in (M30.8) satisfies
\[
 \left|\int_{|x|>R}e^{t_0x}\varphi_\nu(x)\,dx\right|
       \le C' H_\nu^B e^{-(L-|t_0|)R}.
 \]
Choose \(R=C_0+C_1\log H_\nu\), with
fixed constants large enough that
this bound is at most \(1/2\).
The integral over \([-R,R]\) has
absolute value at least \(1/2\).
Cauchy--Schwarz gives
\[
 \frac12\le
       (2R)^{1/2}e^{|t_0|R}\|\varphi_\nu\|_2 .
 \]
Since \(R\ll1+\log H_\nu\), it follows
that, for fixed \(c,C>0\),
\[
                   \|\varphi_\nu\|_2^2\ge cH_\nu^{-C}.        \tag{M30.9}
\]
This argument also covers \(t_0=0\).

For a radial function write
\(\varphi^*(x)=\overline{\varphi(-x)}\), and define
\[
              b_\nu
                =\frac{\varphi_\nu*\varphi_\nu^*}
                             {\|\varphi_\nu\|_2^2}.           \tag{M30.10}
\]
Its first convolution factor is in
\(\mathcal I_{a,\nu}\), and the second
is in \(\mathcal E\). Thus
\[
          b_\nu\in\mathcal I_{a,\nu},
          \qquad b_\nu(0)=1,\qquad
          p_{N,j}(b_\nu)\le C_{N,j}H_\nu^{C_{N,j}}.            \tag{M30.11}
\]
The value at zero is exactly the
squared \(L^2\) norm in (M30.10).
For the seminorm bound, use
\[
 p_{N,j}(f*g)
   \le p_{N,j}(f)\int_{\mathbf R}e^{N|y|}|g(y)|\,dy
   \le 2p_{N,j}(f)p_{N+1,0}(g)
 \]
and then (M30.6), (M30.9).
No stability of \(\mathcal I_a\)
under the involution for arbitrary
\(a\) is needed for this membership:
only its first factor and its
convolution-ideal property are used.

\emph{Summing over the compact character family.}
For \(g\in\mathcal S(C_K^1)\) at the
chosen compact level, write
\[
       g(c)=\sum_{\nu\in\Lambda}g_\nu\beta_\nu(c)^{-1},
       \qquad
       g_\nu=\int_{C_K^1}g(c)\beta_\nu(c)\,dc .
 \]
The coefficients decrease faster
than every power of \(H_\nu\).
Define, in the norm-splitting coordinates,
\[
 (B_ag)(c\,\tau(e^x))
       =\sum_{\nu\in\Lambda}
                   g_\nu\beta_\nu(c)^{-1}b_\nu(x).
                                                               \tag{M30.12}
\]
Each compact derivative costs a
fixed power of \(H_\nu\), and
(M30.11) bounds every weighted
radial derivative by a fixed power.
The rapid Fourier coefficients
therefore make (M30.12) absolutely
convergent in every seminorm of
the elementary \(E_K\) level.
The required source coefficient
bound uses finitely many compact
derivatives, so \(B_a\) is
continuous at that level.

The functions \(b_\nu\) were
defined from the character,
its fixed norm extension and
(M30.5), independently of any
choice of generators used to
prove the bounds. Thus the
definitions agree when levels
are enlarged. The elementary
inductive-limit definition of
the Bruhat smooth topology
now gives the asserted global
continuous map.
Every finite partial sum lies
in \(\mathcal I_a\), and this
ideal is closed in \(E_K\);
hence the whole sum lies in it.
At \(x=0\), (M30.11) gives
\(rB_ag=g\), with convergence
in the compact smooth topology.
The character formula also
shows that \(B_a\) commutes
with compact translations.
This proves (M30.2).

Equations (M30.3) follow from
\(rB_a=\operatorname{id}\).
If \(f\in E_K^0\), then \(P_af=f\).
If \(f\in\mathcal I_a\), then
\(P_af\in\mathcal I_a\cap E_K^0\).
The maps induced by inclusion
and \(P_a\) on the two quotient
spaces are therefore inverse.
Both are continuous by the
defining quotient topologies,
which proves (M30.4).

\emph{Compact test density and extension of equalities.}
For \(f\in E_K^0\), choose smooth
radial cutoffs \(\kappa_R(x)\)
equal to one on \([-R,R]\)
and zero outside \([-2R,2R]\).
The functions \(\kappa_R f\)
are in \(\mathcal D^0(C_K)\):
their compact-variable level
is unchanged, their norm
support is compact, and
\(\kappa_R(0)f(c)=0\).
They converge to \(f\) in
every weighted derivative
seminorm. Indeed the
differences and the derivatives
of the cutoffs are supported
where \(|x|\ge R\); a bound
at one larger exponential
weight gives an additional
factor tending to zero.
This proves density.

Suppose \(T_1,T_2\) are
continuous on \(E_K\), vanish
on \(\mathcal I_a\), and
agree on \(\mathcal D^0(C_K)\).
By density they agree on
\(E_K^0\). For every \(f\in E_K\),
\[
            T_i(f)=T_i(P_af),
 \]
since \(f-P_af\in\mathcal I_a\).
They are therefore equal
on all of \(E_K\). The
reverse implication is immediate.
\end{proof}

\paragraph{An exact restricted-trace criterion.}
The splitting gives a useful precise
test of the residual source claim.
For each spectral pair \((\chi,\rho)\)
choose real weights
\[
                   0\le w_{\chi,\rho}\le m_{\chi,\rho}.
 \]
Define
\[
 \begin{split}
 \tau_{\mathrm{full}}(f)
    &=\sum_{\chi,\rho}m_{\chi,\rho}
              \widehat f(\beta_\chi|\cdot|^{\,\rho-a}),\\
 \tau_w(f)
    &=\sum_{\chi,\rho}w_{\chi,\rho}
              \widehat f(\beta_\chi|\cdot|^{\,\rho-a}).
 \end{split}                                                    \tag{M30.13}
\]
These are continuous on \(E_K\) and
annihilate \(\mathcal I_a\),
unconditionally. To verify convergence,
at each compact level the unconditional
Jensen estimate (M28.13) counts
the nontrivial zeros in each unit
height window by
\(O(\log(H_\nu+|n|))\).
Their real parts are in the fixed
strip \([0,1]\). Integration by
parts in the compact variables
and in the norm variable bounds
their Mellin values by arbitrary
negative powers of \(H_\nu+|n|\),
uniformly on bounded sets at that
level. In the norm variable the
integrals of derivatives are
controlled by finitely many
exponential-Schwartz seminorms;
compact radial support is not
needed. The same summation as
in MT-M28 proves absolute
convergence and continuity.
Vanishing on \(\mathcal I_a\)
is the full zero-jet kernel
identity from MT-M24/MT-M27.

Consequently
\[
 \left.\tau_{\mathrm{full}}\right|_{\mathcal D^0(C_K)}
       =\left.\tau_w\right|_{\mathcal D^0(C_K)}
       \quad\Longleftrightarrow\quad
             w_{\chi,\rho}=m_{\chi,\rho}
                     \text{ for every }(\chi,\rho).           \tag{M30.14}
\]
For the nontrivial implication,
the splitting extends the equality
to \(E_K\). Fix one spectral
pair and use the finite-jet
interpolation of MT-M18 in
that single character sector
to get transform value one
at that zero and zero at
all other zeros, with every
unwanted higher jet zero.
Evaluating the difference of
the two functionals gives
\(m_{\chi,\rho}-w_{\chi,\rho}\).
It must vanish. The reverse
implication follows directly
from the sums.
Equivalently, any omitted
zero or lowered multiplicity
has a witness already in
\(\mathcal D^0(C_K)\):
apply \(P_a\) to this
finite-spectral interpolant,
then approximate by the
compact radial cutoffs above.
Continuity keeps the nonzero
difference for all sufficiently
large cutoffs.

In particular, at \(a=1/2\)
choose \(w_{\chi,\rho}=m_{\chi,\rho}\)
inside a specified open
centered strip and zero outside.
Equality of the full and
strip-restricted functionals
on tests vanishing at norm
one is equivalent to the
absence of every omitted
zero. Requiring this for
all positive strip widths
is exactly the Hecke
zero-location statement.
The test restriction does
not weaken that criterion.
This equivalence supplies
no proof of either side.

\paragraph{What this changes in the source comparison.}
Masullo \S1.4.6, original PDF p.~30,
restricts the proposed geometric
trace calculation to
\(\varphi(C_K^1)=0\);
\S1.7.2, pp.~34--35,
notes that a literal convolution
square \(h*h^*\) satisfying
this condition is zero.
That last assertion is correct:
at the identity
\[
                     (h*h^*)(1)=\|h\|_2^2.
 \]
The new splitting shows that
it is not an obstruction
at the corrected arithmetic
quotient level. The class
of \(h*h^*\) has the
representative \(P_{1/2}(h*h^*)\)
vanishing on \(C_K^1\), and
every continuous quotient
trace takes the same value
on these two representatives.
The new representative need
not itself be a convolution
square or have any pointwise
positivity property.

Let \(W_{\mathrm{raw},1/2}\)
be the classical raw Weil
distribution in the centering
convention of MT-M24. Its
spectral version is
\[
 W_{\mathrm{spec}}(f)
   =\widehat f(|\cdot|^{1/2})
      +\widehat f(|\cdot|^{-1/2})
      -W_{\mathrm{raw},1/2}(f)
   =\tau_{\mathrm{full}}(f).
\]
An exact classical trace
locator is Meyer,
\emph{On a representation of
the idele class group related
to primes and zeros of
\(L\)-functions},
Theorem~5.8, author pp.~43--44,
with the spectral identification
of Theorem~5.11, p.~47.
The complete supplied proof
of Theorem~5.8 and its
normalization were read in
this run; its upstream
trace-theoretic lemmas
remain cited classical inputs.
Here its continuation to
\(E_K\) can be defined by
the absolutely convergent
full spectral sum (M30.13),
so no unproved geometric
kernel continuity is inserted.

Theorem~12.5.1.1 in the
Masullo draft, original
pp.~519--520, asserts that
the same geometric expression
equals the trace with only
the zeros inside the chosen
strip, on the restricted
compact tests. These original
pages and \S1.4.6 were
inspected as page images
for MT-M29. By (M30.14),
the remaining gate is
exactly that geometric
identity independently
of zero location, rather
than admissibility of
positivity tests or a
missing extension from
the restricted test space.
The uniform ideal-valued
splitting is unconditional;
the value of the full
functional on convolution
squares is nonnegative
under \(A_D\) in the
Dirichlet family and
under \(M\) in the
general Hecke family,
as in MT-M19/MT-M28.
No positivity of the
corrected representatives
is inferred.

\paragraph{Run and framework locators.}
This is a new proof in
\texttt{GRH\_NIGHT/complexity},
\texttt{work/root/idele\_restriction\_splitting.tex}.
It advances the arithmetic
quotient comparison from
MT-M24/MT-M27/MT-M28 and
the determinant comparison
MT-M29. Its geometric
framework is the trace
functional on a quotient
by theta relations, beside
the inherited SA23
\texttt{geo-H-lidskii}
and \texttt{geo-H-genus}
rows, and cote~1
scans~79--80,111--112.
It supplies neither an
integral cycle comparison
nor an independent
Gan--Gross--Prasad
period positivity theorem.

\subsection*{MT-S504. The exact relative central condition in the $S_5$ transfer}

Retain the $S_5$ extension and local hypotheses of MT-S501:
complex conjugation and the unramified Frobenius at $5$ both
have cycle type $(12)(34)$. Write $E=L^{A_5}$, let $F$
be a degree-six subfield, and let $H=EF$. Let $\eta$ be the
quadratic character of $E/\mathbb Q$. Use Calegari's labeling
of the five-dimensional character $\rho_5$, whose values on
the classes $1,2A,2B,3A,4A,5A,6A$ are
$(5,1,1,-1,-1,0,1)$. Then the all-prime Artin factorization is
\[
       \zeta_H(s)=\zeta_F(s)L(s,\eta)L(s,\rho_5).
 \tag{S504.1}
\]
Under the conditional input $M_H$, the strictly more focused
sufficient condition
\[
 \operatorname{ord}_{1/2}\zeta_H
       \le\operatorname{ord}_{1/2}\zeta_F
                     +\operatorname{ord}_{1/2}L(s,\eta)
 \tag{S504.2}
\]
implies the same conclusion as MT-S502: every nontrivial
irreducible Artin factor is entire and satisfies GRH, and
the factors are automorphic by Calegari's argument.
In fact equality then holds in (S504.2). The hypothesis
allows zeros in the two other factors; it does not require
$\zeta_H(1/2)\ne0$.

\begin{proof}
Put $m_5=\operatorname{ord}_{1/2}L(s,\rho_5)$, interpreted
as the meromorphic order, so negative values are poles.
Equation (S504.1) shows that (S504.2) is exactly $m_5\le0$.
The factorization and the following precise classical inputs
are Calegari, \emph{The Artin conjecture for some
$S_5$-extensions}, \texttt{arXiv:1112.1152v1}: the comment
after Theorem~1.2, p.~2; Lemmas~4.8--4.9, pp.~10--11;
and the final proof of Theorem~1.2, p.~11. Their relevant
proofs have been read in full. The automorphic construction
in preceding sections remains a cited classical input.

That construction supplies a cuspidal automorphic
$\varpi$ on $\mathrm{GL}_5/\mathbb Q$, an exceptional set
$\Omega$ of primes, and a meromorphic comparison
\[
 \Lambda(s,\varpi)=\Lambda(s,\rho_5)\mu(s),\qquad
 \mu(s)=\mu_B(s)
       \prod_{p\in\Omega}
             \left(\frac{1+p^{-s}}{1-p^{-s}}\right)^2.
 \tag{S504.3}
\]
Here $\mu_B$ is holomorphic and nonzero on $(0,1)$.
Lemma~4.8 says that if $\Omega$ is infinite, $\mu$ has
a pole at a real point in $(0,1)$. Lemma~4.9 says that
finite $\Omega$ forces exact local equality everywhere,
including ramified and archimedean factors. The completed
automorphic numerator in (S504.3) is holomorphic on this
interval.

Brauer induction and $M_H$ imply that $L(s,\rho_5)$ is
holomorphic and nonzero at every real point of $(0,1)$
other than $1/2$: every Hecke factor in its signed
factorization has that property. Its archimedean factors
are finite and nonzero on the interval. Thus a real pole
of $\mu$ could only occur at $1/2$ and would require
$m_5>0$. Our hypothesis $m_5\le0$ rules that out;
a pole of the denominator would instead give a zero
of the comparison. Lemma~4.8 makes $\Omega$ finite,
and Lemma~4.9 proves $\Lambda(s,\varpi)=\Lambda(s,\rho_5)$.
The source's final argument also gives the exact match
for the quadratic twist $\rho_5\otimes\eta$.

The other irreducible dimensions $4$ and $6$ are already
automorphic by Theorem~1.1, as recorded in MT-S501;
the linear factors are classical Dirichlet functions.
This proves all the automorphy and holomorphy assertions.
Now holomorphy makes $m_5\ge0$, so $m_5=0$ and equality
holds in (S504.2). Finally signed Brauer factorization
and $M_H$ locate the nontrivial zeros of every entire
factor on the critical line, exactly as in MT-ART07.
\end{proof}

All $S_5$ irreducibles are realizable over $\mathbb R$;
Fr\"ohlich--Queyrut's real-representation theorem gives
root number $+1$. The meromorphic central-order parity
calculation of MT-ART11 therefore makes $m_5$ even.
Consequently failure of the sufficient condition leaves
only the gate
\[
 \operatorname{ord}_{1/2}\zeta_H
       -\operatorname{ord}_{1/2}\zeta_F
       -\operatorname{ord}_{1/2}L(s,\eta)
                       =m_5\ge2.
 \tag{S504.4}
\]
This is a necessary condition for this route to remain
unsettled, not a claim that a positive central order
prevents automorphy. It removes base-factor central zeros
from the test exactly; no independence or noncoincidence
of zeros is assumed.

Geometrically the virtual permutation identity in
(S504.1) isolates a single rank-five Artin motive from
the arithmetic finite covers. It sharpens the
MT-S502 one-value gate in the cote~1/SGA5 trace-divisor
programme, and the $S_5$ portion of Q-ARTIN-POS.
It leaves the actual relative central order as an
arithmetic input; neither designated GRH theorem
evaluates it.

\subsection*{MT-S505. Real-axis Artin holomorphy in the specified $S_5$ family}

Retain the extension and local hypotheses of MT-S501 and the
notation of MT-S504. Conditional on $M_H$, every irreducible
Artin factor of this extension is holomorphic at every real
point of $(0,1)$, including the centre. In particular
\[
 \operatorname{ord}_{1/2}\zeta_H
  \ge \operatorname{ord}_{1/2}\zeta_F
                  +\operatorname{ord}_{1/2}L(s,\eta).
 \tag{S505.1}
\]
The difference in (S505.1) is a nonnegative even integer.
If it is zero, every irreducible factor is automorphic
and satisfies Artin holomorphy and GRH, by MT-S504.
Thus no additional central nonvanishing hypothesis is
needed for the real-axis holomorphy assertion itself.

\begin{proof}
Write $m_5=\operatorname{ord}_{1/2}L(s,\rho_5)$,
allowing a negative order for a meromorphic pole.
If $m_5<0$, the sufficient condition $m_5\le0$
of MT-S504 holds. That theorem proves automorphy
and forces $m_5=0$, a contradiction. Hence $m_5\ge0$.
Meromorphicity then means that $L(s,\rho_5)$
is holomorphic at $s=1/2$. This argument does
not assume its central value is nonzero.

For completeness the other rank-five factor
is also holomorphic at the centre for an
unconditional elementary reason. Put
$J=\operatorname{Gal}(L/F)$, of order $20$.
Its Sylow $5$-subgroup $P$ is normal, since
its number divides $4$ and is $1$ modulo $5$.
The normalizer of $P$ in $S_5$ has order $20$;
thus $J$ is that normalizer. In the standard
affine action on $\mathbb F_5$, its elements
have cycle types
\[
       1\text{ once},\quad 5A\text{ four times},\quad
       2B\text{ five times},\quad 4A\text{ ten times}.
\]
Using Calegari's character table, arXiv:1112.1152v1,
p.~3, the average over these elements is $1$
for the trivial character and for $\rho_5\otimes\eta$,
and $0$ for $\eta,\rho_4,\rho_4\otimes\eta,\rho_5,\rho_6$.
For example the two rank-five averages are
\[
 \frac{5+4\cdot0+5\cdot1+10\cdot(-1)}{20}=0,
 \qquad
 \frac{5+4\cdot0+5\cdot1+10\cdot1}{20}=1.
\]
Frobenius reciprocity therefore proves the
permutation identity and its all-prime
Artin factorization
\[
 \operatorname{Ind}_{J}^{S_5}\mathbf1
          =\mathbf1\oplus(\rho_5\otimes\eta),\qquad
 \zeta_F(s)=\zeta_{\mathbb Q}(s)L(s,\rho_5\otimes\eta).
 \tag{S505.2}
\]
For $0<s<1$, the alternating series
\[
 A(s)=\sum_{n\ge1}(-1)^{n-1}n^{-s}
     =\sum_{j\ge1}\bigl((2j-1)^{-s}-(2j)^{-s}\bigr)
\]
converges and is strictly positive. Analytic continuation
of the elementary identity from $s>1$ gives
$A(s)=(1-2^{1-s})\zeta_{\mathbb Q}(s)$ on this interval.
The first factor is negative, so
$\zeta_{\mathbb Q}(s)<0$ throughout $(0,1)$.
Thus (S505.2) expresses the twisted rank-five
factor as a holomorphic quotient on the whole
interval, without any GRH assumption.

The factors of dimensions $4$ and $6$ are
entire by the cited Calegari Theorem~1.1,
and the linear factors are holomorphic
on $(0,1)$. Finally for $\rho_5$ at a
real point other than $1/2$, signed Brauer
induction and $M_H$ give a product and
quotient of nonzero holomorphic Hecke values,
so it is holomorphic there too.
This proves the real-axis assertion for
every irreducible factor.

Taking central orders in (S504.1) gives
(S505.1) with difference $m_5$. The real
realizability/root-number-$+1$ and
meromorphic functional-equation parity
argument already checked in MT-S504
makes $m_5$ even. The equality case is
exactly the sufficient condition of
MT-S504. All conclusions follow.
\end{proof}

This settles the real-axis portion of the
Artin pole gate for this precise $S_5$
family. Possible critical-line poles away
from the real axis are not ruled out in
the remaining $m_5\ge2$ case. Geometrically,
(S505.2) is the degree-six permutation
Artin motive of the finite cover, while
(S505.1) is an effectivity statement for
the central coefficient of the relative
rank-five divisor. It concerns that
coefficient, not the effectivity of the
whole Artin divisor required by
Q-ARTIN-POS and the cote~1/SGA5 trace programme.

\subsection*{MT-M31. Local Euler differences: actual range, Hilbert closure,
and the Schanuel dependency}

The distinction in this result is between the range of a bounded
operator and the closure of that range.  It occurs before any
GRH or positivity question.

\paragraph{The precise source assertion.}
Masullo, Proposition 11.2.0.1(2), original PDF pp.~484--485,
asserts that the extension of $\mathrm{PV}_{S_0}$ to the
zero-line Hilbert completion has closed range when
$s\geq 2|S_0|$.  Equation (11.2.0.3), p.~484, identifies
the closure of the smooth image with that actual operator range.
The proof on p.~484 gives, for $S_0=\{v\}$, the exact formula
\[
  P_v\Phi_F(f)=\Phi_F(P_vf),\qquad
  P_vf(x)=f(x)-f(\pi_v^{-1}x).
  \tag{M31.1}
\]
There is no additional norm-character multiplier in this formula.
On p.~485 the proof identifies the range with an intersection
of kernels of continuous coefficient functionals.  The argument
below proves that this identifies the \emph{closure}, and that
the actual range is not closed even for one prime over $\mathbf Q$.
The original images of both pages have been read, including
the multiplier and the range/closure equation.

\paragraph{General analytic formulation.}
For $s\geq0$ put
\[
 H_s=L^2(\mathbf R,w_s(x)\,dx),\qquad
 w_s(x)=(1+x^2)^{s/2},\qquad
 (U_hf)(x)=f(x-h).
 \tag{M31.2}
\]
This is precisely the radial form of the source norm
(9.1.2.3), original p.~399, with the $a=0$ completion
of (10.1.1.3), p.~434; see MT-M24--M25 for the actual
theta-image identification.  Fix a nonempty finite list
of $h_v>0$ and $a_v\in\mathbf C$ with $|a_v|=1$, and set
\[
 P=\prod_v(1-a_vU_{h_v}),\qquad
 D(z)=\prod_v(1-a_ve^{h_vz}).
 \tag{M31.3}
\]
Repeated factors are allowed.  Let $\mathcal T\subset\mathbf R$
be the locally finite set of $t$ with $D(it)=0$, and let
$m_t$ be its multiplicity.  Define
\[
 q_s=\max\{0,\lceil(s-1)/2\rceil\},\qquad
 J_{t,j}(f)=\int_{\mathbf R}f(x)x^je^{itx}\,dx.
 \tag{M31.4}
\]
The following assertions hold without any conjectural input:
\begin{enumerate}
\item $P$ is a bounded injective operator on $H_s$, and its
actual range is not closed, for every $s\geq0$.
\item Its exact closed range replacement is
\[
 \overline{P H_s}^{\,H_s}
 =\bigcap_{t\in\mathcal T}\ \bigcap_{0\leq j<\min(m_t,q_s)}
            \ker J_{t,j}.
 \tag{M31.5}
\]
Every displayed functional is continuous.  Thus for $0\leq s\leq1$
the actual range is dense and proper; for $s>1$ its closure
is proper and the range is a proper dense subspace of that closure.
\item If $\mathcal A$ is any dense linear subspace of $H_s$, then
\[
 \overline{P\mathcal A}^{\,H_s}
       =\overline{PH_s}^{\,H_s}.
 \tag{M31.6}
\]
In particular the corrected quotient is the Hilbert quotient by
(M31.5).  It is not the non-Hausdorff topological quotient
$H_s/PH_s$.
\end{enumerate}

\paragraph{Boundedness and injectivity.}
The inequality
\[
 1+(x+y)^2\leq2(1+x^2)(1+y^2)
 \tag{M31.7}
\]
shows that $\|U_h\|_{H_s\to H_s}\leq
2^{s/4}(1+h^2)^{s/4}$.  It gives the analogous estimate
on $H_{-s}=L^2(w_s^{-1}dx)$ as well.  Each factor and
their product are bounded.  Since $w_s\geq1$, $H_s$
embeds in ordinary $L^2(\mathbf R)$.  The ordinary Fourier
transform of $Pf$ is the transform of $f$ multiplied by
\[
            \prod_v(1-a_ve^{-ih_v\xi}).
\]
Its real zeros form a locally finite set, of Lebesgue measure
zero.  Consequently $Pf=0$ implies $f=0$ in $L^2$,
hence in $H_s$.

\paragraph{A quantitative failure of closed range.}
Choose one root $t_0$ of one of the factors of $D(it)$,
and choose nonzero $\psi\in C_c^\infty((1,2))$.
For $R\geq1$ put
\[
             f_R(x)=e^{-it_0x}\psi(x/R).
 \tag{M31.8}
\]
Changing variables gives constants $c,C>0$, independent of $R$,
such that
\[
      cR^{s+1}\leq\|f_R\|_{H_s}^2\leq CR^{s+1}.
 \tag{M31.9}
\]
For a factor with $a_ve^{it_0h_v}=1$,
\[
 (1-a_vU_{h_v})f_R(x)
   =e^{-it_0x}\{\psi(x/R)-\psi((x-h_v)/R)\}.
 \tag{M31.10}
\]
The mean-value bound, with support contained in an interval
of length $O(R)$ and $|x|=O(R)$, gives
\[
 \|(1-a_vU_{h_v})f_R\|_{H_s}
      \leq C_vR^{(s-1)/2}.
\]
The remaining finitely many factors are bounded independently
of $R$.  Since they commute,
\[
       \frac{\|Pf_R\|_{H_s}}{\|f_R\|_{H_s}}=O(R^{-1}).
 \tag{M31.11}
\]
An injective bounded operator between Hilbert spaces with
closed range has bounded inverse onto that range, by the
bounded-inverse theorem, and hence is bounded below.  Equation
(M31.11) contradicts this.  It proves the first assertion.

\paragraph{The exact closure.}
Use the continuous bilinear duality
\[
       H_s\times H_{-s}\longrightarrow\mathbf C,
       \qquad(f,g)\longmapsto\int f(x)g(x)\,dx.
 \tag{M31.12}
\]
Complex conjugation converts it to the usual Hilbert duality;
the bilinear version keeps the Mellin convention explicit.
The annihilator $\mathcal N$ of $PH_s$ consists precisely of
the $g\in H_{-s}$ satisfying
\[
          P^{\mathrm t}g=
             \prod_v(1-a_vU_{-h_v})g=0.
 \tag{M31.13}
\]
Each $g\in H_{-s}$ is a tempered distribution: Cauchy--Schwarz
bounds its pairing with a Schwartz test by
$\|g\|_{H_{-s}}\|\phi\|_{H_s}$.  With Fourier convention
$\widehat g(\xi)=\int g(x)e^{-i\xi x}dx$, (M31.13) is
\[
                         D(i\xi)\widehat g(\xi)=0.
 \tag{M31.14}
\]
It follows that $\widehat g$ is supported on $\mathcal T$.
Near an isolated $t\in\mathcal T$, a distribution killed
by the smooth function $D(i\xi)$ is a linear combination
of $\delta_t,\ldots,\delta_t^{(m_t-1)}$.  To verify the
order bound, divide by the nonvanishing smooth factor in
$D(i\xi)=(\xi-t)^{m_t}b(\xi)$ and use
$(\xi-t)\delta_t^{(j)}=-j\delta_t^{(j-1)}$.
The classical point-support distribution theorem,
\emph{Expository Notes on Distribution Theory}, AMSC/MATH~673,
Fall~2018, Theorem~2.1, p.~4 (University of Maryland),
then gives exactly this finite list.

Here localization is bounded in the required Hilbert norm.
For any Schwartz kernel $k$, (M31.7) and Minkowski give
\[
 \|k*g\|_{H_{-s}}
 \leq C_s\left(\int |k(y)|(1+|y|)^{s/2}dy\right)
                 \|g\|_{H_{-s}}.
 \tag{M31.15}
\]
Multiplying $\widehat g$ by a compactly supported smooth
function supported near one $t$ is therefore bounded on
$H_{-s}$.  The resulting inverse transform is a polynomial
times $e^{itx}$.  A nonzero polynomial of degree $j$
times this exponential belongs to $H_{-s}$ exactly when
\[
                              2j<s-1.
 \tag{M31.16}
\]
Indeed its squared absolute value is comparable to $|x|^{2j}$
for all sufficiently large $|x|$; the weighted integral
then has tail comparable to $\int_1^\infty x^{2j-s}dx$.
In particular equality in (M31.16) is not allowed.
This argument isolates each frequency first, so cancellation
among distinct frequencies cannot evade the cutoff.

For completeness the finite exponential-polynomial sums are
dense in $\mathcal N$ in its actual $H_{-s}$ norm.
Choose $\eta\in C_c^\infty(\mathbf R)$ equal to one near zero,
and put $\eta_R(\xi)=\eta(\xi/R)$, $R\geq1$.
Its inverse Fourier kernel has the form $k_R(x)=Rk(Rx)$,
with integral one.  Equation (M31.15) gives uniformly bounded
operators because
\[
 \int |k_R(y)|(1+|y|)^{s/2}dy
   =\int |k(u)|(1+|u|/R)^{s/2}du
   \leq\int |k(u)|(1+|u|)^{s/2}du.
\]
Translations are strongly continuous on $H_{-s}$, first on
$C_c^\infty$ and then by density and (M31.7).
Writing $k_R*g-g$ as the integral of
$k(u)(U_{u/R}g-g)$ proves convergence to zero in $H_{-s}$
by dominated convergence, using the same polynomial bound.
Each $\eta_R\widehat g$ has finite support in $\mathcal T$,
so its inverse transform is a finite sum of the terms
$x^je^{itx}$ with $j<m_t$ and (M31.16).
Conversely all these terms satisfy (M31.13), and their
closed span therefore equals $\mathcal N$.

Taking the annihilator once more proves (M31.5).
Cauchy--Schwarz proves continuity of every $J_{t,j}$ in that
formula using (M31.16).  For $s>1$, at least one root exists
and $J_{t,0}$ is a nonzero functional, so the closure is proper.
For $s\leq1$ no term survives and the closure is all of $H_s$.
Nonclosedness gives all remaining strictness assertions.
Finally boundedness of $P$ and density of $\mathcal A$ prove
(M31.6) directly by approximating each vector before applying $P$.

\paragraph{Application to the literal adelic assertion.}
It is enough to take $F=\mathbf Q$, $S_0=\{p\}$, and
the trivial character of the compact norm-one group.
Choose a norm splitting $C_{\mathbf Q}\simeq
C_{\mathbf Q}^1\times\mathbf R$ with radial coordinate
$x=\log|g|$, and normalize the compact Haar measure to one.
On this reducing sector (M31.1) is
$f(x)\mapsto f(x)-f(x+\log p)$.
Replacing $x$ by $-x$ gives exactly $1-U_{\log p}$ on
(M31.2).  Reflection is isometric because $w_s$ is even.

The ambient operator on the full adelic Hilbert space is
injective as well.  If $\Phi(g)=\Phi(\pi_p^{-1}g)$
almost everywhere, decompose $C_{\mathbf Q}$ into the
disjoint norm shells
$\{k\log p\leq\log|g|<(k+1)\log p\}$.
Translation by $\pi_p^{-1}$ maps successive shells onto
one another and preserves Haar measure.  The unweighted
$L^2$ mass is consequently equal on all shells.  Since
$H_s\subset L^2$ for $s\geq0$, it must be zero.
Thus the restriction to any closed invariant Hilbert subspace
is also injective.  If its range were closed, that restriction
would be bounded below on every sector it contains.

The actual zero-line completion of the trivial-character
theta image contains the whole radial $H_s$ sector.
This follows from MT-M24 and the exact MT-M25 kernel formula
at $a=0$, since the pole-removed completed zeta function has
no nontrivial zero on $\Re z=0$.  The latter is the classical
zero-free boundary theorem and functional equation, with no
RH premise.  Thus the normalized vectors from (M31.8)
belong to the actual completion, and (M31.11) contradicts
Proposition 11.2.0.1(2), including its stated range $s\geq2$.
This proves a counterexample in the specified space, not only
in an unrelated radial model.

By (M31.6), the left side of source (11.2.0.3) is precisely
the closure in (M31.5); its right side, read as the actual
extended-operator range, is a proper dense subspace of it.
Replacing the right side by its closure is the necessary
correction.  No conclusion about the truth of $M_H$ follows
from this counterexample.

\paragraph{What the local-pole input actually assumes.}
The following is a targeted dependency trace, not an audit
of the entire draft.
\begin{enumerate}
\item Theorem 7.1.3.2(1), original pp.~324--325, says that
two unramified finite places of distinct residue
characteristics have at most one common Euler pole.
If $s,s'$ were distinct common poles, then
$(s-s')\log q_v$ and $(s-s')\log q_w$ would both be
nonzero integral multiples of $2\pi i$, forcing
$\log q_v/\log q_w\in\mathbf Q$, impossible for powers
of distinct rational primes.  This part is unconditional.
\item Part (3b) of that theorem explicitly assumes Schanuel
for its sharper coincidence exclusion for a character
which is an algebraic character times a norm power.
The proof on p.~325 uses the four-exponentials consequence
of Schanuel.  Its weaker three-residue-characteristic
exclusion uses the unconditional six-exponentials theorem.
The source does not present the four-exponentials assertion
as an unconditional input in that argument.
\item For one rational prime $p$, the pole multiplicity of
$\prod_{v\mid p}L_v(s,\chi_v)$ is at most the number
of factors, hence at most $[F:\mathbf Q]$.
Each unramified Euler factor has simple poles, and ramified
factors contribute none.  This elementary bound does not
use Schanuel or coincidence control between different primes.
Definition 7.3.1.4, p.~343, and Theorem 7.3.1.6,
pp.~344--350, organize the construction by precisely these
$p$-primary multiplicities.  The supplied proof of the latter
was read as a dependency map; its invoked structure and
interpolation theorems are not certified here.
\item The source's zero-line coefficient norm is
$\int_{\mathbf R}|x|^{2j}(1+x^2)^{-s/2}dx$,
displayed in Proposition 11.2.0.4, p.~486.
The strict condition is (M31.16), not a conjectural
arithmetic statement.  If $j\leq [F:\mathbf Q]-1$
and $s\geq2[F:\mathbf Q]$, it holds strictly.
Thus the stated working degree $2[F:\mathbf Q]$
retains every $p$-primary jet without Schanuel.
The source's general boundary cutoff (11.2.0.6), p.~485,
has the same equality problem already corrected by MT-M25;
it does not affect this sufficient even-integer degree.
\item Theorem 11.2.2.1 and Remark 11.2.2.2,
pp.~490--491, use $s\geq2[F:\mathbf Q]$ and
$N(\varepsilon_p)\leq[F:\mathbf Q]$ to assert a
trace with full local multiplicity.  The statement and its
complete supplied proof contain no Schanuel premise.
Their Hilbert quotient and coefficient construction instead
depend on the actual/closed-range identification above.
The correct local quotient can be constructed using
(M31.5); the further circle model and local trace are
treated separately in MT-M32.
\end{enumerate}

\paragraph{The precise downstream gate.}
The dependency is explicit in the text, rather than inferred
from proximity of the statements.  Theorem 11.2.0.2,
pp.~485--486, uses the space asserted closed in
(11.2.0.3); the definition of the orbital Hilbert quotient
at pp.~489--490 and Theorem 11.2.2.1 then use this space.
Theorem 11.3.2.1, pp.~495--496, identifies a localized
geometric kernel with the local orbital action and invokes
the finite-rank trace approximation (11.2.2.12) from
Theorem 11.2.2.1.  Finally the proof of Theorem 12.5.1.1,
p.~520, explicitly invokes Theorem 11.3.2.1 for the
contributions away from the chosen finite set of places,
and Theorem 12.4.1.1 for its remaining local contribution.
The supplied statements and complete displayed proof of
11.3.2.1 were read for this trace.  The full proof of
12.4.1.1 is not claimed audited here.

Consequently the local closure correction preserves a
meaningful local orbital quotient, but does not by itself
prove the geometric-kernel identification in 11.3.2.1
or the narrow-strip trace identity of 12.5.1.1.
The Schanuel refinement belongs to a different branch:
the introductory Expectation 1.7.1, original pp.~37--38,
labels the sharper zero-width multiplicity consequences
as expectations not fully worked out.  No hidden Schanuel
assumption has been established in the positive-width
geometric trace route by this targeted dependency reading.
Its remaining issue is the stated analytic comparison,
not the elementary $p$-primary pole bound.

\subsection*{MT-M32. The corrected primary Euler quotient has a Riesz basis
and a local Poisson trace}

This result proves the local trace after the closure correction
of MT-M31.  It does not identify the global geometric kernel
with this local action.

\paragraph{Statement and conventions.}
Retain $H_s,w_s,U_h$ from MT-M31.  Fix $h>0$, positive
integers $d_1,\ldots,d_b$, and unit complex numbers
$a_1,\ldots,a_b$.  Define
\[
 Q(w)=\prod_{v=1}^{b}(1-a_vw^{d_v}),\qquad
 P=Q(U_h),\qquad
 Z_s=H_s/\overline{PH_s}^{\,H_s}.
 \tag{M32.1}
\]
Let $\lambda_1,\ldots,\lambda_L$ be the distinct roots
of $Q$, all on the unit circle, with multiplicities
$r_1,\ldots,r_L$.  Choose $0\leq\theta_\ell<2\pi$
such that $\lambda_\ell=e^{i\theta_\ell}$, and put
\[
 q_s=\max\{0,\lceil(s-1)/2\rceil\},\qquad
 r'_\ell=\min(r_\ell,q_s),\qquad
 t_{\ell,k}=\frac{\theta_\ell+2\pi k}{h}.
 \tag{M32.2}
\]
All roots and multiplicities here come from a finite Euler
polynomial; no Hecke GRH or Schanuel premise occurs.

The dual of $Z_s$, using the bilinear integral pairing of
MT-M31, has the Riesz basis
\[
 h^{-1/2}x^j e^{it_{\ell,k}x},
 \quad 1\leq\ell\leq L,\quad k\in\mathbf Z,
                  \quad0\leq j<r'_\ell.
 \tag{M32.3}
\]
When every $r'_\ell=0$, the quotient is zero and the basis
is empty.  The corresponding coefficient map on $Z_s$
is a bounded isomorphism onto the $\ell^2$ space of the
indices in (M32.3); its inverse supplies a biorthogonal
Riesz basis of the quotient.

For $f\in C_c^\infty(\mathbf R)$, or more generally for
an exponential-Schwartz test, let
\[
 C_f u(x)=\int_{\mathbf R}f(y)u(x-y)\,dy,
 \qquad M_f(z)=\int_{\mathbf R}f(y)e^{zy}\,dy.
 \tag{M32.4}
\]
The induced operator on $Z_s$ is trace class and satisfies
\[
 \operatorname{Tr}_{Z_s} C_f
      =\sum_{\ell=1}^{L}\sum_{k\in\mathbf Z}
                           r'_\ell M_f(it_{\ell,k})
      =h\sum_{n\in\mathbf Z}f(nh)
                          \sum_{\ell=1}^{L}r'_\ell\lambda_\ell^n.
 \tag{M32.5}
\]
Both expressions converge absolutely.  The map from tests
to trace-class operators is continuous in trace norm, with
uniform tails on bounded sets of tests at a fixed compact
support, and likewise for bounded sets in exponential
Schwartz space.

If $s>2\max_\ell r_\ell-1$, every local jet survives,
and the formula simplifies to
\[
 \operatorname{Tr}_{Z_s}C_f
       =\sum_{v=1}^{b}d_vh
                     \sum_{n\in\mathbf Z}a_v^{-n}f(nd_vh).
 \tag{M32.6}
\]
The strict condition is satisfied in particular by
$s\geq2b$.  For a primary number-field Euler product,
$h=\log p$, $d_v=f_v$ is the residue degree, and
$b\leq[F:\mathbf Q]$; hence the source's sufficient
degree $s\geq2[F:\mathbf Q]$ retains all multiplicities.

\paragraph{A finite matrix-weight model.}
If no index survives the cutoff, MT-M31 gives $Z_s=0$
and all the assertions are immediate.  Assume henceforth
that at least one index survives.
Let $\mathcal N_s\subset H_{-s}$ be the annihilator of
$PH_s$.  By the fully proved formula in MT-M31,
it is the closed span of the functions in (M32.3).
We construct a bounded isomorphism
\[
 \mathcal V_s:=\bigoplus_{\ell=1}^{L}
                     \bigoplus_{j=0}^{r'_\ell-1}L^2(0,h)
                  \longrightarrow\mathcal N_s.
 \tag{M32.7}
\]
For a coefficient vector $c=(c_{\ell,j})$, extend each
$c_{\ell,j}$ periodically with period $h$ and set
\[
 (Ac)(x)=\sum_{\ell,j}
       x^j e^{i\theta_\ell x/h}c_{\ell,j}(x\bmod h).
 \tag{M32.8}
\]
Endpoint representatives have measure zero and make no
difference to this definition.  Write $x=u+nh$ with
$0\leq u<h$, $n\in\mathbf Z$, and define
\[
 v_{\ell,j}(u,n)=(u+nh)^j
                         e^{i\theta_\ell(u/h+n)}.
 \tag{M32.9}
\]
The finite Hermitian matrix
\[
 G_{(\ell,j),(m,k)}(u)
   =\sum_{n\in\mathbf Z}
       \frac{\overline{v_{\ell,j}(u,n)}v_{m,k}(u,n)}
                  {(1+(u+nh)^2)^{s/2}}
 \tag{M32.10}
\]
converges absolutely and uniformly on $0\leq u\leq h$.
Indeed $j,k<q_s$ imply $j+k-s<-1$, so the tails
are bounded by a fixed summable power of $1+|n|$.
Thus $G(u)$ is continuous.  Tonelli, applied first to
finite coefficient sums and then to their squared absolute
values, gives the exact norm identity
\[
           \|Ac\|_{H_{-s}}^2
                =\int_0^h c(u)^*G(u)c(u)\,du.
 \tag{M32.11}
\]

The matrix $G(u)$ is positive definite for every fixed $u$.
If a vector $z$ has $z^*G(u)z=0$, then
\[
       \sum_{\ell,j}z_{\ell,j}
          e^{i\theta_\ell u/h}(u+nh)^j\lambda_\ell^n=0
                       \quad\text{for every }n\in\mathbf Z.
 \tag{M32.12}
\]
Exponential-polynomial sequences with distinct nonzero
$\lambda_\ell$ are linearly independent.
Here is a direct verification.  For one index $\ell$, apply
the finite-difference operator
$\prod_{m\ne\ell}(S-\lambda_m)^{r'_m}$, where
$(Sa)_n=a_{n+1}$, to (M32.12).  Every other component
vanishes.  If the $\ell$-component had a nonzero polynomial
of degree $j$, the resulting component would have the
same degree and a nonzero leading coefficient, multiplied
by $\prod_{m\ne\ell}(\lambda_\ell-\lambda_m)^{r'_m}$.
Applying $(S-\lambda_\ell)^j$ would then give a nonzero
constant multiple of $\lambda_\ell^n$, a contradiction.
Thus all the polynomials vanish.  The powers $(u+nh)^j$
are linearly independent polynomials in $n$, so $z=0$.

Continuity on the compact interval now gives constants
$0<c\leq C<\infty$, depending on this fixed local product
and $s$, with
\[
                  cI\leq G(u)\leq CI\quad(0\leq u\leq h).
 \tag{M32.13}
\]
Consequently (M32.8) is bounded and bounded below, and
its image is closed.  On each summand of $\mathcal V_s$,
the Fourier orthonormal basis
$h^{-1/2}e^{2\pi iku/h}$ maps to the corresponding
function in (M32.3).  Its image therefore contains a dense
subspace of $\mathcal N_s$ by MT-M31.  Conversely, finite
Fourier vectors are dense in $\mathcal V_s$ and their
images belong to the closed space $\mathcal N_s$, so
$A\mathcal V_s\subset\mathcal N_s$.  Its closed
image is exactly $\mathcal N_s$.  This proves (M32.7)
and the asserted Riesz basis, in the actual weighted norm.
No global Riesz assertion for a union of different rational
prime lattices is being made.

The bilinear integral pairing identifies $\mathcal N_s$
isometrically with the continuous dual of $Z_s$.
Taking the dual of the bounded isomorphism from the
standard $\ell^2$ basis onto (M32.3), and using reflexivity
of the Hilbert quotient, proves that the coefficient map
\[
 [u]\longmapsto
 \left(h^{-1/2}\int_{\mathbf R}u(x)x^je^{it_{\ell,k}x}\,dx
                                      \right)_{\ell,k,j}
 \tag{M32.14}
\]
is a bounded bijection with bounded inverse onto $\ell^2$.
This also proves the existence and uniform boundedness
of the biorthogonal quotient vectors.  It is stronger
than merely knowing that the individual coefficient
functionals are continuous.

\paragraph{The exact scalar quotient norm for one untwisted factor.}
For $P=1-U_h$ and $s>1$, the preceding model has a
particularly concrete primal form.  Put
\[
 W(u)=\sum_{n\in\mathbf Z}(1+(u+nh)^2)^{-s/2},
 \qquad
 (Rf)(u)=\sum_{n\in\mathbf Z}f(u+nh),\quad 0\leq u<h.
 \tag{M32.14a}
\]
The function $W$ is continuous and bounded above and
below by positive constants.  For almost every $u$,
Cauchy--Schwarz proves absolute convergence of $Rf(u)$
and
\[
 |Rf(u)|^2\leq W(u)
         \sum_{n\in\mathbf Z}|f(u+nh)|^2w_s(u+nh).
\]
Thus $R:H_s\to L^2((0,h),W(u)^{-1}du)$ is a contraction.
It is onto: a prescribed $F$ in the target has lift
\[
 f(u+nh)=\frac{F(u)}{W(u)w_s(u+nh)},
 \qquad
 \|f\|_{H_s}^2=\int_0^h\frac{|F(u)|^2}{W(u)}du.
 \tag{M32.14b}
\]
The coefficient $J_{2\pi k/h,0}(f)$ is the corresponding
Fourier coefficient of $Rf$ (with the positive-exponent
convention).  Completeness of Fourier series and MT-M31
give $\ker R=\overline{(1-U_h)H_s}$.  Therefore
\[
 H_s/\overline{(1-U_h)H_s}
        \simeq L^2((0,h),W(u)^{-1}du)
 \tag{M32.14c}
\]
is an isometry with its displayed quotient norm.
This is an explicit circle quotient even though the
actual difference-operator range is not closed.

\paragraph{The actual convolution and its trace norm.}
The polynomial-weight translation estimate of MT-M31
shows that $C_f$ is bounded on $H_s$ and $H_{-s}$.
It commutes with $P$ and preserves the closure of its range,
so it acts on $Z_s$.  On the dual functions one computes
directly
\[
 C_f^{\mathrm t}(x^je^{itx})
  =\int f(y)(x+y)^je^{it(x+y)}dy
  =\sum_{a=0}^{j}\binom ja M_f^{(a)}(it)
                                  x^{j-a}e^{itx}.
 \tag{M32.15}
\]
The same identity, paired with $u$, describes the quotient
operator in the coordinate map (M32.14).  It is the
direct sum over $(\ell,k)$ of triangular matrices of
size $r'_\ell$, with diagonal entry $M_f(it_{\ell,k})$.
Their sizes are bounded by the degree of $Q$, which is
fixed in this theorem.

For each integer $N\geq0$ and each derivative order
$0\leq a<\max r'_\ell$, integration by parts gives
\[
 |M_f^{(a)}(it)|\leq C_{N,a}(f)(1+|t|)^{-N}.
 \tag{M32.16}
\]
The constants are bounded by finitely many test seminorms:
on a fixed compact support they are controlled by finitely
many derivatives of $f$, and on exponential-Schwartz
space by finitely many of its defining seminorms.
For a finite $r\times r$ matrix its trace norm is at most
$r$ times its operator norm.  Hence the sum of the trace
norms of these blocks is finite, by (M32.16) with $N>1$.
Conjugating by the bounded isomorphism (M32.14) multiplies
the trace-norm estimate by at most the product of the two
isomorphism norms.  This proves trace class and its stated
continuity.

The block action was obtained on every coordinate of an
actual bounded convolution operator.  Thus the trace-class
sum equals that actual operator, first on finite coordinate
vectors and then by density.  Finite partial sums have
trace $\sum_{\ell,k}r'_\ell M_f(it_{\ell,k})$ over the
selected blocks.  Trace-norm convergence proves the first
identity in (M32.5).  The estimates also give uniform
tails on the stated bounded sets.  In particular no
unjustified passage from a Markushevich system to a
trace expansion is used.

\paragraph{Poisson summation with the normalization fixed.}
For a single $\theta$, periodize
$f(x)e^{i\theta x/h}$ with period $h$.
It is a smooth periodic function whose Fourier coefficients
are
\[
        h^{-1}M_f\bigl(i(\theta-2\pi k)/h\bigr).
\]
Its absolutely convergent Fourier series at zero gives
\[
 \sum_{k\in\mathbf Z}M_f\bigl(i(\theta+2\pi k)/h\bigr)
       =h\sum_{n\in\mathbf Z}f(nh)e^{i\theta n}.
 \tag{M32.17}
\]
The argument applies to exponential-Schwartz tests too,
since their periodization and all derivatives converge
uniformly.  Applying (M32.17) to each of the finitely
many roots proves the second equality in (M32.5).

If $r'_\ell=r_\ell$, count the roots with their
multiplicities in the original factorization of $Q$.
For a factor $1-a_vw^{d_v}$, its roots satisfy
$\lambda^{d_v}=a_v^{-1}$ and their power sum is
\[
 \sum_{\lambda^{d_v}=a_v^{-1}}\lambda^n
   =\begin{cases}
       d_va_v^{-n/d_v},&d_v\mid n,\\
       0,&d_v\nmid n.
     \end{cases}
 \tag{M32.18}
\]
This holds for positive, zero and negative integers $n$,
by the finite geometric-series identity on roots of unity.
Substituting in (M32.5) gives (M32.6).

\paragraph{Adelic normalization and geometric meaning.}
Here is an explicit conversion to a primary number-field
Euler product, with no hidden sign convention.
Choose a norm splitting
$C_F\simeq C_F^1\times\mathbf R$ and Haar measure
$dk\,dx$ with $\int_{C_F^1}dk=1$.
Fix a compact character $\beta$, extended to be trivial
on the chosen norm section, and write vectors in its
sector as $\Phi(k,x)=\beta(k)^{-1}u(x)$.
For a finite place $v\mid p$, write the uniformizer idele
as
\[
             \pi_v=(k_v,-f_v\log p).
\]
The literal operator in (M31.1) becomes
$u(x)-\beta(k_v)u(x+f_v\log p)$.
After the isometric reflection $x\mapsto-x$ this is
the factor in (M32.1) with
\[
              h=\log p,\quad d_v=f_v,
                         \quad a_v=\beta(k_v).
 \tag{M32.19}
\]
For a Bruhat test $\varphi\in\mathcal D(C_F)$, put
\[
 b_\beta(y)=\int_{C_F^1}\varphi(k,y)\beta(k)\,dk,
 \qquad
 \varphi_\beta(k,y)=\beta(k)^{-1}b_\beta(y).
 \tag{M32.20}
\]
The actual convolution on this sector is $C_{b_\beta}$;
after reflection its test is $f(y)=b_\beta(-y)$.
Consequently, when all primary jets survive, (M32.6) reads
\[
     \operatorname{Tr} A_\beta(\varphi)
       =\sum_{v\mid p}\log q_v
                          \sum_{n\in\mathbf Z}\varphi_\beta(\pi_v^n).
 \tag{M32.21}
\]
Only finitely many $n$ contribute for a compact radial
support.  If $\beta$ is unramified at $v$,
$\varphi_\beta$ is constant on the local unit group.
For any multiplicative Haar measure $\mu_v$, the
corresponding local expression is equivalently
\[
 \frac{\log q_v}{\mu_v(\mathcal O_{F_v}^{\times})}
          \int_{F_v^{\times}}\varphi_\beta(g_v)\,d\mu_v(g_v).
 \tag{M32.22}
\]
This specifies exactly where the local-unit volume
normalization enters.  The formulas do not silently
replace an integral by a sum while retaining its volume factor.

The finite matrix $G(u)$ is a concrete Hilbert model for
the finite local Euler polynomial $Q$, with polynomial
jet degrees retaining its nilpotent extensions.  Fourier
series on the norm circle supply the infinite lattice of
Euler poles.  The degree bound used for the local trace
therefore has a direct periodic-function explanation.
It supplies the corrected local quotient and trace intended
by Masullo Theorem 11.2.2.1, original pp.~490--491,
after MT-M31's closure correction and the strict jet cutoff.

The constants in (M32.13) are for a fixed prime product,
character and degree.  No uniform comparison as primes,
characters or global geometric truncations vary is asserted.
In particular this result does not show that the
localized kernel in Masullo Theorem 11.3.2.1,
pp.~495--496, acts on this corrected quotient in the
claimed way, nor that its nuclear limits commute with
the global shell degeneration.  That precise comparison
remains the subsequent geometric-trace gate.

\subsection*{MT-M33. Finite Euler quotients disappear in a separated strip}

This is an unconditional obstruction to a bounded module
identification, not an obstruction to every possible
cutoff or geometric-kernel calculation.

\paragraph{The actual norm and the centered Euler operator.}
For a compact nonempty interval $I=[u,v]$ and $s\geq0$ put
\[
 H_{I,s}=L^2\left(\mathbf R,
       \max(e^{2ux},e^{2vx})(1+x^2)^{s/2}\,dx\right),
 \qquad (U_yf)(x)=f(x-y).
 \tag{M33.1}
\]
The compact norm-one group may be included with its
volume-one Haar measure; the same estimates hold.
Formula (M33.1) is exactly the fixed-character source norm
(9.2.2.31), original PDF p.~416.  The supremum over
subintervals in that definition is attained at the
endpoints of this compact $I$.

Let $T_af(x)=e^{ax}f(x)$.  The source completion
(10.1.1.3), original p.~434, is the closure of
$T_a\Theta$ in (M33.1), where $\Theta$ is its smooth
theta image.  Both original pages were rendered and
read.  The equivariance convention stated just above
that definition agrees with
\[
 U_yT_a=T_a(e^{-ay}U_y),\qquad
            T_aU_yT_a^{-1}=e^{ay}U_y.
 \tag{M33.2}
\]
These identities fix the sign of the centering twist.

In a compact character sector, write a uniformizer idele
as $(k_v,-h_v)$, with $h_v=\log q_v>0$, and put
$b_v=\beta(k_v)$, so $|b_v|=1$.
The uncentered difference of MT-M31 is
\[
 P_{v,0}=1-b_vU_{-h_v}.
\]
Transporting this literal operator by $T_a$ gives
\[
 P_{v,a}=T_aP_{v,0}T_a^{-1}
                 =1-b_ve^{-ah_v}U_{-h_v}.
 \tag{M33.3}
\]
Its Mellin multiplier is
$1-b_ve^{-(a+z)h_v}$, whose zeros have real part $-a$.
The operator is not the unscaled uncentered difference
when $a\ne0$.

\paragraph{The inverse with a quantitative convergent bound.}
If $a+u>0$, (M33.3) is a bounded automorphism of
$H_{I,s}$, with inverse
\[
 P_{v,a}^{-1}
       =\sum_{n=0}^{\infty}b_v^ne^{-anh_v}U_{-nh_v},
 \tag{M33.4}
\]
where convergence is in operator norm.  More precisely,
\[
 \|P_{v,a}^{-1}\|
 \leq 2^{s/4}\sum_{n=0}^{\infty}
       (1+n^2h_v^2)^{s/4}e^{-(a+u)nh_v}<\infty.
 \tag{M33.5}
\]
If $a+v<0$, the inverse instead is
\[
 P_{v,a}^{-1}
      =-\sum_{n=1}^{\infty}b_v^{-n}e^{anh_v}U_{nh_v},
 \tag{M33.6}
\]
with bound obtained by replacing the final exponential
in (M33.5) by $e^{(a+v)nh_v}$ and starting at $n=1$.
Thus every finite product
$P_{S,a}=\prod_{v\in S}P_{v,a}$ is a bounded
automorphism whenever $0\notin a+I$.

\begin{proof}
For every $y\in\mathbf R$, the endpoint exponential
weight and the polynomial inequality from MT-M31 give
\[
 \|U_y\|_{H_{I,s}\to H_{I,s}}
 \leq2^{s/4}(1+y^2)^{s/4}
                      e^{\max(uy,vy)}.
 \tag{M33.7}
\]
Indeed, after the change of variables $x=t+y$, each
endpoint weight acquires its factor $e^{2uy}$ or
$e^{2vy}$, and
$1+(t+y)^2\leq2(1+t^2)(1+y^2)$.
For $B_v=b_ve^{-ah_v}U_{-h_v}$ this proves
\[
 \|B_v^n\|\leq2^{s/4}(1+n^2h_v^2)^{s/4}
                                      e^{-(a+u)nh_v}.
\]
If $a+u>0$, the series is summable and $B_v^n\to0$
in operator norm.  The telescoping identity
$(1-B_v)\sum_{n=0}^{N}B_v^n=1-B_v^{N+1}$,
and its commuting counterpart, prove (M33.4).
If $a+v<0$, apply the same argument to
$B_v^{-1}=b_v^{-1}e^{ah_v}U_{h_v}$; the identity
$(1-B_v)^{-1}=-\sum_{n\geq1}B_v^{-n}$ proves
(M33.6).  The factors and their inverses commute,
so their finite products are automorphisms too.
\end{proof}

\paragraph{All closed invariant subquotients have the same property.}
Let $Y_0\subset Y\subset H_{I,s}$ be closed subspaces
invariant under every $U_y$, $y\in\mathbf R$.
Then $P_{S,a}$ and its inverse preserve $Y_0$ and $Y$:
each term of (M33.4) or (M33.6) preserves them and
they are closed under the norm limit.
Consequently $P_{S,a}$ is an automorphism on $Y/Y_0$
with its quotient norm.  It is enough to assume
invariance under the particular shifts appearing in
the inverse series and their products.

In the actual arithmetic setting,
first work in one fixed compact character sector and put
\[
          J_{a,I,s}=\overline{T_a\Theta}^{\,H_{I,s}}
 \tag{M33.8}
\]
This space is invariant under all norm translations by (M33.2),
since $\Theta$ is translation invariant.
Therefore, without GRH,
\[
 P_{S,a}J_{a,I,s}=J_{a,I,s},\qquad
 \overline{T_aP_{S,0}\Theta}^{\,H_{I,s}}=J_{a,I,s}
                   \quad(0\notin a+I).
 \tag{M33.9}
\]
For the second equality, use
$T_aP_{S,0}=P_{S,a}T_a$, density in (M33.8), and
the fact that an automorphism maps a dense subspace
onto a dense subspace of its whole target.
In particular
\[
       J_{a,I,s}/\overline{T_aP_{S,0}\Theta}=0.
 \tag{M33.10}
\]
For $a=1/2$ and $I=[-\alpha,\alpha]$ this applies
to every $0\leq\alpha<1/2$, including the line
$\alpha=0$.  This concerns the finite Euler quotient,
not the quotient by the full global theta ideal:
the latter need not vanish.

The full compact-group assertion has the same proof:
replace each scalar $b_v$ by the corresponding compact
unitary translation.  Its integer powers have norm one,
so the bounds (M33.5)--(M33.7) are uniform in the compact
character.  For that version the closed subspaces must
also be invariant under these compact translations and
their inverses.  The full theta closure has this
invariance.  Norm-translation invariance alone is the
hypothesis for the fixed-character statement above.

\paragraph{The exact intertwining obstruction.}
Let $Z_{S,s}$ be the corrected nonzero local quotient of
MT-M32, in the original unreflected coordinates:
\[
       Z_{S,s}=H_s/\overline{P_{S,0}H_s}.
 \tag{M33.11}
\]
For instance one prime and $s>1$ gives the nonzero
scalar circle model (M32.14c).
Ordinary norm translations act on this quotient.
Give the same underlying Hilbert space its centered
action
\[
                   V_y=e^{-ay}U_y.
 \tag{M33.12}
\]
Then the polynomial in these operators
\[
       \prod_{v\in S}(1-b_ve^{-ah_v}V_{-h_v})
       =\prod_{v\in S}(1-b_vU_{-h_v})
       =0\quad\hbox{on }Z_{S,s}.
 \tag{M33.13}
\]
The last equality follows from its definition as a
quotient by the closure of the range, and does not
require the actual range to be closed.

Assume $0\notin a+I$, and let $W=Y/Y_0$ be any
closed invariant subquotient as above, with its
inherited norm-translation action.  Every bounded
linear map
\[
       T:(Z_{S,s},V)\longrightarrow W
       \quad\hbox{or}\quad
       T:W\longrightarrow(Z_{S,s},V)
 \tag{M33.14}
\]
which intertwines each norm translation is zero.
For the first direction, applying the same operator
polynomial to the intertwining relation gives
$P_{S,a}T=0$, while $P_{S,a}$ is injective on $W$.
For the second direction it gives $TP_{S,a}=0$,
while $P_{S,a}$ is surjective on $W$.
This proves both assertions.  The twist in (M33.12)
is essential: comparing the untwisted circle action
would compare a different polynomial and would not
prove this statement.

\paragraph{Location of the remaining source comparison.}
The source itself separates the spaces before attempting
the later comparison.  Theorem 11.1.2.1, p.~480,
is explicitly a theorem on the $0$-line.  The local
orbital quotient (11.2.2.1), p.~489, is also explicitly
defined with $a=0$, and Theorem 11.2.2.1,
pp.~490--491, gives its local trace.
In contrast, the absolute and invariant Paley--Wiener
constructions at pp.~487--489 allow a general center
$a$, while Theorem 11.3.2.1, pp.~495--496, sets
$a=1/2$ and permits $I=[-\alpha,\alpha]$ for
$\alpha\geq0$.  Its proof invokes the zero-line
carving theorem and the local orbital trace after
cutoffs and shell degeneration.  The exact norm and
centering convention at pp.~416 and 434 have been
visually checked for this comparison; the cited
transition statements were read in the supplied text.

Equations (M33.9)--(M33.14) show that the needed transition
cannot be a nonzero bounded identification of the
centered local circle module with a closed invariant
subquotient of the narrow Hilbert space.  Any valid
geometric calculation must instead justify the actual
non-equivariant cutoff operators and their limits.
Neither the local Poisson identity in MT-M32 nor the
conditional full narrow-strip spectral trace in MT-M28
supplies that justification.  No conclusion that the
Hecke-GRH endpoint is false is drawn here.

In the cote 1 trace programme (scans 79--80 and 111--112),
this locates the issue in transport from a local
closed-point Euler module to the global trace space.
The local Poisson calculation survives; the separated
strip changes the relevant translation spectrum, so
the geometric comparison needs its own cutoff-limit
argument.

\subsection*{MT-M34. Finite cutoffs erase a genuine Euler condition;
the exact cost of the escaping correction}

This result concerns the literal cutoff-strictness assertion in
Masullo's Theorem 11.1.2.1(2), original PDF p.~480.
It remains a counterexample after replacing every actual
Euler-difference range by its Hilbert closure as in MT-M31.
The uncut Euler quotient and its local trace survive, as
proved in MT-M32; the finite cutoff does not preserve that
quotient.

\paragraph{Statement.}
Fix $s>1$ and $h>0$, and retain the notation
\[
 w_s(x)=(1+x^2)^{s/2},\qquad
 H_s=L^2(\mathbf R,w_s(x)\,dx),\qquad
 U_hf(x)=f(x-h).
 \tag{M34.1}
\]
Let
\[
 V=\overline{(1-U_h)H_s}^{\,H_s},\qquad
 (Rf)(u)=\sum_{n\in\mathbf Z}f(u+nh),\quad0\le u<h.
 \tag{M34.2}
\]
The periodization is a bounded surjection onto $L^2(0,h)$,
and $V=\ker R$ is a proper closed subspace, by the exact
weighted quotient model (M32.14a)--(M32.14c).
If a measurable set $\Omega\subset\mathbf R$ has a
complement containing an interval of length $h$, up to
sets of measure zero, then its indicator multiplier satisfies
\[
               M_\Omega V=M_\Omega H_s.
 \tag{M34.3}
\]
This is equality of the actual images, not merely equality
of their closures.  In particular it holds for every $L>0$
and both
\[
 \Omega_{+,L}=[-L,\infty),\qquad
 \Omega_{0,L}=[-L,L].
 \tag{M34.4}
\]
The corresponding projections converge strongly to the
identity as $L\to\infty$, although $V\ne H_s$.

More generally, if $\rho\in L^\infty(\mathbf R)$ vanishes
almost everywhere on an interval of length $h$, then
\[
                   M_\rho V=M_\rho H_s.
 \tag{M34.5}
\]
This includes every bounded compactly supported cutoff.
In (M34.5) the common image need not be closed.

There is also an exact quantitative form of the obstruction.
Let $g\in H_s$ be compactly supported, put $F=Rg$, and
suppose $F\ne0$.  For either set in (M34.4), take $L$
large enough that $g$ is supported in $\Omega$.  Write
$E=\mathbf R\setminus\Omega$ and
\[
 W_E(u)=\sum_{\substack{n\in\mathbf Z\\u+nh\in E}}
                   \frac1{w_s(u+nh)}.
 \tag{M34.6}
\]
Then the minimum norm of a lift in the proper closed
subspace $V$ with prescribed cutoff $g$ is exactly
\[
 \min_{\substack{f\in V\\M_\Omega f=g}}\|f\|_{H_s}^2
       =\|g\|_{H_s}^2+
                      \int_0^h\frac{|F(u)|^2}{W_E(u)}\,du.
 \tag{M34.7}
\]
In particular, as $L\to\infty$,
\[
 \begin{split}
 \min_{\substack{f\in V\\M_{\Omega_{+,L}}f=g}}
                 \|f\|_{H_s}^2
  &=\|g\|_{H_s}^2+
        h(s-1)L^{s-1}\|F\|_{L^2(0,h)}^2(1+o(1)),\\
 \min_{\substack{f\in V\\M_{\Omega_{0,L}}f=g}}
                 \|f\|_{H_s}^2
  &=\|g\|_{H_s}^2+
       \frac{h(s-1)}2L^{s-1}\|F\|_{L^2(0,h)}^2(1+o(1)).
 \end{split}
 \tag{M34.8}
\]
Thus the norms themselves grow as $L^{(s-1)/2}$,
with leading constants $\sqrt{h(s-1)}\,\|F\|_2$
and $\sqrt{h(s-1)/2}\,\|F\|_2$, respectively.
For the source cutoff $L=\log\Lambda$ these are exact
leading constants in powers of $\log\Lambda$.

\paragraph{Repair outside the cutoff.}
For completeness, the scalar model of MT-M32 gives
\[
 W(u)=\sum_{n\in\mathbf Z}w_s(u+nh)^{-1},\qquad
 |Rf(u)|^2\le W(u)\sum_n|f(u+nh)|^2w_s(u+nh).
 \tag{M34.9}
\]
The positive continuous periodic function $W$ is bounded
above and below.  Therefore $R:H_s\to L^2(0,h)$ is
bounded.  The lift $f(u+nh)=F(u)/(W(u)w_s(u+nh))$
shows surjectivity.  The exact closure identity $V=\ker R$
is the already proved one in (M32.14c), including all
continuous Fourier coefficient constraints.  This also
shows that $V$ is proper.

Choose a half-open interval $I_0=[c,c+h)$ contained in
$\mathbf R\setminus\Omega$ up to a null set.  For almost
every $u\in[0,h)$ there is a unique integer $n(u)$ with
$u+n(u)h\in I_0$.  Define a lift $L_{I_0}$ by
\[
 (L_{I_0}F)(u+n(u)h)=F(u),\qquad
 L_{I_0}F=0\quad\hbox{off }I_0.
 \tag{M34.10}
\]
The piecewise translation $u\mapsto u+n(u)h$ preserves
Lebesgue measure.  Hence
\[
 \|L_{I_0}F\|_{H_s}^2
       \le\bigl(\sup_{x\in I_0}w_s(x)\bigr)\|F\|_2^2,
 \qquad R L_{I_0}F=F.
 \tag{M34.11}
\]
For a prescribed $g\in M_\Omega H_s$, the vector
\[
                        f=g-L_{I_0}Rg
 \tag{M34.12}
\]
belongs to $V$ and satisfies $M_\Omega f=g$.  This proves
(M34.3), and also supplies a bounded right inverse for each
fixed cutoff.  For (M34.5), start with any $f\in H_s$ and
replace it by $f-L_{I_0}Rf$ on an interval where $\rho=0$.
Its $M_\rho$ image is unchanged.  The reverse inclusions
in both image equalities follow from $V\subset H_s$.

For the two sets in (M34.4), the complements contain such
an interval for every finite $L$.  Their indicators increase
pointwise to one.  Dominated convergence applied to
$|f|^2w_s$ proves strong convergence of their multipliers
to the identity.  This proves all the qualitative assertions.

\paragraph{The minimum escaping correction.}
For a feasible $f$ in (M34.7), put $k=f-g$.  Then $k$ is
supported in $E$, $Rk=-F$, and the disjoint supports give
\[
                       \|f\|_{H_s}^2=\|g\|_{H_s}^2+
                                                    \|k\|_{H_s}^2.
 \tag{M34.13}
\]
For each fixed $L$, the measurable function $W_E$ is
positive and bounded above and below by positive constants
on $[0,h)$: the upper bound follows from $W_E\le W$;
for a lower bound, choose one full length-$h$ interval
in $E$ and use its bounded weight.  No continuity across
the endpoints modulo $h$ is needed.

Fiberwise Cauchy--Schwarz gives
\[
 |F(u)|^2\le W_E(u)
         \sum_{u+nh\in E}|k(u+nh)|^2w_s(u+nh).
 \tag{M34.14}
\]
After integration this is the lower bound in (M34.7).
Equality is attained by the measurable function
\[
 k(u+nh)=
 \begin{cases}
 -\dfrac{F(u)}{W_E(u)w_s(u+nh)},&u+nh\in E,\\[6pt]
 0,&u+nh\notin E.
 \end{cases}
 \tag{M34.15}
\]
The lower bound for $W_E$ proves that this function belongs
to $H_s$.  Its absolutely convergent periodization equals
$-F$ and its squared norm is the displayed integral in
(M34.7).  Equality in Cauchy--Schwarz on each fiber, or
strict convexity of the Hilbert norm, shows that this
minimum lift is unique up to equality almost everywhere.

Set $\varphi(t)=(1+t^2)^{-s/2}$ for $t\ge0$.
It is decreasing and
\[
 \int_L^\infty\varphi(t)\,dt
           =\frac{L^{1-s}}{s-1}(1+O(L^{-2})).
 \tag{M34.16}
\]
Indeed $\varphi(t)=t^{-s}(1+O(t^{-2}))$ uniformly for
$t\ge L\ge1$, and both powers can be integrated.
The decreasing-function comparison between a sum on a
lattice of step $h$ and its integral shows, uniformly in
the residue $u\in[0,h)$, that each one-sided tail satisfies
\[
 \sum_{u+nh>L}\varphi(u+nh)
   =\frac1h\int_L^\infty\varphi(t)\,dt+O(\varphi(L)),
 \tag{M34.17}
\]
with the same estimate for $u+nh<-L$ after reflection.
For clarity, the first point beyond $L$ lies between $L$
and $L+h$; comparing the decreasing sum with integrals
starting at that point incurs at most one summand, and
moving the integration endpoint back to $L$ incurs at
most $h\varphi(L)$.  These bounds are uniform in the
residue and also cover endpoint equalities.
Consequently
\[
 \begin{split}
 W_{(-\infty,-L)}(u)
    &=\frac{L^{1-s}}{h(s-1)}(1+o(1)),\\
 W_{\mathbf R\setminus[-L,L]}(u)
    &=\frac{2L^{1-s}}{h(s-1)}(1+o(1)),
 \end{split}
 \tag{M34.18}
\]
uniformly in $u$.  The uniform positive leading terms
allow reciprocation and integration against $|F|^2$.
Equations (M34.7) and (M34.18) prove (M34.8).

\paragraph{A precise missing bounded-lift hypothesis.}
Here is a useful general form of the failed inference.
Let $V$ be a closed subspace of a Hilbert space $H$, and
let orthogonal projections $E_n$ converge strongly to the
identity.  If a fixed $g\in H$ has lifts $v_n\in V$ with
\[
                 E_nv_n=E_ng,\qquad
                            \sup_n\|v_n\|<\infty,
 \tag{M34.19}
\]
then $g\in V$.  To prove this without a subsequence
argument, let $z\in V^\perp$.  Since
$\langle v_n-g,E_nz\rangle=0$,
\[
 |\langle g,z\rangle|
  =|\langle v_n-g,z\rangle|
  \le(\sup_n\|v_n\|+\|g\|)\,\|(1-E_n)z\|
          \longrightarrow0.
 \tag{M34.20}
\]
Thus $g\perp V^\perp$ and $g\in V$.  Equation (M34.8)
proves that this bounded-lift condition fails quantitatively
in the present cutoff problem whenever $Rg\ne0$.

\paragraph{Actual source specialization and exact refutation.}
Take $F=\mathbf Q$, the trivial continuous character of
$C_{\mathbf Q}$, $S=\{\infty\}$, and any rational prime
$p$.  Then $\operatorname{Pic}(\mathbf Z)=1$, the character
is unramified outside $S$, and the source index set
(11.1.2.1), p.~480, includes
\[
                         S_i=\varnothing,
                         \qquad S_{i'}=\{p\}.
 \tag{M34.21}
\]
These are distinct with $S_i\subset S_{i'}$; no place
of $S_i$ divides $p$, while a place of $S_{i'}$ does.
Choose any $s\ge2=2[\mathbf Q:\mathbf Q]$.
Thus every hypothesis of Theorem 11.1.2.1(2) is satisfied.

Normalize Haar measure on $C_{\mathbf Q}^1$ to have
mass one and use $x=\log|g|$ in its trivial sector.
The zero-line Hilbert space is exactly $H_s$.
Let $\Theta$ be the actual smooth theta image in this
sector.  MT-M24 and MT-M25 identify its closure with
$H_s$: the pole-removed completed rational zeta function
has no zero on $\Re z=0$, using the classical zero-free
boundary $\Re z=1$ and its functional equation.
This use of the boundary zero-free theorem requires no
GRH.  The source's one-prime operator is
$P_p=1-U_{-h}$ with $h=\log p$, as checked in MT-M31.
Since
\[
                    1-U_{-h}=-U_{-h}(1-U_h)
 \tag{M34.22}
\]
and $\ker R$ is translation invariant,
$\overline{P_pH_s}=V$; density and boundedness also give
$\overline{P_p\Theta}=V$.

Original Definition 9.3.3.3 and formulas
(9.3.3.4)--(9.3.3.5), p.~428, say that for $\Lambda>1$
the two cutoffs in Theorem 11.1.2.1 are exactly
$M_{\Omega_{+,L}}$ and $M_{\Omega_{0,L}}$ with
$L=\log\Lambda$.  The source defines its cutoff spaces
in (11.1.2.3), p.~480, as the Hilbert closures of the
cutoff smooth images.  For either multiplier $M_\Omega$,
boundedness and (M34.3) therefore imply
\[
 \overline{M_\Omega P_p\Theta}^{\,H_s}
       =\overline{M_\Omega V}^{\,H_s}
       =M_\Omega H_s
       =\overline{M_\Omega\Theta}^{\,H_s}.
 \tag{M34.23}
\]
The common subspace is closed because $M_\Omega$ is an
orthogonal projection.  Thus the source inclusion
(11.1.2.6) is equality for \emph{every} finite $\Lambda>1$,
for both stipulated types of cutoff.  In particular, for
the non-decreasing sequence $\Lambda_k=e^{k+1}$, no
subsequence has the strict inclusion asserted in
Theorem 11.1.2.1(2).  The uncut inclusion (11.1.2.5)
is nevertheless strict, since $V\ne H_s$.

The source's proof on p.~482, immediately after
(11.1.2.16), infers non-strictness of the uncut inclusion
from equality at each cutoff and strong convergence of
the cutoff multipliers to the identity.  Equations
(M34.3)--(M34.8) disprove exactly that inference and
identify the missing uniform control of its lifts.
The source later invokes Theorem 11.1.2.1 in the proof
of its geometric trace Theorem 11.3.2.1, pp.~495--496;
the main Theorem 12.5.1.1, p.~520, in turn invokes
Theorem 11.3.2.1.  This is a targeted dependency trace,
not a proof audit of the remaining kernel arguments.

\paragraph{Input and placement.}
The new cutoff and norm arguments above are elementary
and unconditional; the exact Euler closure and actual
theta-density identification are explicit antecedents
MT-M31--MT-M32 and MT-M24--MT-M25.  Original PDF
pp.~428, 480 and 482 were rendered and visually read
to check the definitions, full strictness hypotheses and
the disputed implication.  MT-M34 refutes this literal
intermediate assertion and proves its quantitative
replacement.  It does not assert that the conditional
Hecke-GRH endpoint $M_H$ is false, or that a different
geometric trace construction is impossible.

For cote 1, scans 79--80 and 111--112, this distinguishes
the trace of a genuine local Euler quotient from a
trace attributed to its finite radial cutoffs.  The
quotient supplies the local closed-point contribution;
the finite cutoff images alone carry no such quotient
condition.  A comparison of these traces therefore
requires quantitative control beyond strong convergence.

\subsection*{MT-M35. Finite Euler products become dense after cutoff;
prime exhaustion and cutoff do not commute}

MT-M34 proves actual cutoff-image equality for one Euler
factor.  The following different argument proves equality
of \emph{closed} cutoff images for every nonzero finite
translation polynomial, including Euler products with
incommensurate logarithms.  It does not assert actual-image
surjectivity for these general products.

\paragraph{Finite-translation statement.}
Let $s\ge0$, and set
\[
 H_s=L^2(\mathbf R,(1+x^2)^{s/2}\,dx),\qquad
 U_tf(x)=f(x-t).
 \tag{M35.1}
\]
Let $P$ be a nonzero finite translation polynomial.  Combine
coincident shifts first, delete zero coefficients, and write
\[
                P=\sum_{j=0}^{m}c_jU_{t_j},\qquad
          t_0<t_1<\cdots<t_m,\quad c_j\ne0.
 \tag{M35.2}
\]
Every term is bounded on $H_s$ and $H_{-s}$, by the
polynomial-weight translation bound (M25.7).
Let $\rho\in L^\infty(\mathbf R)$ have essential support
contained in some half-line, in either orientation.  Then
\[
 \overline{M_\rho P H_s}^{\,H_s}
                 =\overline{M_\rho H_s}^{\,H_s}.
 \tag{M35.3}
\]
If $\Theta\subset H_s$ is any dense linear subspace,
the left side is also $\overline{M_\rho P\Theta}$.
For an indicator $\rho=\mathbf1_\Omega$, with $\Omega$
essentially contained in a half-line, this says
\[
 \overline{M_\Omega P\Theta}^{\,H_s}
      =\overline{M_\Omega\overline{P H_s}}^{\,H_s}
      =M_\Omega H_s.
 \tag{M35.4}
\]
In particular it applies to both source cutoffs
$\Omega_{+,L}=[-L,\infty)$ and $\Omega_{0,L}=[-L,L]$.

\paragraph{Endpoint propagation for the weighted dual.}
We first prove the elementary unique-continuation fact
needed here.  Suppose a measurable locally integrable
function $g$ satisfies, almost everywhere,
\[
                 \sum_{j=0}^{m}d_jg(x-\tau_j)=0,
 \quad \tau_0<\cdots<\tau_m,\quad d_j\ne0,
 \tag{M35.5}
\]
and is supported in a half-line.  Then $g=0$ almost
everywhere.  The assertion for $m=0$ is immediate.
For $m\ge1$, translate the equation so that $\tau_0=0$,
and write $H=\tau_m>0$.

If $g=0$ almost everywhere on $(-\infty,A)$, let
$\varepsilon=\tau_1>0$.  For almost every
$x\in[A,A+\varepsilon)$ all the terms with $j\ge1$
evaluate $g$ below $A$.  Since $d_0\ne0$, equation
(M35.5) gives $g(x)=0$ there.  If by induction $g=0$
below $A+n\varepsilon$, the same equation on
$[A+n\varepsilon,A+(n+1)\varepsilon)$ has every
$x-\tau_j<A+n\varepsilon$ for $j\ge1$, and gives
zero on the next interval.  Countably many such steps
prove $g=0$ on the whole line.

If $g=0$ almost everywhere on $(B,\infty)$, put
$\varepsilon'=H-\tau_{m-1}>0$ and evaluate (M35.5)
at $x=y+H$.  It becomes
\[
 d_mg(y)=-\sum_{j<m}d_jg(y+H-\tau_j).
 \tag{M35.6}
\]
For almost every $y\in(B-\varepsilon',B]$, all
arguments on the right exceed $B$.  Induction in the
opposite direction, on consecutive intervals of length
$\varepsilon'$, proves the result.  Finite translations
of the null exceptional sets remain null at each step,
and their countable union is still null.  Combining the
coincident shifts before selecting the two endpoint gaps
is essential to the proof as stated.

We now use the continuous \emph{bilinear} pairing
\[
            \langle f,g\rangle=\int_{\mathbf R}f(x)g(x)\,dx
             \quad (f\in H_s,\ g\in H_{-s}).
 \tag{M35.7}
\]
It identifies $H_{-s}$ with the continuous linear dual
of $H_s$, with complex conjugation absorbed in the
representing function.  Consequently
\[
              P^{\mathrm t}=\sum_{j=0}^{m}c_jU_{-t_j},
              \qquad M_\rho^{\mathrm t}=M_\rho.
 \tag{M35.8}
\]
There is no coefficient conjugation in this convention.
All functions in $H_{-s}$ are locally integrable, since
the weight is bounded above and below on compact sets.

If $g\in H_{-s}$ annihilates $M_\rho PH_s$, then
$P^{\mathrm t}(\rho g)=0$ in $H_{-s}$, hence almost
everywhere.  The function $\rho g$ lies in $H_{-s}$ and
has essential support in a half-line.  Applying
(M35.5)--(M35.6), with the reflected and sorted shifts
from (M35.8), gives $\rho g=0$.  Thus $g$ annihilates
$M_\rho H_s$.  Conversely every annihilator of the
latter image plainly annihilates the former.  Equality
of continuous annihilators gives equality of the closed
linear subspaces, proving (M35.3).  This last assertion
also follows directly from Hilbert orthogonal
complements after converting the bilinear dual to the
usual conjugate-linear Riesz representation.

Boundedness of $M_\rho P$ and density of $\Theta$ prove
the dense-domain assertion.  For an indicator multiplier,
$M_\Omega$ is an orthogonal projection, so its image is
closed.  Boundedness again gives
$\overline{M_\Omega PH_s}=
\overline{M_\Omega\overline{PH_s}}$.
This proves (M35.4).

\paragraph{Finite Euler products and the cofinal cutoff intersection.}
For a finite set $T$ of rational primes, define
\[
                P_T=\prod_{p\in T}(1-U_{-\log p}),
                   \qquad P_\varnothing=1.
 \tag{M35.9}
\]
These are nonzero translation polynomials.  More generally
any product $\prod_v(1-a_vU_{h_v})$, where $h_v>0$ and
$a_v\ne0$, has nonzero constant and largest-shift
coefficients; combining coincident interior shifts cannot
cancel those endpoints.  Reflecting all shifts gives the
same conclusion for (M35.9).  No rational dependence or
independence of the shifts is used in (M35.3).

Let $\Theta$ be the actual trivial compact-character
theta image for $\mathbf Q$, and take $s>1$.
MT-M24--MT-M25, as used in MT-M34, give
$\overline\Theta^{\,H_s}=H_s$ unconditionally from the
classical zero-free boundary of the completed zeta function.
Put
\[
 V_T=\overline{P_T\Theta}^{\,H_s}
                 =\overline{P_TH_s}^{\,H_s}.
 \tag{M35.10}
\]
If $T\subset T'$, the factorization
$P_{T'}=P_TP_{T'\setminus T}$ gives $V_{T'}\subset V_T$.
For either source cutoff and every finite $T$, equation
(M35.4) gives
\[
                  \overline{M_\Omega P_T\Theta}^{\,H_s}
                                  =M_\Omega H_s.
 \tag{M35.11}
\]
Thus for any increasing cofinal sequence $T_k$ of finite
sets of rational primes,
\[
 \bigcap_{k\ge0}
       \overline{M_\Omega P_{T_k}\Theta}^{\,H_s}
                                   =M_\Omega H_s.
 \tag{M35.12}
\]
Every finite stage is the displayed full cutoff space;
there is no exchange of an image with an intersection
in this argument.  Any quotient between two of these
closed cutoff stages is zero.

\paragraph{The uncut cofinal intersection is zero.}
In contrast, for the same $s>1$,
\[
                         \bigcap_{k\ge0}V_{T_k}=\{0\}.
 \tag{M35.13}
\]
Indeed $\int_{\mathbf R}(1+x^2)^{-s/2}dx<\infty$,
so Cauchy--Schwarz gives a continuous embedding
$H_s\subset L^1(\mathbf R)$.  Each functional
\[
                  J_t(f)=\int_{\mathbf R}f(x)e^{itx}\,dx
 \tag{M35.14}
\]
is continuous on $H_s$, and $t\mapsto J_t(f)$ is
continuous by dominated convergence.  The norm-translation
convention gives
\[
       J_t(U_{-\log p}f)=e^{-it\log p}J_t(f).
 \tag{M35.15}
\]
Hence if $p\in T$, every $f\in V_T$ satisfies
$J_{2\pi n/\log p}(f)=0$ for all $n\in\mathbf Z$.
If $f$ belongs to the intersection in (M35.13), cofinality
makes these vanishings valid for every rational prime $p$.

The union of these frequency grids is dense in
$\mathbf R$.  For any fixed $t$, choose primes
$p_j\to\infty$ and choose an integer $n_j$ nearest to
$t\log p_j/(2\pi)$; then
\[
       \left|t-\frac{2\pi n_j}{\log p_j}\right|
                                  \le\frac\pi{\log p_j}\longrightarrow0.
 \tag{M35.16}
\]
Continuity therefore implies $J_t(f)=0$ for every real
$t$.  To recall the uniqueness argument in this setting,
Fourier inversion for a smooth compactly supported test
$\phi$ and Fubini give $\int f\phi=0$: the Fourier
transform of $\phi$ is integrable and $f\in L^1$.
Thus the distribution defined by $f$ vanishes, and
$f=0$ almost everywhere.  This proves (M35.13).

Equations (M35.12)--(M35.13) locate the precise failure
of a cutoff/prime-exhaustion interchange.  For every
fixed finite $L$, the cofinal cutoff intersection is
the full cutoff space; these spaces increase with $L$
and their union is dense in $H_s$.  The uncut cofinal
intersection is zero.  The result concerns subspaces
and their specified closures; it does not identify
the limits of any separately defined kernel operators.

\paragraph{Comparison with the literal Schwartz image in
Theorem 11.1.2.1(3).}
The source's original PDF p.~480 defines the cutoff
spaces in (11.1.2.3) exactly as the closures in (M35.11).
Its cofinal set-theoretic limit (11.1.2.4) is therefore
(M35.12), for $F=\mathbf Q$, $S=\{\infty\}$ and the
trivial character.  Its assertion (3), however, is an
identification with the closure of the specific Schwartz
image (11.1.2.7).  That separate assertion is \emph{not}
refuted by the present example: its stated right side is
also the full cutoff space, as we now check directly.

Every idele class of $\mathbf Q$ has a representative
with positive real component $t=|g|$ and all finite
components in $\mathbf Z_p^\times$: multiply by a rational
number to cancel the finitely many nonzero valuations,
and choose its sign to make the real component positive.
For $S=\{\infty\}$, the sum in (11.1.2.7), before
its cutoff, is then
\[
 \sum_{\xi\in\mathbf Q^\times}
   \left(f_\infty\otimes\prod_p\mathbf1_{\mathbf Z_p^\times}\right)
                   (\xi g)
                          =f_\infty(t)+f_\infty(-t).
 \tag{M35.17}
\]
Indeed only $\xi=1$ and $\xi=-1$ are units at every
finite prime.  This function is in the trivial compact
sector.  The source conditions in this case are
\[
 f_\infty\in\mathcal S(\mathbf R),\qquad
                  f_\infty(0)=0,\qquad
                  \widehat f_\infty(0)=\int_{\mathbf R}f_\infty=0.
 \tag{M35.18}
\]

Fix $L>0$ and take $u\in C_c^\infty(\Omega^\circ)$
for either $\Omega=\Omega_{+,L}$ or $\Omega_{0,L}$.
For $t\ne0$ define
\[
                    f_0(t)=\tfrac12u(\log|t|),
                    \qquad f_0(0)=0.
 \tag{M35.19}
\]
This is an even smooth compactly supported function
whose support is separated from zero, so
$f_0(e^x)+f_0(-e^x)=u(x)$ and $f_0(0)=0$.
Choose an even smooth compactly supported function $r$
with support in $0<|t|<e^{-L}$ and with $\int r\ne0$.
Then
\[
                    f_\infty=f_0-
                            \frac{\int f_0}{\int r}\,r
 \tag{M35.20}
\]
satisfies both conditions in (M35.18).  The repair term
in (M35.17) is supported in $x<-L$, outside either
cutoff.  Consequently the source's cutoff Schwartz
image contains $u$.  Smooth compactly supported
functions in $\Omega^\circ$ are dense in $M_\Omega H_s$:
truncate away from the endpoints and infinity, and then
use ordinary smooth approximation where the weight is
bounded above and below.  Thus the closure of the
right side in (11.1.2.7) is precisely $M_\Omega H_s$,
as claimed here.  No further refutation of part (3)
follows from this rational example.

\paragraph{Input and placement.}
The endpoint-propagation and cutoff-density theorems
are unconditional Hilbert-space deductions, with no
individual-zero or global-Riesz-basis assumption.
The actual theta-density input is the already proved
MT-M24--MT-M25.  Original PDF pp.~428, 480 and 482
were visually checked in MT-M34; the displayed formula
(11.1.2.7) and all its source hypotheses are retained
in the explicit comparison above.  MT-M35 strengthens
the corrected cutoff analysis without declaring the
conditional Hecke-GRH endpoint $M_H$ false or replacing
the remaining geometric-kernel comparison by a bounded
equivariant map.

\subsection*{MT-M36. The literal section average is zero;
normalized compact projections and genuine idelic cutoffs}

This is a targeted measure and normalization check of
Masullo's formula (11.3.2.2), original PDF p.~495.
The formula uses ordinary indicator functions and the
ordinary Haar integral on the idelic fundamental domain.
It does not display a Dirac distribution or a rescaled
limit of neighborhoods.  Under precisely that
interpretation its value is identically zero.
We prove this, and then give two legitimate replacements
with their exact scopes and trace properties.

\paragraph{The source's measure conventions.}
Write $\mathbb I_F=\mathrm{GL}_1(\mathbf A_F)$,
$\mathbb I_F^1=\ker|\cdot|_F$, and let
$\tau:\mathbf R_{>0}\to\mathbb I_F$ be a continuous
homomorphic section of the norm.  The source defines
\[
 E_F^1=\bigcup_{i=1}^{h_F}
              g_i(E_0\times\widehat{\mathcal O}_F^\times)
                         \subset\mathbb I_F^1
 \tag{M36.1}
\]
in Definition A.2.3.6, (A.2.3.14), original p.~549,
and explicitly gives this set the subspace topology.
It is not assigned the discrete topology.  Definition
A.3.2.1, pp.~559--560, uses Haar measures on ideles
and their class groups; the counting measure is on
$F^\times$.  The same distinction is explicit in
the compact-character integration of \S2.2.5.
The normalization (A.3.2.1), p.~560, gives volume one
on $C_{\mathbf Q}^1$.  The general measure discussion
in \S12.2, p.~504, is ordinary completed Radon/Lebesgue
integration of measurable functions.

The displayed two-variable function (11.3.2.2) is
therefore, with its suppressed Haar differential written,
\[
 B_0(x,y)=\sum_{\xi\in F^\times}\int_{E_F^1}
    \mathbf1_{\tau(\mathbf R_{>0})}(\xi gx)
    \mathbf1_{\tau(\mathbf R_{>0})}(\xi gy)\,d\mu_1(g).
 \tag{M36.2}
\]
The tensor symbol in the source just separates the two
named variables of these indicator functions.
For every $x,y\in\mathbb I_F$,
\[
                              B_0(x,y)=0.
 \tag{M36.3}
\]

\begin{proof}
Haar measure on $\mathbb I_F^1$ has no atoms.
Indeed this group contains an infinite compact subgroup:
embed $\mathcal O_{F_v}^\times$ at any one finite place
and use $1$ at every other place.  If one singleton had
positive Haar measure $c$, every singleton would have
the same measure, and $N$ distinct points in this compact
subgroup would have measure $Nc$ for every $N$.
This contradicts finiteness of a Radon measure on a
compact set.  Consequently every singleton in $E_F^1$
is null as well.  Equivalently, the quotient Haar
measure pulled back to a fundamental domain is atomless;
the discrete quotient by $F^\times$ does not change this
fact.  For $F=\mathbf Q$ this is particularly explicit:
$E_{\mathbf Q}^1=\{1\}\times\widehat{\mathbf Z}^\times$,
with its usual compact, atomless, volume-one Haar measure.

The section image is closed, hence Borel: it is the
fixed-point set of the continuous map
$z\mapsto\tau(|z|_F)$ on the Hausdorff group
$\mathbb I_F$.  Fix $\xi$ and $x$.  If
$\xi gx=\tau(t)$ with $g\in E_F^1$, the product formula
and $|g|_F=1$ force $t=|x|_F$.  Hence the only possible
value of $g$ is
\[
                       g=\xi^{-1}\tau(|x|_F)x^{-1}.
 \tag{M36.4}
\]
The first indicator in the integrand of (M36.2) is
supported on at most this one point.  Its product with
the second indicator has integral zero.  The field
$F$ is countable, so (M36.2) is a countable sum of zeros.
This proves (M36.3), without any interchange of
conditionally convergent expressions.
\end{proof}

\paragraph{The precise normalization failure.}
The source's indentation condition (12.2.1.2), original
p.~505, requires an orbit integral equal to one.
For a discrete subgroup $H$ this reads
$\sum_{\xi\in H}b(\xi x)=1$, with $b\ge0$.
Its two-variable lift is the averaged product
\[
                    B_b(x,y)=\sum_{\xi\in H}b(\xi x)b(\xi y),
 \tag{M36.5}
\]
as in (12.2.2.1), p.~508.  Nonnegativity and reindexing
give the necessary orbit identity
\[
 \sum_{\eta\in H}B_b(x,\eta x)
    =\left(\sum_{\xi\in H}b(\xi x)\right)^2=1.
 \tag{M36.6}
\]
Thus the zero function in (M36.3) cannot provide this
normalized lift.  The source's actual lifting rule
(12.2.2.2), p.~508, is multiplication of the pulled-back
kernel by the two-variable factor.  Inserting (M36.2)
into that rule produces the zero kernel.

This also gives a concrete nonzero test for the literal
interface in Theorem 11.3.2.1.  Take $F=\mathbf Q$,
the trivial character, $S=\{\infty\}$, $s=2$, and any
allowed interval $I$.  Fix a rational prime $p$ and
choose a nonnegative $\psi\in C_c^\infty((0,\infty))$
supported in a small neighborhood of $\log p$, containing
no other element of $\{n\log p:n\in\mathbf Z\}$, with
$\psi(\log p)>0$.  The test
\[
                             \varphi(g)=\psi(\log|g|)
 \tag{M36.7}
\]
is compactly supported on $C_{\mathbf Q}$, is invariant
under every finite local unit group, and satisfies all
the displayed ramification and support hypotheses.
For a local uniformizer $\pi_p$, whose idele norm is
$p^{-1}$, the right side of (11.3.2.1) is
\[
       \frac{\log p}{\mu_p(\mathbf Z_p^\times)}
                           \psi(\log p)>0,
 \tag{M36.8}
\]
because only the power $\pi_p^{-1}$ contributes.
Under the stated multiplication rule the left side
with (M36.2) is zero.  This refutes the literal
normalization/interface, not a reinterpretation using
a separately specified distributional regularization.

\paragraph{A normalized finite compact-level replacement.}
Let $K$ be an infinite compact metrizable profinite
abelian group with probability Haar measure; in the
actual rational example $K=C_{\mathbf Q}^1\simeq
\widehat{\mathbf Z}^\times$.  For an open subgroup
$U\subset K$, put $m_U=\mu(U)=1/[K:U]$ and define
\[
 a_U(k)=\frac{\mathbf1_U(k)}{m_U},\qquad
 B_U(k,l)=\int_Ka_U(rk)a_U(rl)\,d\mu(r).
 \tag{M36.9}
\]
Then
\[
 B_U(k,l)=\frac{\mathbf1_U(kl^{-1})}{m_U}
         =\sum_{\chi\in\widehat{K/U}}\chi(k)\overline{\chi(l)}.
 \tag{M36.10}
\]
The first equality follows because the two cosets of
$U$ in the integrand are either disjoint or identical;
in the latter case their intersection has measure $m_U$.
The second is the elementary character orthogonality
identity on the finite abelian group $K/U$.

Consequently the integral operator with kernel $B_U$
is exactly the orthogonal projection
\[
      (E_Uf)(k)=\frac1{m_U}\int_U f(ku)\,d\mu(u)
                    \quad\hbox{on }L^2(K).
 \tag{M36.11}
\]
Its range is the space of $U$-invariant functions,
and
\[
 \operatorname{rank}E_U=\operatorname{Tr}E_U=m_U^{-1},
       \quad\|E_U\|=1,\quad
       \|B_U\|_{L^2(K\times K)}^2=m_U^{-1}.
 \tag{M36.12}
\]
All these assertions follow either from (M36.10) or
from the orthonormal basis of normalized coset indicators.

If the neighborhood indicators are left unnormalized,
their averaged product is instead
\[
 \int_K\mathbf1_U(rk)\mathbf1_U(rl)\,d\mu(r)
       =m_U\mathbf1_U(kl^{-1})=m_U^2 B_U(k,l).
 \tag{M36.13}
\]
Its operator norm is $m_U^2$ and its trace is $m_U$.
Along a nested sequence of open subgroups $U_j$ forming
a neighborhood base at $1$, one has $m_{U_j}\to0$,
because their intersection is the Haar-null singleton.
Thus these unnormalized operators tend to zero even
in trace norm.  The normalized $E_{U_j}$ instead
converge strongly to the identity.  To prove the latter,
locally constant functions are dense in $L^2(K)$:
continuous functions are dense, and uniform continuity
approximates each by functions constant on the cosets
of a sufficiently small open subgroup.  Every such
function is fixed by all sufficiently late $E_{U_j}$,
and the projections have norm one.  Density proves
strong convergence.  There is no operator-norm
convergence to the identity, since each range is finite
dimensional whereas $L^2(K)$ is infinite dimensional.

These projections are a genuine normalized regularization
of the compact section.  Their diverging diagonal values
and Hilbert--Schmidt norms in (M36.12) also show why they
cannot be replaced by an ordinary point-indicator kernel.

\paragraph{Trace compatibility on the actual Hilbert modules.}
The same averaging defines orthogonal projections $E_U$
for any strongly continuous unitary representation of
$K$ on a Hilbert space $\mathcal H$.  They converge
strongly to the identity as $U$ shrinks to $1$:
\[
 \|E_Uv-v\|
       \le\sup_{u\in U}\|R(u)v-v\|\longrightarrow0.
 \tag{M36.14}
\]
This includes the compact action on the weighted idelic
Hilbert spaces and their closed invariant quotients:
the norm weight depends only on the radial variable.
If $T$ is trace class, then
\[
 \|E_UT-T\|_1\to0,\qquad
 \|TE_U-T\|_1\to0,\qquad
             \operatorname{Tr}(E_UT)\to\operatorname{Tr}(T).
 \tag{M36.15}
\]
For a complete proof, write a singular-value expansion
$T=\sum_ns_n u_n\otimes v_n^*$ with
$\sum_ns_n<\infty$ and unit vectors.  The first
trace-norm difference is at most
$\sum_ns_n\|(E_U-1)u_n\|$, which tends to zero by
(M36.14) and domination by $2s_n$.  For the second
use self-adjointness of $E_U$ and the same argument
on $v_n$.  Continuity of the trace gives the last claim.

There is exact stabilization for a test already at
finite compact level.  If a convolution test $\varphi$
is invariant under an open subgroup $U_0\subset K$,
its convolution operator $A_\varphi$ on any of these
modules satisfies
\[
        E_U A_\varphi=A_\varphi E_U=A_\varphi
                          \qquad(U\subset U_0).
 \tag{M36.16}
\]
Indeed translation of the test by any $u\in U_0$
does not change its convolution operator; integration
over $U$ gives both identities, using commutativity.
Thus if $A_\varphi$ is trace class its trace is preserved
exactly by all these finite-level projections.  On the
arithmetic quotients that trace-class hypothesis is
available unconditionally in the wide strips of MT-M27
and, with its explicit Hecke-GRH premise, in the narrow
positive-width strips of MT-M28.  Formula (M36.16) does
not supply that analytic hypothesis by itself.
For general number fields the compact group can have
connected archimedean factors; the finite-rank assertions
(M36.10)--(M36.12) here are expressly for the profinite
group, not for an untruncated archimedean character family.

\paragraph{A genuine discrete-idelic indentation and its orbital allocation.}
There is a different, also explicit, correction for
the discrete $\mathbf Q^\times$-normalization in
(M36.5)--(M36.6).  Use the standard section and set
\[
 D=\tau(\mathbf R_{>0})\,\widehat{\mathbf Z}^\times
       \subset\mathbb I_{\mathbf Q},\qquad b=\mathbf1_D.
 \tag{M36.17}
\]
The set $D$ is an open subgroup, hence also closed.
Every idele has a unique expression $q d$ with
$q\in\mathbf Q^\times$ and $d\in D$: cancel its
finite valuations by a rational number and then fix
the sign of its real component.  Uniqueness follows
from $\mathbf Q^\times\cap D=\{1\}$.
The map $D\to C_{\mathbf Q}$ is a homeomorphism;
openness follows from the quotient map and the
openness of $D$.  Therefore $b$ is continuous and
nonnegative, satisfies $\sum_qb(qx)=1$, and its support
meets the inverse image of any compact subset of
$C_{\mathbf Q}$ in a compact set.  These are exactly
the three indentation conditions of Lemma 12.2.1.1.
No auxiliary connected-domain assumption is needed
for this direct verification.

For $x\in D$, its two-variable lift has
\[
                   B_b(x,qx)=\mathbf1_{q=1}.
 \tag{M36.18}
\]
Thus it provides the normalized total orbital integral,
but places all of that diagonal mass in the identity
translate.  This allocation is not forced by the
normalization.  For a rational prime $p$ and
$0\le t\le1$, another valid indentation is
\[
                    b_t=(1-t)\mathbf1_D+t\mathbf1_{pD}.
 \tag{M36.19}
\]
Its support is a union of two clopen translates of
$D$, it has the same compactness property, and
$\sum_qb_t(qx)=1$.  Direct substitution, for $x\in D$,
gives
\[
 B_{b_t}(x,qx)=
 \begin{cases}
       (1-t)^2+t^2,&q=1,\\
       t(1-t),&q=p\text{ or }p^{-1},\\
       0,&\text{otherwise}.
 \end{cases}
 \tag{M36.20}
\]
The sum of these values is one.  Consequently, when
an integrable quotient diagonal is pulled back and
multiplied by $B_{b_t}$, its total orbital integral is
unchanged, while its individual rational-translate
contributions vary with $t$.  This follows directly
from (M36.20), since $x$ and $qx$ represent the same
class; no trace-limit interchange is used.

The normalized compact projection and the genuine
discrete-idelic indentation solve different problems.
Neither calculation identifies the source's desired
prime-by-prime geometric mass with the local Poisson
trace of MT-M32.  A correct comparison must specify
its normalized lift and prove the required orbital
allocation and limiting identity.  Positivity or the
conditional Hecke-GRH endpoint does not follow from
the measure repair alone.

\paragraph{Audit scope.}
Original PDF pp.~495, 505, 508, 549 and 560 were
rendered and visually read.  The definitions and
normalization text on pp.~504--506 and 559--560 were
also read from the supplied mathematical text.
This is not an audit of all the kernel theorems in
Chapter 12.  The literal failure (M36.3) and all the
replacement constructions above are unconditional.
The endpoint $M_H$ remains conditional throughout.

For cote 1, scans 79--80 and 111--112, the distinction
is between a normalized lift of a quotient trace and
its claimed local closed-point contributions.  The
former is available in the two explicit forms above;
the latter still requires a canonical allocation and
a proved geometric comparison.

\subsection*{MT-M37. Compatible nonnegative indentations cannot carry
the external prime shells}

The normalized replacements in MT-M36 do not by themselves
establish an orbital decomposition.  Here is a stronger
obstruction to one particular repair: a \emph{single}
nonnegative indentation normalized both globally and on
the initial local subgroup must vanish outside that subgroup.
This is exactly the simultaneous normalization offered in
Masullo's Proposition 12.2.2.4, original pp.~512--513.
An exact defect identity quantifies the obstruction
when the local normalization is relaxed.

\paragraph{The two normalizations force support.}
Let $G$ be an abelian group, $H\subset G$ a countable
subgroup, and $G_0\subset G$ a subgroup such that
\[
                 G=HG_0,\qquad H_0=H\cap G_0.
 \tag{M37.1}
\]
Suppose $b:G\to\mathbf R_{\ge0}$ satisfies
\[
 \sum_{h\in H}b(hx)=1\quad(x\in G),\qquad
 \sum_{h\in H_0}b(hx)=1\quad(x\in G_0).
 \tag{M37.2}
\]
Then
\[
                       b(y)=0\qquad(y\notin G_0).
 \tag{M37.3}
\]
For its averaged product
\[
                     B_b(x,y)=\sum_{h\in H}b(hx)b(hy),
 \tag{M37.4}
\]
one has
\[
                  B_b(x,y)=0\qquad(x^{-1}y\notin G_0).
 \tag{M37.5}
\]
In particular, if $G_0\subset G_k\subset G_{k+1}$ and
$\alpha_k=\mathbf1_{G_k}$, then
\[
       (\alpha_{k+1}(y)-\alpha_k(y))B_b(x,y)=0
                                      \quad(x\in G_0,\ y\in G).
 \tag{M37.6}
\]

\begin{proof}
For $x\in G_0$, subtract the second identity in (M37.2)
from the first.  Nonnegativity gives
\[
                       \sum_{h\in H\setminus H_0}b(hx)=0,
 \tag{M37.7}
\]
so every term is zero.  Any $y\notin G_0$ has the form
$hx$ with $h\in H$ and $x\in G_0$ by (M37.1).
That $h$ cannot belong to $H_0$, so (M37.7) proves
(M37.3).  If a term in (M37.4) is nonzero, both
$hx$ and $hy$ belong to $G_0$; hence $x^{-1}y\in G_0$.
This proves (M37.5).  The series (M37.4) is well-defined:
each factor is at most one by (M37.2), so its nonnegative
sum is at most $\sum_hb(hx)=1$.
Finally, for $x\in G_0$, either $y\notin G_0$ and
(M37.5) applies, or $y\in G_0\subset G_k$ and the
indicator difference is zero.  This proves (M37.6).
\end{proof}

The use of $hx$ instead of $h^{-1}x$ makes no difference
to the normalization, since inversion permutes each
subgroup.  The nonnegativity assumption is indispensable
to the termwise implication in (M37.7); it is part of
the source's definition of an indentation.

\paragraph{Exact defect and external orbital mass.}
Keep (M37.1) and the global normalization in (M37.2),
but do not impose the local normalization.  For $x\in G_0$
and a coset $c\in H/H_0$, define
\[
 A_c(x)=\sum_{h\in c}b(hx),\qquad
 \delta(x)=1-A_{H_0}(x).
 \tag{M37.8}
\]
Then $A_c\ge0$, $\sum_cA_c=1$, and $0\le\delta\le1$.
The total mass assigned to external translates is
\[
 \begin{split}
 W_{\mathrm{ext}}(x)
   &:=\sum_{\eta\in H\setminus H_0}B_b(x,\eta x)\\
   &=1-\sum_{c\in H/H_0}A_c(x)^2.
 \end{split}
 \tag{M37.9}
\]
In particular
\[
       2\delta(x)(1-\delta(x))
          \le W_{\mathrm{ext}}(x)
          \le2\delta(x)-\delta(x)^2\le2\delta(x).
 \tag{M37.10}
\]

\begin{proof}
Tonelli for nonnegative countable sums gives
\[
 \sum_{\eta\in H}B_b(x,\eta x)
    =\sum_{h\in H}\sum_{\eta\in H}b(hx)b(h\eta x)=1.
 \tag{M37.11}
\]
If $\eta$ ranges only over $H_0$, then $h\eta$ runs
over the same coset of $H_0$ as $h$.  Summing coset
by coset gives
\[
                   \sum_{\eta\in H_0}B_b(x,\eta x)
                              =\sum_c A_c(x)^2.
 \tag{M37.12}
\]
Subtracting proves (M37.9).  The identity coset has
mass $1-\delta$ and the other cosets have total mass
$\delta$, so
\[
       0\le\sum_{c\ne H_0}A_c^2\le\delta^2.
 \tag{M37.13}
\]
Substitution into (M37.9) gives both bounds in (M37.10).
\end{proof}

Here is the precise consequence for varying indentations.
Suppose $G$ carries a Haar measure, $b$ is measurable,
$D\subset G_0$ is a
measurable fundamental domain for $H$, and
$d\in L^1(D)$ is a prescribed quotient diagonal.
For any measurable selection of external translates,
its integral after multiplication by $B_b$ has absolute
value at most
\[
                    2\int_D\delta(x)|d(x)|\,d\mu_G(x).
 \tag{M37.14}
\]
Indeed the selected nonnegative weights are bounded by
$W_{\mathrm{ext}}$, and the quotient diagonal is the
same at $x$ and $\eta x$.  Equation (M37.10) proves
the bound; absolute convergence follows from
$W_{\mathrm{ext}}\le1$ and $d\in L^1$.
Thus if a sequence $b_j$ has defects $\delta_j\to0$
almost everywhere on $D$, every such external integral
tends to zero for the fixed diagonal, by dominated
convergence.  If the diagonals also vary, the same
conclusion follows either from
\[
       \|\delta_j\|_\infty\to0,\quad
                       \sup_j\|d_j\|_1<\infty,
 \tag{M37.15}
\]
or from $d_j\to d$ in $L^1$ and
$\delta_j\to0$ almost everywhere.  In the latter case
\[
 \int_D\delta_j|d_j|
       \le\|d_j-d\|_1+\int_D\delta_j|d|\longrightarrow0.
 \tag{M37.16}
\]
These are explicit sufficient conditions, not an
assertion that a varying source kernel meets them.

\paragraph{A trace-norm criterion for the varying diagonals.}
The second alternative has a useful operator formulation.
Fix a Hilbert realization $\mathcal H=L^2(D,\mu)$.
There is a canonical linear contraction
\[
 \mathcal D:\mathcal S_1(\mathcal H)\longrightarrow L^1(D,\mu),
 \qquad
 \operatorname{tr}(M_fT)=\int_D f\,\mathcal D(T)\,d\mu
                   \quad(f\in L^\infty(D,\mu)),
 \tag{M37.20}
\]
where $M_f$ denotes multiplication and $\mathcal S_1$ the
trace-class operators.  In particular
$\int_D\mathcal D(T)\,d\mu=\operatorname{tr}T$.
This is the trace diagonal; no pointwise restriction of an
arbitrary measurable two-variable kernel is being asserted.

\begin{proof}
Write a singular-value expansion
$T=\sum_n s_n\,u_n\otimes\overline{v_n}$, with
$\sum_n s_n=\|T\|_1$ and unit vectors $u_n,v_n$.
Set $\mathcal D(T)=\sum_n s_nu_n\overline{v_n}$.
Cauchy--Schwarz gives
$\|u_n\overline{v_n}\|_{L^1}\le1$, so the series converges
in $L^1$ with norm at most $\|T\|_1$.
For every bounded $f$, the rank-one trace formula and
trace-norm convergence give (M37.20).
Bounded measurable functions separate $L^1$ functions;
therefore (M37.20) makes the definition independent of
the expansion and proves linearity.  Applying the same
bound to $T_1-T_2$ proves contraction.
\end{proof}

Consequently, if $\|T_j-T\|_1\to0$ in this fixed realization,
$d_j=\mathcal D(T_j)$, and $\delta_j\to0$ almost everywhere,
every selected external orbital contribution in (M37.14)
tends to zero, with the explicit bound
\[
 2\|T_j-T\|_1+
       2\int_D\delta_j\,|\mathcal D(T)|\,d\mu
                         \longrightarrow0.
 \tag{M37.21}
\]
Indeed $\|d_j-\mathcal D(T)\|_{L^1}\le\|T_j-T\|_1$,
and the second term tends to zero by dominated convergence.
The assertion also applies to a fixed closed Hilbert subspace,
or to a closed subquotient, realized isometrically by $J$ in
this $L^2$ space: replace $T_j$ by $JT_jJ^*$, which preserves
trace norms and traces because it adds only zero singular
values.  If $w$ is measurable with $0<w<\infty$ almost
everywhere, a weighted $L^2(D,w\mu)$ realization can first
be conjugated by the unitary $Uu=\sqrt w\,u$ into $L^2(D,\mu)$.
The trace identity gives
$\mathcal D_\mu(UTU^{-1})=w\mathcal D_{w\mu}(T)$,
so (M37.21) uses the density for the fixed reference measure
after that conjugation.
This formulation requires that the proposed geometric diagonal
has actually been identified with the fixed trace diagonal.
Trace class alone does not identify it with a chosen version
of a measurable kernel on a null diagonal.

\paragraph{The source's compatible-indentation construction.}
Original Proposition 12.2.2.4, p.~512, assumes a second
countable locally compact abelian $G$, a discrete $H$
with counting measure, an increasing exhaustive family
of open subgroups $G_k$, and $H_k=H\cap G_k$.
Among its stated hypotheses are a fundamental domain
$D\subset G_0$ and the analogous fundamental-domain
property for each $G_k/H_k$.  In particular
$G=HD\subset HG_0$, which gives (M37.1).
Its asserted indentation $b$ is global and has
$\alpha_kb$ an $H_k$-indentation on $G_k$ for every
$k$.  At $k=0$ this is precisely the second
normalization in (M37.2).  The first is the global
indentation condition (12.2.1.2), p.~505.

Therefore every nonnegative indentation with the
asserted simultaneous normalization satisfies
(M37.3)--(M37.6).  This conclusion uses only the
two normalizations and $G=HG_0$; none of the additional
differentiation or domain-regularity hypotheses is
invoked.  The source's own construction on p.~513
explicitly extends an $H_0$-indentation on $G_0$
by zero to $G$, consistent with (M37.3).
In particular its tail integrand (12.2.2.26), evaluated
at $x\in D\subset G_0$ and $y=h^{-1}x$ outside $G_k$,
is already zero for each $k$.  The proposition's
zero-tail limit is not refuted by this observation.

\paragraph{The actual rational external-prime shells.}
Take
\[
 G=\mathbb I_{\mathbf Q},\quad H=\mathbf Q^\times,
 \quad G_0=\mathbf R^\times\times
                           \widehat{\mathbf Z}^\times,
 \quad H_0=\{1,-1\}.
 \tag{M37.17}
\]
Every idele can have its finite valuations canceled
by a rational number, so $G=HG_0$ directly.
Let $T_0=\varnothing$ and let $T_k$ exhaust all rational
primes, adding one prime at a time.  Define
\[
 G_k=\mathbf R^\times\times
         \prod_{p\in T_k}\mathbf Q_p^\times\times
         \prod_{p\notin T_k}\mathbf Z_p^\times.
 \tag{M37.18}
\]
These are open subgroups; adding one prime has quotient
$\mathbf Z$ by its valuation.  The standard section
$\tau(t)$ and every $g\in E_{\mathbf Q}^1=
\{1\}\times\widehat{\mathbf Z}^\times$ belong to $G_0$.
For \emph{any} single nonnegative $b$ satisfying the
global $\mathbf Q^\times$ normalization and the local
$\{1,-1\}$ normalization on $G_0$, (M37.6) gives
\[
 (\alpha_{k+1}-\alpha_k)(\xi g\tau(t))
                  B_b(\tau(t),\xi g\tau(t))=0
 \quad(\xi\in\mathbf Q^\times,\ g\in E_{\mathbf Q}^1,\ t>0).
 \tag{M37.19}
\]
This is pointwise, before integrating or taking any
kernel limit.  Multiplying by an additional radial
cutoff or by a pulled-back measurable kernel does
not change the vanishing wherever that kernel is
defined.  Thus no such single compatible indentation
can turn the source's external shell in (11.3.2.1)
into its nonzero local Poisson value, exhibited by
the explicit test (M36.7)--(M36.8).

The dependency distinction is precise.  The proof of
Theorem 11.3.2.1 on original p.~495 actually cites
Proposition 12.2.2.2 for its indented-kernel trace
interpretation, and on p.~496 claims an approximation
by the local orbital finite-rank kernels.  It does
\emph{not} there explicitly cite Proposition 12.2.2.4.
The latter is the later available construction of a
single indentation compatible with the whole local
filtration.  MT-M37 proves that substituting that
compatible construction for the invalid literal
section average of MT-M36 cannot supply the asserted
external-shell trace.  It does not invent a source
citation or refute the zero-tail statement itself.

\paragraph{What a varying-level repair must address.}
The family $b_t$ of MT-M36 drops the local normalization
at the initial subgroup and can redistribute mass to
the $p$ and $p^{-1}$ translates; its defect is exactly
$\delta=t$ on the standard domain, and it attains the
lower bound $2t(1-t)$ in (M37.10).  Thus there is no
obstruction here to all nonnegative global indentations.
The obstruction is to imposing the two normalizations
on the \emph{same} function.

A level-dependent family could instead normalize only
at its current subgroup and need not have zero defect
at the original $G_0$.  Such a replacement must specify
the resulting orbital weights and justify its kernel
and trace limits; (M37.14)--(M37.16) give concrete
tests for any claim that the initial defect disappears.
No boundedness, $L^1$ convergence, or invariant orbital
allocation is inferred from the existence of the
individual indentations.  The Hecke-GRH endpoint
$M_H$ remains conditional.

\paragraph{Audit scope.}
Original pp.~512--513 were freshly rendered and read
in full for the hypotheses and extension-by-zero
construction.  Original pp.~495, 505 and 508 were
visually checked in MT-M36, and the complete supplied
proof text on pp.~495--496 was read as a dependency
map.  The new support and defect statements above
are fully proved here and require no GRH input.


\subsection*{MT-S506. Dirichlet support of possible poles for every $S_5$ extension}

Let $L/\mathbb Q$ be a Galois extension with group $G=S_5$,
let $E=L^{A_5}$, and let $\eta$ be its quadratic character.
No restriction on complex conjugation or on ramification at $5$
is imposed in this subsection.  Use the irreducible representations
\[
 \begin{gathered}
 \mathbf1,\quad\eta,\quad V,\quad V\eta,\quad W,\quad W\eta,\quad X,\\
 V=\mathbf C^5/\mathbf C\mathbf1,\quad
 W=\operatorname{Sym}^2V-\mathbf1-V,\quad X=\bigwedge^2V.
 \end{gathered}
 \tag{S506.1}
\]
The notation for $W$ first denotes a virtual representation;
the character calculation below proves that it is an actual
irreducible representation of dimension $5$.  These are
Calegari's $\rho_4,\rho_5,\rho_6$, respectively, and $X\eta=X$.
All $L$-functions here include the finite factors at ramified primes.

For $z$ in the open critical strip $0<\Re z<1$, put
\[
 (a,b,c,d,x,y,e)=\operatorname{ord}_{s=z}
 (\zeta_{\mathbb Q},L_\eta,L_V,L_{V\eta},L_W,L_{W\eta},L_X).
 \tag{S506.2}
\]
An order is positive at a zero and negative at a pole.
Unconditionally $L_X$ is an entire Hecke $L$-function of a
nontrivial finite-order character over a degree-six field.
Consequently $a,b,e\ge0$, and the pole orders satisfy
\[
 \begin{aligned}
 \max(0,-c)&\le\min(a,e),&
 \max(0,-d)&\le\min(b,e),\\
 \max(0,-x)&\le b,&
 \max(0,-y)&\le a.
 \end{aligned}
 \tag{S506.3}
\]
In particular every possible pole of an irreducible Artin factor
in this strip lies in
\[
          Z(\zeta_{\mathbb Q})\ \cup\ Z(L_\eta).
 \tag{S506.4}
\]
Under the designated input $A_D$, every such pole therefore lies
on $\Re z=1/2$.  This conclusion uses only rational Dirichlet GRH.
It neither asserts that these poles exist nor proves their absence,
and it gives no general zero-location theorem for the other factors.

\begin{proof}
The complete character table, in the order (S506.1), is
\[
\begin{array}{c|r|rrrrrrr}
\text{cycle type}&\text{class size}&\mathbf1&\eta&V&V\eta&W&W\eta&X\\\hline
1^5&1&1&1&4&4&5&5&6\\
2\,1^3&10&1&-1&2&-2&1&-1&0\\
2^2\,1&15&1&1&0&0&1&1&-2\\
3\,1^2&20&1&1&1&1&-1&-1&0\\
3\,2&20&1&-1&-1&1&1&-1&0\\
4\,1&30&1&-1&0&0&-1&1&0\\
5&24&1&1&-1&-1&0&0&1
\end{array}
 \tag{S506.5}
\]
Indeed $\chi_V(g)=\#\operatorname{Fix}(g)-1$ and
\[
 \chi_W(g)=\frac{\chi_V(g)^2+\chi_V(g^2)}2-1-\chi_V(g),\qquad
 \chi_X(g)=\frac{\chi_V(g)^2-\chi_V(g^2)}2.
\]
These formulas give every displayed entry.  Weighted inner products
of the seven columns are $\delta_{ij}$: multiplication by the
class-size diagonal matrix gives $120I$ for the Gram matrix.
A virtual character of norm one is the positive or negative of
an irreducible character; its positive value at the identity fixes
the positive sign.  Thus all seven columns are actual irreducibles.
Their squared dimensions sum to $120$, so the list is complete.
This also proves $X\eta=X$ directly, since $X$ vanishes on odd
permutations.

Let $J=S_4$ be a point stabilizer and let
$F_{20}=\operatorname{AGL}_1(\mathbb F_5)\subset S_5$ be the
normalizer of a $5$-cycle.  Write an affine element as
$t\mapsto ut+v$, and let $\tau(u)$ be a primitive quartic character
of $\mathbb F_5^\times$, inflated to $F_{20}$.  The sign restricted
to $F_{20}$ is $\tau^2$.  Let
$T=S_2\times S_3$ be a two-subset stabilizer and let
$\alpha$ be sign on its $S_2$ factor and trivial on its $S_3$ factor.
Frobenius reciprocity and (S506.5) give
\[
\begin{array}{c|l}
\text{induced character}&\text{decomposition in }G\\\hline
\operatorname{Ind}_{J}^{G}\mathbf1&\mathbf1+V\\
\operatorname{Ind}_{J}^{G}(\eta|_J)&\eta+V\eta\\
\operatorname{Ind}_{F_{20}}^{G}\mathbf1&\mathbf1+W\eta\\
\operatorname{Ind}_{F_{20}}^{G}(\eta|_{F_{20}})&\eta+W\\
\operatorname{Ind}_{F_{20}}^{G}\tau&X\\
\operatorname{Ind}_{T}^{G}\alpha&V+X\\
\operatorname{Ind}_{T}^{G}(\alpha\eta|_T)&V\eta+X.
\end{array}
 \tag{S506.6}
\]
Here is an explicit check of the potentially sensitive quartic
identity.  The affine group has the identity, four nonidentity
translations of cycle type $5$, five elements of type $2^2\,1$,
and ten of type $4\,1$.  The first two sets have multiplier $1$,
the third has multiplier $-1$, and the last two multiplier classes
have $\tau$-values $i$ and $-i$.  For $X$ the weighted inner-product
numerator is $6+4+10=20$.  For every other column the numerator
is zero, because the two quartic multiplier classes cancel and
the remaining two multiplier contributions cancel.  Division by
$|F_{20}|=20$ proves the identity.  The other rows follow by the
same finite averaging formula
\[
 \bigl\langle\operatorname{Ind}_{H}^{G}\lambda,R\bigr\rangle_G
       =\frac1{|H|}\sum_{h\in H}\overline{\lambda(h)}\chi_R(h).
 \tag{S506.7}
\]
For the two $T$ rows, $\alpha$ and $\alpha\eta$ are respectively
sign on the $S_2$ and on the $S_3$ factor, which fixes the twist
convention.

For a linear character $\lambda$ of any subgroup $H\subset G$,
Artin induction and class field theory identify
$L(s,\operatorname{Ind}_H^G\lambda)$ with its Hecke $L$-function
over $L^H$.  It is holomorphic in the open critical strip,
including when $\lambda$ is trivial, in which case it is a
Dedekind zeta function.  For $\tau\ne\mathbf1$ it is entire.
Taking orders in (S506.6) yields
\[
 \begin{gathered}
 a+c\ge0,\quad b+d\ge0,\quad a+y\ge0,\quad b+x\ge0,\\
 e\ge0,\quad c+e\ge0,\quad d+e\ge0.
 \end{gathered}
\]
The orders exist as integers because Brauer induction gives
meromorphic continuation of every Artin factor.  Since $a,b,e$
are nonnegative, these inequalities prove (S506.3)--(S506.4).
The last assertion is precisely the Dirichlet zero-location
input $A_D$ applied to the trivial and the quadratic characters.
\end{proof}

The classical all-prime induction and multiplication identities
are Prasad--Yogananda, \emph{A Report on Artin's holomorphy
conjecture}, Section~2.1, Properties~1--2, PDF p.~3;
Hecke entireness is Section~3.1, pp.~4--5; meromorphicity is
Theorem~3.3, p.~5.  The character labels agree with Calegari,
\emph{The Artin conjecture for some $S_5$-extensions},
author-hosted publisher PDF p.~4, Section~2 character table.
The character calculations here are proved explicitly; the
classical analytic inputs are cited, not re-proved.

\subsection*{MT-S507. The complete monomial cone and the remaining five-dimensional gate}

Retain (S506.1)--(S506.2), and assume additionally that
$L_V$ and $L_{V\eta}$ are holomorphic at $z$, that is, $c,d\ge0$.
This is only a pointwise hypothesis for the following cone theorem.
It holds at every point for the Calegari family of MT-S501:
complex conjugation has cycle type $(12)(34)$, $5$ is unramified,
and $\operatorname{Frob}_5$ has the same cycle type.
Calegari's Theorem~1.1, author-hosted publisher PDF p.~2, supplies
this four-dimensional holomorphy unconditionally.  No GRH input
is needed for the cone theorem or its numerical bounds.

At $z$, all tests
\[
 \operatorname{ord}_z L\bigl(s,\operatorname{Ind}_H^G\lambda\bigr)
             \ge0\qquad(H\subset G,\ \dim\lambda=1)
 \tag{S507.1}
\]
are equivalent, after imposing $c,d\ge0$, to the following twelve
linear inequalities:
\[
 \begin{gathered}
      a,b,c,d,e\ge0,\\
 g_1=x+y\ge0,\quad g_2=b+x\ge0,\quad g_3=a+y\ge0,\\
 g_4=c+x+e\ge0,\quad g_5=d+y+e\ge0,\\
 g_6=a+c+x\ge0,\quad g_7=b+d+y\ge0.
 \end{gathered}
 \tag{S507.2}
\]
The assertion is an exact description of this cone, not a claim
that all of its integer points occur as arithmetic order vectors.
In particular $c,d\ge0$ are extra holomorphy inputs, not deductions
from the monomial tests alone.

Define the relative and total Dedekind orders
\[
 q=\operatorname{ord}_z\frac{\zeta_L}{\zeta_E}
       =4(c+d)+5(x+y)+6e,
 \qquad
 r=\operatorname{ord}_z\zeta_L=a+b+q.
 \tag{S507.3}
\]
If $L_W$ has a pole of order $v>0$ at $z$, then
\[
 \begin{gathered}
 x=-v,\quad y\ge v,\quad b\ge v,\quad c+e\ge v,
       \quad a+c\ge v,\\
 v\le\min\{b,c+e,a+c\},\qquad q\ge4v,\qquad r\ge5v.
 \end{gathered}
 \tag{S507.4}
\]
For a pole of $L_{W\eta}$, interchange
$(a,b),(c,d),(x,y)$.  Both five-dimensional factors cannot have
a pole at the same point.  In particular $q\le3$ implies that
all irreducible factors are holomorphic at $z$, and every possible
pole order is at most $\lfloor q/4\rfloor$.
In the branch $c=0$ of (S507.4),
\[
        a,b,e,y\ge v,\qquad q\ge6v,\qquad r\ge8v.
 \tag{S507.5}
\]
The analogous assertion with $d=0$ holds for $W\eta$.
All these numerical constants are sharp for the exact cone
(S507.2), including the extra branch condition in (S507.5).

Let $Z(f)$ denote its zero set within the open critical strip.
For this holomorphy family the corresponding precise necessary
simultaneous-zero conditions are
\[
\begin{aligned}
 \operatorname{Pole}(L_W)&\subseteq
 Z(L_\eta)\cap\bigl(Z(L_V)\cup
             (Z(\zeta_{\mathbb Q})\cap Z(L_X))\bigr),\\
 \operatorname{Pole}(L_{W\eta})&\subseteq
 Z(\zeta_{\mathbb Q})\cap\bigl(Z(L_{V\eta})\cup
             (Z(L_\eta)\cap Z(L_X))\bigr).
\end{aligned}
 \tag{S507.6}
\]
Under $A_D$ these possible poles lie on the critical line by
MT-S506.  Under $M_H$ and the Calegari local hypotheses they
must additionally be nonreal, by MT-S505.  Neither hypothesis
by itself supplies the missing noncoincidence statement in
(S507.6), and no simplicity of any zero is asserted here.

\begin{proof}[Proof of the exact cone]
Besides the rows already established in (S506.6), use the
subgroup $A_4$ fixing the fifth letter and a nontrivial
character $\nu$ of its quotient $A_4/V_4=C_3$; the two-subset
stabilizer $T$; and
\[
  D_8=\langle r_0,s_0\rangle\subset S_4,\qquad
  r_0=(1234),\quad s_0=(13),\qquad
  \lambda(r_0)=-1,\quad\lambda(s_0)=1.
\]
These values respect $r_0^4=s_0^2=1$ and
$s_0r_0s_0=r_0^{-1}$, so they define a linear character.
Formula (S506.7) and the table give
\[
 \begin{gathered}
 \operatorname{Ind}_{A_4}^{G}\nu=W+W\eta,\qquad
 \operatorname{Ind}_{D_8}^{G}\lambda=V+W+X,\\
 \operatorname{Ind}_{T}^{G}\mathbf1=\mathbf1+V+W.
 \end{gathered}
 \tag{S507.7}
\]
Twisting by $\eta$ gives the other two needed rows, since
$X\eta=X$.  For the $A_4$ identity, its identity and three
involutions contribute $5+3$ to each rank-five inner product;
the eight $3$-cycles split into two sets of four with character
values $\omega,\omega^2$, so their weighted contribution is
$-4(\omega+\omega^2)=4$.  Division by $12$ gives one.
The other inner products vanish.  For the $D_8$ identity,
its elements have cycle types $1$, two $4$-cycles, three double
transpositions, and two transpositions.  Their $\lambda$-weights
sum, within those types, to $1,-2,-1,2$, respectively.
Against the seven columns this yields
$(0,0,1,0,1,0,1)$.  For $T$, the permutation action on the ten
two-subsets has character $1+\chi_V+\chi_W$ by direct counting
of invariant two-subsets.  This proves every row of (S507.7).
Thus all seven $g_i$ are monomial tests.  The tests with
$H=G$ give $a,b\ge0$, and the quartic row gives $e\ge0$.
Together with the stated $c,d\ge0$ premise this proves the
necessity of (S507.2).

For completeness, the next two tables give every distinct induced
linear-character multiplicity vector, together with a nonnegative
integer decomposition of its order functional into (S507.2).
The columns labeled $a,\ldots,e$ are the coefficients of that
functional, not values of the order variables.  Thus each last
entry can be verified by substituting the definitions of $g_i$.
\[
\begin{array}{r|rrrrrrr|l}
 i&a&b&c&d&x&y&e&\text{nonnegative decomposition}\\\hline
1 & 0 & 0 & 0 & 0 & 0 & 0 & 1 & e \\
2 & 0 & 0 & 0 & 0 & 0 & 0 & 2 & 2e \\
3 & 0 & 0 & 0 & 0 & 1 & 1 & 0 & g_1 \\
4 & 0 & 0 & 0 & 1 & 0 & 0 & 1 & d+e \\
5 & 0 & 0 & 0 & 1 & 0 & 1 & 1 & g_5 \\
6 & 0 & 0 & 0 & 1 & 1 & 1 & 1 & g_1+d+e \\
7 & 0 & 0 & 1 & 0 & 0 & 0 & 1 & c+e \\
8 & 0 & 0 & 1 & 0 & 1 & 0 & 1 & g_4 \\
9 & 0 & 0 & 1 & 0 & 1 & 1 & 1 & g_1+c+e \\
10 & 0 & 0 & 1 & 1 & 0 & 0 & 2 & c+d+2e \\
11 & 0 & 0 & 1 & 1 & 1 & 1 & 1 & g_1+c+d+e \\
12 & 0 & 0 & 1 & 1 & 1 & 1 & 2 & g_4+g_5 \\
13 & 0 & 0 & 1 & 1 & 2 & 2 & 2 & g_1+g_4+g_5 \\
14 & 0 & 0 & 2 & 2 & 2 & 2 & 4 & 2g_4+2g_5 \\
15 & 0 & 1 & 0 & 0 & 0 & 0 & 0 & b \\
16 & 0 & 1 & 0 & 0 & 1 & 0 & 0 & g_2 \\
17 & 0 & 1 & 0 & 1 & 0 & 0 & 0 & b+d \\
18 & 0 & 1 & 0 & 1 & 0 & 1 & 0 & g_7 \\
19 & 0 & 1 & 0 & 1 & 1 & 1 & 0 & g_1+b+d \\
20 & 0 & 1 & 0 & 2 & 0 & 1 & 1 & g_7+d+e \\
21 & 0 & 1 & 0 & 2 & 1 & 2 & 1 & g_1+g_7+d+e \\
22 & 0 & 1 & 1 & 1 & 0 & 1 & 1 & g_7+c+e \\
\end{array}
\]

\[
\begin{array}{r|rrrrrrr|l}
 i&a&b&c&d&x&y&e&\text{nonnegative decomposition}\\\hline
23 & 0 & 1 & 1 & 1 & 2 & 1 & 1 & g_1+g_4+b+d \\
24 & 0 & 1 & 1 & 3 & 2 & 3 & 3 & g_1+g_4+g_5+g_7+d+e \\
25 & 1 & 0 & 0 & 0 & 0 & 0 & 0 & a \\
26 & 1 & 0 & 0 & 0 & 0 & 1 & 0 & g_3 \\
27 & 1 & 0 & 1 & 0 & 0 & 0 & 0 & a+c \\
28 & 1 & 0 & 1 & 0 & 1 & 0 & 0 & g_6 \\
29 & 1 & 0 & 1 & 0 & 1 & 1 & 0 & g_1+a+c \\
30 & 1 & 0 & 1 & 1 & 1 & 0 & 1 & g_6+d+e \\
31 & 1 & 0 & 1 & 1 & 1 & 2 & 1 & g_1+g_5+a+c \\
32 & 1 & 0 & 2 & 0 & 1 & 0 & 1 & g_6+c+e \\
33 & 1 & 0 & 2 & 0 & 2 & 1 & 1 & g_1+g_6+c+e \\
34 & 1 & 0 & 3 & 1 & 3 & 2 & 3 & g_1+g_4+g_5+g_6+c+e \\
35 & 1 & 1 & 0 & 0 & 0 & 0 & 0 & a+b \\
36 & 1 & 1 & 0 & 0 & 1 & 1 & 0 & g_2+g_3 \\
37 & 1 & 1 & 0 & 0 & 1 & 1 & 2 & g_2+g_3+2e \\
38 & 1 & 1 & 1 & 1 & 0 & 0 & 0 & a+b+c+d \\
39 & 1 & 1 & 1 & 1 & 1 & 1 & 0 & g_6+g_7 \\
40 & 1 & 1 & 1 & 1 & 2 & 2 & 0 & g_1+g_6+g_7 \\
41 & 1 & 1 & 2 & 2 & 1 & 1 & 2 & g_6+g_7+c+d+2e \\
42 & 1 & 1 & 2 & 2 & 3 & 3 & 2 & g_1+g_4+g_5+g_6+g_7 \\
43 & 1 & 1 & 4 & 4 & 5 & 5 & 6 & g_1+3g_4+3g_5+g_6+g_7 \\
\end{array}
\]

Here is an exhaustive finite algorithm for these tables, with no
numerical optimization premise.  List all $120$ permutations of
five letters and their exact composition table.  Begin with the
trivial subgroup; repeatedly adjoin every individual permutation
to every subgroup found, compute the resulting generated subgroup,
and discard duplicates.  This finds every subgroup: given a
subgroup $H$, successively adjoin elements of $H$ until its
finite generating list has been exhausted.  Each such step appears
in this search.  It yields $156$ subgroups and $19$ conjugacy
classes; conjugate subgroups give the same induced patterns.
For each representative, retain a generating list.

The exponent of $S_5$ divides $60$, so every linear character of
one of its subgroups has values among the $60$th roots of unity.
Assign to each generator all phases $j\in\mathbb Z/60$ allowed
by its order.  Propagate those phases through its Cayley graph,
reject inconsistent assignments, and check
$j(hh')=j(h)+j(h')$ for all pairs.  This both finds every linear
character and verifies that each retained assignment is one.
Compute (S506.7) in the exact ring
$\mathbb Q[Z]/(\Phi_{60}(Z))$, using
$\overline{\lambda(h)}=Z^{-j(h)}$; reduce every sum modulo
$\Phi_{60}$ and divide by $|H|$.  The resulting multiplicities
are the nonnegative integers listed above.  As checks, the number
of characters found for each $H$ is $|H/[H,H]|$, and the
weighted dimension of every row is $120/|H|$.
Discarding repeated vectors gives exactly the $43$ displayed rows.
All steps use finite permutations, integer arithmetic and exact
cyclotomic reduction.

The full executable enumeration and its output are
\texttt{work/checks/s5\_hecke\_cone.py} and
\texttt{s5\_hecke\_cone.json}; their exploratory floating-point
optimizer is not used in this proof.  The separate exact verifier
\texttt{work/checks/s5\_cone\_decomposition\_verification.py}
and its JSON output reconstruct all $43$ rows as nonnegative
integer combinations of the twelve generators and verify the
sharpness witnesses below.  Its input SHA-256 is
\begin{center}\small\ttfamily
0b1ab9c1b47311477e3e113cced1e568a4c6fc417c8077875e87e7d8637ea917
\end{center}
The exhaustive permutation/character portion was re-executed
without its optimizer, and the displayed decompositions were
checked in integer arithmetic.  This is a transparent finite
computation with its algorithm and complete output supplied;
it is not a proof-assistant certification.

Every order vector satisfying (S507.2) pairs nonnegatively
with every displayed row.  Exhaustiveness therefore proves
sufficiency for (S507.1).  This establishes the exact cone.
\end{proof}

\begin{proof}[Proof of the pole bounds, sharpness and support condition]
The regular character of $G$ has multiplicities equal to
irreducible dimensions, whereas the character for the quadratic
field $E$ is $\mathbf1+\eta$.  All-prime Artin factorization
therefore gives (S507.3), including at ramified primes.
If $x=-v<0$, the inequalities $g_1,g_2,g_4,g_6\ge0$
give the first line of (S507.4).  Since $d,e,x+y\ge0$,
\[
 q=4(c+e)+4d+5(x+y)+2e\ge4(c+e)\ge4v,
 \qquad r=a+b+q\ge5v.
\]
If $c=0$, $g_4,g_6\ge0$ force $e,a\ge v$, so
$q\ge6e\ge6v$ and $r\ge a+b+q\ge8v$.
Twisting interchanges $a,b$, $c,d$ and $x,y$, giving every
symmetric assertion.  Since $g_1\ge0$, $x,y$ cannot both be
negative.  The other five irreducible factors are already
holomorphic at $z$ under the stated hypotheses.  A remaining
pole has integral order $v\ge1$ and $q\ge4v$; this proves both
the $q\le3$ holomorphy criterion and the floor bound.

The integer vector
\[
        (a,b,c,d,x,y,e)=(0,1,1,0,-1,1,0)
 \tag{S507.8}
\]
satisfies (S507.2): the seven $g_i$ equal
$(0,0,1,0,1,0,2)$ and the five required coordinates are
nonnegative.  It has a pole order $v=1$, $q=4$ and $r=5$.
Its sign twist provides the symmetric witness.  For the
$c=0$ branch, the vector
\[
        (a,b,c,d,x,y,e)=(1,1,0,0,-1,1,1)
 \tag{S507.9}
\]
has $g_i=(0,0,2,0,2,0,2)$, $q=6$ and $r=8$.
Multiplying these vectors by any positive integer gives every
pole order $v$.  By the exact cone, all $43$ monomial tests are
nonnegative at these witnesses.  This proves sharpness within
the specified positivity information.  It makes no claim that
these integer vectors arise from number fields or analytic
$L$-functions.

Finally $x<0$ forces $b>0$.  If $c>0$ the first branch of
(S507.6) holds.  If $c=0$, then $a,e>0$, which is exactly
its other branch.  Since $c\ge0$, these cases exhaust all
possibilities.  The sign twist proves the second inclusion.
Under $A_D$, (S506.4) confines the poles to the critical line.
Under $M_H$ and the Calegari family hypotheses, MT-S505 excludes
every real point of the open critical strip, including the
centre.  This leaves only the nonreal points described above.
\end{proof}

\paragraph{Input separation and the geometric gate.}
MT-S506 is a new reduction using classical character theory,
Artin induction and Hecke analyticity; only its last pole-location
assertion invokes $A_D$.  MT-S507 additionally uses the
four-dimensional holomorphy premise, supplied in the indicated
family by Calegari's theorem.  Its stronger exclusion of real
poles uses the separately proved, $M_H$-conditional MT-S505.
The finite cone does not establish noncoincidence, pole absence,
automorphy of the rank-five factors, or a new proof of GRH.

On the arithmetic base with the ramified primes removed, these
representations are the Artin motives of finite covers and their
direct summands.  Induction corresponds to finite pushforward
from a cover; (S507.1) records positivity of the zero orders
for all one-dimensional induced summands.  The witness (S507.8)
shows precisely why that information, even supplemented by the
rank-four holomorphy input, does not force the order divisor of
each irreducible summand to be effective.  The additional
all-prime induction identities are compatible with inertia
invariants at the removed primes.  Thus (S507.6), with the order
penalties (S507.4)--(S507.5), is a smaller explicit interface for
the Artin-positivity gate, rather than an identification of
zero-independent positive cohomology.  The attached run is
\texttt{GRH\_NIGHT/complexity}; the inherited interfaces are
MT-ART07 and MT-S501/MT-S504/MT-S505, while the cone certificate
and the deductions in MT-S506/MT-S507 are new in this run.

\subsection*{MT-S508. One quadratic Hecke input and the exact automorphic central-order gate}

Retain the Calegari family of MT-S501: $L/\mathbb Q$ has group
$S_5$, complex conjugation has type $(12)(34)$, and the extension
is unramified at $5$ with Frobenius of the same type.  Let
$E=L^{A_5}$, let $F=L^{F_{20}}$ have degree $6$, and put $H=EF$.
Write $\eta$ for the quadratic character of $E/\mathbb Q$ and
$\eta_F$ for its restriction to $\operatorname{Gal}(L/F)$.
Then $H/F$ is quadratic, and its nontrivial Hecke character is
$\eta_F$.  Define the \emph{finite} $L$-functions
\[
 C(s)=L_{\mathbb Q}(s,\eta),\qquad
 B(s)=L_F(s,\eta_F)=\frac{\zeta_H(s)}{\zeta_F(s)}.
 \tag{S508.1}
\]
Both $B$ and $C$ are entire.  The all-prime induction identity
of MT-S506 gives
\[
             B(s)=C(s)L_W(s),\qquad
             \zeta_F(s)=\zeta_{\mathbb Q}(s)L_{W\eta}(s).
 \tag{S508.2}
\]
Let $\varpi$ be the particular cuspidal representation of
$\mathrm{GL}_5/\mathbb Q$ supplied by Calegari's construction,
and put $A(s)=L_f(s,\varpi)$.  This is its standard finite
$L$-function, which is entire.  Fix this comparison representation;
the statements below concern its exact transfer, not an arbitrary
choice of an automorphic representation of the same degree.
Here matching at every place means agreement of the standard
local $L$-factors, including the archimedean factor, in the
precise sense of the cited comparison theorem.

The only additional zero-location input needed here is
\[
  (\mathcal R_B)\qquad
  B(\sigma)\ne0\quad
  \bigl(0<\sigma<1,\ \sigma\ne\tfrac12,\ \sigma\in\mathbb R\bigr).
 \tag{S508.3}
\]
It concerns one quadratic Hecke character over one degree-six
field and only its real zeros.  GRH for this single Hecke
$L$-function implies $(\mathcal R_B)$; universal $M_H$ is a
sufficient input but is not required.  No Dirichlet GRH assumption
is needed for the central-order and automorphy assertions below.

Put
\[
 \beta=\operatorname{ord}_{1/2} B,\qquad
 b=\operatorname{ord}_{1/2} C,\qquad
 n=\operatorname{ord}_{1/2} A.
 \tag{S508.4}
\]
These are nonnegative integers.  Under $(\mathcal R_B)$ one has
\[
             \beta\ge b+n,
       \qquad \operatorname{ord}_{1/2}L_W=\beta-b\ge n\ge0.
 \tag{S508.5}
\]
Moreover the following are equivalent:
\[
 \begin{gathered}
       \beta\le b+n,\\
       \beta=b+n,\\
       \text{Calegari's exceptional set }\Omega\text{ is finite},\\
       \text{the specified }\varpi\text{ matches }W
             \text{ at every place}.
 \end{gathered}
 \tag{S508.6}
\]
Whenever these conditions hold, $\varpi\otimes\eta$ also matches
$W\eta$ at every place.  Thus every irreducible Artin factor of
$L/\mathbb Q$ is automorphic and every nontrivial one is entire.
If the comparison has not closed, its remaining gate is the
strict positive discrepancy
\[
                        \beta-b-n\ge1.
 \tag{S508.7}
\]
The older sufficient inequality $\beta\le b$ is a special case
of (S508.6).  The new condition retains possible central zeros
of the known automorphic numerator.

\paragraph{The exact classical comparison and its normalization.}
The classical construction and local comparison are Calegari,
\emph{The Artin conjecture for some $S_5$-extensions},
author-hosted publisher PDF: Theorem~1.1 p.~2, construction of
$\varpi$ p.~10, Theorem~4.7 p.~12, Lemma~4.8 pp.~13--14,
Landau Lemma~4.9 p.~14, Lemma~4.10 pp.~14--15, and the final
Theorem~1.2 proof p.~15.  This is the publisher numbering;
the last exact-comparison lemma is numbered 4.9 in the
arXiv v1 locators used in MT-S504.

We use the actual archimedean-completed quotient
\[
 R(s)=\frac{G_{\varpi}(s)A(s)}{G_W(s)L_W(s)},
 \tag{S508.8}
\]
where $G$ denotes the archimedean factor.  Including the usual
positive conductor powers would only multiply this by a
nowhere-zero entire function and would not change any order
or real-pole argument.  With the convention that a finite local
Euler factor is $P_p(p^{-s})^{-1}$, its correct factorization is
\[
 \begin{aligned}
 R(s)&=R_{\mathcal B}(s)
      \prod_{p\in\Omega}
             \left(\frac{1+p^{-s}}{1-p^{-s}}\right)^2,\\
 R_{\mathcal B}(s)&=
       \frac{G_{\varpi}(s)}{G_W(s)}
       \prod_{p\in\mathcal B}
             \frac{P_{W,p}(p^{-s})}{P_{\varpi,p}(p^{-s})}.
 \end{aligned}
 \tag{S508.9}
\]
Here $\mathcal B$ is the finite set of exceptional bad primes,
chosen disjoint from $\Omega$.  The source's displayed
$\mu_{\mathcal B}$ on publisher p.~13 reverses the two finite
polynomials relative to its subsequent reciprocal-Euler-factor
convention.  Equation (S508.9) fixes the direction directly from
the definition (S508.8); the source's holomorphic, nonzero
real-interval property is unaffected by this finite inversion.
In particular at $p\in\Omega$ its explicit polynomials are
\[
 P_{\varpi,p}(T)=(1+T^2)(1-T)^3,\qquad
 P_{W,p}(T)=(1+T^2)(1-T)(1+T)^2,
\]
which give exactly the displayed positive Euler quotient.

The cited comparison supplies two precise implications:
if $\Omega$ is infinite then $R$ has a pole at some real
$\sigma\in(0,1)$; if $\Omega$ is finite then the local factors
agree at all places, including the bad and archimedean ones.
The first uses its nonnegative Euler-series coefficients and
Landau's theorem; the second uses the cited finite-place
comparison with functional equations and twists.  The source's
final proof applies the same comparison to the quadratic twist.
We retain these classical results as inputs, not a new proof
of their automorphic construction.

The ratio $G_{\varpi}/G_W$ is holomorphic and nonzero on $(0,1)$.
For clarity, the Artin factor can be checked directly here.
Complex conjugation has trace $1$ on $W$, so its $+1,-1$
eigenspaces have dimensions $3,2$ and
\[
 G_W(s)=\Gamma_{\mathbb R}(s)^3\Gamma_{\mathbb R}(s+1)^2
       =\Gamma_{\mathbb R}(s)\Gamma_{\mathbb C}(s)^2.
 \tag{S508.10}
\]
Also $\eta$ is even, so $G_C(s)=\Gamma_{\mathbb R}(s)$.
The field $F$ has signature $(2,2)$, since the degree-six
permutation character $\mathbf1+W\eta$ takes value $2$ on
complex conjugation.  The real places split in $H/F$, because
$E$ is real.  Thus
$G_B(s)=\Gamma_{\mathbb R}(s)^2\Gamma_{\mathbb C}(s)^2
=G_C(s)G_W(s)$, consistently with (S508.2).
Here $\Gamma_{\mathbb R}(s)=\pi^{-s/2}\Gamma(s/2)$ and
$\Gamma_{\mathbb C}(s)=2(2\pi)^{-s}\Gamma(s)$.
All these factors are finite and nonzero at every real
$0<s<1$.  The source verifies that $G_{\varpi}$ likewise
uses standard real/complex gamma factors there.

\begin{proof}[Proof of the one-Hecke theorem]
First, the subgroup $F_{20}$ contains odd permutations, so
its intersection with $A_5$ has index two.  It follows that
$H=EF$ is the quadratic extension of $F$ cut out by the
nontrivial restriction $\eta_F$.  Artin induction, including
ramified factors, gives (S508.1)--(S508.2).

Substituting (S508.2) into the actual quotient yields
\[
          R(s)=U(s)\,\frac{A(s)C(s)}{B(s)},\qquad
          U(s)=\frac{G_{\varpi}(s)}{G_W(s)}.
 \tag{S508.11}
\]
For any real $0<\sigma<1$, the function $U$ is holomorphic
and nonzero.  Therefore
\[
 \operatorname{ord}_{\sigma}R
   =\operatorname{ord}_{\sigma}A
     +\operatorname{ord}_{\sigma}C
     -\operatorname{ord}_{\sigma}B.
 \tag{S508.12}
\]
This formula explicitly accounts for zeros of the automorphic
numerator; discarding them would give only the older sufficient
condition.

Under $(\mathcal R_B)$, the denominator in (S508.11) is nonzero
at every real point of $(0,1)$ other than the centre.
The other factors are holomorphic there, even if $C$ has a
zero.  Thus the only possible real pole of $R$ is at $1/2$,
where its order is $n+b-\beta$.
If $\Omega$ is infinite, the classical dichotomy forces a pole
there, hence $\beta>b+n$.  If $\Omega$ is finite, the exact
comparison gives $R\equiv1$ in the archimedean normalization
of (S508.8), hence $\beta=b+n$.  These two cases exhaust the
possibilities and prove (S508.5)--(S508.7).

The reverse implication in the last equivalence is also literal:
if all local factors of the specified comparison agree, then
no prime can have the two distinct polynomials displayed above,
so its exceptional set is empty.  Conversely, finite $\Omega$
forces exact local equality by the cited comparison.
For the quadratic twist, equality outside finitely many places
persists after twisting; the same exact-comparison argument,
as explicitly applied on source p.~15, identifies
$\varpi\otimes\eta$ with $W\eta$ at every place.  No second
Hecke zero-location hypothesis is needed for this automorphy
conclusion.  The dimensions $4$ and $6$ are already automorphic
in the specified family by Theorem~1.1, and the two linear
characters are classical.  The claimed holomorphy follows.
\end{proof}

\paragraph{The positive Euler quotient controls a whole half-plane.}
The same single real-zero hypothesis $(\mathcal R_B)$ implies
that $R$ is holomorphic and nonzero on $\Re s>1/2$.
In that half-plane it gives the exact divisor relation
\[
        \operatorname{ord}_z B
             =\operatorname{ord}_z A+\operatorname{ord}_z C.
 \tag{S508.13}
\]
In particular $L_W$ is holomorphic there and has the same
zero orders as $A$.  The poles of $L_W$ in the open critical
strip can only occur on its critical line, away from the real
axis.  This assertion requires neither $A_D$ nor full GRH for $B$.

\begin{proof}
Put
\[
 D_\Omega(s)=\prod_{p\in\Omega}
             \left(\frac{1+p^{-s}}{1-p^{-s}}\right)^2.
\]
This is the positive-coefficient Dirichlet series in the proof
of source Lemmas~4.8--4.9.  It is meromorphic on the plane
because $D_\Omega=R/R_{\mathcal B}$.  If $\Omega$ is infinite,
that proof identifies an abscissa of absolute convergence
$\gamma\in(0,1)$ and a pole at $\gamma$; the convergence at $1$
uses the stronger prime-sum bound proved in Theorem~4.7, not
Dirichlet density zero alone.  Since $R_{\mathcal B}$ is nonzero
and holomorphic at every real point of $(0,1)$, this is also a
pole of $R$.  Under $(\mathcal R_B)$, (S508.11) shows that this
pole can only be at $1/2$.  Hence $\gamma=1/2$.
If $\Omega$ is finite, the finite Euler product instead
converges for $\Re s>0$.

For every real $\sigma>1/2$ the nonnegative Dirichlet series
therefore converges, and its prime coefficients are $4$.
It follows that $\sum_{p\in\Omega}p^{-\sigma}<\infty$.
The logarithms of the Euler factors consequently converge
absolutely and locally uniformly on $\Re s>1/2$.
Each individual factor is nonzero in $\Re s>0$, so the
infinite product is holomorphic and nonzero on $\Re s>1/2$.
This also covers an empty or finite exceptional set.

The gamma quotient in $R_{\mathcal B}$ is holomorphic and
nonzero on all of $\Re s>0$: its standard gamma factors have
no poles or zeros there.  The finite polynomial factors are
also nonzero there, since their roots have modulus one and
$|p^{-s}|<1$.  Thus $R_{\mathcal B}$, and consequently $R$,
are holomorphic and nonzero on $\Re s>1/2$.
The same is true of $U$.  Taking orders in (S508.11) proves
(S508.13), and $L_W=UA/R$ proves the asserted holomorphy
and equality of zero orders.  The meromorphic Artin functional
equation for the selfdual $W$ reflects poles in
$0<\Re s<1/2$ to poles in $1/2<\Re s<1$;
its gamma and conductor factors are nonzero in this open strip.
Thus it excludes poles in the left half as well.  The centre
is already holomorphic by (S508.5).
\end{proof}

\paragraph{A single Hecke GRH also transfers to the known automorphic numerator.}
Assume now the full GRH for $B$, which implies
$(\mathcal R_B)$.  Then both $A=L_f(s,\varpi)$ and $C=L(s,\eta)$
satisfy GRH, whether or not the central equality (S508.6) holds.
Moreover the meromorphic divisor of $L_W=B/C$ in the open
critical strip is supported on the critical line.  This last
statement allows poles on that line; it does not prove Artin
holomorphy there.

\begin{proof}
For $1/2<\Re z<1$, GRH for $B$ says that
$\operatorname{ord}_z B=0$.  The two terms on the right of
(S508.13) are nonnegative, since $A$ and $C$ are holomorphic.
They are therefore both zero.  Their usual completed functional
equations then exclude zeros in $0<\Re z<1/2$ as well.
For the cuspidal unitary $\varpi$, the functional equation
pairs its standard $L$-function with that of its contragredient,
which is the complex conjugate under $s\mapsto\bar s$;
thus a zero is reflected across $\Re s=1/2$.
All gamma factors are nonzero in the open strip.  This proves
GRH for $A$ and $C$.  The factorization $L_W=B/C$ now proves
its critical-line divisor support.  If the central equality
also holds, the already proved exact transfer makes $L_W$
entire, so the divisor assertion is its ordinary GRH.
\end{proof}

This is specifically a \emph{one-$B$} transfer.  It is not a
claim that an $M_H$-conditional zero-location statement for
$\varpi$ was previously unavailable through its known base-change
product.  No GRH for $L_{W\eta},L_V,L_{V\eta}$ or $L_X$ follows
from the argument just given.  Once exact transfer is proved,
GRH for $\zeta_F$ separately gives GRH for $L_{W\eta}$ from
(S508.2), and universal $M_H$ is one sufficient way to supply
the other zero-location inputs.

\paragraph{The quadratic-twist expression for the same central gate.}
Put
\[
 A_\eta(s)=L_f(s,\varpi\otimes\eta),\qquad
 n_\eta=\operatorname{ord}_{1/2}A_\eta,
 \qquad f=\operatorname{ord}_{1/2}\zeta_F,
 \qquad k=\beta-b-n.
\]
Under $(\mathcal R_B)$ these nonnegative orders satisfy
\[
             k=n_\eta-f,\qquad 0\le k\le n_\eta.
 \tag{S508.14}
\]
Consequently $n_\eta\le f$ is equivalent to the same exact
transfer condition (S508.6).  In particular the additional
nonvanishing $A_\eta(1/2)\ne0$ suffices.  Neither that
nonvanishing nor this inequality is asserted unconditionally.

\begin{proof}
Every prime in $\Omega$ has Frobenius type $4\,1$, so
$\eta(p)=-1$.  Twisting its two displayed local polynomials
replaces $T$ by $-T$ and inverts their quotient.  Hence the
actual twisted completed ratio is
\[
 R_\eta(s)=\frac{G_{\varpi\eta}(s)A_\eta(s)}
                     {G_{W\eta}(s)L_{W\eta}(s)}
          =R_{\mathcal B,\eta}(s)D_\Omega(s)^{-1}.
 \tag{S508.15}
\]
The finite and archimedean factor $R_{\mathcal B,\eta}$ is
holomorphic and nonzero at every real $0<s<1$ before any
finiteness assertion about $\Omega$: $\eta$ restricts trivially
to $E$, so $(\varpi\otimes\eta)_E=\varpi_E$ has the same
known finite-image base change, and the local parameter argument
of source Lemma~4.8 still gives root-of-unity finite Euler
parameters after twisting.  In particular $\eta$ is even,
so it does not change the archimedean factors.  At the centre $R_{\mathcal B}$ and
$R_{\mathcal B,\eta}$ have order zero, and (S508.12) gives
$\operatorname{ord}_{1/2}D_\Omega=-k$.  Thus
$\operatorname{ord}_{1/2}R_\eta=k$.
On the other hand $\zeta_{\mathbb Q}(1/2)\ne0$, so (S508.2)
gives $\operatorname{ord}_{1/2}L_{W\eta}=f$.
Taking orders in the first quotient of (S508.15) gives
$k=n_\eta-f$.  Since $f\ge0$ and $k\ge0$ by (S508.5),
(S508.14) follows.  Equality $k=0$ is exactly (S508.6).
\end{proof}

\paragraph{Real-axis holomorphy needs only the one real-zero input.}
Under $(\mathcal R_B)$ every irreducible Artin factor in this
family is holomorphic at every real point of $(0,1)$.
For $W$, the half-plane argument and its functional equation
exclude every noncentral real pole, while (S508.5) handles
the centre.  For $W\eta$, (S508.2) and the elementary inequality
$\zeta_{\mathbb Q}(\sigma)<0$ for $0<\sigma<1$ prove
holomorphy on that real interval without any GRH input.
The other nontrivial factors are already entire by the
classical theorem.  Thus the earlier $A_D+(\mathcal R_B)$
corollary strengthens to $(\mathcal R_B)$ alone after the
positive Euler quotient is used.  This remains a conditional
statement about one Hecke function, not a verification of
$(\mathcal R_B)$.

\paragraph{The precise residual exceptional-prime law.}
If $k>0$, then $\Omega$ is infinite and, as $\epsilon\downarrow0$,
\[
 \sum_{p\in\Omega}p^{-1/2-\epsilon}
          =\frac{k}{4}\log\frac1\epsilon+O(1).
 \tag{S508.16}
\]
Thus an unclosed transfer forces an explicit positive logarithmic
singularity of the exceptional-prime harmonic sum, with the
same integer $k$ as the central excess.  This is an Abelian
statement at the Dirichlet-series boundary; no asymptotic
formula for the unweighted count of these primes is inferred.

\begin{proof}
By (S508.9), $D_\Omega$ has a pole of order $k$ at $1/2$.
It is positive for real $s>1/2$, so its leading Laurent
coefficient there is positive.  Consequently
\[
       \log D_\Omega(1/2+\epsilon)
             =k\log(1/\epsilon)+O(1).
\]
In the absolutely convergent half-plane the Euler logarithm is
\[
 \log D_\Omega(s)=4\sum_{p\in\Omega}p^{-s}+E(s),\qquad
 E(s)=4\sum_{p\in\Omega}\sum_{j\ge1}
              \frac{p^{-(2j+1)s}}{2j+1}.
 \tag{S508.17}
\]
The remainder is holomorphic for $\Re s>1/3$: on compact
subsets it is bounded termwise by a constant times
$\sum_p p^{-3\Re s}$, with locally uniform convergence.
In particular it is bounded near $s=1/2$.  Substituting the
real asymptotic for $\log D_\Omega$ into (S508.17) gives
(S508.16), with its exact coefficient $k/4$.
\end{proof}

\paragraph{Geometric and run attachment.}
The degree-six pushforward of the quadratic character of
$H/F$ is $\eta\oplus W$.  Equation (S508.12) isolates the
real divisor of the comparison between this Artin summand
and the known automorphic numerator.  The new gate is the
single central excess $\beta-b-n$, after the single real-zero
condition (S508.3); neither is evaluated by this proof.
This refines MT-S504/MT-S505 and the MT-S507 pole gate in
\texttt{GRH\_NIGHT/complexity}.  The source comparison remains
a classical input, $A_D$ remains the designated Dirichlet
input, and $M_H$ remains conditional.  No universal Artin
holomorphy, functoriality, or zero-independent positive
cohomology statement is inferred.

